\documentclass{article}

% AGENTS: We're targeting arXiv only now. The ICML conditionals are vestigial—don't maintain them.
\newif\ifarXiv
\arXivtrue

% Common packages
\usepackage{physics, mathtools, amssymb, amsmath, amsthm}
\usepackage{booktabs}
\usepackage{graphicx}
\usepackage{placeins}
\usepackage{xcolor}
\usepackage{enumitem}

\ifarXiv
  % arXiv-specific: biblatex and custom margins
  \usepackage[style=authoryear,natbib=true]{biblatex}
  \addbibresource{bibliographyml.bib}
  \usepackage[margin=1.2in]{geometry}
  \usepackage{hyperref}
\else
  % ICML-specific: natbib (loaded by icml2026) and ICML formatting
  \usepackage{hyperref}
  \usepackage[preprint]{icml2026}  % Change to [accepted] for camera-ready
\fi
\usepackage{tikz}
\usepackage{cleveref}

% Custom label for non-standard enumerate items (e.g., primed labels)
\makeatletter
\newcommand{\customlabel}[2]{%
  \phantomsection\begingroup\def\@currentlabel{#2}\label{#1}\endgroup}
\makeatother

\input{bregman-macros}
\DeclareMathOperator{\KL}{KL}
\DeclareMathOperator{\Cov}{Cov}
\let\laplacian\relax
\DeclareMathOperator{\laplacian}{\Delta}
\newcommand{\model}{\mathcal{M}}
\newcommand{\vmodel}{\mathcal{V}}
\newcommand{\F}{\mathcal{F}}
\newcommand{\G}{\mathcal{G}}
\newcommand{\hPn}{\frac{1}{n}\sum_{i=1}^n}

% Patch proof environment so \paragraph works inside it.
% amsthm's proof is a trivlist; we redefine \paragraph locally to use \item.
\usepackage{etoolbox}
\makeatletter
\newcommand{\proofparagraph}[1]{%
  \par\medskip\noindent\textbf{#1}\hspace{0.5em}\ignorespaces}
\AtBeginEnvironment{proof}{\let\paragraph\proofparagraph}
\makeatother

\newtheorem{theorem}{Theorem}[section]
\newtheorem{lemma}[theorem]{Lemma}
\newtheorem{corollary}[theorem]{Corollary}
\newtheorem{proposition}[theorem]{Proposition}
\newtheorem{definition}{Definition}[section]
\newtheorem{remark}{Remark}[section]
\newtheorem{assumption}{Assumption}[section]

%%% Variant theorem numbering: "Corollary X.YQ" and "Corollary X.YE"
%%% Inspired by amsmath's subequations environment.
\newcommand{\setupvariants}{%        % call once before the first variant
  \stepcounter{theorem}%
  \xdef\variantbase{\arabic{theorem}}%
  \addtocounter{theorem}{-1}%
}
\newcommand{\theoremvariant}[1]{%    % call before each \begin{corollary} etc.
  \setcounter{theorem}{\numexpr\variantbase-1\relax}%
  \renewcommand{\thetheorem}{\thesection.\variantbase\textup{#1}}%
}
\newcommand{\resettheorem}{%         % call after each \end{corollary} etc.
  \renewcommand{\thetheorem}{\thesection.\arabic{theorem}}%
}

% Paper-specific notation
\newcommand{\calS}{\mathcal{S}}
\newcommand{\hgammazero}{\hat\gamma_0}
\newcommand{\hL}{\hat{L}}
\newcommand{\tL}{\tilde{L}}
\newcommand{\xs}{x^\conj}             % for duality proof
\newcommand{\ys}{y^\conj}
\newcommand{\zs}{z^\conj}

% ICML running title (short form for headers)
\ifarXiv\else
  \icmltitlerunning{Balancing Weights}
\fi

\begin{document}

\ifarXiv
  % arXiv title block
  \title{Balancing Weights: When and Why They Work}
  \author{David A. Hirshberg}
  \maketitle

  \noindent\textit{Error bounds for balancing weights and Riesz representers that work when you undersmooth.}
\else
  % ICML title block
  \twocolumn[
    \icmltitle{Balancing Weights: When and Why They Work}

    \begin{icmlauthorlist}
      \icmlauthor{David A. Hirshberg}{todo}
    \end{icmlauthorlist}

    \icmlaffiliation{todo}{Affiliation TBD}

    \icmlcorrespondingauthor{David A. Hirshberg}{email@example.edu}

    \icmlkeywords{balancing weights, Riesz representers, undersmoothing, method of moments, Bregman divergence}

    \vskip 0.3in
  ]
  \printAffiliationsAndNotice{}
\fi

\begin{abstract}
High- and infinite-dimensional generalizations of the method of moments are pervasive in modern machine learning. In causal inference, balancing weights and Riesz representers are often estimated by enforcing moment balance conditions; in empirical likelihood, probability distributions are estimated via analogous moment constraints. We study a broad class of \textbf{regularized method-of-moments estimators} defined by approximately enforcing moment conditions uniformly over all functions in a seminorm ball. This formulation blends classical method-of-moments and minimum-distance ideas, induces invariance along directions in the null space of the seminorm, and naturally accommodates rich, potentially infinite-dimensional moment classes.

We show that these estimators admit a dual characterization as generalized linear models minimizing a regularized \textbf{Bregman loss}, where the index is determined by the seminorm geometry and the link function is induced by the dispersion regularizer. This perspective provides a unified view of balancing weights, Riesz representer estimation, and related procedures.

These estimators are typically tuned to prioritize control of moment imbalance over dispersion in the primal, a choice that is well suited to downstream estimators in which the resulting weights or representers are averaged and variance is attenuated while bias is not. \textbf{This deliberate undersmoothing manifests as weak regularization in the dual.} In this regime, many Bregman losses grow rapidly away from the population optimum, placing the problem outside the scope of standard concentration-based analyses. We address this challenge by establishing \textbf{uniform lower bounds} for the resulting Bregman objectives, adapting Mendelson's stable lower bounds approach to this setting. Using these, we establish nonasymptotic convergence rates for both the weights and the corresponding maximal moment imbalance in terms of localized complexity fixed points.
\end{abstract}

\ifarXiv
  \noindent\textbf{Keywords:} balancing weights, Riesz representers, undersmoothing, method of moments, Bregman divergence, stable lower bounds.
\fi

\ \\
\section{Introduction}
\label{sec:intro}

This paper establishes nonasymptotic error bounds
for a broad class of regularized method-of-moments estimators---the
class that includes balancing weights,
Riesz representer estimation,
and empirical likelihood.
We show that the bias, convergence rate,
and asymptotic variance
are all governed by a single quantity:
the \emph{Bregman-offset complexity},
a localized complexity measure
that generalizes the offset Rademacher complexity
of \citet{liang2015learning}
from quadratic to arbitrary convex losses.
Our bounds apply to models defined
via seminorms,
which induce a mix of exact and inexact balance,
and to essentially arbitrary dispersion penalties
beyond the standard quadratic and entropic cases.

The key technical challenge is
that these estimators are typically undersmoothed---tuned
to prioritize imbalance control
over dispersion---and
weak regularization translates
to dual losses that grow rapidly.
When individual terms in the sample average
can be exponentially large
even though the population average is moderate,
standard concentration arguments fail.
We address this by extending
the stable lower bounds approach
of \citet{mendelson2021extending}
from sums of squares
to sample-averaged Bregman divergences,
obtaining uniform lower bounds
that hold without concentration.

Different dispersion penalties---quadratic,
positive-part quadratic,
and entropy-like---each
have their own analytical challenges:
exponential growth, dead zones, nonlinear links.
Our framework handles all three,
showing that the choice between them
is a modeling decision, not a statistical one.
We illustrate the approach by deriving
concrete bounds for Sobolev models
and reproducing kernel Hilbert spaces.

\subsection{Estimating differentiable functionals}
\label{sec:estimation-problem}

We want to estimate a functional
$\psi = \P \psiZ(\mu)$
of an unknown function $\mu$.
\begin{definition}\label{def:differentiable}
Call $\P\psiZ$ \emph{differentiable} at $m$
on the tangent space $\mathcal{H} \subseteq L_2(\P)$
if there exist a $Z$-measurable linear \emph{derivative functional} $\dpsiZ$
and a function $\griesz \in \mathcal{H}$
such that, for each $h \in \mathcal{H}$,
\begin{equation}\label{eq:riesz-rep}
\P\psiZ(m + th)
  = \P\psiZ(m) + t \cdot \dpsiZ(h) + o(t) % o(t) is standard definitional usage here
  \qqtext{where} \dpsiZ(h) = \P \griesz(Z) h(Z).
\end{equation}
We call $\griesz$ the \emph{Riesz representer} of the derivative $\P\dpsiZ$ on $\mathcal{H}$.
\end{definition}
Both the derivative $\dpsiZ^m$ and the
representer $\gamma_{\dot\psi^m}$ depend on the
point of linearization $m$;
we drop the superscript where $m$ is
irrelevant (linear functionals)
or implied by context.
When $\mathcal{H}$ is closed in $L_2(\P)$
and $\dpsiZ$ is $L_2$-continuous,
the Riesz representation theorem
gives $\griesz \in \mathcal{H}$.
For further discussion,
see \Cref{sec:riesz-existence}.

\paragraph{Observations.}
We observe independent but not necessarily identically
distributed pairs $(Y_1, Z_1) \ldots (Y_n, Z_n)$
where $Y_i$ is a scalar outcome
and $Z_i$ is in an arbitrary set $\mathcal{Z}$.
$\hP$ is the empirical measure, i.e.
\begin{equation}\label{eq:empirical-product}
\hP f(Y,Z)
  = \frac{1}{n} \sum_{i=1}^n f(Y_i, Z_i).
\end{equation}
We define its population counterpart $\P$
using an independent copy $(\bar Y_1, \bar Z_1) \ldots (\bar Y_n, \bar Z_n)$ of our sample.
\begin{equation}\label{eq:P-def}
\P f(Y,Z)
  = \E\qty{\frac{1}{n} \sum_{i=1}^n
      f(\bar Y_i, \bar Z_i)
        \;\middle|\; (Y_1, Z_1) \ldots (Y_n, Z_n)}.
\end{equation}

$Y$ is a proxy for $\mu$
in the sense that $\mu$ is the $\P$-projection
of $Y$ onto $\mathcal{H}$:
\begin{equation}\label{eq:orthogonality}
\P h(Z) \{Y - \mu(Z)\} = 0
  \qqtext{for all} h \in \mathcal{H}.
\end{equation}
When $\mathcal{H}$ is all bounded functions of $Z$,
$\mu(Z) = \E(Y \mid Z)$ is the conditional mean
and $Y - \mu(Z)$ is orthogonal to every
function of $Z$.
More generally, $\mu$ need only be a projection
onto $\mathcal{H}$;
Example~\ref{ex:average-linear-regression}
illustrates a case where
$\mu(Z) \neq \E(Y \mid Z)$.

\paragraph{The one-step estimator.}
Given an estimator $\hat\mu$ of $\mu$,
the plug-in $\hP \psiZ(\hat\mu)$
is generally biased.
The \emph{one-step estimator} corrects it
with a weighted average of the
proxy residuals $Y - \hat\mu(Z)$.
\begin{equation}\label{eq:one-step}
\hat\psi
  = \hP \psiZ(\hat\mu)
  + \hP \hgamma(Z) \{ Y - \hat\mu(Z) \}
  \qqtext{where hopefully}
  \hP \hgamma(Z) f(Z) \approx \dpsiZ(f)
  \text{ for } f = \mu - \hat\mu.
\end{equation}
If our hopes bear out,
$\hat\psi \approx \hP\psiZ(\hat\mu) + \dpsiZ(\mu - \hat\mu)$,
the first-order approximation
to the target $\psi$.
The weights $\hgamma$ that we're hoping for
approximate the functional's action---at
least on $f = \mu - \hat\mu$---as
its Riesz representer $\griesz$ does.

\subsubsection{Example: Estimating a Treatment-Specific Mean}
\label{ex:treatment-specific-mean}
Consider a study with selection on observables.
We observe independent triples $(W_i,X_i,Y_i)$ for a
treatment indicator $W_i$, covariate vector $X_i$, and
outcome $Y_i$ and take $Z_i=(W_i,X_i)$ to be the
treatment/covariate pair.
The nuisance is $\mu(w,x)=\E(Y \mid W=w,X=x)$
and the tangent space $\mathcal{H}$ is all bounded
functions of $(W,X)$.
The derivative functional $\dot\psi^w_Z(m)=m(w,X)$
picks out treatment $w$, so
$\E \dot\psi^w_Z(m)=\psi^w(m)$, and the estimand is
\begin{equation}\label{eq:tsm-estimand}
\psi^w(\mu) = \E \mu(w,X).
\end{equation}
Here $\E(Y \mid Z) = \mu(Z)$, so $Y$ is a valid proxy.
The Riesz representer is the inverse propensity weight
$\gamma_{\dot\psi^w}(W,X) = \ind{W=w}/\pi_w(X)$
for the propensity score
\begin{equation}\label{eq:tsm-riesz}
\pi_w(x) := \P(W=w \mid X=x).
\end{equation}
It solves
$\E \gamma_{\dot\psi^w}(W,X) m(W,X) = \E m(w,X)$.
This justifies IPW estimation: by iterated expectations,
$\E \gamma_{\dot\psi^w}(Z) Y
  = \E \gamma_{\dot\psi^w}(Z) \mu(Z) = \psi^w(\mu)$.
Other estimators like AIPW and TMLE also involve
these weights, either explicitly or as components
of their bias-correction terms.

\subsubsection{Example: Average Linear Regression Function}
\label{ex:average-linear-regression}
In the same setting, now let $W$ be a continuously-valued
treatment and consider the slope coefficient from the
regression of $Y$ on $W$ in the subpopulation with
covariates $X = x$,
\begin{equation}\label{eq:cond-slope}
b_0(x) = \frac{\Cov(Y, W \mid X=x)}
              {\Var(W \mid X=x)}.
\end{equation}
The estimand is the average coefficient
$\psi(\mu) = \E b_0(X)$
\citep{graham2022semiparametrically}.
The nuisance is the conditional linear predictor
$\mu(W,X) = a_0(X) + W \cdot b_0(X)$,
and the tangent space $\mathcal{H}$ is the space
of functions linear in $W$:
$h(W,X) = \alpha(X) + W \cdot c(X)$.
When $W$ is binary, $b_0(X)$ is the
conditional average treatment effect
and $\psi(\mu)$ is the average treatment effect.
When $W$ is continuous,
$\mu$ is a projection of $Y$ onto $\mathcal{H}$
rather than the conditional mean:
$\mu(Z) \neq \E(Y \mid Z)$ in general.

\subsubsection{Example: Variance of the Conditional Average Treatment Effect}
\label{ex:vcate}
In the same setting as Example~\ref{ex:treatment-specific-mean},
the \emph{conditional average treatment effect}
$\tau(X) = \mu(1,X) - \mu(0,X)$
measures how the treatment effect varies with covariates.
The \emph{variance of the CATE}
\begin{equation}\label{eq:vcate}
V(\mu) = \Var(\tau(X))
       = \E[\tau(X)^2] - (\E\tau(X))^2
\end{equation}
quantifies treatment effect heterogeneity.
We assume $V > 0$ throughout
(constant effects are a degenerate case
with a faster-than-parametric rate).
The squared ATE $(\E\tau)^2$
is a function of the linear functional
from Example~\ref{ex:treatment-specific-mean};
the new content is the nonlinear component
$\E[\tau(X)^2]$.
Taking $\psiZ(m) = (m(1,X) - m(0,X))^2$
gives $\P\psiZ(\mu) = \E[\tau(X)^2]$.
The derivative at $\mu$ in direction $h$
\begin{equation}\label{eq:vcate-derivative}
\dpsiZ^\mu(h)
  = 2\,\E\bigl[
      \tau(X)\,
      (h(1,X) - h(0,X))
    \bigr]
\end{equation}
depends on $\mu$ through the CATE,
and the Riesz representer is
\begin{equation}\label{eq:vcate-riesz}
\griesz^\mu(W,X)
  = 2\tau(X)
    \qty\Big{\frac{\ind{W=1}}{\pi_1(X)}
         - \frac{\ind{W=0}}{\pi_0(X)}}.
\end{equation}
This is the ATE weight
modulated by the CATE:
the treatment-specific mean asks
``what is the average effect?'';
the variance of the CATE asks
``how does the effect vary?'',
and the representer
upweights individuals
whose treatment effect is large.

\subsection{Estimating Riesz Representers}
\label{sec:estimating-riesz}

We solve an optimization problem that encompasses many estimation problems, from empirical likelihood to estimation of a covariate-balancing propensity score in causal inference.
\begin{equation}
\label{eq:primal-moment-matching}
\begin{aligned}
&\hgamma := \argmin_{g} L_{\hP}(g) \qfor
L_{\hP}(g) =  \zeta\{b(g)\} + \hP (\chiZ \comp g)(Z) \\
&\quad \text{where} \ b(g) := \sup_{m \in B_\rho} \hP g(Z)\, m(Z) - \dotpsiZ(m), \\
&\quad \chiZ \ \text{is convex for each $Z$}, \ \zeta \ \text{is nondecreasing and convex, and} \ B_\rho = \{m : \rho(m) \le 1\}.
\end{aligned}
\end{equation}
The first term is a function of our \emph{maximal imbalance}
$b(g)$---the
largest in-sample violation of the moment conditions
defining our Riesz representer
for functions in our \emph{model} $B_\rho$.
The second term penalizes dispersion of the weights;
$\zeta$ controls the tradeoff.
It's convenient to analyze the estimator
through an equivalent dual problem,
which characterizes the weights
as following a generalized linear model.%
\footnote{We use the term GLM in the expansive sense:
a model of the form
$\ell(\calS) := \{ \ell \comp f \ : f \in \calS \}$
where $\ell$ is called the \emph{link}
and $\calS$---a linear subspace of functions---is
called the \emph{index class}.
This includes the case that $\calS$ is the space
of linear-in-covariates functions
$\{ f(z)=z^T \beta \ : \ \beta \in \R^p \}$
and, more generally, the case that
$\calS=\{ f \ : \ \rho(f) < \infty \}$
for any seminorm $\rho$.}
$\hgamma = \dchisZ \comp \hat\phi$.
The corresponding \emph{index} $\hat\phi$
is chosen from
the index class $V_\rho := \{f : \rho(f) < \infty\}$
and the \emph{link} $\dchisZ$
is the derivative
of our dispersion penalty's convex conjugate.
The index solves this dual:
\begin{equation}
\begin{aligned}
\label{eq:dual-problem}
\hat\phi &:= \argmin_{f} L^\conj_{\hP}(f) \qfor
L^\conj_{\hP}(f) := \hP (\chisZ \comp f)(Z) - \dotpsiZ(f) + (\zetas \comp \rho)(f) \\
&\quad \text{where} \ \chisZ \ \text{and} \ \zetas \ \text{are the convex conjugates of} \ \chiZ \ \text{and} \ \zeta.
\end{aligned}
\end{equation}
The problem may appear more familiar
when we look at the loss difference
rather than the loss itself, i.e.
\begin{equation}
\label{eq:loss-difference}
\begin{aligned}
L^\conj_{\hP}(f) - L^\conj_{\hP}(\tilde\phi)
&=  \hP D_{\chisZ}\{f(Z), \tilde\phi(Z)\} + \hP (\dchisZ \comp \tilde\phi)(Z) \cdot \{f(Z) - \tilde\phi(Z)\} - \dotpsiZ( f - \tilde\phi ) \\
&\qquad + (\zetas \comp \rho)(f) - (\zetas \comp \rho)(\tilde\phi)%
\footnote{When $\zetas$ is non-differentiable, $D_{\zetas \comp \rho}$ is not uniquely defined. We use a subgradient-based definition that satisfies the key identity \eqref{eq:loss-difference-with-foc}; see \Cref{sec:rate-of-convergence} for details.}
\\
&\qquad\text{for}\quad D_{\chisZ}(b, a) = \chisZ(b) - \chisZ(a) - \dchisZ(a)(b - a).
\end{aligned}
\end{equation}
Here $f \comp g$ denotes the composition $(f \comp g)(z)=f\{g(z)\}$
and $\dchisZ$ and $\ddchisZ$ denote the first and second derivatives of $\chisZ$.
The \emph{Bregman divergence} $D_{\chisZ}$ is the error
of the first-order Taylor approximation to $\chisZ$,
which is approximately quadratic near $\tilde\phi$:
$D_{\chisZ}(f, \tilde\phi) \approx (\ddchisZ \comp \tilde\phi)(f - \tilde\phi)^2/2$
for $f \approx \tilde\phi$.

\paragraph{The Index-Scale Limit.}
The dual solution $\hat\phi$ converges
to a deterministic limit $\tilde\phi$.
Because $\dchisZ$ is the link,
the second term in the loss difference
is a mean-zero `noise-like process.'
But to get a non-vacuous result
when $\rho(\griesz)=\infty$,
we take $\tilde\phi$ to be the minimizer
of the population loss $L^\conj_{\P}$,
which is a \emph{regularized Bregman projection}
of $\griesz$,
and centers the same term `almost for free.'
We will be more precise
in \Cref{sec:results} below.
Until then, there is little harm
in thinking of $\dchisZ \comp \tilde\phi$
as $\griesz$,
which happens to satisfy $\rho(\griesz)<\infty$.

\paragraph{The Weight Scale.}
What we set out to estimate were the weights $\hgamma$,
not the indices $\hat\phi$.
Bregman duality relates our error
in estimating the index
to an error in estimating
the corresponding weights.
\begin{equation}\label{eq:bregman-duality}
D_{\chisZ}(\hat\phi, \tilde\phi) = D_{\chiZ}(\tilde\gamma, \hgamma)
\qfor \tilde\gamma = \dchisZ \comp \tilde\phi.
\end{equation}
For the quadratic dispersion penalty,
the mapping is simple:
the weight is the index
and averaged Bregman divergence on both scales
is the sample two-norm of the difference.
For the commonly-used entropy penalty,
the situation is more complex---our
weight convergence is in
sample Kullback--Leibler divergence,
$\smash{\KL_{\hP}(\tilde\gamma, \hgamma)
:= \hP\, \tilde\gamma(Z) \log\{\tilde\gamma(Z)/\hgamma(Z)\}
- \hP\{\tilde\gamma(Z) - \hgamma(Z)\}}$.
We summarize the moving parts
in the following table,
writing $h := \hat\phi - \tilde\phi$
for the index-scale error.
\begin{center}
\begin{tabular}{l c c c c c}
\toprule
& $\chiZ$ & $\chisZ$ & link $\dchisZ$
& index divergence
& weight divergence \\
\midrule
quadratic & $\gamma^2/2$ & $\phi^2/2$ & $\phi$
& $\|h\|_{\LtPn}^2/2$
& $\|\tilde\gamma - \hgamma\|_{\LtPn}^2/2$ \\
entropy & $\gamma \log \gamma - \gamma$ & $e^\phi$ & $e^\phi$
& $\hP\, e^{\tilde\phi(Z)}\{e^{h(Z)} - 1 - h(Z)\}$
& $\KL_{\hP}(\tilde\gamma, \hgamma)$ \\
\bottomrule
\end{tabular}
\end{center}
With the $\LtPn$ convergence
we get in the quadratic case,
we can show the estimator's variance
converges to the optimal variance
given only that the Riesz representer
is square-integrable---an
overlap condition that is necessary
for a regular estimator to exist.
But convergence in KL divergence,
as we get in the entropy case,
is substantially weaker.
Nonetheless, we show that
with a mild increase in overlap,
we can arrive at the same conclusion
with smoothness assumptions similar
to the ones we need for the analysis
of bias and weight convergence
(\Cref{sec:variance}).

\subsection{Modeling considerations}
\label{sec:modeling}

The triple $(\rho, \chi, \zeta)$ encodes modeling choices about smoothness, weight structure, and tuning. Of these, $\rho$ is the most consequential and $\zeta$ the least.

\paragraph{Choosing $\rho$: seminorms and balance.}
The choice of $\rho$ determines how balance is enforced. Using a seminorm $\rho$ unifies two common formulations.
\begin{enumerate}
\item When $\rho$ is dichotomous---zero on some vector space $K$ and infinite elsewhere---balance is exact: $B_\rho = K$ and we require $\hP g(Z)m(Z) - \dotpsiZ(m)=0$ for all $m \in K$. Here we either use a subspace of dimension $n$---one moment condition per weight---as in classical method of moments, or we use a lower-dimensional one, necessitating the use of some other criterion like our \emph{dispersion penalty} $\hP (\chiZ \comp g)(Z)$ to act as a tiebreaker. Concrete realizations of the latter include empirical likelihood \citep{owen2001empirical}, entropy balancing \citep{hainmueller2012entropy}, inverse probability tilting \citep{graham2012inverse}, and the covariate balancing propensity score \citep{imai2014covariate}. Nonparametric extensions use sieves---a sequence of subspaces $K_n$ growing with $n$---typically chosen on the basis of approximation theory \citep[e.g.,][]{chan2016globally}.
\item When $\rho$ is a norm---so its kernel is trivial---balance is approximate: minimizing $L_{\hP}$ asks for a compromise between imbalance and dispersion \citep{zubizarreta2015stable, kallus2020generalized}. The ball can even be infinite-dimensional, allowing fully nonparametric estimation without sieves.
\end{enumerate}
A seminorm with nontrivial kernel combines both: functions in $\ker\rho$ are balanced exactly, while the quotient $\calS_\rho / \ker\rho$ is balanced approximately.

Many common constraints on weights can be understood in terms of kernels of seminorms. Exact balance for $k \in \ker\rho$ gives $\hP \hgamma(Z) k(Z) = \dotpsiZ(k)$. When the kernel contains the constant function, this gives $\hP \hgamma(Z) = \psiZ(1)$, so in Example~\ref{ex:treatment-specific-mean} our weights average to one---a constraint often imposed via post-hoc normalization. In Example~\ref{ex:treatment-specific-mean}, weights that average to one give an estimator that is equivariant to shifts in the outcome---the well-known H\'ajek estimator. The seminorm perspective is the key to expressing desiderata like this in a way that generalizes across estimands; we develop this in our worked examples (\Cref{sec:worked-examples}). We can derive a seminorm with any prescribed kernel from a norm $\norm{\cdot}$ via minimization: $\rho(m) := \inf_{k \in K}\norm{m - k}$.

\paragraph{Choosing $\chi$: dispersion penalties and effective domain.}
The effective domain of $\chiZ$---the set where it is finite---constrains which weights are allowed. Different problems call for different constraints. In Example~\ref{ex:treatment-specific-mean}, the Riesz representer $\griesz$ is zero for treatment groups $W\neq w$ and nonnegative generally, and we might want to constrain our weights $\hgamma$ to have these properties. The resulting weights come from a logit model for the propensity score, fit by balancing rather than maximum likelihood. When used in combination with a translation-invariant seminorm, so our weights average to one, the weighted average $\hP \hgamma(Z) Y$ will be a convex combination of the treated group outcomes $\{ Y_i \ : \ W_i = w\}$. In practice it will almost always be within their range; the constraints make this a guarantee rather than an observation.

One choice is to start with the quadratic penalty $\chiZ(\gamma)=\gamma^2/2$ and impose the constraints by adding a term that blows up to infinity when they are violated.
\begin{equation}
\chiZ(\gamma) :=
\begin{cases}
  \gamma^2/2 + \mathcal{I}(\gamma \ge 0) & \text{ if } \  W(Z)=w \\
  \mathcal{I}(\gamma=0)                  & \text{ otherwise. }
\end{cases}
\qwhere \mathcal{I}(A)=\begin{cases} 0 & \text{if } A \text{ holds} \\ +\infty & \text{otherwise.}\end{cases}
\end{equation}
Because $\dchis(\phi) = \phi_+ := \max(\phi,0)$, this approach, proposed by \citet{zubizarreta2015stable}, gives weights $\hgamma(Z_i) = \hat\phi(Z_i)_+$. That is, it gives us inverse propensity weights based on a generalized linear model for the propensity score with reciprocal positive part link: $\hat \pi_w(\cdot) = \ell_\chi \comp \hat \phi$ for $\ell(f)=1/f_+$.

Another option, with effective domain $\R_{\ge 1} := \{ \gamma : \gamma \ge 1 \}$ matching the natural range of inverse propensity weights $\griesz\{(W,X)\} = \ind{W=w}/\pi_w(X) \ge 1$, is this entropy-like dispersion penalty:
\begin{equation}
\chiZ(\gamma) =  \begin{cases}
    (\gamma-1)\{\log(\gamma-1) -  1\} + \mathcal{I}(\gamma - 1 \ge 0) & \text{if} \ W(Z)=w \\
    \mathcal{I}(\gamma=0)                                         & \text{otherwise}
\end{cases}
\end{equation}
With this penalty, $\hgamma(w,\cdot) = \dchisZ(\hat\phi) = 1 + e^{\hat\phi}$,
so the propensity score is $\smash{1/\hgamma(w,\cdot) = 1/(1+e^{\hat\phi})}$---a logit model.
Entropy balancing, as introduced by \citet{hainmueller2012entropy}, takes a similar approach to estimating the ATT. In that context, the Riesz representer is the odds $\smash{\griesz(-1,\cdot) = \pi_1/(1-\pi_{1})}$ for control units. In our dual formulation, $\hgamma(-1,\cdot) = e^{\hat\phi}$, so the propensity is $\smash{e^{\hat\phi}/(1+e^{\hat\phi})}$---again a logit model.

Given the same effective domain, the choice of $\chi$ is a second-order modeling decision. With nonparametric index classes, we can choose an index that `inverts the link' to give an equivalent model. For Sobolev or H\"older seminorms, when the two links have similar curvature, the reparametrized index $f_b$ has similar regularity to the original $f_a$---$\rho(f_b) \approx \rho(f_a)$---because
\[
\dchis_a \comp f_a = \dchis_b \comp f_b \qfor f_b=(\dchis_b)^{-1} \comp \dchis_a \comp f_a
\qqtext{with} \dot f_b = \frac{\ddchis_a \comp f_a}{\ddchis_b \comp f_b} \cdot \dot f_a.
\]
As we will see in \Cref{sec:statistical-considerations}, the choice is often similarly inconsequential for the statistical behavior.

\paragraph{Choosing $\zeta$.}
\label{sec:zeta}
The choice of $\zeta$ is the least consequential. Different choices of $\zeta$ are equivalent, up to tuning parameters, to constraints on $\rho$ via Lagrange multipliers. If tuning is done by cross-validation or another automated procedure%
\footnote{See \citet{cui2024selective} for automated selection of tuning parameters via cross-validation in semiparametric settings.}%
, the choice matters only insofar as grids in different parameterizations correspond to different grids in the constrained form---ideally not at all and in practice relatively little.

\section{Context}
\label{sec:context}
We organize the context around the three ingredients
of the inference program.
The one-step estimator's error decomposes as
\begin{equation}\label{eq:decomposition}
\begin{aligned}
\hat\psi - \psi
&= \P\psiZ(\hat\mu) + \dpsiZ(\mu - \hat\mu) - \psi
  && \text{(linearization error)} \\
&\quad+ \hP \hgamma(Z)\{\mu(Z) - \hat\mu(Z)\}
       - \dpsiZ(\mu - \hat\mu)
  && \text{(imbalance)} \\
&\quad+ \hP \hgamma(Z)\varepsilon
       + (\hP - \P)\psiZ(\hat\mu) \qfor \varepsilon = Y - \mu(Z).
  && \text{(noise average)}
\end{aligned}
\end{equation}
What we hope for is a $t$-statistic
$(\hat\psi - \psi)/\hat\sigma$
that is approximately standard normal.
Essentially, this requires three ingredients:
negligible linearization error,
negligible imbalance (bias control),
and a well-behaved noise average (weight convergence).
For each, we describe the problem
and the existing approaches,
leading to the tools our contribution builds on.

\subsection{Negligible linearization error}
\label{sec:linearization}
The linearization error
$\P\psiZ(\hat\mu) + \dpsiZ(\mu - \hat\mu) - \psi$
measures how well the first-order approximation
$\psi + \dpsiZ(\mu - \hat\mu)$ captures
the full nonlinearity of $\psiZ$.
In our examples, $\psiZ$ is linear
and this term vanishes exactly.
For nonlinear functionals,
it is controlled by smoothness of the functional
and the rate at which $\hat\mu \to \mu$.

\subsection{Imbalance and bias control}
\label{sec:imbalance-bias}
The imbalance measures how well
the weights $\hgamma$ approximate
the action of $\dpsiZ$---at least on
the error $\mu - \hat\mu$.
Controlling it---making
the bias $b(\hgamma)$ negligible---is
the central analytical challenge;
once it is negligible,
asymptotic normality follows
from design-conditional arguments
\citep{benmichael2021balancing,pham2023stable}.

\paragraph{Maximal imbalance.}
Because $(\hat\mu - \mu)/\rho(\hat\mu - \mu) \in B_\rho$,
a uniform bound on the maximal imbalance
$\sup_{m \in B_\rho}
  |\hP \hgamma(Z) m(Z) - \dotpsiZ(m)|$
controls the imbalance on $\hat\mu - \mu$
up to the factor $\rho(\hat\mu - \mu)$.
A natural approach is to bound
$b(\hgamma)$ directly
by controlling imbalance
over a function class.
Early work in this direction
\citep{athey2018approximate,kallus2020generalized}
bounded the maximal imbalance
at rate $n^{-1/2}$
over the full model,
then required convergence
of the outcome estimator
$\hat\mu$ in the model norm
so that the product
of imbalance and model error
is negligible.
\citet{hirshberg2018augmented}
refined this
using the equicontinuity
of the empirical process
to achieve imbalance $\ll n^{-1/2}$
on a localized set,
replacing model-norm convergence
with $L_2$ consistency.
Both approaches, however,
rely on a comparison
to the population-optimal weights,
and therefore require
convergence of $\hgamma$.
This limits the degree
of undersmoothing allowed.
\citet{kong2025asymptotics}
removed this limitation
in the quadratic case
$\chis = x^2/2$
by bounding the bias
via the Donoho--Liu modulus of continuity
over the full norm ball $B_\rho$,
without weight convergence.
Whether this extends
beyond the quadratic case---where
the modulus involves a Bregman divergence,
not a sum of squares---is not obvious.

\paragraph{Nonuniform bias bounds.}
When both the weights and outcome model converge,
the one-step estimator's remainder,
after centering at the population Riesz representer,
decomposes as
$\hP(\hgamma - \griesz)(\hat\mu - \mu)$
plus a linearization term.
By Cauchy--Schwarz,
this is bounded by
$\|\hgamma - \griesz\|_{L_2(\hP)}
\cdot \|\hat\mu - \mu\|_{L_2(\hP)}$:
a product of $L_2$ estimation errors
for the weights and outcome model,
each converging at its natural rate.
When this product is $o_P(n^{-1/2})$,
the bias is negligible without undersmoothing.
When the weights and outcome model
are estimated by projection---least
squares onto a basis,
or regularized regression
in an RKHS---the
first-order conditions
make the estimation errors
orthogonal to the fitted subspace,
sharpening the bound further.
This mechanism appears
in finite-dimensional models
\citep{vermeulen2015bias,avagyan2017honest,tan2010bounded},
RKHS settings
\citep{hirshberg2018augmented},
and sparse models
where the active set
plays the role of the basis
\citep{bradic2019sparsity,bradic2019minimax}.
These approaches exploit the specific structure
of $\hat\mu$ to get sharper bounds,
but do not apply when
the structure of $\hat\mu$ is unknown
or when weight convergence fails.

\subsection{Well-behaved noise average}
\label{sec:noise-convergence}
When we have enough of the first two ingredients,
i.e.\ negligible imbalance and linearization error,
what's left is the noise average
$\hP \hgamma(Z) \varepsilon$---the numerator of a $t$-statistic.
The condition for asymptotic normality is weak\footnotemark:
it's sufficient that our weights $\hgamma(Z_i)$ are infinitesimal---
that no single weight dominates
the sum of squares $\sum_i \hgamma^2(Z_i)$---when
we use the natural variance estimate
\begin{equation}\label{eq:variance-estimate}
\hat\sigma^2
  = \frac{1}{n}\hP \qty[\hgamma(Z)\qty{Y - \hat\mu^\sigma(Z)}
    + \psiZ(\hat\mu^\sigma) - \hat\psi]^2.
\end{equation}%
\footnotetext{This applies when the weights
are estimated without the use of $Y_1 \ldots Y_n$---``design-based''
weights \citep{benmichael2021balancing}---as
in Examples~\ref{ex:treatment-specific-mean}--\ref{ex:average-linear-regression}.
More precisely,
the form of weight convergence required is
$\hP \hgamma^2 \nu / \P \tilde\gamma^2 \nu \to 1$
for $\nu(Z) = \Var(Y \mid Z)$
and population regularized weights
$\tilde\gamma = \dchis \comp \tilde\phi$.
Whether the resulting variance
equals the semiparametric efficiency bound
depends on whether
$\P\tilde\gamma^2\nu \to \P\griesz^2\nu$,
which requires $\griesz \in L_2(\P)$
and $\tilde\gamma \to \griesz$ in $\LtP$
(\Cref{cor:efficient-variance}).}

Bregman convergence ensures
$\hP \hgamma^2 / \P \tilde\gamma^2 \to 1$
(\Cref{lem:variance-equiv}),
so the self-normalized CLT holds
with variance proportional to
$\P\tilde\gamma^2\nu$---always finite.
Whether this variance is \emph{optimal}
depends on whether
$\P\tilde\gamma^2 / \P \griesz^2 \to 1$---whether
the regularized weights' sum of squares
converges to that of the Riesz representer.
In the quadratic case,
this follows from the Cauchy--Schwarz inequality
and square-integrability of the Riesz representer, as
\begin{equation*}
\begin{aligned}
\P \tilde\gamma^2 - \P \griesz^2
&= \P(\tilde\gamma - \griesz)^2
+ 2\P \griesz(\tilde\gamma - \griesz) \\
&\le \|\tilde\gamma - \griesz\|_{\LtP}^2
+ 2\|\griesz\|_{\LtP}\,
\|\tilde\gamma - \griesz\|_{\LtP}
\to 0 \qqtext{as} \tilde\gamma \to \griesz
  \text{ in } \LtP.
\end{aligned}
\end{equation*}

\paragraph{Why this is hard.}
In the entropy case,
the whole pipeline is more involved.
Establishing convergence
in Bregman divergence itself is harder.
The exponential link amplifies
large index values---for
$\chisZ(\phi) = e^\phi$,
the Bregman divergence is
$D_{\chisZ}(\tilde\phi+h, \tilde\phi) = e^{\tilde\phi}\{e^{h} - 1 - h\}$,
which grows like $e^h$ for large $h$.
Undersmoothing allows
$\rho(h)$ as large as $\sqrt{n}$
for the index error $h = \hat\phi - \tilde\phi$.
Errors like this can be `spiky:'
$h(Z_i)$ can be large in some coordinates
even when it's on a small $\LtP$-sphere
and for these coordinates Bregman divergence will be exponentially large.
And even with Bregman convergence in hand,
converting from KL divergence
to the sum-of-squares ratio
requires additional conditions and additional work---the
Cauchy--Schwarz argument above
has no immediate analogue when the link is nonlinear
(\Cref{sec:variance}).

But as \citet{mendelson2015learning} observed,
we don't need to know
how large the Bregman divergence is---only
\emph{how small it isn't}.
Rather than concentration,
we need a \emph{uniform lower bound}:
a guarantee that
$\hP D_{\chisZ}\{f(Z), \tilde\phi(Z)\}$
exceeds some threshold
for all $f$ on the sphere.
The machinery for obtaining
uniform lower bounds of this kind
comes from two streams of work
in nonasymptotic learning theory.

\paragraph{Stable lower bounds.}
The small-ball method
of \citet{mendelson2015learning}
replaces concentration inequalities
with anticoncentration:
to bound the excess risk of a learner,
it suffices to show that
the loss increments are bounded below
on most coordinates,
even when they are unbounded above.
\citet{mendelson2021extending}
developed this into a general framework
of \emph{stable lower bounds}:
uniform lower bounds
on sample-averaged sums of squares
that survive dropping
a small fraction of coordinates.
These bounds give convergence rates
for empirical risk minimizers
under heavy-tailed losses,
where classical concentration arguments
--- which bound the supremum
of the empirical process ---
break down.
The existing framework handles sums of squares.
For general dispersions $\chis$, however,
the loss difference at the optimizer
produces the Bregman divergence
$D_{\chisZ}\{\hat\phi(Z), \tilde\phi(Z)\}$,
not a sum of squares,
and working directly
in the divergence's natural geometry
avoids an $L_2$--Bregman correspondence
that breaks down
when the link function is nonlinear.

\paragraph{Offset complexity.}
\citet{liang2015learning}
introduced offset Rademacher complexity
as a localization device
for learning with square loss.
The key idea is that
the sample-averaged quadratic penalty
$\hP h^2$
can replace the population penalty
$\P h^2$ used in
classical local Rademacher complexity
\citep{bartlett2005local},
because the offset --- the
difference between sample and population
averages --- is itself controlled
by the empirical process.
For convex function classes,
this yields fast rates
without a preliminary estimation step.
The natural generalization replaces
the quadratic penalty
with a Bregman divergence,
recovering the offset Rademacher complexity
when $\chis = x^2/2$.
\citet{vijaykumar2021localization}
extends the offset approach
to general losses
satisfying a uniform convexity condition
that encompasses both
exponential concavity
and self-concordance,
recovering optimal rates
for the $p$-loss ($1 < p < \infty$)
and for logistic regression.

\subsection{The quadratic baseline}
\label{sec:quadratic-baseline}

The augmented minimax linear estimator (AMLE) and automatic debiased machine learning (Auto-DML) carry out the full inference program in the quadratic case \citep{hirshberg2018augmented,chernozhukov2022automatic}. Both are instances of the one-step estimator \eqref{eq:one-step}. The AMLE's minimax weights and the lasso-minimum-distance weights (LMD) \citep{chernozhukov2022locally} often used in Auto-DML solve the primal and dual problems \eqref{eq:primal-moment-matching}--\eqref{eq:dual-problem} using quadratic dispersion $\chi_Z(g) = g^2/2$. The corresponding link function is the identity $\dchis_Z(f)=f$, so the weights are $\hgamma=\hat\phi$ where
\begin{equation}
\begin{aligned}
\hgamma &:= \argmin_g L_{\hP}(g) \qfor L_{\hP}(g) := \frac{1}{2\eta} \qty{\sup_{m \in B_\rho} \hP g(Z) m(Z) - \dotpsiZ(m)}^2 + \hP \frac{g(Z)^2}{2}, \\
\hat\phi &:= \argmin_f L^\conj_{\hP}(f) \qfor L^\conj_{\hP}(f) := \hP \frac{f(Z)^2}{2} - \dotpsiZ(f) + \frac{\eta \rho(f)^2}{2}.
\end{aligned}
\end{equation}
The methods differ in the way imbalance
is scaled relative to dispersion.
AMLE uses the quadratic scaling $\zeta(\rho)=\rho^2/(2\eta)$%
\footnote{In AMLE's notation, $\eta = n/(2\sigma^2)$.} while LMD uses a constraint-inducing $0/\infty$ penalty $\zeta = \mathcal{I}(\cdot \le s)$ for some $s \ge 0$. In the dual, the latter translates to the scaled seminorm penalty $\zetas \comp \rho = s^{-1}\rho$. The `lasso' in the name comes from the use of the seminorm
\begin{equation}
\label{eq:lasso-seminorm}
\rho(f) =
\begin{cases}
  \norm{\beta}_1 &\qif f(\cdot)=\sum_{j=0}^p \beta_j b_j(\cdot) \\
  \infty         &\qif \text{otherwise}
\end{cases}
\ \text{ for some set of features } b_0 \ldots b_p.
\end{equation}

The choice of quadratic $\chi$ here---or at least a dispersion penalty whose effective domain is all of $\R$---is crucial. It's  what allows them to be universal. And in the same vein, universal recommendations for $\rho$ must be generic: without knowing the problem, you can't tailor a seminorm to model $f=\mu-\hat\mu$. More tailored methods often specify other choices
that (i) constrain $\hgamma$ to a set
that $\griesz$ is either in or close to, e.g. positive weights averaging to one in Example~\ref{ex:treatment-specific-mean}
and (ii) give the resulting estimator $\hat\psi$ desirable \emph{boundedness} and \emph{equivariance} properties. These choices arise in each of our worked examples (\Cref{sec:worked-examples}).

\paragraph{Minimax linear estimation.}
When $\psiZ$ is linear, there is usually little harm in using a simpler Riesz-weighting (RW) estimator: the one-step estimator \eqref{eq:one-step} with $\hat\mu = 0$:
\[\hat\psi_{\text{RW}} = \hP \hgamma(Z) Y.\]
Classical results \citep{donoho1994statistical} establish that the minimax linear estimator---the AMLE for $\hat\mu=0$---is nearly minimax optimal for estimating linear functionals under convex constraints
like $\hat\mu \in B_\rho$. And recent work \citep{kong2025asymptotics} shows that RW is asymptotically efficient --- i.e. exactly optimal in a local asymptotic minimax sense --- under essentially the same assumptions under which \citet{hirshberg2018augmented} showed efficiency of the AMLE.

The circumstances these results do not cover are ones in which $\mu$ is believed to satisfy some \emph{non-convex} constraint and $B_\rho$ is a smaller convex set. The classic example is $\mu$ satisfying a sparsity constraint $\hat\mu(\cdot)=\sum_j \beta_j b_j(\cdot)$ for $\norm{\beta}_{\ell_0} \le s_0$, which --- if the matrix of features satisfies a restricted eigenvalue condition ---
implies that a lasso estimator's error function $\hat\mu-\mu$ has coefficients with small $\ell_1$ norm, i.e. is in a seminorm ball $s_1 B_\rho$ of small radius $s_1$ for $\rho$ as in \eqref{eq:lasso-seminorm}.  In this context, these theoretical results are uninformative and empirical evidence favors augmented approaches like the AMLE \citep{athey2018approximate}.

\section{Contribution}
\label{sec:contribution}

We carry out the semiparametric inference program---bias
control, weight convergence,
and asymptotic efficiency---for
one-step estimators using balancing weights.
Our results cover a large class
of seminorm-based models
and dispersion penalties
including all popular options in the literature.
We do so in a way that decouples
questions of bias and weight convergence,
permitting a clear understanding
of the effects of undersmoothing.
The Bregman-offset complexity emerges
as the single crucial parameter;
we develop tools to bound it
in RKHS models.
We work through the application
of this approach to estimating
the average treatment effect
and average linear regression coefficient,
the latter involving a nontrivial conditioning structure
rarely discussed in the literature
on balancing approaches.

% --- old contribution paragraphs (saved) ---
% \paragraph{Strong duality}
% (\Cref{prop:duality}).
% The primal problem~\eqref{eq:primal-moment-matching}
% and dual problem~\eqref{eq:dual-problem}
% have the same value,
% and their solutions are linked
% by $\hgamma = \dchis \comp \hat\phi$,
% for any dispersion $\chis$
% and penalty $\zetas \comp \rho$
% satisfying mild regularity conditions.
%
% \paragraph{Bias bounds}
% (\Cref{thm:bias-decomp}).
% The worst-case bias at the estimated weights
% is controlled, on both sides,
% by the offset Rademacher complexity.
% No weight convergence is needed.
% In the undersmoothing limit,
% the offset complexity sharpens
% to the local Rademacher complexity
% of $B_\rho$ at $\griesz$.
%
% \paragraph{Rates and variance}
% (\Cref{thm:rate-main},
% Lemma~\ref{lem:variance-equiv}).
% Uniform lower bounds on sample-averaged
% Bregman divergences establish
% convergence of the dual solution
% $\hat\phi \to \tilde\phi$
% at explicit rates
% for general convex losses.
% A comparison lemma
% (\Cref{lem:variance-equiv})
% shows that convergence
% in Bregman divergence
% is strong enough to stabilize the variance:
% $\hP D_{\chisZ}\{\hat\phi(Z), \tilde\phi(Z)\} \to 0$
% implies
% $\hP \hgamma^2 / \hP \griesz^2 \to 1$.
% For entropy-like dispersions,
% this requires an interpolation argument
% to control the nonlinear link.
%
% \paragraph{Implications.}
% [old version had treatment-specific means instead of
%  average linear regression coefficients]
% --- end old ---

\paragraph{Bregman-offset complexity.}
The offset complexity
\begin{equation}\label{eq:offset-complexity-intro}
\Xi(s) = \sup_{h \in B_\rho}
  (\hP - \P)\{\dotpsiZ(h) - \tilde\gamma(Z)\, h(Z)\}
  - \frac{1}{s}\, \hP\, D_{\chisZ}\{\tilde\phi(Z) + s\,h(Z),\, \tilde\phi(Z)\}
\end{equation}
measures worst-case imbalance
over $B_\rho$,
penalized by the Bregman divergence
from $\tilde\phi$ at scale~$s$.
For the quadratic $\chis = x^2/2$,
the penalty is $s\hP\, h(Z)^2/2$.
For general $\chis$,
it is the sample-averaged Bregman divergence
$\hP\, D_{\chisZ}\{\tilde\phi + sh,\, \tilde\phi\}$.
In either case,
the offset complexity controls
both the bias and the rate of convergence
at different scales~$s$.
To make it useful,
we need the sample penalty
to remain comparable
to its population counterpart
uniformly over $h \in B_\rho$.
Stable lower bounds---the
machinery of \citet{mendelson2015learning,mendelson2021extending}
for controlling sample averages from below---were
developed for exactly this purpose.
Their existing results apply to sums of squares;
we extend them
to sample-averaged Bregman divergences,
where the asymmetry of exponential-type links
makes the divergence
fundamentally different from a squared norm.
The extension handles all three
dispersion families,
including the positive-part quadratic
$\chis(\phi) = \phi_+^2/2$,
whose non-smooth Bregman divergence
requires a separate accounting argument.
For RKHS models,
the resulting bounds take a concrete form
in terms of eigenvalues
of the curvature-weighted integral operator.

\paragraph{Bias}
(\Cref{thm:bias-decomp}).
In Donoho--Liu theory
\citep{donoho1991geometrizing,donoho1994statistical},
the minimax bias under quadratic dispersion
$\chiZ = \gamma^2/2$
is characterized in terms of
the modulus of continuity
$\omega(t) = \sup_{h \in B_\rho}
  \hP \dotpsiZ(h) - t \cdot \hP \chiZ(h)$.
The bias at the optimal weights is
$b(\hgamma_t) = \omega(t) - t\,\omega'(t)$,
a characterization that holds for convex
dispersion penalties $\chiZ$ generally.
Generalizing an approach taken by \citet{kong2025asymptotics}
in the quadratic case, we bound this expression
in terms of the offset complexity:
the bias in our primal problem~\eqref{eq:primal-moment-matching}
satisfies
\[
b(\hgamma) \lesssim \Xi\{\rho(\tilde\phi)\}.
\]
Because convergence of the weights
does not enter the argument, we get useful
bias bounds for even the most undersmoothed
estimators.

\paragraph{Rates and variance}
(\Cref{thm:rate-main},
\Cref{prop:duality},
\Cref{lem:variance-equiv}).
Strong duality
(\Cref{prop:duality})
links the primal and dual solutions
by $\hgamma = \dchis \comp \hat\phi$
for any dispersion and penalty
satisfying mild regularity conditions---significantly
broader than existing duality results.
The same offset complexity
that controls the bias
also determines the rate of convergence
$\hat\phi \to \tilde\phi$
(\Cref{thm:rate-main}).
A comparison lemma
(\Cref{lem:variance-equiv})
shows that convergence
in Bregman divergence---for
entropy-like dispersions,
via an interpolation argument
that controls the nonlinear link---is
strong enough to stabilize the variance,
giving
\[
\hP D_{\chisZ}\{\hat\phi(Z), \tilde\phi(Z)\} \to 0
\quad\implies\quad
\hP \hgamma^2 / \P \tilde\gamma^2 \to 1.
\]

\paragraph{Implications.}
We carry two examples---average
treatment effects
and average linear regression coefficients
\citep{graham2022semiparametrically}---from
formulation through
bias bounds, convergence rates,
and asymptotic normality,
discussing the choice of dispersion
and penalty in each context.
There is a folk wisdom
that entropy penalties are preferable
because they enforce positivity of the weights,
and a competing folk wisdom
that the positivity-constrained quadratic
is better behaved
because the link $\gamma = \phi_+$ is Lipschitz.
Each has its weirdness:
entropy weights blow up exponentially
in the dual parameter,
and positivity-constrained weights
have a junction region
where the Bregman curvature vanishes.
Our results show that both are fine---the
stable lower bounds approach handles each,
and the choice between them
is a modeling decision, not a statistical one.
All of this is done in the independent
but not identically distributed setting
that our examples require.
Have it your way.

\section{Abstract Results}
\label{sec:results}
\label{sec:statistical-considerations}

\subsection{Organizing principle}
\label{sec:organizing-principle}

The statistical properties of the estimator $\hat\phi$
are governed by a single notion:
the local complexity of the seminorm ball $B_\rho$
near $\tilde\phi$,
measured in the curvature-weighted $L_2\{\ddchis(\tilde\phi)\}$ geometry.
This complexity determines
the worst-case imbalance
at the estimated weights (the bias)
and the rate at which $\hat\phi$ converges to $\tilde\phi$.
The rate, in turn,
determines whether the convergence is strong enough
to stabilize $\hP \hgamma^2$.
We described the two key quantities---offset complexity for the bias
and stable lower bounds for the rate---in
\Cref{sec:noise-convergence};
here we state the formal results.

\subsection{Setup}
\label{sec:setup}

We will show that our estimator $\hat\phi$ converges to the minimizer of the corresponding population loss $L^\conj_{\P}$, or equivalently the \emph{regularized Bregman projection} of $\dotpsiZ$.
\begin{equation}
\label{eq:regularized-bregman-projection}
\tilde\phi
= \argmin_{f} L^\conj_{\P}(f) \quad \text{where} \quad L^\conj_{\P}(f) = \P D_{\chisZ}\{f(Z), \phi_{P\dot\psi}(Z)\} + \zetas\{\rho(f)\}.
\end{equation}
Here $\phi_{P\dot\psi}$ is the dual index of the Riesz representer: $\dchis \comp \phi_{P\dot\psi} = \griesz$.
The first order conditions for this projection tell us that we can center the noise term by adding a penalty subgradient term to our loss.
\begin{equation}
\begin{aligned}
\label{eq:loss-difference-with-foc}
L^\conj_{\hP}(f) - L^\conj_{\hP}(\tilde\phi)
&= \hP D_{\chisZ}\{f(Z), \tilde\phi(Z)\} \\
&+ (\hP-\P) (\dchisZ \comp \tilde\phi)(Z)\{f(Z) - \tilde\phi(Z)\} - \dotpsiZ(f - \tilde\phi) \\
&+ D_{\zeta \comp \rho}(f, \tilde \phi).
\end{aligned}
\end{equation}

This decomposition separates the statistical question---how close is $\hat\phi$ to $\tilde\phi$?---from the deterministic approximation question of how well $\phi_{P\dot\psi}$ can be approximated by functions with bounded seminorm.

\paragraph{Scope.}
\label{sec:scope}
We restrict our attention to dispersion penalties with a simple form of dependence on $Z$:
in terms of some \emph{base penalty} $\chi(\cdot)$, $\chiZ(\gamma) = \chi\{W(\gamma-a)\}$ where $W = w(Z) \in \{ \pm 1\}$ is a sign-valued function of $Z$ and $a \in\R$ is constant. This suffices for the choices we considered in ``Choosing $\chi$,'' where the sign flip $W \in \{\pm 1\}$ accounts for the treatment arm. And it's a simple form of dependence because, although sign flips and translation affect the conjugate and its first derivative, it does not affect curvature or Bregman divergence:
\[
\chisZ(\phi) = a\phi + \chis(W\phi), \quad \dchisZ(\phi) = a + W\dchis(W\phi), \quad \text{and} \quad \ddchisZ(\phi) = \ddchis(W\phi).
\]
The last of these is convenient because it allows us to approximate the $\chisZ$-Bregman divergence in terms of the curvature of the base penalty.
\[
D_{\chisZ}(\phi+h,\phi) = \chisZ(\phi+h) - \chisZ(\phi) - \dchisZ(\phi)h \approx \tfrac{1}{2}\ddchisZ(\phi)h^2 = \tfrac{1}{2}\ddchis(W\phi)h^2.
\]
Write $\|h\|_{\chis}^2 = \P\, \ddchis(\tilde\phi)\, h^2$ for the curvature-weighted $L_2$ norm at $\tilde\phi$.

For the base penalty $\chi$, we consider two settings.
\begin{description}
\item[Setting Q (quadratic family).]
The base penalty is $\chis(\phi) = \phi^2/2$ (quadratic)
or $\chis(\phi) = \phi_+^2/2$ (positive-part quadratic).
Curvature is bounded: $\ddchis \le 1$.
For the positive-part case, $\chis$ is not twice-differentiable at zero;
notationally we let $\ddchis(\phi) := \ind{\phi \ge 0}$,
one particular subgradient of $\dchis$.
Weight convergence follows directly from Bregman convergence,
and no interpolation inequality is needed for the variance.
Boundedness instead comes from a seminorm embedding:
if $\ker(\rho)$ is finite-dimensional,
decomposing $h = h_0 + k$
with $h_0 \in \ker(\rho)^\perp$ and $k \in \ker(\rho)$ gives
\begin{equation}\label{eq:seminorm-embedding}
\|h\|_{L_\infty}
  \le \kappa_{\mathrm{emb}}\, \rho(h) + C_K\, \|h\|_{\LtP},
\end{equation}
where $\kappa_{\mathrm{emb}}$ is the embedding constant
for $\inf_{k \in \ker(\rho)} \|h - k\|_{L_\infty}
  \le \kappa_{\mathrm{emb}}\, \rho(h)$
and $C_K$ is the $\LtP$-to-$L_\infty$
norm equivalence constant
on the finite-dimensional kernel.
For Sobolev seminorms with $s > d/2$,
$\kappa_{\mathrm{emb}}$ is the Sobolev embedding constant;
$C_K$ depends on $\dim(\ker\rho)$ and $\pmin$.
\item[Setting E (exponential family).]
The base penalty is \emph{entropy-like}:
the log-derivative $(\log \ddchis)'$ lies in $[\alpha, \beta]$
(Level~2: pseudo-self-concordance and curvature growth;
cf.\ generalized self-concordance in \citet{bach2010self})
and $(\log \dchis)'$ lies in $[\alpha_1, \beta_1]$
(Level~1: link--curvature ratio).
Level~2 implies convexity of $\ddchis$;
Level~1 fixes the ratio between link and curvature.
Entropy, with $\chis(\phi) = e^{\tau\phi}/\tau$,
is the canonical example
($\alpha = \alpha_1 = \beta = \beta_1 = \tau$).
The exponential link $\gamma = e^{\tau\phi}$
makes variance unstable under undersmoothing;
bias and variance results
require an interpolation inequality
relating the sup-norm to the seminorm
and a curvature-weighted $L_2$ norm,
\begin{equation}
\label{eq:interpolation}
\|h\|_{L_\infty} \le \kappa\, \rho(h)^\theta\, \|h\|_{L_2\{\ddchis(\tilde\phi)\}}^{1-\theta}
\end{equation}
for some exponent $\theta \in (0,1)$,
constant $\kappa$,
and all $h \in B_\rho$ orthogonal to $\ker(\rho)$.
For Sobolev seminorms
$\rho = |\cdot|_{\dot H^s}$ with $s > d/2$,
this is the Gagliardo--Nirenberg inequality
with $\theta = d/(2s)$
(\Cref{lem:sobolev-interpolation}).
\end{description}
These settings restrict the rate and variance results;
the bias bound (\Cref{sec:bias})
applies to any convex $\chi$.
The challenges differ between them.
For entropy-like dispersions,
the exponential link makes the Bregman divergence's level sets
teardrop-shaped:
$D_{\chis}(\tilde\phi + h, \tilde\phi)$
is exponentially large in some coordinates
and small in others,
so standard concentration fails
and we need lower bounds
that tolerate this asymmetry.
For the positive-part quadratic,
the curvature $\ddchis(\phi) = \ind{\phi \ge 0}$
vanishes in the dead zone $\{\tilde\phi < 0\}$,
where the Bregman divergence grows only linearly
rather than quadratically.

\subsection{Duality}
\label{sec:duality-main}

The dual relationship $\hgamma = \dchis \comp \hat\phi$ between the solutions of \eqref{eq:primal-moment-matching} and \eqref{eq:dual-problem} has been established in various special cases. The following result, proven in \Cref{sec:duality}, establishes it for the full generality of our formulation.

\begin{proposition}[Strong Duality]
\label{prop:duality}
Suppose $\rho$ is a seminorm on functions $f:\mathcal{Z} \to \R$ with finite-dimensional kernel $K$;
$\hP$ is a measure supported on finitely many points $Z_1 \ldots Z_n \in \mathcal{Z}$; $\chi:\R \to \R \cup \{ +\infty \}$ is strictly convex on its effective domain; $\zeta:\R \to \R \cup \{ +\infty \}$ is convex, nondecreasing, and coercive; and the point evaluation functionals $m \to m(Z_i)$ and $\dot\psi_{Z_i}$ are continuous on the quotient space $\calS/K$. Under a qualification condition---there exists $\gamma$ with $L_{\hP}(\gamma) < \infty$ at which $\zeta$ is continuous---then $L_{\hP}$ has a unique minimizer $\hgamma$ and \smash{$\dchis(\hat\phi_j) \to \hgamma$} for any sequence \smash{$(\hat\phi_j)$} along which $L^\conj_{\hP}$ converges to its minimum.
\end{proposition}
The characterization of $\hgamma$ via a minimizing sequence $\hat\phi_j$ rather than a minimizer $\hat\phi$ saves verifying that $L^\conj_{\hP}$ attains its minimum. The conditions are otherwise weak: aside from convexity, the qualification condition asks only that no function in $B_\rho$ have arbitrarily large imbalance---except for kernel functions, which have $\rho(k)=0$ irrespective of scale and must be balanced exactly.

\subsection{Bias}
\label{sec:bias}

The bias of the one-step estimator is the worst-case imbalance $b(\hgamma) = \sup_{f \in B_\rho} |\hP \hgamma(Z) f(Z) - \dotpsiZ(f)|$ at the estimated weights. This is a purely primal quantity: it depends on the weights $\hgamma$ themselves, not on how close they are to any target in dual space. In particular, bounding the bias requires no weight convergence.

The key quantity is the offset complexity $\Xi(s)$ defined in~\eqref{eq:offset-complexity-intro}, which measures how much the centered empirical process can push the optimizer away from $\tilde\phi$, offset by the sample-averaged Bregman divergence. For large increments, the superlinearity of $D_{\chisZ}$ makes the Bregman term dominate, localizing the supremum.

\begin{theorem}[Bias bound]
\label{thm:bias-decomp}
Let $\chi$ and $\zeta$ satisfy the assumptions of \Cref{sec:duality-main}. Fix $s > 0$ with $2s \cdot \rho(\tilde\phi) < 1$, where $\tilde\phi$ is the population minimizer of the dual loss \eqref{eq:dual-problem}. Let $b = b(\hgamma)$ denote the worst-case bias at the estimated weights $\hgamma = \dchis \comp \hat\phi$. Then $|b| \le 2\Xi(s)$.
\end{theorem}

\paragraph{Undersmoothing limit.} As regularization vanishes, $\rho(\tilde\phi) \to \infty$ and the constraint $s \cdot \rho(\tilde\phi) < 1/2$ forces $s \to 0$. In this limit, $\tilde\gamma = \dchis \comp \tilde\phi \to \griesz$ (centering at the true target) and the superlinearity of $D_{\chis}$ makes the Bregman offset dominate for all nonzero increments---only infinitesimal perturbations contribute. The offset complexity converges to the local Rademacher complexity of $B_\rho$ at $\griesz$: the complexity of unit-seminorm increments around the true target. There is no harm in letting $\rho(\tilde\phi) \to \infty$ arbitrarily fast: this only strengthens the localization.

\subsection{Variance}
\label{sec:variance}

The asymptotic variance of the one-step estimator is $(1/n)\hP \hgamma^2 v$ for $v(z) = \Var(Y \mid Z=z)$ (see~\eqref{eq:decomposition}). Because the estimated weights $\hgamma = \dchis \comp \hat\phi$ are random, we need to (i) establish that $\hat\phi$ converges to $\tilde\phi$ and (ii) show that this convergence is strong enough to replace $\hgamma$ with the deterministic population weights $\tilde\gamma = \dchis \comp \tilde\phi$ in the variance.

For the positive-part quadratic ($\gamma = \phi_+$), Bregman convergence gives $L_2$ weight convergence directly: $D_{\chis}(\hat\phi, \tilde\phi) = (\hat\phi_+ - \tilde\phi_+)^2/2$, so task (ii) is immediate from the rate. For exponential family dispersions ($\gamma = e^{\tau\phi}$), the nonlinear link amplifies dual-space errors, and an interpolation inequality is needed to control the resulting $e^{2\tau h}$ terms. We address both tasks below: weight convergence via uniform Bregman lower bounds, and variance equivalence via a deterministic comparison lemma.

The argument has three steps.
First, we establish a stable lower bound
on the quadratic approximation
$\frac{1}{2}\hP\, \ddchis(\tilde\phi)\, h^2$
and lift it to the sample-averaged Bregman divergence
via Jensen's inequality---this
is what lets us work
without concentration.
Second, we combine this lower bound
with Rademacher complexity control
to get a rate of convergence
for $\hat\phi \to \tilde\phi$
(\Cref{thm:rate-main}).
Third, we show that Bregman convergence
is strong enough to stabilize the variance
(\Cref{lem:variance-equiv}).
The first step is the new technical ingredient;
the third is a deterministic consequence.

\subsubsection{Stable lower bounds}

The key tool for task~(i) is the \emph{stable lower bound} of \citet{mendelson2021extending}, discussed in \Cref{sec:uniform-lower-bounds}: a guarantee that sample-averaged sums of squares remain close to their population counterparts even after dropping a small set of coordinates. Stability allows us to derive uniform bounds from bounds on a finite net: nonnegativity of the Bregman divergence lets us drop coordinates where off-net approximation error is hard to control, and stability ensures this doesn't cost much. The formal definition and sufficient conditions (bounded and $L_p$ cases) are given in \Cref{sec:uniform-lower-bounds}.

We adapt this to our setting by working with the quadratic approximation to the Bregman divergence, $\frac{1}{2}\ddchis(\tilde\phi)h^2 \approx D_{\chisZ}(\tilde\phi + h, \tilde\phi)$, using stable lower bounds on the curvature-weighted functions $\tilde{\mathcal{H}} := \{ \ddchis(\tilde\phi)^{1/2}h \ : \ h \in \mathcal{H} \}$.

\paragraph{When $\ddchis$ is convex,} Jensen's inequality implies that the quadratic approximation is a lower bound (\Cref{lem:quadratic-bregman}). Quadratic homogeneity allows us to work around the poor control on $\rho(f)$ that we get when we undersmooth: we can prove uniform lower bounds on $\rho$'s unit ball $B_\rho$, then scale up by a factor of $s^2$ to get a uniform lower bound on $\frac{1}{2}\hP \ddchis(\phi)h^2$ and therefore on our sample-averaged Bregman divergence. For entropy-like dispersions $\gamma = e^{\tau\phi}$, this requires only that sample averages not be too small, so the exponential growth of individual terms is not an obstacle.

\paragraph{When $\chis(\phi)=\phi_+^2/2$,}
the approximation is exact away from
a junction region near zero:
$D_\chis(\phi+h, \phi) = h^2/2$
when $h \ge -\phi$,
but growth slows to linear for $h < -\phi$.
\Cref{lem:two-pile} handles this
by transferring a stable lower bound
from the active zone
to the full Bregman divergence,
at the cost of degrading the tolerance $\xi$
by a dead-zone term
that depends on how much mass
$\tilde\phi$ puts near zero.

\subsubsection{Rate of convergence}

Combining stable lower bounds with Rademacher complexity control yields the following rate. For positive numbers $r, s, \nu$ with $s^2 \ge \nu r^2$, define the \emph{increment class}
\[
H_{r,s} = \{h : \|h\|_{\chis} = r,\; D_{\zetas \comp \rho}(\tilde\phi + h, \tilde\phi) \le s^2\}.
\]

\begin{theorem}[Rate of Convergence]
\label{thm:rate-main}
Let $\xi \in (0,1)$. Suppose:
\begin{enumerate}
\item[(i)] Every $\tilde h = \sqrt{\ddchis(\tilde\phi)} \cdot h$ with $h \in H_{r,s}$ satisfies a $(\xi/2, \ell, k)$ stable lower bound.
\item[(ii)] For some convex increasing $F: \R_{\ge 0} \to \R_{\ge 0}$, these complexity bounds are satisfied:
\begin{align}
&\E \sup_{h \in H_{r,s}} \qty| \hP \sigma_i \sqrt{\ddchis(\tilde\phi)} h(Z_i) | \le c \xi \sqrt{\nu} r \min\qty(\sqrt{\ell/n}, \sqrt{k/n}), \label{eq:rademacher-complexity} \\
&\E F\qty[ 2 \sup_{h \in H_{r,s}} \qty| \hP \sigma_i \qty{(\dchisZ \comp \tilde\phi)(Z)h(Z) - \dotpsiZ(h)} | ] \le \delta \cdot F(\nu r^2). \label{eq:F-complexity}
\end{align}
\end{enumerate}
Then with probability at least $1 - \delta - 2e^{-c\min(\ell,k)}$,
\[
\|\hat\phi - \tilde\phi\|_{\chis} < r \qqtext{and} \hP D_{\chisZ}(\hat\phi, \tilde\phi) + D_{\zetas \comp \rho}(\hat\phi, \tilde\phi) \le \nu r^2.
\]
Here $\sigma_1 \ldots \sigma_n$ is an independent Rademacher sequence and $c$ is a universal constant.
\end{theorem}
For the positive-part quadratic $\chis(\phi) = \phi_+^2/2$, condition~\ref{cond:shell-i} is obtained via \Cref{lem:two-pile}, which transfers an active-zone SLB to the full Bregman divergence at a degraded tolerance $\xi$ depending on $M$ and the mass $\P(\tilde\phi \in [0,\epsilon])$ near the dead-zone boundary.

The proof combines \Cref{lem:generic-rate} (rate via offset complexity) with the Bregman lower bounds of \Cref{sec:uniform-lower-bounds}; see \Cref{sec:rate-of-convergence}.

\paragraph{Sobolev models.}
For the Sobolev seminorm $\rho(g) = \|(-\Delta)^{s/2} g\|_{L_2}$ with $s > d/2$ and quadratic regularization $\zetas(\rho) = \eta \rho^2$, the \emph{effective dimension}
\[
d_\eta = \sum_j \frac{\lambda_j}{\max(\lambda_j, \eta)} \qqtext{for} \lambda_j = j^{-2s/d}
\]
governs the rate. The interpolation inequality (\Cref{lem:sobolev-interpolation}) gives $L_\infty$ control on the seminorm ball, converting to a stable lower bound via \Cref{lem:slb-bounded}; eigenvalue estimates on the curvature-weighted integral operator give the Rademacher bounds.

\begin{corollary}[Sobolev Models with Quadratic Regularization]
\label{cor:sobolev-main}
For the Sobolev seminorm $\rho(g) = \|(-\Delta)^{s/2} g\|_{L_2}$ with $s > d/2$ and quadratic regularization $\zetas(\rho) = \eta\rho^2$, let $X$ have density $p \ge \pmin > 0$ on $[0,1]^d$. Suppose $\underline{\ddchis} \le \ddchis(\tilde\phi) \le \overline{\ddchis}$ a.s.\ and that
\begin{equation}
\label{eq:sobolev-sample-size}
n\pmin^{1-\theta}\eta^\theta
  \ge C_0\, \frac{\overline{\ddchis}}{\underline{\ddchis}}\,
  \log\frac{2}{\delta}
\qqtext{for} \theta = \frac{d}{2s}.
\end{equation}
Then with probability at least $1 - \delta$,
\[
\|\hat\phi - \tilde\phi\|_{\chis}^2 \le \frac{C (\overline{\ddchis}/\underline{\ddchis}) d_\eta}{\delta n}
\qqtext{and}
\hP D_{\chis}(\hat\phi, \tilde\phi) + \eta\rho^2(\hat\phi - \tilde\phi) \le \frac{C (\overline{\ddchis}/\underline{\ddchis}) d_\eta}{\delta n}.
\]
\end{corollary}
The curvature ratio $\overline{\ddchis}/\underline{\ddchis}$ is a condition number: it equals $1$ for the quadratic, and is bounded under pseudo-self-concordance when $\tilde\phi$ is bounded. The sample size condition~\eqref{eq:sobolev-sample-size} asks that the effective sample size after interpolation, $n\pmin^{1-\theta}\eta^\theta$, dominate the curvature ratio.

\subsubsection{Variance equivalence}

With weight convergence established, it remains to show that convergence in Bregman divergence is strong enough to preserve the variance: that the estimated variance $\hP \hgamma^2 v$ can be replaced by its population counterpart $\hP \tilde\gamma^2 v$. For the quadratic family, this is automatic: $D_{\chis}(\hat\phi, \tilde\phi) = (\hat\phi - \tilde\phi)^2/2$ (or $(\hat\phi_+ - \tilde\phi_+)^2/2$), so Bregman convergence is $L_2$ weight convergence. For entropy-like dispersions, the nonlinear link $\hgamma = e^{\tau\hat\phi}$ amplifies dual-space errors, and an interpolation bound $\|\hat\phi - \tilde\phi\|_{L_\infty} \to 0$ is needed in addition to Bregman convergence to ensure $\hP \hgamma^2 v / \hP \tilde\gamma^2 v \to 1$. \Cref{lem:variance-equiv} and \Cref{cor:variance-interpolation} in \Cref{sec:variance-equiv} give deterministic conditions under which this holds. This is the only place where the Level~1 bound ($\alpha_1\, \dchis \le \ddchis \le \beta_1\, \dchis$) is needed: it bridges Bregman convergence in $\ddchis$-space to weight convergence in $\dchis$-space.

\section{Worked Examples}
\label{sec:worked-examples}

We illustrate the framework on three applications
from \Cref{sec:intro},
chosen to highlight different aspects
of the argument.
The treatment-specific mean (\Cref{sec:example-tsm})
is the cleanest case:
all three dispersions give the same rate,
and the differences are structural---how
Bregman convergence relates to weight convergence,
and when interpolation is needed
to bridge the gap.
Average linear regression (\Cref{sec:example-regression})
introduces misspecification:
the nuisance is a projection,
not the conditional mean,
and orthogonality of the Riesz representer
is what keeps the bias from dominating.
The variance of the CATE (\Cref{sec:example-vcate})
introduces nonlinearity:
the functional is quadratic in the nuisance,
so conditions~\ref{cond:an-F} and~\ref{cond:an-G}---vacuous
for the linear examples---require
convergence of the outcome estimator.

\subsection{Treatment-specific mean}
\label{sec:example-tsm}

Suppose we observe $(Z_i,Y_i)$ with $Z_i=(W_i,X_i)$ where $W_i$ is a treatment indicator, $X_i$ is a covariate vector, and $Y_i$ is the response. The \emph{treatment-specific mean} $\tau_w = \E \mu(w,Z_i)$ for $\mu(z)=\E(Y_i \mid Z_i=z)$ can be written as $\E \dotpsiZ(\mu)$ where $\dotpsiZ(m) = m(w, X)$. Taking $\P=\hP$ in \eqref{eq:primal-moment-matching}, we seek weights $g(Z)$ that \emph{balance} functions of $X$ across treatment groups. The inverse propensity weights $\gipwstar(Z) := \ind{W=w}/\pi_w(X)$ for $\pi_w(x) = \P(W=w \mid X=x)$ achieve perfect population balance, so \eqref{eq:primal-moment-matching} estimates them via moment matching:
\[
\P \gipwstar(Z)m(Z) = \P \dotpsiZ(m) \qqtext{so we ask that}
\hP \hgamma(Z_i)m(Z_i) \approx \hP \dotpsiZ(m).
\]
The estimator error decomposes as
\[
\begin{aligned}
\hP \hgamma(Z) Y - \tau_w
&= \underbrace{\hP \hgamma(Z)\mu(Z) - \P \dotpsiZ(\mu)}_{\text{in-sample bias}} + \underbrace{\hP \hgamma(Z)\varepsilon + (\hP-\P)\dotpsiZ(\mu)}_{\text{averaged mean-zero noise}}
\end{aligned}
\]
for $\varepsilon=Y - \E(Y \mid Z)$. Our results bound the in-sample bias term via the offset complexity (\Cref{sec:bias}), without requiring weight convergence.

\paragraph{Headline rates.}
Under Sobolev smoothness $s > d/2$
with quadratic regularization
$\zetas(\rho) = \eta\rho^2$
and propensity overlap $\pi_w(x) \ge \pi_{\min} > 0$,
all three dispersions give the same leading rate
of convergence for $\hat\phi \to \tilde\phi$:
\[
\hP\, D_{\chis}(\hat\phi, \tilde\phi)
  + \eta\,\rho^2(\hat\phi - \tilde\phi)
  \le \frac{C\, d_\eta}{\delta\, n}
\]
with probability $1 - \delta - 2e^{-c\, n\, \pmin^{1-\theta}\eta^\theta}$,
where $d_\eta \asymp \eta^{-\theta}$
is the effective dimension
and $\theta = d/(2s)$.
The bias bound is
\[
|b(\hgamma)|
  \le C\, n^{-1/2}\,
    (n^{1/2}\,\tilde\rho)^{-\theta/(2-\theta)}
  + 4\eta\,\rho(\tilde\phi),
\]
which vanishes faster than $n^{-1/2}$
under overlap
(because $\tilde\rho$ is bounded and $\eta \to 0$).
Asymptotic normality follows from
\Cref{thm:asymptotic-normality}
once we verify infinitesimal weights
and variance equivalence---both
of which hold under overlap
for all three dispersions
(\Cref{sec:app-tsm}).
The rates are identical
because the Bregman divergence is locally quadratic:
$D_{\chis}(\tilde\phi + h, \tilde\phi)
  \approx \frac{1}{2}\ddchis(\tilde\phi)\,h^2$,
and the curvature $\ddchis(\tilde\phi)$
enters only through the ratio
$\overline{\ddchis}/\underline{\ddchis}$,
which is bounded under overlap
for all three families.

\paragraph{What varies across dispersions.}
The dispersions differ not in rate
but in structure:
what the weights look like,
how Bregman convergence relates
to weight convergence,
and what additional conditions are needed.
\begin{description}
\item[Quadratic] ($\chis(\phi) = \phi^2/2$,
weights $\hgamma = \hat\phi$).
Simplest case.
Bregman convergence \emph{is} $L_2$ weight convergence:
$D_{\chis}(\hat\phi, \tilde\phi)
  = (\hat\phi - \tilde\phi)^2/2$.
No interpolation needed.
But the weights are unconstrained---they
can be negative,
which is unnatural for inverse probability weights.
\item[Positive-part quadratic]
($\chis(\phi) = \phi_+^2/2$,
weights $\hgamma = \hat\phi_+$).
Weights are nonnegative by construction.
Bregman convergence still gives
$L_2$ weight convergence directly,
because $D_{\chis}(\hat\phi, \tilde\phi)
  \ge (\hat\phi_+ - \tilde\phi_+)^2/2$.
The additional cost is a tolerance degradation
from the dead zone:
the two-pile analysis (\Cref{lem:two-pile})
degrades $\xi$ by a term depending on $M$
and the mass $\P(\tilde\phi \in [0,\epsilon])$.
\textcolor{red}{Under overlap with Sobolev regularization,
\Cref{prop:anticonc-sobolev}
gives $\omega(\epsilon) \to 0$,
so the dead-zone cost vanishes.}%
% STALE: prop:anticonc-sobolev no longer requires overlap or Sobolev
\item[Entropy-like]
($\chis(\phi) = e^{\tau\phi}/\tau$,
weights $\hgamma = e^{\tau\hat\phi}$).
Weights are automatically positive
and the link $\gamma = e^{\tau\phi}$
is a logistic model for propensity scores---the
most common choice in practice.
But Bregman convergence does \emph{not}
imply $L_2$ weight convergence:
$D_{\chis}(\tilde\phi+h, \tilde\phi)
  = (e^{\tau h} - 1 - \tau h)\,
    e^{\tau\tilde\phi}/\tau$
grows exponentially,
so small Bregman divergence does not preclude
large pointwise errors
$|\hat\phi(Z_i) - \tilde\phi(Z_i)|$
at individual coordinates.
Variance equivalence
(\Cref{cor:variance-interpolation})
circumvents this gap:
we show $\hP \hgamma^2 \nu \to \hP \tilde\gamma^2 \nu$
directly via an interpolation argument,
without passing through $L_2$.
This is the only case
where the Level~1 bound
($\alpha_1\,\dchis \le \ddchis \le \beta_1\,\dchis$)
and the interpolation inequality~\eqref{eq:interpolation}
are needed.
\end{description}

\paragraph{What the link costs you.}
The headline rates are identical under overlap,
so the dispersion choice might seem cosmetic.
It is not.
\Cref{fig:weight-fan} shows
estimated weight functions $\hat\gamma(x)$
across twelve values of $\eta$
for the propensity score
$\pi(x) = \tfrac{1}{2}x^{0.7}$
with Sobolev regularization
($s = 1$, $d = 1$,
$J = 200$ cosine basis,
$n = 2000$).
At this overlap exponent
the inverse propensity weights satisfy
$\griesz = 1/\pi_w \in L_p$ for $p < 10/7$,
so the quadratic CLT applies
without any overlap condition
(\Cref{sec:app-tsm}).
The entropy CLT requires
$\log(1/\pi) \in H^1$,
which fails:
the derivative $\alpha/x$ is not square-integrable.

\begin{figure}[t]
\centering
\begin{minipage}[t]{0.48\textwidth}
\centering
\includegraphics[width=\textwidth]{scratch/plots/viz-fan-Q.pdf}
\end{minipage}
\hfill
\begin{minipage}[t]{0.48\textwidth}
\centering
\includegraphics[width=\textwidth]{scratch/plots/viz-fan-E.pdf}
\end{minipage}

\medskip

\includegraphics[width=0.48\textwidth]{scratch/plots/viz-density.pdf}
\caption{Estimated weight functions $\hat\gamma(x)$
for quadratic (left, blue) and entropy (right, red)
dispersions across twelve values of $\eta$
from $10^{-0.5}$ (lightest) to $10^{-4.5}$ (heaviest).
Dashed: oracle $1/\pi(x) = 2/x^{0.7}$.
Bottom: covariate densities by treatment group.
Propensity $\pi(x) = \tfrac{1}{2}x^{0.7}$,
$\mu(x) = \sin(2\pi x)$,
$n = 2000$, $H^1$ cosine sieve with $J = 200$.
The quadratic fan spreads gradually
(max weight $3$ to $20$);
the entropy fan shoots past the frame
at small $\eta$
(max weight $4$ to $300$).}
\label{fig:weight-fan}
\end{figure}

\FloatBarrier

For the quadratic dispersion,
$\hat\gamma = \hat\phi$
grows linearly in the regularized index.
As $\eta$ decreases
across four orders of magnitude,
the weight function
approaches the oracle gradually:
the most aggressive curve reaches
$\hat\gamma(0) \approx 20$.
For the entropy dispersion,
$\hat\gamma = e^{\tau\hat\phi}$
amplifies the same index exponentially.
At $\eta = 0.002$
the entropy weights reach~$13$
while quadratic reaches~$8$;
at $\eta = 3 \times 10^{-5}$,
entropy reaches~$300$
while quadratic reaches~$20$.
Most of the entropy fan
lies above the frame.

The exponential link
is doing exactly what it should:
approximating $1/\pi$
through $e^{\tau\hat\phi}$
rather than $\hat\phi$ directly.
But the amplification means
the range of $\eta$ values
producing stable weights is narrower.
The quadratic estimator tolerates
a four-order-of-magnitude sweep
with bounded weights;
the entropy estimator
produces stable weights
over roughly two orders.

Despite this,
when each method is tuned to its own oracle---the
$\eta$ minimizing RMSE
over Monte Carlo replications
with oracle outcome regression
$\hat\mu = \mu$---the
gap closes:
at $n = 2000$,
the oracle RMSE is $0.042$ for quadratic
and $0.043$ for entropy
($3\%$ gap);
at $n = 5000$,
both achieve $0.026$.
The entropy estimator compensates
by choosing heavier regularization,
taming the exponential weights
at the cost of more smoothing.

\paragraph{Choosing a dispersion.}
The quadratic dispersion
is the most forgiving default.
Its weights grow linearly in the index,
the theory requires
the mildest conditions---condition~\ref{cond:an-C}
is automatic
because Bregman divergence $=$ $L_2$ distance---and
the tuning window is wide.
The positive-part quadratic adds nonnegativity
at negligible theoretical cost:
the dead-zone tolerance degrades
but does not change the rate.
The entropy dispersion is natural
when the weights model
a logistic propensity score
and positivity is essential.
The guarantees hold
when $\log(1/\pi)$
has the penalty's native smoothness,
but the exponential link
narrows the tuning window
and requires stronger structural conditions---the
interpolation inequality for variance equivalence
and pseudo-self-concordance for bias control.
Practitioners using entropy-type losses need not switch:
where the theory applies,
the CLT is valid
and oracle-tuned performance
is essentially identical to quadratic.
But there is less room for tuning error.

\subsection{Average linear regression}
\label{sec:example-regression}

The nuisance $\mu(W,X) = a_0(X) + W \cdot b_0(X)$
is a conditional linear predictor---a projection
of $\E(Y \mid W,X)$ onto functions linear in $W$,
not the conditional mean itself.
The residual $\mu_\perp(Z) = \E(Y \mid Z) - \mu(Z)$
is nonzero in general,
so this is a misspecified setting.

Misspecification would be fatal
without the orthogonality from
the Riesz representer.
Here $\griesz(W,X) = \{W - \E(W \mid X)\}/\Var(W \mid X)$,
which is linear in $W$
and therefore lies in the tangent space $\mathcal{H}$.
The projection property gives
$\P\griesz\, \mu_\perp = 0$,
killing the leading bias term.
The remaining misspecification contribution
is controlled by the regularization:
$|\P\tilde\phi\, \mu_\perp| \le C\eta$,
which vanishes as $\eta \to 0$.

Only the quadratic dispersion is natural here:
the weights $\hgamma = \hat\phi$
estimate a contrast $\{W - \E(W \mid X)\}/\Var(W \mid X)$
that takes both signs,
so positivity constraints would be inappropriate.
Entropy-like dispersions produce weights
$\hgamma = e^{\tau\hat\phi} > 0$
or $\hgamma = -e^{-\tau\hat\phi} < 0$;
neither can represent a contrast
that changes sign with $W$.
Asymptotic normality follows
under the same Sobolev conditions as the
treatment-specific mean,
with the additional requirement
$\sqrt{n}\,\eta \to 0$
to ensure the regularization-induced bias
from misspecification is negligible.
See \Cref{sec:app-regression}.

\subsection{Variance of the CATE}
\label{sec:example-vcate}

The variance of the CATE
$V = \Var(\tau(X))$
of Example~\ref{ex:vcate}
is the first nonlinear functional
to which we apply the framework.
The nonlinear component
$\E[\tau(X)^2]$
has $\psiZ(m) = (m(1,X) - m(0,X))^2$;
the variance
$V = \E[\tau^2] - (\E\tau)^2$
subtracts the squared ATE,
a linear functional
handled by \Cref{sec:example-tsm}.

The derivative at $\mu$
has Riesz representer~\eqref{eq:vcate-riesz},
the ATE weight modulated by the CATE.
We estimate it by a single penalized projection:
the moment-matching condition~\eqref{eq:primal-moment-matching}
asks that $\hgamma$ balance functions of $(W,X)$
against $\dpsiZ^{\hat\mu}$,
the derivative formed from the outcome estimate $\hat\mu$.
Because $\griesz^\mu$ depends on $\mu$,
we solve this at $\hat\mu$
rather than at the unknown $\mu$.

\paragraph{Headline rates.}
Under Sobolev smoothness $s > d/2$,
overlap $\pi_w(x) \ge \pi_{\min} > 0$,
and heterogeneous effects $V > 0$,
the quadratic dispersion gives the same
rate of Bregman convergence
as the treatment-specific mean.
Only the quadratic dispersion is natural here:
the representer $\griesz^\mu$
takes both signs
(positive for treated, negative for control
when $\tau > 0$),
so sign-definite links
are inappropriate---the same
situation as for average linear regression.
The representer is bounded:
$\|\griesz^\mu\|_{L_\infty}
  \le 2\|\tau\|_{L_\infty}/\pi_{\min}$.

\paragraph{The linearization error.}
Because $\psiZ$ is quadratic in $m$,
the linearization remainder is exact:
\begin{equation}\label{eq:vcate-linearization}
R_{\mathrm{lin}}
  = \P\psiZ(\hat\mu)
    + \dpsiZ^\mu(\mu - \hat\mu)
    - \P\psiZ(\mu)
  = \|\hat\tau - \tau\|_{\LtP}^2.
\end{equation}
Condition~\ref{cond:an-G} requires this to be
$o_p(n^{-1/2})$,
i.e.\ $\hat\mu$ converges
faster than $n^{-1/4}$ in $L_2$.
Under Sobolev smoothness $s > d/2$,
nonparametric regression achieves
$\|\hat\mu - \mu\|_{\LtP}^2
  \lesssim n^{-2s/(2s+d)}$,
which is $o(n^{-1/2})$
when $s > d/2$:
the same smoothness condition
that drives the weight analysis.

\paragraph{What changes.}
Conditions~\ref{cond:an-A}--\ref{cond:an-C}---imbalance,
infinitesimal weights,
and variance equivalence---are
verified as in
\Cref{sec:example-tsm},
because the penalized projection
for $\griesz^\mu$
has the same analytical structure
as for the TSM.
What is new:
\begin{description}
\item[Conditions~\ref{cond:an-F} and~\ref{cond:an-G}]
are no longer vacuous.
The linearization error~\eqref{eq:vcate-linearization}
and functional stability
$(\hP - \P)[\hat\tau^2 - \tau^2]
  = o_p(n^{-1/2})$
both require convergence
of the outcome estimator.
Under Sobolev smoothness,
both hold automatically.
\item[The representer depends on $\mu$.]
The penalized projection targets
$\dpsiZ^{\hat\mu}$,
not $\dpsiZ^\mu$.
The difference
$\dpsiZ^{\hat\mu}(h) - \dpsiZ^\mu(h)
  = 2\,\E[(\hat\tau - \tau)(h(1,X) - h(0,X))]$
is controlled by
$\|\hat\tau - \tau\|_{\LtP}$,
which vanishes under the same
rate condition.
\item[Denominator anchoring.]
For the treatment-specific mean,
the first-order condition~\eqref{eq:foc-quadratic} at $h = 1$
gives $\P\tilde\phi = p_w$
and Jensen anchors $\P\tilde\phi^2 \ge p_w^2$
uniformly in $\eta$.
Here, $\dpsiZ^\mu(1) = 0$---the
derivative kills constant functions---so
the anchor comes instead
from the convergence
$\tilde\phi \to \griesz^\mu$
under $\eta \to 0$,
giving $\P\tilde\phi^2 \to \P(\griesz^\mu)^2 > 0$.
This is not an additional condition:
the tuning window for~\ref{cond:an-A}--\ref{cond:an-C}
already requires $\eta \to 0$.
\end{description}

See \Cref{sec:app-vcate}
for the full verification.

\ifarXiv
  \printbibliography
\else
  \bibliography{bibliographyml}
  \bibliographystyle{icml2026}
\fi


\appendix

\section{Preliminaries}
\label{sec:general-framework}

\subsection{Empirical process tools}
\label{sec:empirical-process-tools}

We collect standard empirical process inequalities
for averages of independent random variables,
stated in terms of the empirical measure
$\hP f = (1/n)\sum_{i=1}^n f(Z_i)$
and its population counterpart
$\P f = (1/n)\sum_{i=1}^n \E f(\bar Z_i)$.

\begin{lemma}[Variance of suprema; Example~3.6 of \citealt{boucheron2013concentration}]
\label{lem:efron-stein}
Let $Z_1 \ldots Z_n$ be independent
(not necessarily identically distributed).
For any class $\F$
of measurable functions,
$\Var \sup_{f \in \F} (\hP - \P) f
  \le (2/n) \sup_{f \in \F} \P f^2$.
\end{lemma}

\begin{lemma}[Symmetrization; Lemma~6.3 of \citealt{ledoux2013probability}]
\label{lem:symmetrization}
Let $Z_1 \ldots Z_n$ be independent
and $\sigma_1 \ldots \sigma_n$ be independent Rademacher
variables independent of the data.
For any class $\F$
of bounded functions
and any convex increasing $\Psi : \R_+ \to \R$,
\[
\E \Psi\qty{\sup_{f \in \F} |(\hP - \P) f|}
  \le \E \Psi\qty{
    2 \sup_{f \in \F}
    \qty|\hP \sigma f|}.
\]
\end{lemma}

\begin{lemma}[Contraction; Theorem~4.12 of \citealt{ledoux2013probability}]
\label{lem:contraction}
Let $Z_1 \ldots Z_n$ be independent
and $\varphi : \R \to \R$ be $L$-Lipschitz
with $\varphi(0) = 0$.
Then
\[
\E \sup_{f \in \F}
  \qty|\hP \sigma\, (\varphi \comp f)|
\le 2L \cdot
  \E \sup_{f \in \F}
  \qty|\hP \sigma f|.
\]
\end{lemma}

\begin{lemma}[Talagrand concentration]
\label{lem:talagrand}
Let $Z_1 \ldots Z_n$ be independent (not necessarily identically distributed).
Let $\F$ be a countable class of functions
with $\|f\|_\infty \le b$ for all $f \in \F$,
and write $W = \sup_{f \in \F} (\hP - \P)f$.
For all $t \ge 0$,
\[
\P\qty(W \ge \E W + t)
  \le \exp\qty{-\frac{nt^2}{2(V + 2b\E W + bt)}}
\qqtext{where}
V = \sup_{f \in \F} \P(f - \P f)^2.
\]
\end{lemma}

\begin{proof}
Theorem~3.3.16 of \citet{gine2016mathematical}
gives, for $S = \sup_f \sum_i \{f(Z_i) - \E f(Z_i)\}$,
the log-MGF bound
\[
\log \E e^{u(S - \E S)}
  \le \frac{u^2 \bar{V}_n}{2(1 - u b)}
\]
for $u \in [0, 1/b)$,
where $\bar{V}_n = nV + 2b\E S$.
Because $W = S/n$,
we have $\P(W \ge \E W + t) = \P(S \ge \E S + nt)$.
Markov's inequality applied to $e^{u(S - \E S)}$ gives
\[
\P(S \ge \E S + nt)
  \le e^{-u nt}\, \E e^{u(S - \E S)}
  \le \exp\qty{-u nt + \frac{u^2 \bar{V}_n}{2(1 - u b)}}
\]
for any $u \in [0, 1/b)$.
Take $u = nt / (\bar{V}_n + bnt)$.
This satisfies $1 - ub = \bar{V}_n/(\bar{V}_n + bnt)$
and $u(\bar{V}_n + bnt) = nt$,
so the log-MGF term is
\[
\frac{u^2 \bar{V}_n}{2(1-ub)}
  = \frac{u^2 \bar{V}_n}{2}
    \cdot \frac{\bar{V}_n + bnt}{\bar{V}_n}
  = \frac{u^2(\bar{V}_n + bnt)}{2}
  = \frac{u \cdot nt}{2}.
\]
The exponent at $u$ is therefore
$-u nt + u nt/2 = -u nt/2 = -n^2 t^2/\{2(\bar{V}_n + bnt)\}$.
Substituting $\bar{V}_n + bnt = nV + 2b\E S + bnt = n(V + 2b\E W + bt)$
gives the stated bound.
\end{proof}

\subsection{Existence of the Riesz representer}
\label{sec:riesz-existence}

The conditions in \Cref{def:differentiable}
are that the derivative functional
$h \mapsto \P\dpsiZ(h)$ be $L_2(\P)$-continuous---that
is, $|\P\dpsiZ(h)| \le C\, \norm{h}_{L_2(\P)}$ for some constant $C$---and
that $\mathcal{H}$ be closed in $L_2(\P)$.
When both hold,
the Riesz representation theorem
\citep[Theorem~6.15]{folland1999real}
gives a unique $\griesz \in \mathcal{H}$
satisfying~\eqref{eq:riesz-rep}.

\subsection{Basic bounds for entropy-like dispersions}
\label{sec:e-type-basics}

Suppose $\chis$ is three times continuously differentiable
and entropy-like:
\begin{equation}
\label{eq:e-type}
\alpha\, \ddchis \le \dddchis \le \beta\, \ddchis
\qquad\text{(Level 2)},
\qquad
\alpha_1\, \dchis \le \ddchis \le \beta_1\, \dchis
\qquad\text{(Level 1)}.
\end{equation}
Level~2 strengthens pseudo-self-concordance
($|\dddchis| \le \beta\, \ddchis$)
by requiring curvature growth:
$\dddchis \ge \alpha\, \ddchis$ with $\alpha > 0$,
so that the log-curvature $(\log \ddchis)'$
lies in $[\alpha, \beta]$ rather than $[-\beta, \beta]$.
Level~1 relates the link $\dchis$
to the curvature $\ddchis$.
All results in this paper use Level~2 only,
except \Cref{lem:variance-equiv},
which also requires Level~1
to convert Bregman convergence
(in $\ddchis$-space)
to weight convergence
(in $\dchis$-space).
For entropy ($\chis(\phi) = e^{\tau\phi}/\tau$),
all consecutive derivative ratios
equal $\tau$,
so $\alpha = \alpha_1 = \beta = \beta_1 = \tau$.

\begin{lemma}[Consequences of entropy-like structure]
\label{lem:e-type-bounds}
Suppose $\abs{h},\abs{h'} \le M$ and let $\Omega = e^{\beta M}$. If $\chis$ satisfies Equation~\ref{eq:e-type},
the following identities hold.
\begin{alignat}{2}
e^{-\beta|h|}\, \ddchis(\phi)
  &\le \ddchis(\phi + h)
  \le e^{\beta|h|}\, \ddchis(\phi)
  &&\qquad \textup{(Curvature ratio)}
  \label{eq:curvature-ratio-psc}\\
\frac{e^{-\beta|h|}}{2}\, \ddchis(\phi)\, h^2
  &\le D_{\chis}(\phi+h, \phi)
  \le \frac{e^{\beta|h|}}{2}\, \ddchis(\phi)\, h^2
  &&\qquad \textup{(Bregman sandwich)}
  \label{eq:bregman-sandwich}\\
|\dchis(\phi+h) - \dchis(\phi+h')|
  &\le \Omega\, \ddchis(\phi)\, |h - h'|
  &&\qquad \textup{(Derivative Lipschitz)}
  \label{eq:derivative-lipschitz}\\
\P\, \ddchis(\phi)\, h^2
  &\le 2\Omega\, \P\, D_{\chis}(\phi+h, \phi)
  &&\qquad \textup{(Weighted $L_2$)}
  \label{eq:weighted-l2-from-bregman}\\
D_{\chis}(\phi + v, \phi)
  &\ge \frac{\ddchis(\phi)}{\alpha^2}\,(e^{\alpha v} - \alpha v - 1)
  &&\qquad \textup{(Curvature growth)}
  \label{eq:curvature-growth-bregman}
\end{alignat}
where the last inequality holds for $v \ge 0$.
\end{lemma}

\begin{proof}
\paragraph{Curvature ratio~\eqref{eq:curvature-ratio-psc}.}
Pseudo-self-concordance gives
$|(\log \ddchis)'| = |\dddchis/\ddchis| \le \beta$.
Integrating from $\phi$ to $\phi + h$,
\[
\abs{\log \frac{\ddchis(\phi + h)}{\ddchis(\phi)}}
  = \abs{\int_\phi^{\phi+h} (\log \ddchis)'(s)\, ds}
  \le \beta |h|.
\]
Exponentiating gives~\eqref{eq:curvature-ratio-psc}.

\paragraph{Bregman sandwich~\eqref{eq:bregman-sandwich}.}
By the fundamental theorem of calculus,
$D_{\chis}(\phi+h, \phi)
  = \int_0^h \{\dchis(\phi + s) - \dchis(\phi)\}\, ds$.
Applying it again to
$\dchis(\phi + s) - \dchis(\phi)
  = \int_0^s \ddchis(\phi + t)\, dt$
gives
\begin{equation}
\label{eq:bregman-ftc}
D_{\chis}(\phi+h, \phi)
  = \int_0^h \int_0^s \ddchis(\phi + t)\, dt\, ds.
\end{equation}
Because $|t| \le |h|$,
the curvature ratio~\eqref{eq:curvature-ratio-psc}
gives
$e^{-\beta|h|}\, \ddchis(\phi)
  \le \ddchis(\phi + t)
  \le e^{\beta|h|}\, \ddchis(\phi)$,
so
\[
e^{-\beta|h|}\, \ddchis(\phi)
  \int_0^h \int_0^s dt\, ds
\le D_{\chis}(\phi+h, \phi)
\le e^{\beta|h|}\, \ddchis(\phi)
  \int_0^h \int_0^s dt\, ds.
\]
Evaluating
$\int_0^h \int_0^s dt\, ds = h^2/2$
gives the Bregman sandwich~\eqref{eq:bregman-sandwich}.

\paragraph{Derivative Lipschitz~\eqref{eq:derivative-lipschitz}.}
By the fundamental theorem of calculus,
the difference $\dchis(\phi+h) - \dchis(\phi+h')$
equals the integral of the curvature between $h'$ and $h$:
\[
|\dchis(\phi+h) - \dchis(\phi+h')|
  = \abs{\int_{h'}^h \ddchis(\phi + u)\, du}
  \le \sup_{|u| \le M} \ddchis(\phi + u) \cdot |h - h'|.
\]
The supremum ranges over all $|u| \le M$,
and the curvature ratio~\eqref{eq:curvature-ratio-psc} gives
$\ddchis(\phi + u) \le e^{\beta|u|}\, \ddchis(\phi)
  \le e^{\beta M}\, \ddchis(\phi) = \Omega\, \ddchis(\phi)$,
yielding the derivative Lipschitz bound~\eqref{eq:derivative-lipschitz}.

\paragraph{Weighted $L_2$~\eqref{eq:weighted-l2-from-bregman}.}
The lower bound in the Bregman sandwich~\eqref{eq:bregman-sandwich}
gives
$\ddchis(\phi)\, h^2
  \le 2 e^{\beta|h|}\, D_{\chis}(\phi+h, \phi)$.
Because $|h| \le M$,
the exponential satisfies $e^{\beta|h|} \le e^{\beta M} = \Omega$,
so $\ddchis(\phi)\, h^2 \le 2\Omega\, D_{\chis}(\phi+h, \phi)$.
Integrating both sides against~$\P$ gives the weighted $L_2$ bound~\eqref{eq:weighted-l2-from-bregman}.

\paragraph{Curvature growth~\eqref{eq:curvature-growth-bregman}.}
Integrating $(\log \ddchis)' \ge \alpha$
from $\phi$ to $\phi + t$
gives
$\ddchis(\phi + t) \ge \ddchis(\phi)\, e^{\alpha t}$
for $t \ge 0$.
Substituting into the double integral~\eqref{eq:bregman-ftc},
\[
D_{\chis}(\phi+v, \phi)
  = \int_0^v \int_0^s \ddchis(\phi + t)\, dt\, ds
  \ge \ddchis(\phi) \int_0^v \int_0^s e^{\alpha t}\, dt\, ds
  = \ddchis(\phi) \int_0^v
      \frac{e^{\alpha s} - 1}{\alpha}\, ds
  = \frac{\ddchis(\phi)}{\alpha^2}(e^{\alpha v} - \alpha v - 1). \qedhere
\]
\end{proof}

\begin{lemma}[Quadratic lower bound for Bregman divergences]
\label{lem:quadratic-bregman}
Let $\chi:\R \to \R$ be convex and twice continuously differentiable with $\chi''$ convex. Then
\begin{equation}
\label{eq:bregman-quadratic}
D_\chi(f+h, f) \ge \frac{1}{2}\chi''\qty(f + \frac{h}{3})h^2.
\end{equation}
\end{lemma}

\begin{proof}
We work with the double-integral characterization of Bregman divergence \eqref{eq:bregman-ftc}:
\[
D_\chi(f+h,f)
  = \int_0^1 \{\chi'(f + th) - \chi'(f)\}\,h\,dt
  = \int_0^1 \int_0^t \chi''(f + uh)\,du\,dt \cdot h^2.
\]
The double integral ranges over the triangle
$\{(u,t) : 0 \le u \le t \le 1\}$,
which has area $1/2$
and barycenter $\E U = 1/3$.
Writing the integral as $\frac{1}{2}\E\chi''(f + Uh)\,h^2$
where $U$ is uniform on this triangle,
Jensen's inequality gives
$\frac{1}{2}\E\chi''(f + Uh)
  \ge \frac{1}{2}\chi''(f + h\E U)
  = \frac{1}{2}\chi''(f + h/3)$.
\end{proof}

\begin{lemma}[PSC discount]
\label{lem:psc-discount}
Suppose $\chis$ satisfies
the pseudo-self-concordance condition~\eqref{eq:e-type}
with constant $\beta$,
and the interpolation inequality~\eqref{eq:interpolation}
holds with exponent $\theta$ and constant $\kappa$.
For $h$ with $\rho(h) \le s$,
write $r = \norm{h}_{\chis}$ and
\begin{equation}
\label{eq:psc-discount}
B_s(r) = \frac{\beta\kappa\, s^\theta\, r^{1-\theta}}{3}.
\end{equation}
Then
\begin{equation}
\label{eq:psc-discount-bound}
\P\, D_{\chis}(\tilde\phi + h, \tilde\phi)
  \ge \tfrac{1}{2}\, e^{-B_s(r)}\, r^2.
\end{equation}
\end{lemma}

The map $r \mapsto e^{-B_s(r)}\, r$
is non-decreasing when $(1-\theta)\, B_s(r) \le 1$,
because
\begin{equation}
\label{eq:psc-monotonicity}
[e^{-B_s(r)}\, r]'
  = e^{-B_s(r)}\{1 - (1-\theta)\, B_s(r)\} \ge 0.
\end{equation}
This condition controls the critical radius arguments
in the offset complexity bounds below.

\begin{proof}
\Cref{lem:quadratic-bregman} applied coordinatewise gives
\begin{equation}
\label{eq:psc-discount-step1}
D_{\chis_Z}(\tilde\phi + h, \tilde\phi)
  \ge \frac{1}{2}\,\ddchis(\tilde\phi + h/3)\, h^2.
\end{equation}
The interpolation inequality~\eqref{eq:interpolation}
and the constraint $\rho(h) \le s$ give
$\norm{h}_{L_\infty}
  \le \kappa\,\rho(h)^\theta\,\norm{h}_{\chis}^{1-\theta}
  \le \kappa\, s^\theta\, r^{1-\theta}$.
The curvature ratio bound~\eqref{eq:curvature-ratio-psc}
then gives
\begin{equation}
\label{eq:psc-discount-pointwise}
\ddchis(\tilde\phi + h/3)
  \ge e^{-\beta\norm{h}_{L_\infty}/3}\,\ddchis(\tilde\phi)
  \ge e^{-B_s(r)}\,\ddchis(\tilde\phi).
\end{equation}
Substituting the curvature bound~\eqref{eq:psc-discount-pointwise}
into the pointwise lower bound~\eqref{eq:psc-discount-step1}
and integrating against $\P$ gives
\[
\P\, D_{\chis}(\tilde\phi + h, \tilde\phi)
  \ge \frac{1}{2}\, e^{-B_s(r)}\,
    \P[\ddchis(\tilde\phi)\, h^2]
  = \frac{1}{2}\, e^{-B_s(r)}\, r^2.
  \qedhere
\]
\end{proof}

For Sobolev classes, the Gagliardo--Nirenberg inequality provides $L_\infty$ control from a Sobolev seminorm and $L_2$ norm.

\begin{lemma}[Sobolev interpolation]
\label{lem:sobolev-interpolation}
Let $\mathcal{X} \subseteq \R^d$ be bounded with Lipschitz boundary, $\mu$ be Lebesgue measure,
and let $s > d/2$. There exists a constant $C = C(s, d, \mathcal{X})$ such that for all $h$,
\[
\|h\|_{L_\infty} \le C \rho_s(h)^\theta \|h\|_{L_2(\mu)}^{1-\theta} \quad \text{where } \ \theta = \frac{d}{2s} \ \text{ and } \ \rho_s(h) = \|(-\Delta)^{s/2} h\|_{L_2(\mu)}.
\]
\end{lemma}

If $\P$ has density $p \ge \pmin > 0$ with respect to $\mu$, then $\|h\|_{L_2(\mu)} \le \pmin^{-1/2} \|h\|_{\LtP}$, giving the interpolation inequality~\eqref{eq:interpolation} with $\kappa = C\pmin^{-(1-\theta)/2}$:
\begin{equation}
\label{eq:sobolev-interpolation}
\|h\|_{L_\infty} \le C \pmin^{-(1-\theta)/2} \rho_s(h)^\theta \|h\|_{\LtP}^{1-\theta}.
\end{equation}

\subsection{Population-averaged Bregman lower bounds}

The results in this subsection
bound the population-averaged Bregman divergence
$\P D_{\chis}(\tilde\phi + h, \tilde\phi)$
from below
for perturbations $h$ with $\rho(h) \le s$,
where $\tilde\phi$ minimizes the population dual loss $L^\conj_{\P}$,
$V = (\ker\rho)^{\perp}$ is the orthogonal complement
of the kernel of $\rho$ in $\LtP$,
and $\norm{h}_{\chis} = (\P[\ddchis(\tilde\phi)\, h^2])^{1/2}$
is the curvature-weighted $L_2$ norm.

\begin{lemma}[Jensen lower bound for curvature growth]
\label{lem:jensen-lower}
Suppose $\chis$ is convex with
the curvature growth condition
$\dddchis \ge \alpha\, \ddchis$
for some $\alpha > 0$,
the interpolation inequality~\eqref{eq:interpolation}
holds with exponent $\theta$ and constant $\kappa$,
and $0 < \underline{\ddchis} \le \ddchis(\tilde\phi) \le \overline{\ddchis}$
under~$\P$.
Let $h \in V$
with $\P h = 0$, $\rho(h) \le s$,
and $\norm{h}_{\chis} = r$.
Write
\[
c_1 = \frac{\underline{\ddchis}}
  {2\overline{\ddchis}^2\, \kappa},
\qquad
c_0 = \frac{\underline{\ddchis}}
  {8\overline{\ddchis}\, \kappa^2}.
\]
\begin{enumerate}
\item[\textup{(i)}]
  \textup{(Quadratic floor.)}
  \begin{equation}
  \label{eq:jensen-quadratic-floor}
  \P\, D_{\chis}(\tilde\phi + h, \tilde\phi)
    \ge c_J\, \frac{r^{2+4\theta}}{s^{4\theta}},
  \qquad
  c_J = \frac{\underline{\ddchis}^3}
    {16\overline{\ddchis}^5\, \kappa^4}.
  \end{equation}

\item[\textup{(ii)}]
  \textup{(Exponential growth.)}
  If $\alpha\, c_1\, r^{1+\theta}/s^\theta \ge 2$, then
  \begin{equation}
  \label{eq:jensen-bregman-lower}
  \P\, D_{\chis}(\tilde\phi + h, \tilde\phi)
    \ge \frac{c_0}{\alpha^2}\,
      \frac{r^{2\theta}}{s^{2\theta}}\,
      e^{\alpha\, c_1\, r^{1+\theta}/s^\theta}.
  \end{equation}
\end{enumerate}
\end{lemma}

\begin{proof}
The idea: because $\P h = 0$,
the positive part of $h$
has substantial mass,
which Jensen's inequality converts
into exponential growth
of the Bregman divergence.

Write $w = \ddchis(\tilde\phi)$
and $S_{\pm} = \{h \gtrless 0\}$.

\paragraph{Positive-part mass.}
Because $\P h = 0$,
the positive and negative parts balance:
$\int_{S_+} h\, d\P = \int_{S_-} |h|\, d\P$.
The interpolation inequality~\eqref{eq:interpolation}
and $\rho(h) \le s$ give
$\norm{h}_{L_\infty}
  \le \kappa\, s^\theta\, r^{1-\theta}$.
Because $w \le \overline{\ddchis}$,
\begin{equation}
\label{eq:f-plus-mass}
r^2
  = \int h^2\, w\, d\P
  \le \norm{h}_\infty \int |h|\, w\, d\P
  \le \overline{\ddchis}\, \kappa\, s^\theta\, r^{1-\theta}
    \cdot 2\int_{S_+} h\, d\P,
\end{equation}
giving
\begin{equation}
\label{eq:positive-mass-lower}
\int_{S_+} h\, d\P
  \ge \frac{r^{1+\theta}}{2\overline{\ddchis}\, \kappa\, s^\theta}.
\end{equation}

\paragraph{Jensen's inequality.}
Drop $S_-$ (where $D_{\chis} \ge 0$)
and apply
\Cref{lem:e-type-bounds}(v)
on $S_+$ (where $h \ge 0$):
\[
\P\, D_{\chis}(\tilde\phi + h, \tilde\phi)
  \ge \frac{1}{\alpha^2} \int_{S_+}
    w\, (e^{\alpha h} - \alpha h - 1)\, d\P.
\]
The function $v \mapsto e^{\alpha v} - \alpha v - 1$
is convex, so Jensen's inequality
on $S_+$ with the probability measure
$w\, d\P / W_+$ (where $W_+ = \int_{S_+} w\, d\P$)
gives
\[
\P\, D_{\chis}(\tilde\phi + h, \tilde\phi)
  \ge \frac{W_+}{\alpha^2}\,
    (e^{\alpha \bar v} - \alpha \bar v - 1),
\]
where $\bar v = \int_{S_+} h\, w\, d\P / W_+$.

\paragraph{The two regimes.}
The weighted mean satisfies
$\bar v \ge c_1\, r^{1+\theta}/s^\theta$
by~\eqref{eq:positive-mass-lower},
$w \ge \underline{\ddchis}$,
and $W_+ \le \overline{\ddchis}$.

\paragraph{Quadratic floor~\eqref{eq:jensen-quadratic-floor}.}
The bound $e^v - v - 1 \ge v^2/2$ gives
\[
\frac{W_+}{\alpha^2}(e^{\alpha \bar v} - \alpha \bar v - 1)
  \ge \frac{W_+}{2}\, \bar v^2
  \ge \frac{\underline{\ddchis}\, \P(S_+)}{2}\,
    c_1^2\, \frac{r^{2+2\theta}}{s^{2\theta}}.
\]
The mass $\P(S_+) \ge r^{2\theta} / (2\overline{\ddchis}\, \kappa^2\, s^{2\theta})$
follows from~\eqref{eq:positive-mass-lower}
and $\norm{h}_\infty \le \kappa\, s^\theta\, r^{1-\theta}$,
giving~\eqref{eq:jensen-quadratic-floor}.

\paragraph{Exponential growth~\eqref{eq:jensen-bregman-lower}.}
When $\alpha\, c_1\, r^{1+\theta}/s^\theta \ge 2$,
the bound $e^v - v - 1 \ge e^v/4$ gives
\[
\frac{W_+}{\alpha^2}(e^{\alpha \bar v} - \alpha \bar v - 1)
  \ge \frac{W_+}{4\alpha^2}\,
    e^{\alpha\, c_1\, r^{1+\theta}/s^\theta}.
\]
Because $W_+ \ge \underline{\ddchis}\, \P(S_+)
  \ge 4\, c_0\, r^{2\theta}/s^{2\theta}$,
this gives~\eqref{eq:jensen-bregman-lower}.
\end{proof}


\subsection*{Guide to results}

The appendix is organized by proof technique:
duality (\Cref{sec:duality}),
properties of the regularized projection (\Cref{sec:regularized-projections}),
stable lower bounds for Bregman divergences (\Cref{sec:uniform-lower-bounds}),
bias bounds (\Cref{sec:bias-bound}),
offset complexity (\Cref{sec:offset-complexity}),
rates (\Cref{sec:rate-of-convergence}),
variance equivalence (\Cref{sec:variance-equiv}),
and asymptotic normality (\Cref{sec:asymptotics}).
The following table maps practical questions
to the relevant results.

\medskip

\begin{center}
\setlength{\tabcolsep}{4pt}
\renewcommand{\arraystretch}{1.3}
\begin{tabular}{lll}
\textbf{Question} & \textbf{Result} & \textbf{Section} \\
\hline
Bias at estimated weights & \Cref{thm:bias-decomp} & \S\ref{sec:bias-bound} \\
Rate $\hat\phi \to \tilde\phi$ (general) & \Cref{thm:rate-main} & \S\ref{sec:rate-of-convergence} \\
Rate (Sobolev models) & \Cref{cor:sobolev-main}, \Cref{cor:rkhs} & \S\ref{sec:rkhs} \\
Rate (Q$_+$, Sobolev) & \Cref{cor:sobolev-two-pile} & \S\ref{sec:sobolev-two-pile} \\
Variance equivalence (Setting~E) & \Cref{lem:variance-equiv} & \S\ref{sec:variance-equiv} \\
CLT (abstract conditions) & \Cref{thm:asymptotic-normality} & \S\ref{sec:asymptotics} \\
CLT (Sobolev, quadratic family) & \Cref{thm:general-clt} & \S\ref{sec:asymptotics} \\
Tuning window & \Cref{cor:tuning-window} & \S\ref{sec:asymptotics} \\
\hline
Treatment-specific mean & & \S\ref{sec:app-tsm} \\
Average linear regression & & \S\ref{sec:app-regression} \\
Variance of the CATE & & \S\ref{sec:app-vcate} \\
\end{tabular}
\end{center}

\section{Duality}
\label{sec:duality}

Strong duality and existence of optimal weights require finite-dimensional structure. In this section, weights are vectors $\gamma \in \R^n$ connected to functions via a linear map $L: \calS \to \R^n$. This specialization, together with a general penalty $\zeta$, yields existence and uniqueness of the primal minimizer and a characterization of the optimal weights via any dual minimizer or minimizing sequence.

\begin{proposition}
\label{prop:duality-full}
Let $\calS$ be a vector space with a seminorm $\rho$ with finite dimensional kernel $K$. Furthermore, let $\ell_0 : \calS \to \R$ and $L:\calS \to \R^n$ be linear maps with the property that $\ell_\gamma(\cdot) = \ell_0(\cdot) - \gamma^T L(\cdot)$ is continuous with respect to $\rho$ for all $\gamma \in \R^n$ for which it is identically zero on the kernel $K$.
Let $\zeta: \R_{\ge 0} \to \R_{\ge 0} \cup \{\infty\}$ and $\chi:\R^n \to \R_{\ge 0} \cup \{\infty\}$ be convex lower-semicontinuous functions with $\zeta$ nondecreasing and coercive and $\chi$ coercive and strictly convex on its effective domain. Define
\begin{align*}
p(\gamma) &:= \chis(\gamma) + \zeta\qty{\sup_{m \in \calS : \rho(m) \le 1} \ell_{\gamma}(m)}, \\
d(f) &:= \chis(Lf) - \ell_0(f) + \zetas\{\rho(f)\}.
\end{align*}
If there exists some $\gamma \in \R^n$ for which
$p(\gamma) < \infty$ and $\zeta(x)$ is continuous at the point $x(\gamma)=\sup_{m \in \calS : \rho(m) \le 1}\ell_{\gamma}(m)$, then the following hold.
\begin{enumerate}
    \item $p$ has a unique minimum at a vector $\hgamma \in \R^n$.
    \item $\hgamma = (\dchis \comp L)(\hat\phi)$ for any minimizer $\hat\phi$ of $d$, and $(\dchis \comp L)(f_j) \to \hgamma$ for any minimizing sequence $(f_j)$.
\end{enumerate}
\end{proposition}

\subsection{Application to balancing weights.} This result applies to the problem of estimating the treatment-specific mean $\mu_1 = \E m_1(X)$ using balancing weights. In this setting, we observe treatment-covariate pairs $(W_i, Z_i)$ for $i = 1, \ldots, n$, and we take
\[
L(m) = [m(X_1), \ldots, m(X_n)]^T, \qquad \ell_0(m) = \sum_{i=1}^n W_i m(Z_i).
\]
The primal $p(\gamma)$ penalizes dispersion of the weights $\gamma$ via $\chi$ and maximal imbalance for $m \in \model := \{ m : \rho(m) \le 1 \}$ via $\zeta$.

\paragraph{Continuity.}
To illustrate the continuity condition, we consider two model classes, each with a norm and seminorm variant.
\[
\begin{aligned}
&\textit{Linear model:}
  &&m(x) = \beta_0 + \beta^T x,
  &&&\rho(m) = |\beta_0| + \|\beta\|_1 \text{ (norm)}
  \quad\text{or}\quad
  \rho(m) = \|\beta\|_1 \text{ (seminorm)}.
  \\
&\textit{Sobolev model:}
  && m: \R \to \R,
  &&&\rho^2(m) = \|m\|_{L_2}^2 + \|m^{(s)}\|_{L_2}^2 \text{ (norm)}
  \quad\text{or}\quad
  \rho^2(m) = \|m^{(s)}\|_{L_2}^2 \text{ (seminorm)}.
\end{aligned}
\]
The norm penalizes overall size (the intercept $\beta_0$ or the $L_2$ norm of $m$), while the seminorm does not: constants have zero linear seminorm, and polynomials of degree $< s$ have zero Sobolev seminorm. In the case $s=1$, the Sobolev seminorm gives a model with unpenalized intercept. We can define Sobolev norms for multivariate functions $m: \R^p \to \R$ analogously, with powers $\laplacian^{s/2}$ of the Laplacian playing the role of the $s$th derivative, and multivariate polynomials of order $< s$ have zero seminorm.

\emph{Norm case.} When $\rho$ is a norm, it has an empty kernel and we require continuity for all $\gamma$, i.e.\ that point evaluation functionals $m \mapsto m(Z_i)$ are continuous with respect to $\rho$ for $X_i$ in our sample. What this means is that we are requiring that $m(X_1) \ldots m(X_n)$ be small when $\rho(m)$ is.
\begin{itemize}
\item For our one-norm-constrained linear model, H\"older's inequality tells us that this holds when $\|X_i\|_\infty$ is bounded. And in this context, where we're asking only for the existence of a bound on the particular sample we have, this is essentially vacuous: all this means is that $X_1 \ldots X_n$ are in $\R^p$ rather than an extended version of $\R^p$ that includes infinity.
\item For the Sobolev model, what we are asking is that if a function is typically small and doesn't vary much, it's small at our sample points. This is not vacuous---a `spike function' that is zero virtually everywhere but jumps up to one at a single point $X_i$ can have arbitrarily small Sobolev norm---but the Sobolev embedding theorem tells us it is satisfied when the number of bounded derivatives $s$ exceeds half the covariate dimension $p$ \citep[Theorem 4.12]{adams2003sobolev}.
\end{itemize}
\emph{Seminorm case.} When $\rho$ is a seminorm with nontrivial kernel $K$, the primal forces $\gamma$ to balance the kernel perfectly: we must have $\sum_i (\gamma_i - W_i) k(Z_i) = 0$ for all $k \in K$. For such $\gamma$, the imbalance $\sum_i (\gamma_i - W_i) m(Z_i)$ is invariant to adding elements of $K$ to $m$---a generalization of translation-invariance.
\begin{itemize}
\item For the linear model with $\rho(m) = \|\beta\|_1$, the kernel is constant functions, so perfect balance requires only $\sum_i \gamma_i = \sum_i W_i$. Writing $m(x) = \beta_0 + \beta^T x$, the imbalance becomes $\beta^T \sum_i (\gamma_i - W_i) X_i$, which depends only on the slope vector $\beta$. H\"older's inequality bounds this when $\|X_i\|_\infty$ is bounded---the same condition as the norm case.
\item For the Sobolev model with $\rho^2(m) = \|m^{(s)}\|_{L_2}^2$, the kernel consists of polynomials of degree $< s$, so perfect balance requires $\sum_i (\gamma_i - W_i) q(Z_i) = 0$ for all such $q$. The polynomial contribution to imbalance vanishes, so we can replace $m$ with its remainder $r = m - q$ after subtracting the best-fitting polynomial. The Sobolev embedding theorem bounds $\|r\|_\infty$ in terms of $\rho(m)$, giving the same dimension condition $s > p/2$ as the norm case.
\end{itemize}

\paragraph{Conditions on dispersion penalties.}
A common choice is $\chis(\gamma) = \|\gamma\|^2/2$, which is strictly convex on all of $\R^n$. Sometimes a nonnegativity constraint $\gamma \ge 0$ is added; in our framework, this can be rewritten by taking $\chis(\gamma) = \sum_i \gamma_i^2/2$ for  $\gamma \ge 0$ (coordinatewise) and $+\infty$ otherwise. This $\chis$ is strictly convex on its effective domain---the set where it is finite---which is $\{ \gamma \ge 0\}$. It is also lower semicontinuous: it is continuous on the interior of its effective domain, and at the boundary it jumps \emph{up} to $+\infty$---the `lower semi-' in lower semicontinuity means that jumps upward but not downward are allowed.
In general, strict convexity of $\chis$ on its effective domain implies differentiability of $\chi$ everywhere.

For the constrained quadratic, the conjugate is $\chi(\theta) = \|\theta_+\|^2/2$ with gradient $\nabla \chi(\theta)= \theta_+$, the (coordinatewise) positive part of the vector $\theta$. This gradient maps negative components to zero, reflecting that the primal solution is pinned to zero on the boundary of the effective domain. The entropy penalty $\chis(\gamma) = \sum_i (\gamma_i \log \gamma_i - \gamma_i)$ is similar. The formula is only defined for $\gamma \ge 0$, so when people write it they implicitly mean that it's $+\infty$ elsewhere, and on its effective domain $\{ \gamma \ge 0\}$ it is strictly convex.

\subsection{General linear functionals}
The treatment-specific mean example above uses $\dotpsiZ(m) = m\{(w, X)\}$, but our framework accommodates any linear functional $\dotpsiZ$. For instance, taking $\dotpsiZ(m) = \partial m_w(X)/\partial w |_{w=W}$ targets the average marginal effect of a continuous treatment, and $\dotpsiZ(m) = \int m_1(x) p(x) dx$ for a target density $p$ addresses covariate shift between sample and target populations \citep{hirshberg2018augmented}. In each case, the Riesz representer $\griesz$ satisfying $\P \griesz(Z) m(Z) = \P \dotpsiZ(m)$ plays the role that the inverse propensity weights play in the treatment-specific mean setting.

The continuity conditions generalize accordingly: we require $m \mapsto \dotpsiZ(m)$ to be continuous with respect to $\rho$, or---in the seminorm case---continuous on the quotient $\calS / K$. For our linear model examples, this amounts to asking that $\dotpsiZ$ not place mass on individual points in a way that evades the smoothness constraint. The Sobolev embedding condition $s > p/2$ ensures this for point-evaluation functionals; for smoother functionals like integrals against densities, weaker conditions suffice.

\subsection{Dual Bregman identity}
The Bregman divergences of $\chi$ and $\chis$ are related by
\begin{equation}
\label{eq:dual-bregman-identity}
D_\chi(a, b) = D_{\chis}\{\chi'(b), \chi'(a)\},
\end{equation}
which follows from the identity $(\chis)' = (\chi')^{-1}$. This connects the dual loss $d(g)$, which involves $D_{\chis}$, to the primal interpretation in terms of weights $\gamma = \dchis(g)$.

For the unconstrained quadratic $\chis(\theta) = \|\theta\|^2/2$, the identity is trivial: $\chi = \chis$ and both divergences equal $\|a-b\|^2/2$. For the positive-part quadratic $\chis(\theta) = \|\theta_+\|^2/2$, the conjugate is $\chi(\gamma) = \|\gamma\|^2/2$ for $\gamma \ge 0$ and $+\infty$ otherwise---not self-dual. The dual Bregman divergence satisfies $D_{\chis}(g, \tilde\phi) = (g - \tilde\phi)^2/2$ when both $g, \tilde\phi \ge 0$, and vanishes when both $g, \tilde\phi < 0$. This ``dead zone'' means only the positive region $\{\tilde\phi \ge 0\}$ contributes to the Bregman divergence, leading to effective dimensions restricted to the active region in our rate theorems.

For the entropy penalty $\chis(\theta) = \sum_i (\theta_i \log \theta_i - \theta_i)$, the primal Bregman divergence $D_\chi$ is relative entropy (KL divergence) between the weights. The dual $D_{\chis}$ has exponential growth in $\theta - \tilde\phi$, but \eqref{eq:dual-bregman-identity} shows this corresponds to KL divergence of the implied weights $\dchis(\theta)$. Thus, uniform lower bounds on $\hP D_{\chis}(g, \tilde\phi)$ translate to bounds on how close the implied weights $\dchis(g)$ are to the target $\dchis(\tilde\phi)$, measured in primal Bregman divergence.

\subsection{Proofs}
\label{sec:duality-proofs}

The proof works in an abstract bias-dispersion framework that also underlies the bias analysis (\Cref{sec:bias-bound}). Let $V$ be a vector space with a seminorm $\rho$ and $W$ a space of weights, paired by a bilinear form $\inner{\cdot}{\cdot}: W \times V \to \R$. The function class is the unit ball $B_\rho = \{f \in V : \rho(f) \le 1\}$. The target functional $\psi: V \to \R$ is linear. Let $\chi: W \to \R \cup \{+\infty\}$ be convex, lower-semicontinuous, and strictly convex on its effective domain, with conjugate $\chis(\phi) = \sup_{\gamma \in W} \inner{\gamma}{\phi} - \chi(\gamma)$. For duality, we specialize to $W = \R^n$ and $\inner{\gamma}{f} = \gamma^T L(f)$ for a linear map $L: V \to \R^n$. The finite-dimensionality of $W$ is essential: it gives coercivity of the primal, compactness in the bidual (Banach-Alaoglu), and the Goldstine density argument below.

\begin{lemma}[Existence of minima for coercive convex functions]
\label{lem:coercive-minimum}
A proper, lower-semicontinuous, coercive, strictly convex function
on a reflexive Banach space
has a unique minimizer.
\end{lemma}
\begin{proof}
See \citet[Theorem~2.19]{peypouquet2015convex}.
\end{proof}

\begin{lemma}[Fenchel--Rockafellar duality]
\label{lem:fenchel-rockafellar}
Let $s$ and $r \comp A$ be proper, convex, lower-semicontinuous
functions on a Banach space,
where $A$ is a bounded linear operator.
If $p(x) = s(x) + r(Ax)$ has a finite minimum $\hat p$,
then the dual
$d(z^{**}) = s^*(-A^* z^{**}) + r^*(z^{**})$
also attains its minimum,
$\min d = \hat p$,
and every dual minimizer $\hat z^{**}$
satisfies $-A^* \hat z^{**} \in \partial s(\hat x)$
at the primal minimizer $\hat x$.
\end{lemma}
\begin{proof}
See \citet[Theorem~3.50]{peypouquet2015convex}.
\end{proof}

\begin{lemma}[Goldstine's theorem]
\label{lem:goldstine}
For a Banach space $X$,
the canonical image of the unit ball of $X$
is weak$^*$-dense
in the unit ball of the bidual $X^{**}$.
\end{lemma}
\begin{proof}
See \citet[Theorem~2.6.26]{megginson2012introduction}.
\end{proof}

\begin{proof}[Proof of
Proposition~\ref{prop:duality-full}]
The gradient $\nabla\chis$ is continuous on all of $\R^n$. Strict convexity of $\chi$ on its effective domain implies $\chis$ is differentiable where it is finite, coercivity of $\chi$ ensures $\chis$ is finite everywhere, and a differentiable convex function on $\R^n$ has a continuous gradient.

Because $p$ is a proper, strictly convex, coercive, and lower-semicontinuous on the reflexive space $\R^n$, it has a unique minimizer $\hgamma \in \R^n$ (\Cref{lem:coercive-minimum}). The solution $\hgamma$ must satisfy $\ell_{\hgamma}(k)=0$ for all $k \in K$:
if it did not for some $k \in K$, then $\hat p$ would not be finite, as $\rho(\alpha k) = 0$ for all $\alpha$ and as $\abs{\alpha} \to \infty$, $\abs{\ell_{\hgamma}(\alpha k)} = \abs{\alpha}\abs{\ell_{\hgamma}(k)}\to \infty$ and therefore $\zeta\{\ell_{\hgamma}(\alpha k)\} \to \infty$ because $\zeta$ is coercive. Thus, without changing our solution, we can restrict $\gamma$ to the affine subspace $\Gamma \subseteq \R^n$ of vectors with this property. Letting $L_K \subseteq \R^n$ be the image of the kernel $K$ under $L$ and $L_K^{\perp}\subseteq \R^n$ its orthogonal complement, we can characterize this subspace as $\Gamma := \hgamma_0 + L_K^\perp$ where $\hgamma_0$ is the unique element of $L_K$ for which $\ell_{\hgamma_0}$ is identically zero on $K$.

For all such $\gamma=\hgamma_0 + x$ with $x \in L_K^\perp$, $\ell_{\hgamma_0} - x^T L$
is the sum of two linear functionals on $\calS$ that vanish on the kernel $K$. These functionals, which we can
interpret as functionals on the quotient space $\calS_K := \calS / K$ which has $\rho$ as its norm, are continuous with respect
to $\rho$ by assumption and therefore in our quotient space's dual $\calS_K^*$. And we can
interpret the supremum $\sup_{m \in \calS : \rho(m) \le 1} \ell_\gamma (m)$ in the definition of $p$ as the
operator norm $\norm{\ell_\gamma}_*$ of $\ell_\gamma$ as an element of $\calS_K^*$. Thus, reparameterized in terms of $x \in L_K^\perp$, our primal $p : L_K^\perp \to \R$ is
\begin{equation}
\label{eq:fr-primal}
\begin{aligned}
p(x) &:= s(x) + r(Ax)
&& \text{ where } \\
s(x) &:= \chi\qty( \hgamma_{0} + x )
&&\text{ maps } \quad L_K^\perp \to \R \cup \{+\infty\}, \\
r(y^*) &:= \zeta\qty( \norm{\ell_{\hgamma_0} + y^* }_* )
&& \text{ maps } \calS_K^* \to \R \cup \{+\infty\}, \\
A(x) &:= -x^T L
&& \text{ maps } \quad L_K^\perp \to \calS_K^*.
\end{aligned}
\end{equation}


Because $p(x)$ has a minimum $\hat x = \hgamma - \hgamma_0$,
its Fenchel-Rockafellar dual
\begin{equation}
\label{eq:fr-dual}
    d_{\calS_K^{**}}(z^{**}) := s^{*}(-A^{*} z^{**}) + r^{*}(z^{**}) \quad \text{mapping} \quad \calS_K^{**} \to \R,
\end{equation}
has a minimum, and each argmin $\hat z^{**}$ satisfies the first order optimality condition $-A^{*} \hat z^{**} \in \partial s(\hat x)$ where $A^{*}$ is the adjoint of $A$ and $s^{*}$ and $r^{*}$ the convex conjugates of $s$ and $r$ (\Cref{lem:fenchel-rockafellar}).
We conclude by characterizing $\hat z^{**}$ and its relationship to our primal solution $\hgamma$ more concretely.
We first simplify $r^*$.
\begin{align*}
    r^{*}(z^{**}) &= \sup_{y^* \in \calS_K^*} z^{**}(y^*) - \zeta\qty(\norm{\ell_{\hgamma_0} + y^*}_{*}) \\
                       &= - z^{**}(\ell_{\hgamma_0}) + \sup_{v^* \in \calS_K^*} z^{**}(v^*) - \zeta\qty(\norm{v^*}_{*}) && \text{where } v^* = y^* + \ell_{\hgamma_0} \\
                       &= - z^{**}(\ell_{\hgamma_0}) + \sup_{t \in \R}\sup_{\substack{u^* \in \calS_K^* \\ \norm{u^*}_{*}=1}} t z^{**}(u^*) - \zeta(t) && \text{where } v^* = t u^* \\
                       &= - z^{**}(\ell_{\hgamma_0}) + \sup_{t \in \R}t\norm{z^{**}}_{**} - \zeta(t) && \text{by definition of } \norm{\cdot}_{**} \\
                       &= - z^{**}(\ell_{\hgamma_0}) + \zetas\qty(\norm{z^{**}}_{**}) && \text{by definition of } \zetas.
\end{align*}
Substituting this into our dual problem~\eqref{eq:fr-dual}, we get this equivalent characterization.
\begin{equation}
\label{eq:dual-for-bidual-elements}
    d_{\calS_K^{**}}(z^{**}) = s^{*}(-A^{*} z^{**}) - z^{**}(\ell_{\hgamma_0}) + \zetas\qty(\norm{z^{**}}_{**}).
\end{equation}

\paragraph{The Reflexive Case.} If the quotient space $\calS_K$ is reflexive, then every element $z^{**}$ of the bidual $\calS_K^{**}$ is an evaluation functional
corresponding to some coset $[z] \in \calS_K$ via $z^{**}(y^*) = y^*([z])$ for all $y^* \in \calS_K^*$. Now recall that $A: L_K^\perp \to \calS_K^{*}$
is the linear map defined by $Ax = -x^T L$  and its adjoint $A^{*}: \calS_K^{**} \to L_K^\perp$ therefore satisfies, for any such pair $(z^{**}, [z])$ the identity
\[
\inner{-A^{*} z^{**}}{x} = \inner{z^{**}}{-Ax} = \inner{-Ax}{[z]} \qqtext{ for all } x \in L_K^\perp.
\]
It follows, because $-Ax [z] = x^T L z$ for all $z \in [z]$, that all $z \in [z]$ satisfy $-A^{*} z^{**}=\Pi^\perp L z$; $\Pi^\perp L z$
and $-A^{*} z^{**}$ are vectors in $L_K^\perp$ whose inner products with all vectors $x \in L_K^\perp$ coincide, so they must be the same
vector.  It follows that $\hat z^{**}$ minimizes \eqref{eq:dual-for-bidual-elements} if and only if every element $\hat z$ in the corresponding coset minimizes the function on $\calS$ given by
\begin{equation}
\label{eq:dual-for-primal-elements}
    d_{[\cdot]}(z) = s^{*}(\Pi^\perp L z) - \ell_{\hgamma_0}(z) + \zetas\{\rho(z)\}.
\end{equation}
Furthermore, the first-order optimality condition from \eqref{eq:fr-dual} determines the $L_K^\perp$ component of $L\hat z$. Because $s(x) = \chi(\hat\gamma_0 + x)$, the chain rule gives $\nabla s(x) = \Pi^\perp \nabla\chi(\hat\gamma_0 + x)$: the gradient of $s$ is the $L_K^\perp$ component of the gradient of $\chi$. The first-order condition $\nabla s(\hat x) = -A^*\hat z^{**} = \Pi^\perp L\hat z$ therefore gives
\[
\Pi^\perp L\hat z = \Pi^\perp \nabla\chi(\hgamma).
\]
In intuitive terms, optimizing in the quotient space determines the $L_K^\perp$ component of $L\hat z$ but says nothing about the $L_K$ component. The within-coset optimization handles the $L_K$ component.

We conclude by showing that optimizing over $\calS$ itself does in fact choose the right kernel element.  To do this, we think of minimizing $z$ as a two-stage process. First, we choose the best representative $\hat z_{[z]}$ within each coset $[z] \in \calS_K$,
then we choose the best of those representatives, i.e.\ the best coset.
\[
\inf_{z \in \calS} d(z) = \inf_{[z] \in \calS_K} \inf_{z \in [z]} d(z).
\]
The infimum within each coset is attained at some element $\hat z_{[z]}$ because $d$ is strictly convex and the coset $[z] = z + K$ finite-dimensional and therefore reflexive. And each within-coset minimizer satisfies a first order condition with no $\rho$ term, because $\rho$ is constant within cosets.
\begin{equation}
\label{eq:within-coset-optimality}
(L|_K)^* \nabla \chis(L \hat z_{[z]}) = \ell_0|_K.
\end{equation}%
That is, the gradient $\xi = \nabla \chis(L \hat z_{[z]}) \in \R^n$ satisfies $\xi^T L k = \ell_0(k)$ for all $k \in K$,  i.e. $\nabla \chis(L \hat z_{[z]}) \in \Gamma = \hgamma_0 + L_K^\perp$. If $\hat z$ minimizes $d$ over all of $\calS$, it minimizes $d$ within its coset and therefore satisfies the within-coset first order condition \eqref{eq:within-coset-optimality}, i.e.\ $\nabla \chis(L\hat z) \in \Gamma$: the $L_K$ component $\Pi_K\nabla\chis(L\hat z) = \hat\gamma_0$. We will show that these two first order conditions---$\Pi^\perp L\hat z = \Pi^\perp\phi$ from the quotient and $\Pi_K\nabla\chis(L\hat z) = \hat\gamma_0$ from the within-coset problem---together imply $\nabla\chis(L\hat z) = \hgamma$.

\textit{Determining $\hgamma$.} Let $\phi = \nabla\chi(\hgamma)$ be the index of the primal optimizer, so that $\hgamma = \nabla\chis(\phi)$. We have established that $L\hat z$ and $\phi$ have the same $L_K^\perp$ component: $\Pi^\perp L\hat z = \Pi^\perp\phi$. Write $c = \Pi^\perp\phi$ for this common value and consider the within-coset first order condition at $c$, which asks: for which $w \in L_K$ does $\Pi_K\nabla\chis(w + c) = \hat\gamma_0$? The full minimizer $\hat z$ satisfies this at $w = \Pi_K L\hat z$, and $\phi$ satisfies it at $w = \Pi_K\phi$ because $\nabla\chis(\phi) = \hgamma \in \Gamma$. Because $\chis$ is strictly convex on its effective domain, the function $w \mapsto \chis(w + c) - \hat\gamma_0^T w$ is strictly convex on $L_K$, so its first order condition has at most one solution. Both $\Pi_K L\hat z$ and $\Pi_K\phi$ are solutions, so they are equal: $L\hat z$ and $\phi$ agree on both components, giving $\nabla\chis(L\hat z) = \nabla\chis(\phi) = \hgamma$.

Where $\chis$ is not strictly convex---on a flat of $\chis$ where its gradient is constant---the within-coset first order condition may have multiple solutions $w$, but $\nabla\chis(w + c)$ takes the same value at each, so $\nabla\chis(L\hat z) = \nabla\chis(\phi) = \hgamma$ regardless.

\textit{Summary.}
We have shown that if $\hat\gamma$ minimizes $p$ and $\hat z$ minimizes $d$, we have the primal-dual relationship $\hgamma = \nabla \chis(L\hat z)$. This establishes our claim in the reflexive case or, more generally, in the case that $d$ is minimized at some point $z \in \calS$. We will now extend this to characterize $\hgamma$ in terms of a \emph{minimizing sequence}
$g_1,g_2,\ldots \in \calS$ along which $d(g_j) \to \inf_{g}d(g)$.

\paragraph{The General Case.}
The strategy:
extract a weak$^*$-convergent subsequence
in the bidual
(Goldstine's theorem + Banach--Alaoglu),
show the limit minimizes the bidual objective,
then recover the primal solution $\hgamma$
via the within-coset first-order condition.

We pick up from \eqref{eq:dual-for-bidual-elements}, where we characterized $p$'s Fenchel-Rockafellar dual as a functional $d_{\calS_K^{**}}$ on the bidual $\calS_K^{**}$. Consider any sequence $z_1,z_2,\ldots$ in $\calS$ along which $d$ converges to its infimum. The cosets $[z_1], [z_2], \ldots$ in $\calS_K$ correspond to evaluation functionals $[z_1]^{**}, [z_2]^{**}, \ldots$ in the bidual---elements of the form $z^{**}(y^*) = y^*(z)$ for $z \in \calS_K$. Coercivity of $d$ implies $\rho(z_j)$ is bounded, so these lie in a bounded ball of $\calS_K^{**}$, which is weak$^*$-compact (Banach-Alaoglu), so some subsequence converges weak$^*$ to a limit $\hat z^{**}$.

This limit minimizes $d_{\calS_K^{**}}$. The unit ball of $\calS_K^{**}$ is weak$^*$-compact and $d_{\calS_K^{**}}$ is lower semicontinuous, so a minimizer exists. Goldstine's theorem (\Cref{lem:goldstine}) says evaluation functionals are weak$^*$-dense in the bidual, so this minimizer is a weak$^*$-limit of evaluation functionals; lower semicontinuity then implies their values get arbitrarily close to the minimum. Thus $\inf_\calS d = \inf_{\calS_K^{**}} d_{\calS_K^{**}}$, and $\hat z^{**}$---having value $\le \liminf d(z_j) = \inf_\calS d$---is a minimizer.

The map $A^*$ lands in the finite-dimensional space $L_K^\perp$, so weak$^*$ convergence $[z_j]^{**} \to \hat z^{**}$ implies that the $L_K^\perp$ components $\gamma_\perp^{(j)} = \Pi^\perp L z_j$ converge: these are the arguments to $s^*$ in \eqref{eq:dual-for-bidual-elements}. The limit $\gamma_\perp$ may depend on which weak$^*$-convergent subsequence we extracted above, because $s^*$ need not be strictly convex, so different minimizers $\hat z^{**}$ of $d_{\calS_K^{**}}$ may yield different $\gamma_\perp$ values. However, the primal minimizer $\hgamma$ is unique regardless: the within-coset argument below determines $\hgamma$ from any such $\gamma_\perp$. Where $\chis$ is not strictly convex, $\nabla \chis$ is constant on flats---all $x$ in a flat region satisfy the Fenchel-Young equality $\chis(x) + \chi(y) = x^T y$ at the same $y = \nabla\chis(x)$---so different $\gamma_\perp$ values within a flat region produce the same $\hgamma$. For example, if $\chi(\gamma) = \|\gamma\|^2/2$ for $\gamma \ge 0$ and $+\infty$ otherwise, then $\chis(x) = \|x_+\|^2/2$ where $x_+$ denotes the positive part; the gradient $\nabla\chis(x) = x_+$ is constant at zero for all $x \le 0$, mapping the entire flat region to the boundary $\gamma = 0$ of $\chi$'s effective domain.

\textit{Within-coset problem.} Recall the two-part structure from the reflexive case: we minimize $d$ first within each coset (optimizing the $L_K$ component of $Lz$), then over cosets (optimizing the $L_K^\perp$ component). The quotient-space argument above handled the second part, showing that $\Pi^\perp L z_j \to \gamma_\perp$. What remains is the first part---to show that $\nabla\chis(L z_j) \to \hgamma$.

For each $j$, let $\tilde z_j$ be the element of $[z_j]$ minimizing $d$. It satisfies the first order condition \eqref{eq:within-coset-optimality}: $\nabla\chis(L\tilde z_j) \in \Gamma$. By the same argument as in the reflexive case---both $L\tilde z_j$ and the index $\phi_j$ of the within-coset primal optimizer at $c_j = \Pi^\perp Lz_j$ satisfy the within-coset first order condition at $c_j$---the within-coset uniqueness gives $\nabla\chis(L\tilde z_j) = \hat\gamma(c_j)$, where $\hat\gamma(c)$ denotes the primal optimizer for the within-coset problem at $L_K^\perp$ component $c$. Because $c_j = \Pi^\perp Lz_j$ converges and $\hat\gamma(c)$ depends continuously on $c$ (by continuity of the gradient $\nabla\chis$ and of the strictly convex within-coset minimizer), $\nabla\chis(L\tilde z_j)$ converges to $\hgamma$. So if we can show $\nabla\chis(Lz_j) - \nabla\chis(L\tilde z_j) \to 0$, then
\begin{equation}\label{eq:gradient-chain}
\nabla\chis(Lz_j) \to \nabla\chis(L\tilde z_j) \to \hgamma.
\end{equation}


Along the minimizing sequence $z_j$, the gap $d(z_j) - d(\tilde z_j)$ must go to zero---otherwise we wouldn't be converging to the infimum. Furthermore, this gap is, in fact, the Bregman divergence $D_{\chis}(Lz_j, L\tilde z_j)$: $\rho$ is constant within cosets, and our within-coset first order condition implies $\ell_0(k) = \nabla\chis(L\tilde z_j)^T L k$ for all elements of the kernel including $z_j - \tilde z_j$. Summarizing,
\[
D_{\chis}(Lz_j, L\tilde z_j) = \chis(Lz_j) - \chis(L\tilde z_j) - \nabla\chis(L\tilde z_j)^T(Lz_j - L\tilde z_j) = d(z_j) - d(\tilde z_j) \to 0.
\]
Bregman divergences satisfy the duality identity $D_{\chis}(a,b) = D_{\chi}\{\nabla\chis(b), \nabla\chis(a)\}$ when these gradients exist, so it follows that $D_{\chi}\{\nabla\chis(L\tilde z_j), \nabla\chis(Lz_j)\} \to 0$. For strictly convex and coercive $\chi$, $D_{\chi}(p_j, q_j) \to 0$ with $p_j$ bounded implies $p_j - q_j \to 0$. And the sequence $\nabla\chis(L\tilde z_j)$ converges by~\eqref{eq:gradient-chain}, so $\nabla\chis(Lz_j) - \nabla\chis(L\tilde z_j) \to 0$. This completes our proof.
\end{proof}



\section{Properties of Regularized Projections}
\label{sec:regularized-projections}
The conditions in \Cref{thm:rate-main} involve the population minimizer $\tilde\phi$, a regularized Bregman projection. We need to know that $\tilde\phi$ is close to the population index $\phi$ (Bregman approximation), that the resulting weights don't collapse to zero (curvature control), and---for the positive-part case---that $\tilde\phi$ doesn't linger near the dead zone boundary (density control). This appendix derives these properties from primitives.
The population minimizer solves
\[
\tilde\phi = \argmin_f \P D_{\chis}(f, \phi) + \eta \rho(f)^2.
\]
Here $\phi = (\dchis)^{-1}(\griesz)$ is the Riesz representer's index. We take $\rho = |\cdot|_{\dot H^s}$, the homogeneous Sobolev seminorm, throughout this section. The implied weights are $\tilde\gamma = \dchis(\tilde\phi)$.

\subsection{General Bregman approximation}

\paragraph{Mollification.} Given a function $\phi$ and a \emph{mollifier}---a smooth, compactly supported kernel $\omega \ge 0$ with $\int \omega = 1$---the \emph{mollification} at scale $\epsilon$ is the convolution
\[
\phi_\epsilon(z) = \int \phi(z - u) \omega_\epsilon(u) \, du \qquad \text{where } \omega_\epsilon(u) = \epsilon^{-d} \omega(u/\epsilon).
\]
The mollification $\phi_\epsilon$ is smooth and converges pointwise to $\phi$ as $\epsilon \to 0$ \citep[Section 2.29]{adams2003sobolev}. Three properties are key:
\begin{enumerate}
\item \emph{Convex combination.} Because $\omega_\epsilon \ge 0$ and $\int \omega_\epsilon = 1$, the value $\phi_\epsilon(z)$ is a weighted average of $\phi$ over a neighborhood of $z$. Sets defined by pointwise bounds are preserved: if $a \le \phi \le b$, then $a \le \phi_\epsilon \le b$.
\item \emph{Pointwise convergence.} As $\epsilon \to 0$, $\phi_\epsilon(z) \to \phi(z)$ pointwise. Combined with uniform bounds, dominated convergence gives $\|\phi_\epsilon - \phi\|_{\LtP} \to 0$.
\item \emph{Seminorm control.} The Fourier transform satisfies $\hat\phi_\epsilon(\xi) = \hat\phi(\xi) \hat\omega(\epsilon\xi)$, so high frequencies are damped. For the $\dot H^s$ seminorm $|f|_{\dot H^s} = \|(-\Delta)^{s/2} f\|_{L_2(\R^d)}$, we have $|\phi_\epsilon|_{\dot H^s} \le C_\omega \epsilon^{-s} \|\phi\|_{L_2(\R^d)}$.
\end{enumerate}

The set $\F$ is not a constraint on the optimization---the minimizer $\tilde\phi$ is unconstrained. It is the region where the Bregman divergence is quadratic, and we only need the \emph{competitor} $\phi_\epsilon$ to lie in $\F$. Say $\chis$ satisfies the \emph{quadratic upper bound} with constant $C_\chi$ on $\F$ if
\[
D_{\chis}(f, \phi) \le C_\chi (f - \phi)^2 \quad \text{for all } f, \phi \in \F.
\]
Because $\tilde\phi$ minimizes the population objective, its objective value is at most that of the competitor $\phi_\epsilon \in \F$. Dropping the non-negative $\zetas_\rho(\tilde\phi)$ from the left side of this comparison and applying the quadratic upper bound to the right gives
\begin{equation}\label{eq:approx-competitor}
\P D_{\chis}(\tilde\phi, \phi) \le C_\chi \|\phi_\epsilon - \phi\|_{\LtP}^2 + \zetas_\rho(\phi_\epsilon)
\end{equation}
whenever $\phi \in \F \cap L_2(\R^d)$ and mollification preserves $\F$. Here $\zetas_\rho = \zetas \comp \rho$ is the penalty from the population objective. When the Riesz representer $\griesz = \dchis(\phi)$ is bounded between $c$ and $C$, the natural choice is $\F = \{f : c \le \dchis(f) \le C\}$.

\subsection{Positive-part quadratic}

\paragraph{Setting.} For the positive-part quadratic dispersion $\chis(\phi) = \phi_+^2/2$, the weights are $\gamma = \phi_+$ (zero when $\phi < 0$). The threshold $\tilde\phi = 0$ separates the \emph{active region} ($\tilde\phi \ge 0$, positive weights) from the \emph{dead zone} ($\tilde\phi < 0$, zero weights). The Bregman divergence vanishes in the dead zone, so the two-pile transfer (\Cref{lem:two-pile}) uses the density control function
\begin{equation}\label{eq:density-control}
\omega(\varepsilon) = \P(\tilde\phi \in [0, \varepsilon]).
\end{equation}
The dead-zone cost in~\eqref{eq:two-pile-tolerance} is $\mu = \omega(\epsilon)\, M^2 + r^2 M^2/\epsilon^2$, optimized over $\epsilon > 0$. \Cref{prop:anticonc-linear} gives the linear bound $\omega(\varepsilon) \le C\varepsilon/\sigma$ for models with bounded density; \Cref{prop:anticonc-sobolev} bounds $\omega(\varepsilon)$ by the mass of $\griesz$ near zero plus the Bregman approximation rate.

\begin{lemma}[Anticoncentration from bounded density]
\label{lem:subgaussian-anticonc}
If $X$ has a density $f_X$
with $\|f_X\|_\infty \le B$,
then $\P(X \in [a, a+\varepsilon]) \le B\varepsilon$
for all $a \in \R$ and $\varepsilon > 0$.
\end{lemma}
\begin{proof}
$\P(X \in [a, a+\varepsilon])
  = \int_a^{a+\varepsilon} f_X(x)\, dx
  \le \|f_X\|_\infty\, \varepsilon$.
\end{proof}

\begin{proposition}[Linear models]\label{prop:anticonc-linear}
Assume the linear form $\tilde\phi(Z) = \alpha + \beta^\top Z$ where the projection $\beta^\top Z$ has a density $f$ with $\|f\|_\infty \le C/\sigma$. Then $\omega(\varepsilon) \le C\varepsilon/\sigma$ for all $\varepsilon > 0$.
\end{proposition}

\begin{proof}
\Cref{lem:subgaussian-anticonc} with $B = C/\sigma$ gives $\P(\beta^\top Z \in [a, a+\varepsilon]) \le C\varepsilon/\sigma$ for any interval of width $\varepsilon$.
\end{proof}

\begin{proposition}[Density control via Bregman divergence]\label{prop:anticonc-sobolev}
For any $\varepsilon > 0$,
\begin{equation}\label{eq:density-control-bound}
\omega(\varepsilon)
  \le \P(\griesz \le 2\varepsilon)
  + \frac{2\,\P D_{\chis}(\tilde\phi, \phi)}{\varepsilon^2}.
\end{equation}
\end{proposition}

\begin{proof}
Split the event $\{\tilde\phi \in [0,\varepsilon]\}$
according to whether $\griesz > 2\varepsilon$.
On $\{\griesz \le 2\varepsilon\}$
there is nothing to prove.
On $\{\tilde\phi \in [0,\varepsilon],\, \griesz > 2\varepsilon\}$,
both $\tilde\phi \ge 0$ and $\phi = \griesz > 0$
lie in the active region,
so the Bregman divergence is quadratic:
$D_{\chis}(\tilde\phi, \phi)
  = (\tilde\phi - \griesz)^2/2
  \ge (\griesz - \varepsilon)^2/2
  > \varepsilon^2/2$.
Markov's inequality gives
\[
\P(\tilde\phi \in [0,\varepsilon],\, \griesz > 2\varepsilon)
  \le \frac{\P D_{\chis}(\tilde\phi, \phi)}{\varepsilon^2/2}
  = \frac{2\,\P D_{\chis}(\tilde\phi, \phi)}{\varepsilon^2}.
\qedhere
\]
\end{proof}


\subsection{Entropy}

The entropy dispersion $\chis(\phi) = e^\phi$ has weights $\gamma = \dchis(\phi) = e^\phi$, which are automatically positive but can grow or collapse exponentially as $\phi$ varies. When the Riesz representer has bounded weights $c \le \griesz \le C$, the Riesz representer's index $\phi = \log \griesz$ lies in $[\log c, \log C]$ and the Bregman divergence is locally quadratic in this range. But the regularized projection $\tilde\phi$ can wander outside this range, where the divergence grows exponentially rather than quadratically. The results below address three questions: how close is $\tilde\phi$ to $\phi$ (approximation), how likely are the implied weights to collapse (curvature control), and how large can $\tilde\phi$ be at fixed regularization (strict overlap).

The entropy Bregman divergence at deviation $h$ from $\phi$ is
\[
D_{\chis}(\phi + h, \phi) = e^\phi g(h) = h^2 e^\phi \int_0^1 (1-t) e^{th}\,dt \qqtext{where} g(h) = e^h - 1 - h.
\]
The integrand $(1-t)e^{th}$ lies between $(1-t)e^{\min(0,h)}$ and $(1-t)e^{|h|}$ on $[0,1]$, and $\int_0^1 (1-t)\,dt = 1/2$, giving the two-sided bound
\[
\frac{e^{-|h|}}{2} h^2 \le \frac{D_{\chis}(\phi + h, \phi)}{e^\phi} \le \frac{e^{|h|}}{2} h^2.
\]
When $c \le e^\phi \le C$ and $|h| \le \log(C/c)$, the exponential factors are controlled by $c$ and $C$:
\begin{equation}\label{eq:bregman-sandwich-app}
\frac{c^2}{2C}h^2 \le D_{\chis}(\phi + h, \phi) \le \frac{C^2}{2c}h^2.
\end{equation}

\begin{corollary}[Entropy approximation rate]\label{cor:entropy-approx}
For entropy with $c \le \griesz \le C$, penalty $\zetas_\rho(f) = \frac{\eta}{2} |f|_{\dot H^s}^2$, and $\phi = \log \griesz \in H^{s_\phi}(\R^d)$,
\begin{equation}\label{eq:approx-rate}
\P D_{\chis}(\tilde\phi, \phi) \le \qty{\frac{C^2}{2c} A_\omega^2 |\phi|_{H^{s_\phi}}^2 + \frac{C_\omega^2}{2} \|\phi\|_{L_2(\R^d)}^2} \eta^{s_\phi/(s_\phi+s)},
\end{equation}
where $A_\omega = \sup_{\xi \in \R^d \setminus \{0\}} |1 - \hat\omega(\xi)| \cdot \|\xi\|^{-s_\phi}$ is finite because $\hat\omega(0) = 1$ and $\hat\omega$ is Schwartz.
\end{corollary}

\begin{proof}
The bound~\eqref{eq:approx-competitor} applies with $\F = \{f : \log c \le f \le \log C\}$ and $C_\chi = C^2/2c$. Mollification preserves $\F$ because convex combinations preserve pointwise bounds. Seminorm control gives $\zetas_\rho(\phi_\epsilon) \le \frac{\eta}{2} C_\omega^2 \epsilon^{-2s} \|\phi\|_{L_2(\R^d)}^2$, and the mollification estimate gives $\|\phi_\epsilon - \phi\|_{L_2} \le A_\omega \epsilon^{s_\phi} |\phi|_{H^{s_\phi}}$. Substituting both into~\eqref{eq:approx-competitor} and choosing $\epsilon = \eta^{\frac{1}{2(s_\phi + s)}}$ balances the two terms.
\end{proof}

The sandwich bound shows that $D_{\chis}$ is quadratic when both points lie in $[\log c, \log C]$. But the regularized projection $\tilde\phi$ can wander outside this range: for large deviations $|h| > \log(C/c)$, the exponential tail of $g(h)$ kicks in. The next proposition controls the probability of weight collapse by splitting into these two regimes.

\begin{proposition}\label{prop:curvature-control}
Assume $c \le \griesz \le C$ with $\log(C/c) \ge 2$, and $\delta \ge c^2/C$. Then
\[
\P(\tilde\gamma < \delta) \le \frac{C_0 \, \P D_{\chis}(\tilde\phi, \phi)}{\log(c/\delta)^2}.
\]
The constant $C_0 = 2C/c^2 + \log(C/c)^2/\{c(\log(C/c) - 1)\}$ depends only on $c$ and $C$. The restriction $\delta \ge c^2/C$ ensures $M = \log(c/\delta) \le \log(C/c)$, so the deviation falls in the quadratic regime of the Bregman divergence.
\end{proposition}

\begin{proof}
Write $h = \tilde\phi - \phi$ for the deviation and $G = \{|h| \le \log(C/c)\}$ for the good set.

On $G$, $|h| \le \log(C/c)$, so the general two-sided bound gives $D_{\chis}(\tilde\phi, \phi)/e^\phi \ge e^{-|h|} h^2/2 \ge (c/C) h^2/2$. The base-point factor $e^\phi \ge c$ (a bound on the population index, not on $\tilde\phi$) then gives
\[
D_{\chis}(\tilde\phi, \phi) \ge \frac{c^2}{2C} h^2, \quad\text{so}\quad
\P h^2 \mathbf{1}_G \le \frac{2C}{c^2} \P D_{\chis}(\tilde\phi, \phi).
\]

On $G^c$, $|h| > \log(C/c) \ge 2$. For $h \ge 2$, $e^h/2 \ge 1 + h$ (because $e^h - 2 - 2h$ is positive at $h = 2$ with derivative $e^h - 2 > 0$), so $g(h) = e^h - 1 - h \ge e^h/2 \ge h$; for $h \le -2$, $g(h) \ge |h| - 1$. Either way $g(h) \ge |h| - 1$, so with the base-point bound $e^\phi \ge c$,
\[
D_{\chis}(\tilde\phi, \phi) = e^\phi g(h) \ge c(|h| - 1) \ge c(\log(C/c) - 1), \quad\text{giving}\quad
\P(G^c) \le \frac{\P D_{\chis}(\tilde\phi, \phi)}{c(\log(C/c) - 1)}.
\]

For the event $\tilde\gamma = e^{\tilde\phi} < \delta$: because $\phi = \log \griesz \ge \log c$, the deviation satisfies $|h| > M$ for $M = \log(c/\delta)$. If $M \le \log(C/c)$, Markov's inequality on $G$ and the bound $1 \le \log(C/c)^2/M^2$ on $G^c$ give
\[
\begin{aligned}
\P(|h| > M, \, G) &\le \frac{\P h^2 \mathbf{1}_G}{M^2} \le \frac{2C}{c^2 M^2} \P D_{\chis}(\tilde\phi, \phi), \\
\P(G^c) &\le \frac{\P D_{\chis}(\tilde\phi, \phi)}{c(\log(C/c) - 1)} \le \frac{\log(C/c)^2}{c(\log(C/c) - 1)} \cdot \frac{\P D_{\chis}(\tilde\phi, \phi)}{M^2}.
\end{aligned}
\]
Adding both pieces gives $C_0 = 2C/c^2 + \log(C/c)^2/\{c(\log(C/c) - 1)\}$.
\end{proof}

\begin{proposition}\label{prop:strict-overlap}
Assume $c \le \griesz \le C$, $s > d/2$ (so that $\|f - \P f\|_\infty \le C_{\mathrm{Sob}}\, \rho(f)$). Then the upper bound $\tilde\phi \le \overline{\tilde\phi}$ holds for $\overline{\tilde\phi} = \log C + C_{\mathrm{Sob}}\sqrt{B_0/\eta}$ and $B_0 = \P D_{\chis}(0, \phi)$.
\end{proposition}

\begin{proof}
Comparing to the competitor $f = 0$,
$\P D_{\chis}(\tilde\phi, \phi) + \eta \rho(\tilde\phi)^2
  \le \P D_{\chis}(0, \phi) + \eta\rho(0)^2 = B_0$.
Dropping the nonnegative $\P D_{\chis}(\tilde\phi, \phi)$ gives
\[
\eta \rho(\tilde\phi)^2 \le B_0,
\]
so $\rho(\tilde\phi) \le \sqrt{B_0/\eta}$. Because $\rho$ is a seminorm that kills constants, constant shifts $f \mapsto f + t$ leave the penalty invariant: $\rho(f+t)^2 = \rho(f)^2$. The Bregman term shifts by $\P(e^{f+t} - e^f) - t \P \griesz$, whose derivative at $t = 0$ is $\P e^f - \P \griesz$. Setting this to zero at the minimizer gives the first order condition
\[
\P e^{\tilde\phi} = \P \griesz.
\]
Jensen's inequality then gives $e^{\P \tilde\phi} \le \P e^{\tilde\phi} = \P \griesz \le C$, so $\P \tilde\phi \le \log C$. Sobolev embedding gives $\|f - \P f\|_\infty \le C_{\mathrm{Sob}} \rho(f)$, so
\[
\tilde\phi \le \P\tilde\phi + \|\tilde\phi - \P\tilde\phi\|_\infty \le \log C + C_{\mathrm{Sob}} \rho(\tilde\phi) \le \log C + C_{\mathrm{Sob}} \sqrt{B_0/\eta}.
\]
\end{proof}

A lower bound follows from the same normalization.
From $\P e^{\tilde\phi} = \P\griesz \ge c > 0$
and $\tilde\phi \le \P\tilde\phi + C_{\mathrm{Sob}}\,\rho(\tilde\phi)$,
we get
$c \le \P e^{\tilde\phi}
  \le e^{\P\tilde\phi + C_{\mathrm{Sob}}\,\rho(\tilde\phi)}$,
so $\P\tilde\phi
  \ge \log c - C_{\mathrm{Sob}}\,\rho(\tilde\phi)$.
By the same Sobolev embedding,
$\tilde\phi
  \ge \P\tilde\phi - C_{\mathrm{Sob}}\,\rho(\tilde\phi)
  \ge \log c - 2\,C_{\mathrm{Sob}}\,\sqrt{B_0/\eta}$.
Because $\overline{\tilde\phi} \to \infty$ as $\eta \to 0$, strict overlap fails in the undersmoothed regime. \Cref{prop:curvature-control} provides the asymptotic substitute.


\section{Uniform Lower Bounds for Bregman Divergences}
\label{sec:uniform-lower-bounds}

\subsection{Uniform Lower Bound on a Single Sphere}

For uniform lower bounds on sample-averaged Bregman divergence,
we will use a partial\footnotemark generalization of the \emph{stable lower bounds} approach from \citet{mendelson2021extending}. The heart
of the approach is the following notion of stability.

\footnotetext{\citet{mendelson2021extending} has a layer of complexity that is important for tournament-type estimators, but unnecessary for what we discuss here. It shows that if we subdivide our sample into blocks of approximately equal size, then with high probability, a uniform lower bound is satisfied on almost every block. Here we consider only the 1-block case, but we generalize that special case to apply to
Bregman divergences generally rather than just squared error.}

\begin{definition}[Stable lower bound for nonnegative functions]
\label{def:slb-general}
A nonnegative function $\psi : \mathcal{Z} \to \R_+$
satisfies a \emph{$(\xi, \ell, k)$ stable lower bound on an event~$E$}
if there exists an event $\Omega_0$ with
$\P(\Omega_0) \ge 1 - 2e^{-k}$
such that on $E \cap \Omega_0$,
\[
\frac{1}{n}\sum_{i \not\in J} \psi(Z_i) \ge (1-\xi) \P\psi(Z)
 \ \text{ for all subsets $J \subset \{1,\ldots,n\}$
   of cardinality $|J| \le \ell$.}
\]
When $E$ is the full sample space,
we say $\psi$ satisfies
a $(\xi, \ell, k)$ stable lower bound.
\end{definition}

The following lemma gives a uniform lower bound on a single Bregman sphere, generalizing \citet[Theorem 3.1]{mendelson2021extending}.
The proof strategy is to establish a lower bound
on a finite $\epsilon$-net, then extend to all of $H_r^\circ$.
This requires four ingredients:
stability of the Bregman divergence
under deletion of a few coordinates~\ref{cond:shell-i},
control of the net size~\ref{cond:shell-ii},
control of the approximation error
between $h$ and its nearest net point~\ref{cond:shell-iii},
and Lipschitz control
of the derivative map~\ref{cond:shell-iv}.
\begin{lemma}[Bregman ULB on a single sphere]
\label{lem:bregman-shell}
Let $Z_1 \ldots Z_n$ be independent random variables
with empirical distribution $\hP$ and distribution $\P$.
Let $\chisZ$ be convex for each $Z$.
Let $H \subset \LtP$ be star-shaped around~$0$,
and let $\Q \ll \P$.
Choose any $r > 0$ and define
$H_r^\circ = \{h \in H : \P D_{\chisZ}(f+h, f) = r^2\}$
and $r_2^2 = \sup_{h \in H_r^\circ} \|h\|_{L_2(\Q)}^2$.
Let $\sigma_1 \ldots \sigma_n$ be independent Rademacher variables.

\smallskip

Suppose that the following conditions hold
for $c_0 = 1/8$, $c_1 = 1/64$, and any $\xi \in (0,1)$.
\begin{enumerate}[label=\textup{(\roman*)}]
\item\label{cond:shell-i} \textbf{(Bregman SLB)}
  For every $h \in H_r^\circ$,
  the function $Z \mapsto D_{\chisZ}(f+h, f)(Z)$
  satisfies a $(\xi/2, \ell, k)$ stable lower bound
  on an event~$E$.
\item\label{cond:shell-ii} \textbf{(Metric entropy)}
  $\log N\{H_r^\circ, \epsilon, L_2(\Q)\} \le k/16$
  for $\epsilon = c_0 \xi r$.
\item\label{cond:shell-iii} \textbf{(Local Rademacher)}
  For $\Delta_\epsilon = \{h - h' : h, h' \in H_r^\circ,\;
    \|h - h'\|_{L_2(\Q)} \le \epsilon\}$,
  \[
  \E \sup_{u \in \Delta_\epsilon}
    \abs*{\hP \sigma u(Z)}
    \le c_1 \xi r \sqrt{\ell/n}.
  \]
\item\label{cond:shell-iv} \textbf{(Derivative equicontinuity)}
  For some constant $L > 0$
  and all $h, h' \in H_r^\circ \cup \{0\}$,
  \[
  \P \{ \dchisZ \comp (f+h)
    - \dchisZ \comp (f+h')\}^2
    \le L^2 \norm{h-h'}_{L_2(\Q)}^2.
  \]
\end{enumerate}

\smallskip

Then for any $t > 0$, writing
$\xi_{t,k} = (1+t)\xi + (2 - \xi) e^{-k}$,
\[
  \inf_{h \in H_r^\circ} \hP D_{\chisZ}(f+h, f)
    \ge (1 - \xi_{t,k}) r^2
\]
on the intersection of $E$
and an event of probability at least
$1 - 2e^{-15k/16} - e^{-\ell/16}
  - \frac{20 L^2 r_2^2}{\ell t^2 r^2}$.

\medskip

Conditions \ref{cond:shell-ii} and \ref{cond:shell-iii}
can be replaced by
\begin{itemize}
\item[\textup{(ii$'$)}]\customlabel{cond:shell-ii-prime}{\textup{ii$'$}}
  \textbf{(Global Rademacher)}
  $\E\sup_{h \in H_r^\circ} \abs*{\hP \sigma h(Z)}
    \le c_1 \xi r \sqrt{\ell/n} / 2.$
\item[\textup{(iii$'$)}]\customlabel{cond:shell-iii-prime}{\textup{iii$'$}}
  \textbf{($L_2$ SLB)}
  For every pair $h_1, h_2 \in H_r^\circ$
  with $\|h_1 - h_2\|_{L_2(\Q)} \ge \epsilon_0$,
  $(h_1 - h_2)^2$ satisfies
  a $(\xi/2, \ell, k)$ stable lower bound on~$E$.
\end{itemize}
The natural choice is $d\Q = \ddchis(\tilde\phi)\,d\P$,
under which $L_2(\Q)$ is the curvature-weighted norm
and $r_2 \approx r$.
For the quadratic dispersion, $\ddchis = 1$ and $\Q = \P$.
\end{lemma}

\subsubsection{Alternative Conditions \ref{cond:shell-ii-prime} and \ref{cond:shell-iii-prime} Imply \ref{cond:shell-ii} and \ref{cond:shell-iii}}

\textit{\ref{cond:shell-ii-prime} $\Rightarrow$ \ref{cond:shell-iii}:} Condition~\ref{cond:shell-iii} bounds the Rademacher complexity of the difference class $\Delta_\epsilon = \{h - h' : h, h' \in H_r^\circ,\; \|h-h'\|_{L_2(\Q)} \le \epsilon\}$. Because $\Delta_\epsilon \subset H_r^\circ - H_r^\circ$, the triangle inequality for the Rademacher process gives
\[
\E\sup_{u \in \Delta_\epsilon} \abs*{\hP \sigma u(Z)} \le 2\E\sup_{h \in H_r^\circ} \abs*{\hP \sigma h(Z)} \le 2 \cdot \frac{c_1}{2} \xi r \sqrt{\ell/n} = c_1 \xi r \sqrt{\ell/n},
\]
which is condition~\ref{cond:shell-iii}. The constraint $\|h-h'\| \le \epsilon$ is not used---the global bound already suffices.

\textit{\ref{cond:shell-ii-prime}+\ref{cond:shell-iii-prime} $\Rightarrow$ \ref{cond:shell-ii}:} This is Sudakov minoration \citep[Theorem 4.4]{mendelson2021extending}. The idea: if $H_r^\circ$ had a large $\epsilon$-packing in $L_2(\Q)$, the Rademacher complexity of $H_r^\circ$ would be large---contradicting~\ref{cond:shell-ii-prime}.

Let $M(\epsilon_0)$ be the $\epsilon_0$-packing number of $H_r^\circ$ in $L_2(\Q)$: the largest set of points with pairwise $L_2(\Q)$-distance at least $\epsilon_0$. Take $\epsilon_0 = c_0 \xi r$. Condition~\ref{cond:shell-iii-prime} ensures that, for any two packing points $h_1, h_2$, the pairwise $L_2(\Q)$-separation is preserved in the sample average: after dropping at most $\ell$ coordinates,
\[
\frac{1}{n}\sum_{i \notin J} (h_1 - h_2)^2(Z_i) \ge (1-\xi/2)\|h_1 - h_2\|_{L_2(\Q)}^2 \ge (1-\xi/2)\epsilon_0^2.
\]
We apply \citet[Theorem~4.4]{mendelson2021extending}
contrapositively,
with one block ($n_{\mathrm{block}} = 1$, $m = n$),
separation $\rho = \epsilon_0 = c_0\xi r$,
and the stable lower bound from~(iii$'$).
If $\log M(\epsilon_0) \ge k/16$,
the theorem gives the Rademacher lower bound
\[
\E_\sigma \sup_{h \in H_r^\circ}
  \abs*{\hP \sigma h(Z)}
\ge c_2\, c_0\, \xi r\, \min\!\qty{\sqrt{\ell/n},\, \sqrt{k/n}},
\]
where $c_2$ is an absolute constant
from \citet[Theorem~4.4]{mendelson2021extending}.
When $k \le \ell$,
comparing with the upper bound
$(c_1/2)\,\xi r\,\sqrt{\ell/n}$ from~\ref{cond:shell-ii-prime}
gives $k \le c_1^2 \ell / (4 c_2^2 c_0^2)$.
So $\log M(\epsilon_0) \le k/16$
whenever $k \ge c_1^2 \ell / (4 c_2^2 c_0^2)$.
Because covering number $\le$ packing number
at the same scale,
this gives condition~\ref{cond:shell-ii}.

\subsection{Proof of Lemma~\ref{lem:bregman-shell}}
\label{sec:proof-main}

\paragraph{Outline.}
The proof proceeds in three steps.
\begin{enumerate}
\item Establish the lower bound on a finite $\epsilon$-net via the stable lower bound property.
\item Control the number of ``bad'' coordinates where the approximation error $|h - \pi h|$ exceeds a threshold. Here $\pi$ maps $h$ to its nearest net point.
\item Extend this lower bound from the net to all of $H_r^\circ$.
\begin{enumerate}
  \item We drop bad coordinates. Because we're working with an average of (positive) Bregman divergences, this gives us a lower bound.
  \item On the remaining coordinates, we decompose our Bregman divergence as the sum of a net term $D_{\chisZ}(f+\pi h, f)$, which we bounded in step 1, plus a remainder $\Gamma$.
  \item  We bound the mean and the fluctuations of the remainder $\Gamma$.
\end{enumerate}
\end{enumerate}

\subsubsection{Step 1: Net Argument}

By Condition~\ref{cond:shell-ii}, there exists an $\epsilon$-net $H' \subset H_r^\circ$ with $\epsilon = c_0 \xi r$ satisfying $\log |H'| \le k/16$. For each net point $h' \in H'$, Condition~\ref{cond:shell-i} provides a $(\xi/2, \ell, k)$ stable lower bound. Specifically, there exists a set $J(h') \subset \{1 \ldots n\}$ with $|J(h')| \le \ell$ such that on the intersection of $E$ and an event of probability $1-2\exp(-k)$,
\[
\frac{1}{n} \sum_{i \notin J(h')} D_{\chisZ}(f+h', f)(Z_i) \ge (1-\xi/2) r^2
\]
Applying the union bound over $|H'| \le \exp(k/16)$ net points, the total failure probability is
\[
\exp(k/16) \cdot 2\exp(-k) = 2\exp(-15k/16).
\]
Thus with probability at least $1 - 2\exp(-15k/16)$, the stable lower bound holds simultaneously for all $h' \in H'$.

\subsubsection{Step 2: Truncation}

For each $h \in H_r^\circ$, let $\pi h \in H'$ denote the nearest net point, so $\|h - \pi h\|_{L_2(\Q)} \le \epsilon$. Define the threshold
\[
A = \sqrt{n/\ell} \cdot \xi r
\]
and the set of bad coordinates
\[
J(h) = \{i \in \{1 \ldots n\} : |(h - \pi h)(Z_i)| > A\}.
\]
We claim $|J(h)| \le \ell$ uniformly over $h \in H_r^\circ$. The proof is by contradiction.

\paragraph{Lower bound from large $J(h)$.} Suppose $|J(h)| > \ell$ for some $h$. Consider the truncation $\varphi_A^2(t) = \min(t^2, A^2)$. On each coordinate $i \in J(h)$, we have $|(h - \pi h)(Z_i)| > A$, so $\varphi_A^2\{(h-\pi h)(Z_i)\} = A^2$, so
\begin{equation}\label{eq:truncation-lower}
\hP \varphi_A^2(h - \pi h) \ge \frac{|J(h)|}{n} \cdot A^2 > \frac{\ell}{n} \cdot A^2 = \frac{\ell}{n} \cdot \frac{n}{\ell} \xi^2 r^2 = \xi^2 r^2.
\end{equation}

\paragraph{Upper bound from Rademacher complexity.} We show $\sup_{h \in H_r^\circ} \hP \varphi_A^2(h - \pi h) < \xi^2 r^2$ with high probability. Because $\varphi_A^2(0) = 0$ and $\varphi_A^2$ has Lipschitz constant $2A$, the contraction principle (\Cref{lem:contraction}) gives
\[
\E \sup_{h \in H_r^\circ} \hP \sigma_i \varphi_A^2\{(h - \pi h)(Z_i)\} \le 2A \cdot \E \sup_{u \in \Delta_\epsilon} |\hP \sigma_i u(Z_i)| \le 2A \cdot c_1 \xi r \sqrt{\ell/n} = 2 c_1 \xi^2 r^2.
\]
By symmetrization (\Cref{lem:symmetrization}), $\E \sup_{h \in H_r^\circ} (\hP - \P)\varphi_A^2(h - \pi h) \le 4 c_1 \xi^2 r^2$. Because $\P\varphi_A^2(h - \pi h) \le \P(h - \pi h)^2 \le \epsilon^2 = c_0^2 \xi^2 r^2$, we have
\[
\E\sup_{h \in H_r^\circ} \hP \varphi_A^2(h - \pi h) \le c_0^2 \xi^2 r^2 + 4c_1 \xi^2 r^2 = \{(1/8)^2 + 4(1/64)\}\xi^2r^2 = (5/64)\xi^2 r^2 \qfor c_0 = 1/8, c_1 = 1/64.
\]
We apply \Cref{lem:talagrand} to $W = \sup_{h \in H_r^\circ} (\hP - \P)\varphi_A^2(h - \pi h)$ with bound $b = A^2 = (n/\ell)\xi^2 r^2$. For the variance term, because $\varphi_A^4(t) \le A^2 \varphi_A^2(t)$,
\[
V = \sup_{h \in H_r^\circ} \frac{1}{n}\sum_i \Var \varphi_A^2\{(h-\pi h)(Z_i)\} \le A^2 \cdot \P\varphi_A^2(h-\pi h) \le A^2 \epsilon^2 = c_0^2 (n/\ell) \xi^4 r^4.
\]
Setting $t = \xi^2 r^2/4$ and using $\E W \le 4c_1 \xi^2 r^2$, the parenthesized factor in our probability's exponent's denominator satisfies
\[
V + 2b\E W + bt \le (n/\ell)\xi^4 r^4 \qty(c_0^2 + 8c_1 + 1/4) = (n/\ell)\xi^4 r^4 \cdot \tfrac{25}{64} \qqtext{recalling} c_0 = 1/8, c_1 = 1/64.
\]
 The exponent is at least $nt^2 / [2 \cdot (n/\ell)\xi^4 r^4 \cdot 25/64] = \ell \cdot 64/(16 \cdot 50) > \ell/16$. Thus, on an event of probability at least $1 - \exp(-\ell/16)$:
\begin{equation}\label{eq:truncation-upper}
\sup_{h \in H_r^\circ} \hP \varphi_A^2(h - \pi h) \le \tfrac{5}{64}\xi^2 r^2 + \tfrac{1}{4}\xi^2 r^2 < \xi^2 r^2/2.
\end{equation}

\paragraph{Contradiction.} On this event, the lower bound \eqref{eq:truncation-lower} exceeds the upper bound \eqref{eq:truncation-upper}, so our premise that $|J(h)| > \ell$
must be incorrect. On this event, $|J(h)| \le \ell$ for all $h \in H_r^\circ$.

\subsubsection{Step 3: Extension from net to all $h \in H_r^\circ$}

\paragraph{Setup and decomposition.}

We define an indicator for coordinates on which $\pi h$ is a good approximation to $h$. For every $h \in H_r^\circ$, as shown in the previous step, this is true of all but $\ell$ coordinates.
\[
1_A(h) = 1_{|h - \pi h| \le A} \qand 1_{\neg A}(h)=1-1_{A}(h) \qqtext{satisfy} \sum_i 1_{\neg A}(h)(Z_i) \le \ell \quad \text{for all } h \in H_r^\circ.
\]
Because Bregman divergences are nonnegative, we can drop the bad coordinates from our lower bound.
\begin{equation}\label{eq:drop-J}
\hP D_{\chisZ}(f+h, f) \ge \hP D_{\chisZ}(f+h, f) \cdot 1_A(h).
\end{equation}
This is the key step: we sacrifice the bad coordinates because we cannot control their deviation from the net well and---because our lower bound is stable---it doesn't cost us much. We then decompose our Bregman divergence into a net term and a remainder.
\begin{equation}\label{eq:decomp-Jc}
\begin{aligned}
&\hP D_{\chisZ}(f+h, f) \cdot 1_A(h) = \hP D_{\chisZ}(f+\pi h, f) \cdot 1_A(h) + \hP \Gamma_h \cdot 1_A(h) \\
&\qfor \Gamma_h = D_{\chisZ}(f+h, f) - D_{\chisZ}(f+\pi h, f) \\
&\hphantom{\qfor \hP \Gamma_h} = D_{\chisZ}(f+h, f+\pi h) + (\dchisZ \comp (f+\pi h) - \dchisZ \comp f)(h - \pi h).
\end{aligned}
\end{equation}
The latter expression for $\Gamma_h$ follows from the Bregman law of cosines,
$D_{\chi}(a,c) - D_\chi(b,c) = D_\chi(a,b) + \{\chi'(b) - \chi'(c)\}(a-b)$
for any convex function $\chi$ and points $a\ldots c$, taking $a = f+h$, $b = f+\pi h$, $c = f$.

The number of bad points we've dropped is less than $\ell$, so the stable lower bound from Step~1 applies to our net term:
\begin{equation}\label{eq:net-Jc}
\hP D_{\chisZ}(f+\pi h, f) \cdot 1_A(h) \ge (1 - \xi/2) r^2.
\end{equation}

We will decompose our remainder into a mean and fluctuation, i.e. as
\[
\hP \Gamma_h \cdot 1_A(h) = \P \Gamma_h \cdot 1_A(h) + (\hP-\P) \Gamma_h \cdot 1_A(h),
\]
and lower bound them in the next two sections.

\paragraph{Remainder: mean.}

The key idea:
because $\P \Gamma_h = 0$
(both $h$ and $\pi h$ lie on the same Bregman sphere),
the good-coordinate mean
equals minus the bad-coordinate mean.
We bound the bad-coordinate mean
using the population Bregman divergence
and the SLB failure probability.

Because $h, \pi h \in H_r^\circ$, $\P D_{\chisZ}(f+h,f) = \P D_{\chisZ}(f+\pi h,f) = r^2$ and therefore $\P \Gamma_h = 0$. It follows that the mean on our good coordinates must be equal and opposite the mean on our bad coordinates:
\[
-\P \Gamma_h \cdot 1_A(h) = \P \Gamma_h \cdot 1_{\neg A}(h).
\]
Thus, to lower bound our good-coordinate mean, we upper bound the bad-coordinate one. And recalling our remainder's definition \eqref{eq:decomp-Jc},
$\Gamma_h \le D_{\chisZ}(f+h, f)$ because $D_{\chisZ}(f+\pi h, f) \ge 0$ and therefore
\[
\P \Gamma_h \cdot 1_{\neg A}(h) \le \P D_{\chisZ}(f+h, f) \cdot 1_{\neg A}(h).
\]
We bound the right side using our stable lower bound for $h \in H_r^\circ$, which implies $\hP D_{\chisZ} (f+h,f) \cdot 1_{A}(h) \ge (1 - \xi/2) r^2$ with probability at least $1 - 2e^{-k}$ and therefore, because it is $\ge 0$ with the complementary probability, that
\[
\P D_{\chisZ}(f+h, f) \cdot 1_{A}(h) \ge (1 - 2e^{-k})(1 - \xi/2) r^2.
\]
It follows, because $\P D_{\chisZ}(f+h, f) \cdot 1_{\neg A}(h) = r^2 - \P D_{\chisZ}(f+h, f) \cdot 1_A(h)$, that
\begin{align*}
\P D_{\chisZ}(f+h, f) \cdot 1_{\neg A}(h) &\le r^2 - (1 - 2e^{-k})(1 - \xi/2) r^2 \\
&= \qty{\frac{\xi}{2} + (2 - \xi) e^{-k}} r^2.
\end{align*}
It follows, given our equal-and-opposite property, that
\begin{equation}\label{eq:mean-bound}
-\P \Gamma_h \cdot 1_A(h) \le \qty{\frac{\xi}{2} + (2 - \xi) e^{-k}} r^2.
\end{equation}

\paragraph{Remainder: fluctuation.}

We bound the variance of our fluctuation term using \Cref{lem:efron-stein}.
\[
\Var \sup_{h \in H_r^\circ} (\hP - \P) \Gamma_h \cdot 1_A(h) \le \frac{2}{n} \sup_{h \in H_r^\circ} \P \{\Gamma_h \cdot 1_A(h)\}^2.
\]
And we can control $\Gamma_h 1_A(h)$  pointwise in terms of  $|h - \pi h|$. Referring to the latter expression for $\Gamma_h$ in \eqref{eq:decomp-Jc}, this is obvious for its second term. For the first,
this follows from the bound $D_{\chis}(a,b) \le D_{\chis}(a,b) + D_{\chis}(b,a) = \{\dchis(a) - \dchis(b)\}(a-b)$. Taking $a = f+h$, $b = f+\pi h$,
\[
D_{\chisZ}(f+h, f+\pi h) \le \{ \dchisZ \comp (f+h) - \dchisZ \comp (f+\pi h)\}(h - \pi h).
\]
Bounding the absolute value of $\Gamma_h$ using this and the triangle inequality, we have
\[
|\Gamma_h| 1_A(h) \le A \qty{\qty|\dchisZ \comp (f+h) - \dchisZ \comp (f+\pi h)| + \qty|\dchisZ \comp (f+\pi h) - \dchisZ \comp f|}.
\]
and, using the elementary bound $(a+b)^2 \le 2a^2 + 2b^2$ and our derivative equicontinuity condition~\ref{cond:shell-iv} with measure $\Q$,
\begin{align*}
\P \{\Gamma_h \cdot 1_A(h)\}^2 &\le 2A^2 \qty{\P\{\dchisZ \comp (f+h) - \dchisZ \comp (f+\pi h)\}^2 + \P\{\dchisZ \comp (f+\pi h) - \dchisZ \comp f\}^2} \\
&\le 2 A^2 L^2 (\epsilon^2 + r_2^2) \le 2A^2 L^2(c_0^2\xi^2 + 1) r_2^2
\qqtext{ where }  A^2 = \xi^2 r^2 n/\ell.
\end{align*}
The first term uses $\|h - \pi h\|_{L_2(\Q)} \le \epsilon$ (the net approximation); the second uses $\|\pi h\|_{L_2(\Q)} \le r_2$. Because $\epsilon = c_0\xi r \le r \le r_2$, the $\epsilon^2$ term is at most $r_2^2$.
For the lower bound we need only control the one-sided fluctuation $G = \sup_{h \in H_r^\circ}\{-(\hP - \P)\Gamma_h \cdot 1_A(h)\}$. We apply Chebyshev with tolerance $t\xi r^2/2$: the Efron-Stein variance bound gives
\begin{equation}\label{eq:fluctuation-cheb}
\begin{aligned}
  \P\qty(G - \E G > t\xi r^2/2)
  &\le  \frac{\frac{2}{n} \times  2A^2 L^2(c_0^2\xi^2 + 1)r_2^2}{(t/2)^2\xi^2 r^4} \\
  &= \frac{16(1+c_0^2\xi^2) L^2 r_2^2}{\ell t^2 r^2} \\
  &\le \frac{20 L^2 r_2^2}{\ell t^2 r^2} \qqtext{recalling that $c_0=1/8$ and $\xi \le 1$}.
\end{aligned}
\end{equation}
By symmetrization (\Cref{lem:symmetrization}),
$\E G
  \le 2\,\E_\sigma \sup_{h \in H_r^\circ}
    |\hP \sigma\, \Gamma_h \cdot 1_A(h)|$.
The contraction principle (\Cref{lem:contraction})
with the $L$-Lipschitz derivative map
reduces this to the Rademacher complexity
of $H_r^\circ$,
which is at most $c_1\,\xi r\,\sqrt{\ell/n}$
by condition~\ref{cond:shell-iii}.
With $A^2 = \xi^2 r^2 n/\ell$,
$\E G \le t\xi r^2/2$.
So $G > t\xi r^2$ implies $G - \E G > t\xi r^2/2$,
and the probability of that event is bounded
by the right side of~\eqref{eq:fluctuation-cheb}.

\subsubsection{Conclusion}

On an event on which our bounds \eqref{eq:drop-J}--\eqref{eq:fluctuation-cheb} hold, all $h \in H_r^\circ$ satisfy
\begin{equation}
\label{eq:extension-bound}
\begin{aligned}
\hP D_{\chisZ}(f+h, f) &\ge \hP D_{\chisZ}(f+\pi h, f) \cdot 1_A(h) + \P \Gamma_h \cdot 1_A(h) + (\hP-\P) \Gamma_h \cdot 1_A(h) \notag \\
&\ge (1 - \xi/2) r^2 - \qty{\frac{\xi}{2} + (2 - \xi) e^{-k}} r^2 - t\xi r^2 \notag \\
&= \qty{1 - (1+t)\xi - (2 - \xi) e^{-k}} r^2.
\end{aligned}
\end{equation}

\paragraph{Probability accounting.} Our failure events have the following probabilities.
\begin{itemize}
\item Net (Step 1): $2\exp(-15k/16)$
\item Truncation (Step 2): $\exp(-\ell/16)$
\item Fluctuation (Step 3): $\dfrac{20 L^2 r_2^2}{\ell t^2 r^2}$
\end{itemize}
By union bound, the total failure probability is at most
\[
2e^{-15k/16} + e^{-\ell/16} + \frac{20 L^2 r_2^2}{\ell t^2 r^2}.
\]
This completes the proof of \Cref{lem:bregman-shell} under Assumptions (i)--(iv).

\subsection{Sufficient Conditions for Stable Lower Bounds}

\citet{mendelson2021extending} gives a number of sufficient conditions for stable lower bounds. In our applications, the increments $h$ are always bounded, so the following result is the workhorse.

\begin{lemma}[Bounded functions]
\label{lem:slb-bounded}
If $|h| \le M$ almost surely, then $h^2$ satisfies a $(\xi, \ell, k)$ stable lower bound with
\[
\ell = c_0 n \xi \frac{\|h\|_{L_2}^2}{M^2} \quad \text{and} \quad k = c_1 n \xi^2 \frac{\|h\|_{L_2}^2}{M^2}.
\]
\end{lemma}

\subsection{Stable Lower Bounds for Bregman Divergences}

The abstract \Cref{lem:bregman-shell} requires
a stable lower bound for each $D_{\chisZ}(\tilde\phi+h,\tilde\phi)$
on the sphere $H_r^\circ$.
This section derives those SLBs
for three dispersion families:
quadratic, positive-part quadratic, and entropy-like.
Each case uses bounded increments ($|h| \le M$)
and \Cref{lem:slb-bounded};
proofs are collected in \Cref{sec:ulb-application-proofs}.

\subsubsection{Quadratic family}

\setupvariants

For the quadratic dispersion $\chis(\phi) = \phi^2/2$,
curvature is constant ($\ddchis = 1$),
so $D_{\chis}(\tilde\phi+h,\tilde\phi) = h^2/2$
and derivative equicontinuity holds with $L = 1$
(the derivative $\dchis(\phi) = \phi$ is linear,
so $\dchis(\tilde\phi + h) - \dchis(\tilde\phi + h') = h - h'$).
The natural measure is $\Q = \P$,
and $r_2^2 = \sup_{h \in H_r^\circ} \|h\|_{L_2}^2 = 2r^2$.

\theoremvariant{Q}
\begin{corollary}[Setting Q: quadratic sphere conditions]
\label{cor:ulb-quadratic}
Let $H_r^\circ = \{h : \P D_{\chis}(\tilde\phi+h,\tilde\phi) = r^2\}$
and suppose $|h| \le M$ for all $h \in H_r^\circ$.
If $\chis(\phi) = \phi^2/2$,
then conditions~\ref{cond:shell-i}, \ref{cond:shell-iii-prime},
and~\ref{cond:shell-iv}
of \Cref{lem:bregman-shell}
hold with $\Q = \P$, $L = 1$, $r_2^2 = 2r^2$,
and SLB parameters
$k = 2c_1\, n\, \xi^2\, r^2/M^2$
and $\ell = 2c_0\, n\, \xi\, r^2/M^2$.
\end{corollary}
\resettheorem

\subsubsection{Positive-part quadratic}

On coordinates where $\tilde\phi \ge 0$,
the positive-part divergence
$D_{\chis}(\tilde\phi+h,\tilde\phi)$
reduces to $h^2/2$---the same as the quadratic case.
The dead zone $\{\tilde\phi < 0\}$ creates coordinates
where $\dchis(\tilde\phi) = 0$
and the Bregman divergence does not see
increments that stay negative.
The new ingredient is the \emph{junction error}
on coordinates where $h$ pushes across
the boundary $\tilde\phi = 0$.
The two-pile decomposition (\Cref{lem:two-pile})
reduces the Bregman SLB
to the quadratic case
at a cost proportional
to the empirical junction error $\hP f_h$.
A separate concentration step
(\Cref{lem:dead-zone-concentration})
bounds $\hP f_h$ uniformly over $H_r$,
producing an event $E$
whose complement has exponentially small probability.

\begin{lemma}[Two-pile SLB transfer]
\label{lem:two-pile}
Let $\chis(x) = x_+^2/2$.
Let $\tilde\phi, h : \mathcal{Z} \to \R$ be measurable
with $|\tilde h| \le M$ almost surely,
where $\tilde h = \ind{\tilde\phi \ge 0} \cdot h$.
Define the junction error
\begin{equation}\label{eq:junction-error}
f_h(Z) = h(Z)^2 \cdot \ind{\tilde\phi(Z) \ge 0,\;\tilde\phi(Z) + h(Z) < 0}.
\end{equation}
Then $D_{\chis}(\tilde\phi+h,\tilde\phi)$
satisfies a $(\xi, \ell, k)$ stable lower bound with
\[
\xi = \xi_0 + \frac{\hP f_h}{r^2},
\]
where $r = \|\tilde h\|_{\LtP}$,
$\xi_0 \in (0,1)$ is a free parameter,
$\ell = c_0\, n\, \xi_0\, r^2/M^2$,
and $k = c_1\, n\, \xi_0^2\, r^2/M^2$.
\end{lemma}

\paragraph{Population bound on the junction error.}
On Pile~2, $|h| > \tilde\phi \ge 0$,
so for any $\epsilon > 0$,
either $\tilde\phi \in [0,\epsilon]$
(the threshold is near the boundary)
or $|h| > \epsilon$
(the perturbation is large).
Splitting,
\begin{align*}
\P f_h
  &= \P\!\qty[h^2 \cdot \ind{\tilde\phi \ge 0,\, \tilde\phi + h < 0}] \\
  &\le \P\!\qty[h^2 \cdot \ind{\tilde\phi \in [0,\epsilon]}]
    + \P\!\qty[h^2 \cdot \ind{|h| > \epsilon}].
\end{align*}
The first term is at most
$M^2 \cdot \P(\tilde\phi \in [0,\epsilon])$
because $|h| \le M$.
For the second, Markov's inequality gives
$\P(|h| > \epsilon) \le r^2/\epsilon^2$,
so the second term is at most
$M^2 \cdot r^2/\epsilon^2$.
Combining,
\begin{equation}\label{eq:junction-population}
\P f_h \le \P(\tilde\phi \in [0,\epsilon]) \cdot M^2
  + \frac{r^2}{\epsilon^2} \cdot M^2 =: \mu.
\end{equation}

\begin{lemma}[Dead-zone concentration]
\label{lem:dead-zone-concentration}
Under the notation of \Cref{lem:two-pile},
let $H_r$ be a class of increments and suppose:
\begin{enumerate}
\item[\textup{(i)}] \textbf{(Boundedness)}
  $|\tilde h| \le M$ almost surely for all $h \in H_r$.
\item[\textup{(ii)}] \textbf{(Rademacher)}
  $\E\sup_{h \in H_r} \abs*{\hP \sigma_i h(Z_i)}
    \le R.$
\item[\textup{(iii)}] \textbf{(Population bound)}
  $\P f_h \le \mu$ for all $h \in H_r$.
\end{enumerate}
Then for any $t > 0$,
there exists an event $E$ with
\[
\P(E^c)
  \le \exp\!\qty(-\frac{n\, t^2}
    {2M^2(\mu + 8MR + t)})
\]
such that on $E$,
$\sup_{h \in H_r} \hP f_h \le \mu + 4MR + t$.
\end{lemma}

Combining \Cref{lem:two-pile,lem:dead-zone-concentration}:
on the event $E$,
the SLB tolerance for all $h \in H_r$ simultaneously is
\begin{equation}\label{eq:two-pile-tolerance}
\xi = \xi_0
  + \frac{\mu + 4MR + t}{r^2}.
\end{equation}
When the density near zero satisfies
a polynomial bound,
the dead-zone cost $\mu/r^2$
vanishes as a power of $r$.

\begin{corollary}[Polynomial density control]\label{cor:two-pile-polynomial}
Under the hypotheses of
\Cref{lem:two-pile,lem:dead-zone-concentration},
if $\P(\tilde\phi \in [0,\epsilon])
  \le C_{\tilde\phi}\, \epsilon^\alpha$
for some $\alpha > 0$,
taking $\epsilon = (r^2/C_{\tilde\phi})^{1/(\alpha+2)}$
in~\eqref{eq:junction-population} gives
$\mu
  = 2\, C_{\tilde\phi}^{2/(\alpha+2)}\,
    r^{2\alpha/(\alpha+2)}\, M^2$.
Equivalently,
$\mu/r^2 = 2\, C_{\tilde\phi}^{2/(\alpha+2)}\,
  r^{-4/(\alpha+2)}\, M^2$.
The linear case $\alpha = 1$ gives
$\mu/r^2 = 2\, C_{\tilde\phi}^{2/3}\, M^2\, r^{-4/3}$.
\end{corollary}

\subsubsection{Entropy-like dispersions}

For dispersions satisfying
the pseudo-self-concordance condition
$|\dddchis| \le \beta\,\ddchis$,
the curvature $\ddchis$ is no longer constant,
so derivative equicontinuity (condition~\ref{cond:shell-iv} of \Cref{lem:bregman-shell})
produces a nontrivial equicontinuity constant $L > 0$.
The growth factor $\Omega = e^{\beta M}$
controls the ratio between the curvature
at a displaced point $\tilde\phi + h$
and at $\tilde\phi$ itself;
it enters the SLB through the Bregman sup-norm
$\overline{D_{\chis}} = \frac{1}{2}\Omega\,\overline{\ddchis}\,M^2$.
When $\ddchis$ is also convex (curvature growth),
\Cref{lem:quadratic-bregman} improves the
ball containment to the effective radius
$r_2^2 = 2\Omega^{1/3}\,r^2$.

\theoremvariant{E}
\begin{corollary}[Setting E: entropy-like sphere conditions]
\label{cor:ulb-entropy}
Suppose $\chis$ is twice continuously differentiable
with $\ddchis$ convex
and $|\dddchis| \le \beta\,\ddchis$.
Let $0 < \underline{\ddchis}
  \le \ddchis\{\tilde\phi(Z)\}
  \le \overline{\ddchis}$ a.s.
Let $H_r^\circ = \{h : \P D_{\chis}(\tilde\phi+h,\tilde\phi) = r^2\}$
and suppose $|h| \le M$ for all $h \in H_r^\circ$.
Write $\Omega = e^{\beta M}$
and $\overline{D_{\chis}}
  = \tfrac{1}{2}\,\Omega\,\overline{\ddchis}\,M^2$.
Then conditions~\ref{cond:shell-i}, \ref{cond:shell-iii-prime},
and~\ref{cond:shell-iv} of \Cref{lem:bregman-shell}
hold with $d\Q = \ddchis(\tilde\phi)\,d\P$,
SLB parameters
\[
k = c_1\, n\, \xi^2\, r^2 / (4\overline{D_{\chis}}),
\qquad
\ell = c_0\, n\, \xi\, r^2 / (2\overline{D_{\chis}}),
\]
and derivative equicontinuity parameters
$L^2 = \Omega^2\,\overline{\ddchis}$,
$r_2^2 = 2\Omega^{1/3}\,r^2$.
\end{corollary}
\resettheorem

\subsubsection{Proofs of application corollaries}
\label{sec:ulb-application-proofs}

\begin{proof}[Proof of Corollary~\ref{cor:ulb-quadratic}]
For $\chis(\phi) = \phi^2/2$,
the Bregman divergence is $D_{\chis}(\tilde\phi+h,\tilde\phi) = h^2/2$,
so $\P D_{\chis} = r^2$ gives $\|h\|_{L_2}^2 = 2r^2$.
\Cref{lem:slb-bounded} applied to $h$
with $|h| \le M$ and $\|h\|_{L_2}^2 = 2r^2$
gives a stable lower bound for $h^2 = 2D_{\chis}$
with parameters proportional to $n\xi^2 r^2/M^2$ and $n\xi r^2/M^2$.
Applying twice (once for the Bregman divergence, once for pairwise differences)
at tolerance $\xi/2$ each and taking a union bound
gives the stated parameters
$k = 2c_1\, n\, \xi^2\, r^2/M^2$
and $\ell = 2c_0\, n\, \xi\, r^2/M^2$,
where $c_0, c_1$ include the factor of $2$ from the tolerance halving.

Because curvature is constant ($\ddchis = 1$),
the derivative $\dchis(\tilde\phi + h) - \dchis(\tilde\phi + h') = h - h'$,
so condition~\ref{cond:shell-iv} of \Cref{lem:bregman-shell}
holds with $L = 1$ and $r_2^2 = 2r^2$.
\end{proof}

\begin{proof}[Proof of Lemma~\ref{lem:two-pile}]
\paragraph{Step 1: Case analysis.}
For $\chis(x) = x_+^2/2$, we have $\dchis(x) = x_+$, so the Bregman divergence is
\[
D_{\chis}(a, b) = \frac{a_+^2}{2} - \frac{b_+^2}{2} - b_+(a - b).
\]
We partition coordinates with $\tilde\phi \ge 0$ into two piles:
\begin{itemize}
\item \textbf{Pile 1} ($\tilde\phi \ge 0$ and $\tilde\phi + h \ge 0$): Here $D = h^2/2$.
\item \textbf{Pile 2} ($\tilde\phi \ge 0$ and $\tilde\phi + h < 0$): Here $|h| > \tilde\phi$ and
\[
D = \tilde\phi|h| - \frac{\tilde\phi^2}{2} = \frac{h^2}{2} - \frac{(|h| - \tilde\phi)^2}{2}.
\]
\end{itemize}
On coordinates with $\tilde\phi < 0$,
the weights vanish ($\tilde\phi_+ = 0$),
so $\tilde h = \dchis(\tilde\phi + h) - \dchis(\tilde\phi) = (\tilde\phi + h)_+ - 0 = (\tilde\phi + h)_+$.
The Bregman divergence
$D = (\tilde\phi + h)_+^2/2 \ge 0$
is nonnegative,
while the quadratic comparator $\tilde h^2/2 = (\tilde\phi + h)_+^2/2 = D$,
so the Bregman divergence matches the quadratic exactly.
These coordinates contribute nonnegatively to both sides
and can be dropped from the lower bound.

\paragraph{Step 2: Decomposition.}
For any $J \subset \{1,\ldots,n\}$, decompose the sample norm on $\{\tilde\phi \ge 0\}$ as
\[
\sum_{i \notin J} \tilde h_i^2 = \underbrace{\sum_{\substack{i \notin J \\ \text{Pile 1}}} h_i^2}_{S_{\mathrm{good}}} + \underbrace{\sum_{\substack{i \notin J \\ \text{Pile 2}}} h_i^2}_{S_{\mathrm{bad}}}.
\]
On Pile~1, $D_i = h_i^2/2 = \tilde h_i^2/2$,
so the Pile~1 contribution is $S_{\mathrm{good}}/2$.
On Pile~2, the identity
$D = h^2/2 - (|h| - \tilde\phi)^2/2$
shows $0 \le D \le h^2/2$,
so $D_i \ge 0$ on Pile~2.
Coordinates with $\tilde\phi < 0$ also have $D_i \ge 0$
as argued above.
Dropping these nonnegative contributions,
\[
\sum_{i \notin J} D_i
  \ge \sum_{\substack{i \notin J \\ \text{Pile 1}}} D_i
  = \frac{S_{\mathrm{good}}}{2}
  = \frac{1}{2}\left(\sum_{i \notin J} \tilde h_i^2 - S_{\mathrm{bad}}\right).
\]

\paragraph{Conclusion.}
Because $|\tilde h| \le M$,
\Cref{lem:slb-bounded} gives
a $(\xi_0, \ell, k)$ stable lower bound for $\tilde h^2$.
Because $f_h \ge 0$
and $S_{\mathrm{bad}}$ sums over $i \notin J$,
$S_{\mathrm{bad}}/n
  \le \frac{1}{n}\sum_{i=1}^n f_h(Z_i)
  = \hP f_h$.
For any $|J| \le \ell$,
\[
\frac{1}{n}\sum_{i \notin J} D_i
  \ge \frac{1}{2}\qty{(1-\xi_0) r^2
    - \hP f_h}
  = \frac{1-\xi}{2}\, r^2. \qedhere
\]
\end{proof}

\begin{proof}[Proof of Lemma~\ref{lem:dead-zone-concentration}]
We apply \Cref{lem:talagrand} to
$W = \sup_{h \in H_r} (\hP - \P) f_h$
with bound $b = M^2$.
For the variance, $f_h^2 \le M^2 f_h$ gives
$V = \sup_h \frac{1}{n}\sum_i \Var f_h(Z_i)
  \le M^2\, \mu$.
For the expectation,
viewing $f_h(Z)$ as a function of $h(Z)$ for fixed $Z$,
it equals $h^2 \cdot \ind{h < -\tilde\phi}$
on $\{\tilde\phi \ge 0\}$
and vanishes elsewhere.
This is $2M$-Lipschitz in $h$ and vanishes at $h = 0$,
so by the contraction principle (\Cref{lem:contraction})
and symmetrization (\Cref{lem:symmetrization}),
$\E W \le 4MR$.
\Cref{lem:talagrand} gives
\[
\P(W \ge 4MR + t)
  \le \exp\!\qty(-\frac{n\, t^2}
    {2M^2(\mu + 8MR + t)}).
\]
Define $E = \{W \le 4MR + t\}$.
Because $\sup_h \hP f_h \le \mu + W$,
on $E$ we have $\sup_h \hP f_h \le \mu + 4MR + t$
uniformly over $h \in H_r$.
\end{proof}

\begin{proof}[Proof of Corollary~\ref{cor:two-pile-polynomial}]
Substituting
$\P(\tilde\phi \in [0,\epsilon])
  \le C_{\tilde\phi}\, \epsilon^\alpha$
into the population bound~\eqref{eq:junction-population}
gives
\[
\mu = C_{\tilde\phi}\, \epsilon^\alpha\, M^2
  + r^2\, M^2 / \epsilon^2.
\]
The two terms are balanced
by the choice
$\epsilon = (r^2/C_{\tilde\phi})^{1/(\alpha+2)}$,
giving $\mu
  = 2\, C_{\tilde\phi}^{2/(\alpha+2)}\,
    r^{2\alpha/(\alpha+2)}\, M^2$.
Equivalently,
$\mu/r^2 = 2\, C_{\tilde\phi}^{2/(\alpha+2)}\,
  r^{-4/(\alpha+2)}\, M^2$,
and the linear case $\alpha = 1$
gives $\mu/r^2 = 2\, C_{\tilde\phi}^{2/3}\, M^2\, r^{-4/3}$.
\end{proof}

\begin{proof}[Proof of Corollary~\ref{cor:ulb-entropy}]
\paragraph{Outline.}
We verify conditions~\ref{cond:shell-i}, \ref{cond:shell-iii-prime}, and~\ref{cond:shell-iv}
of \Cref{lem:bregman-shell}
using the pseudo-self-concordance bounds
from \Cref{lem:e-type-bounds},
then compute the failure probability
from \Cref{thm:bregman-lower}.

\paragraph{Condition~\ref{cond:shell-i}: Bregman SLB.}
The Bregman sandwich~\eqref{eq:bregman-sandwich}
gives $g(Z) = D_{\chisZ}(\tilde\phi+h,\tilde\phi)
  \le \frac{1}{2}\Omega\,\overline{\ddchis}\, h^2$,
so $\|g\|_{L_\infty} \le \overline{D_{\chis}}$.
We apply \Cref{lem:slb-bounded} to $\sqrt{g}$:
$|\sqrt{g}| \le \sqrt{\overline{D_{\chis}}}$
and $\|\sqrt{g}\|_{L_2}^2 = \P g = r^2$,
giving a $(\xi/2, \ell, k)$
stable lower bound for $(\sqrt{g})^2 = g$
with
$k = c_1\, n\, (\xi/2)^2\, r^2/\overline{D_{\chis}}
  = c_1\, n\, \xi^2\, r^2/(4\overline{D_{\chis}})$
and
$\ell = c_0\, n\, (\xi/2)\, r^2/\overline{D_{\chis}}
  = c_0\, n\, \xi\, r^2/(2\overline{D_{\chis}})$.

\paragraph{Condition~\ref{cond:shell-iii-prime}: $L_2$ SLB for differences.}
For $h_1, h_2 \in H_r^\circ$,
boundedness gives $\|h_1-h_2\|_{L_\infty} \le 2M$.
\Cref{lem:slb-bounded}
applied to $h_1 - h_2$
gives a stable lower bound for $(h_1-h_2)^2$
with
$k = c_1\, n\, (\xi/2)^2\,
  \|h_1-h_2\|_{L_2}^2 / (4M^2)$.
Because $\|h_1 - h_2\|_{L_2(\Q)} \ge \epsilon_0$
in condition~\ref{cond:shell-iii-prime}
and $\|h_1 - h_2\|_{L_2}^2 \ge \|h_1 - h_2\|_{L_2(\Q)}^2/\overline{\ddchis}$,
$k \ge c_1\, n\, \xi^2\, \epsilon_0^2/(16M^2\overline{\ddchis})$,
which exceeds $k/16$ from condition~\ref{cond:shell-ii}
when the Rademacher condition holds.

\paragraph{Condition~\ref{cond:shell-iv}: Derivative equicontinuity.}
The derivative Lipschitz
bound~\eqref{eq:derivative-lipschitz}
gives
$|\dchis(\tilde\phi+h) - \dchis(\tilde\phi+h')|
  \le \Omega\, \ddchis(\tilde\phi)\, |h-h'|$.
Squaring and integrating against~$\P$,
and using
$\ddchis(\tilde\phi)^2
  \le \overline{\ddchis}\, \ddchis(\tilde\phi)$
(because $\ddchis(\tilde\phi) \le \overline{\ddchis}$),
\[
\P\{\dchis(\tilde\phi+h)
  - \dchis(\tilde\phi+h')\}^2
  \le \Omega^2\, \overline{\ddchis}\,
    \|h-h'\|_{L_2(\Q)}^2,
\]
so $L^2 = \Omega^2\, \overline{\ddchis}$.

For the ball containment,
\Cref{lem:quadratic-bregman}
and the curvature ratio~\eqref{eq:curvature-ratio-psc}
at the intermediate displacement $h/3$ give
\[
D_{\chisZ}(\tilde\phi+h,\tilde\phi)
  \ge \tfrac{1}{2}\,\Omega^{-1/3}\,\ddchis(\tilde\phi)\,h^2.
\]
The exponent $1/3$ (rather than $1$) comes from
evaluating the curvature ratio at $h/3$
in \Cref{lem:quadratic-bregman},
which is sharper than the sandwich bound
at $h$.
Integrating against~$\P$ gives
$r^2
  \ge \tfrac{1}{2}\,\Omega^{-1/3}\,\|h\|_{L_2(\Q)}^2$,
so $r_2^2
  = \sup_{h \in H_r^\circ}\|h\|_{L_2(\Q)}^2
  \le 2\,\Omega^{1/3}\, r^2$.
\end{proof}

\subsection{Uniform Lower Bound on the Ball}

The sphere bound from \Cref{lem:bregman-shell}
extends to the complement of a Bregman ball
by applying it at each radius in a dyadic grid.
The net and truncation probabilities
sum over $J$ shells;
the fluctuation uses the joint Efron--Stein supremum,
so the variance proxy is a sup over shells, not a sum.

\begin{theorem}[Uniform lower bound for Bregman divergences]
\label{thm:bregman-lower}
Under the hypotheses of \Cref{lem:bregman-shell},
suppose conditions~\ref{cond:shell-i}--\ref{cond:shell-iv}
hold for $H_r^\circ$
at each radius $r \in [\underline{r}, \overline{r}]$,
with parameters $k_r$, $\ell_r$, $L_r$, $r_{2,r}$
depending on $r$.
Then for any $t > 0$,
\[
  \inf_{\substack{h \in H \\
      \underline{r}^2 \le \P D_{\chisZ}(f+h,f) \le \overline{r}^2}}
    \frac{\hP D_{\chisZ}(f+h, f)}
         {\P D_{\chisZ}(f+h, f)}
  \ge 1 - \xi_{t,k_\star}
\]
on the intersection of $E$
and an event of probability at least
$1 - p_{\textup{net}} - p_{\textup{trunc}} - p_{\textup{fluct}}$,
where $k_\star = \inf_{r \in [\underline{r},\overline{r}]} k_r$,
$\xi_{t,k} = (1+t)\xi + (2-\xi)e^{-k}$,
\begin{align*}
p_{\textup{net}}
  &= \sum_{j=0}^{J} 2 e^{-15 k_j/16}, &
p_{\textup{trunc}}
  &= \sum_{j=0}^{J} e^{-\ell_j/16}, &
p_{\textup{fluct}}
  &= \frac{20}{t^2}
     \sup_{r \in [\underline{r},\overline{r}]}
     \frac{L_r^2\, r_{2,r}^2}{\ell_r\, r^2},
\end{align*}
and $r_j = 2^j \underline{r}$,
$J = \lceil\log_2(\overline{r}/\underline{r})\rceil$.
The same result holds with
conditions~\ref{cond:shell-ii-prime}--\ref{cond:shell-iii-prime}
in place of \ref{cond:shell-ii}--\ref{cond:shell-iii}.
\end{theorem}

\begin{proof}
The proof follows that of
\Cref{lem:bregman-shell},
applied at each radius in a dyadic grid.
Write $r_j = 2^j \underline{r}$ for
$j = 0, \ldots, J = \lceil \log_2(\overline{r}/\underline{r}) \rceil$.
Every $h$ with $\P D_{\chisZ}(f+h,f) \in [\underline{r}^2, \overline{r}^2]$
lies on some Bregman sphere $H_{\sqrt{\P D_h}}$;
the dyadic grid covers these
up to a factor of $2$ in radius.

\paragraph{Steps 1--2: net and truncation.}
Apply Steps~1--2 of \Cref{lem:bregman-shell}
at each radius~$r_j$.
The net SLB (Step~1) at radius $r_j$
fails with probability at most $2e^{-15k_j/16}$;
the truncation (Step~2) fails with probability
at most $e^{-\ell_j/16}$.
By union bound, all nets and truncations
hold simultaneously with probability
at least $1 - p_{\textup{net}} - p_{\textup{trunc}}$.

\paragraph{Step 3: joint fluctuation bound.}
Rather than bounding the fluctuation
on each shell separately,
we take a single supremum.
On the event that all truncations hold,
define the normalized fluctuation
\[
  F = \sup_{j=0}^{J}
      \sup_{h \in H_{r_j}^\circ}
      \frac{(\hP - \P)\,
        \Gamma_h^{(j)} \cdot 1_{A_j}(h)}{r_j^2}
\]
where $\Gamma_h^{(j)}$ and $1_{A_j}$
are the remainder and good-coordinate indicator
from the proof of \Cref{lem:bregman-shell}
at radius $r_j$.

Writing
$g_{j,h}(Z)
  = \Gamma_h^{(j)}(Z) \cdot 1_{A_j}(h)(Z) / r_j^2$,
the supremum has the form
$F = \sup_{(j,h)} (\hP - \P)\, g_{j,h}$.
\Cref{lem:efron-stein} gives
\[
  \Var(F) \le \frac{2}{n}
    \sup_{(j,h)} \P\, g_{j,h}^2.
\]
The variance proxy is a supremum over shells,
not a sum:
$F = \sup_{(j,h)}(\hP - \P) g_{j,h}$
is a single supremum
over all shells simultaneously,
so Efron--Stein gives the variance
of the joint supremum
in terms of the worst single-shell
second moment.
No union bound over shells is needed;
the probability bound
for the joint fluctuation
is the same as for any one shell.
The derivative equicontinuity bound
from the proof of \Cref{lem:bregman-shell}
(applied at each radius $r_j$)
gives
$\P\, g_{j,h}^2
  \le 2A_j^2 L_j^2
    (\epsilon_j^2 + r_{2,j}^2) / r_j^4$
where $\epsilon_j = c_0\xi r_j$,
so that
$\Var(F)
  \le 4\xi^2
    \sup_j L_j^2
      (c_0^2\xi^2 + r_{2,j}^2/r_j^2) / \ell_j$.
Chebyshev gives
$\P(F - \E F > t\xi/2)
  \le \Var(F)/(t\xi/2)^2$.
The Rademacher condition~\ref{cond:shell-iii}
bounds $\E F \le t\xi/2$,
so $F > t\xi$ implies $F - \E F > t\xi/2$.
Together,
\[
  \P(F > t\xi)
    \le \frac{20}{t^2}
      \sup_j
      \frac{L_j^2\, r_{2,j}^2}{\ell_j\, r_j^2}
    = p_{\textup{fluct}}.
\]

\paragraph{Conclusion.}
On the event that $F \le t\xi$
and all net and truncation events hold,
the conclusion of \Cref{lem:bregman-shell}
gives $\hP D_{\chisZ}(f+h,f) \ge (1 - \xi_{t,k_j})\, r_j^2$
for every $h \in H_{r_j}^\circ$ and every $j$.
Because $\xi_{t,k}$ decreases in $k$,
the infimum over shells is
$1 - \xi_{t,k_\star}$
with $k_\star = \min_j k_j$.
\end{proof}

Specializing the sphere conditions from
\Cref{cor:ulb-quadratic,cor:ulb-entropy}:

\setupvariants

\theoremvariant{Q}
\begin{corollary}[Setting Q: quadratic ULB]
\label{cor:ulb-quadratic-ball}
Under the hypotheses of \Cref{cor:ulb-quadratic},
\Cref{thm:bregman-lower} gives
\[
p \le C\, \log(\overline{r}/\underline{r})\,
  \exp\!\qty(-c\, n\, \xi^2\, \underline{r}^2/M^2)
  + \frac{C\, M^2}{t^2\, n\, \xi\, \underline{r}^2}.
\]
\end{corollary}
\begin{proof}
Substituting the parameters from \Cref{cor:ulb-quadratic}
into \Cref{thm:bregman-lower}:
the net and truncation probabilities give
$p_{\textup{net}} + p_{\textup{trunc}}
  \le C J\, e^{-c k_\star}$
where $k_\star = 2c_1\, n\, \xi^2\, \underline{r}^2/M^2$
and $J = \lceil\log_2(\overline{r}/\underline{r})\rceil$.
The fluctuation term contributes
$p_{\textup{fluct}}
  \le C\, M^2 / (t^2\, n\, \xi\, \underline{r}^2)$,
using $L^2 r_2^2 / (\ell\, r^2)
  = 2r^2 / (2c_0\, n\, \xi\, r^2/M^2 \cdot r^2)
  = M^2 / (c_0\, n\, \xi\, r^2)$
at the worst-case (smallest) radius.
\end{proof}
\resettheorem

\theoremvariant{E}
\begin{corollary}[Setting E: entropy-like ULB]
\label{cor:ulb-entropy-ball}
Under the hypotheses of \Cref{cor:ulb-entropy},
\Cref{thm:bregman-lower} gives
\[
p \le C\, \log(\overline{r}/\underline{r})\,
  \exp\!\qty(-c\, n\, \xi^2\,
    \underline{r}^2/\overline{D_{\chis}})
  + \frac{C\, \Omega^{10/3}\,
    \overline{\ddchis}^2\, M^2}
    {t^2\, n\, \xi\, \underline{r}^2}.
\]
\end{corollary}
\begin{proof}
Substituting the SLB and equicontinuity parameters
from \Cref{cor:ulb-entropy}
into \Cref{thm:bregman-lower},
the net and truncation probabilities decay as
$p_{\textup{net}} + p_{\textup{trunc}}
  \le C\, J\, e^{-c\, n\, \xi^2\, \underline{r}^2/\overline{D_{\chis}}}$
where $J = \lceil\log_2(\overline{r}/\underline{r})\rceil$.
The fluctuation probability satisfies
\[
p_{\textup{fluct}}
  = \frac{20}{t^2}
    \sup_r \frac{L^2\, r_2^2}{\ell\, r^2}
  = \frac{20}{t^2}\,
    \frac{\Omega^2\, \overline{\ddchis}\,
      2\Omega^{1/3}\, r^2}
    {(c_0\, n\, \xi\, r^2 / (2\overline{D_{\chis}}))\, r^2}
  = \frac{80\,\Omega^{7/3}\,
    \overline{\ddchis}\, \overline{D_{\chis}}}
    {c_0\, t^2\, n\, \xi\, r^2}
  = \frac{40\,\Omega^{10/3}\,
    \overline{\ddchis}^2\, M^2}
    {c_0\, t^2\, n\, \xi\, r^2},
\]
using $\overline{D_{\chis}} = \frac{1}{2}\Omega\,\overline{\ddchis}\,M^2$
in the last step.
This decays as $1/(n r^2)$.
\end{proof}
\resettheorem

\section{Bias Bound via Offset Complexity}
\label{sec:bias-bound}

\subsection{Results}
\label{sec:offset-vs-local}

In this section, we bound the \emph{maximal bias} for the weights $\hgamma$ minimizing our primal problem
\eqref{eq:primal-moment-matching}, i.e.
\[
b(\hgamma) = \sup_{f \in B_\rho}  \dotpsiZ(f) - \hP\, \hgamma(Z) f(Z).
\]
Our bound is phrased in terms of an \emph{offset complexity} of our seminorm's unit ball.

\begin{theorem}[Bias bound]
\label{thm:bias-decomp-restate}
Let $\chi$ and $\zeta$ satisfy the assumptions of \Cref{sec:duality} and $\tilde\phi$
be the minimizer of the population dual loss $L_P^\conj$ \eqref{eq:dual-problem}.  For
any $0 < s < 1/\{2\rho(\tilde\phi)\}$ and $\mu_\rho \in \partial\zetas\{\rho(\tilde\phi)\}$,
\[
\qty|b(\hgamma)| \le 2\overline\xi(s) \qfor \overline\xi(t) = \sup_{f \in B_\rho - t\tilde\phi} (\hP - \P)\dotpsiZ(f) - \tilde\gamma(Z) f(Z) - t \cdot \hP\, D_{\chisZ}(\tilde\phi + f/t, \tilde\phi) + \mu_\rho\qty{1 - t\,\rho(\tilde\phi)}.
\]
\end{theorem}

\subsection{Proofs}
\label{sec:bias-proofs}

We work in the abstract framework
of \Cref{sec:duality-proofs}:
$V$ is the space of centered functions,
$W = V^*$ is its dual,
$\inner{\gamma}{f} = \hP \gamma(Z) f(Z)$
is the sample pairing,
and $B_\rho = \{f : \rho(f) \le 1\}$
is the unit seminorm ball.

\subsubsection*{Proof of Theorem~\ref{thm:bias-decomp}}

The proof works in the abstract bias-dispersion framework introduced in \Cref{sec:duality-proofs}: a vector space $V$ with seminorm $\rho$, a weight space $W$ with bilinear pairing $\inner{\cdot}{\cdot}: W \times V \to \R$, convex $\chi: W \to \R \cup \{+\infty\}$ with conjugate $\chis$, and a linear target functional $\psi: V \to \R$. Unlike duality, the bias analysis does not require $W$ to be finite-dimensional.

\paragraph{The primal problem.} The bias-dispersion tradeoff is governed by
\begin{equation}
\label{eq:bias-primal}
p(\gamma) = \zeta\{b(\gamma)\} + \chi(\gamma) \qwhere b(\gamma) = \sup_{f \in B_\rho} \psi(f) - \inner{\gamma}{f}
\end{equation}
is the worst-case bias at weights $\gamma$. This matches the paper's primal problem~\eqref{eq:primal-moment-matching}: the penalty $\zeta$ controls how imbalance is scaled relative to dispersion $\chi(\gamma)$. For linear $\zeta(b) = b/s$, the primal becomes $p(\gamma) = b(\gamma)/s + \chi(\gamma)$, which is equivalent to minimizing $b(\gamma) + s \cdot \chi(\gamma)$. We work with this linearized form, parametrized by $s > 0$:
\begin{equation}
\label{eq:modulus-primal}
\omega(s) = \inf_\gamma \qty{ b(\gamma) + s \cdot \chi(\gamma) }.
\end{equation}
The parameter $s$ controls the tradeoff: large $s$ penalizes dispersion heavily, forcing weights toward the interior of $\chi$'s effective domain. For general $\zeta$, the solution to \eqref{eq:bias-primal} corresponds to \eqref{eq:modulus-primal} at $s = 1/\zeta'(\hat b)$ where $\hat b$ is the optimal bias.

\paragraph{Assumptions on $\chi$.} Assume $\chi$ is coercive in addition to the conditions in \Cref{sec:duality-proofs}. Coercivity and strict convexity imply $\chis$ is finite everywhere and differentiable with continuous gradient $\nabla\chis$ \citep[Theorem 2.22]{peypouquet2015convex}.

\paragraph{Centering.} Fix a centering point $\phi^* \in V$ and let $\tilde\gamma = \nabla\chis(\phi^*) \in W$ be the corresponding centering direction.

\paragraph{Key objects.} The residual is
\[
M(f) = \psi(f) - \inner{\tilde\gamma}{f},
\]
measuring how far the target functional $\psi$ deviates from the ideal $f \mapsto \inner{\tilde\gamma}{f}$. The offset complexity is
\begin{equation}
\label{eq:abstract-offset-complexity}
\xi(s) = \sup_{\Delta \in B_\rho - s\phi^*} M(\Delta) - s \cdot D_{\chis}(\phi^* + \Delta/s, \phi^*)
\end{equation}
where $D_{\chis}(a, b) = \chis(a) - \chis(b) - \inner{\nabla\chis(b)}{a - b}$ is the Bregman divergence. The offset complexity measures how much the residual $M$ can push the optimizer away from the centering point $\phi^*$, offset by the Bregman divergence penalty.

\paragraph{Application to balancing weights.} The balancing weights problem \eqref{eq:dual-problem} fits the abstract framework with the sample pairing $\inner{\gamma}{f} = \hP \gamma(Z) f(Z)$, matching the primal loss $L_{\hP}$:
\[
\begin{array}{r@{\;=\;}l@{\qquad}l}
\inner{\gamma}{f} & \hP\, \gamma(Z)\, f(Z) & \text{(sample inner product)} \\[3pt]
\chis(\phi) & \hP\, \chisZ\{\phi(Z)\} & \text{(sample conjugate)} \\[3pt]
D_{\chis}(a, b) & \hP\, D_{\chisZ}\{a(Z), b(Z)\} & \text{(sample Bregman)} \\[3pt]
\psi(f) & \dotpsiZ(f) & \text{(target functional)} \\[3pt]
\tilde\gamma & \dchis \comp \tilde\phi & \text{(population weights)}
\end{array}
\]

The abstract framework yields a deterministic lemma bounding bias in terms of offset complexity.

\begin{lemma}[Bias bound]
\label{lem:bias-bound}
Let $\chi$ satisfy the assumptions of \Cref{thm:bias-decomp}. Fix $s > 0$ with $2s \cdot \rho(\phi^*) < 1$. If $\xi(s) \le \nu$, then $|b| \le 2\nu$. If $\nu/2 < \delta \le \xi(s/2)$ for some $\delta > 0$ as well, then $\abs{b} \ge 2\delta - \nu$.
\end{lemma}

The condition $2s \cdot \rho(\phi^*) < 1$ ensures $t\phi^* \in B_\rho$ for all $t \in [s/2, 2s]$, which gives non-negativity $\xi(t) \ge 0$ at the endpoints needed for the secant bounds.

\begin{proof}[Proof of Theorem~\ref{thm:bias-decomp}]
Set $\phi^* = \tilde\phi$ and $\psi = \dotpsiZ$. The centering direction is $\tilde\gamma = \nabla\chis(\tilde\phi) = \dchis \comp \tilde\phi$, and the residual is
\[
M(f) = \hP\, \dotpsiZ(f) - \hP\, \tilde\gamma(Z) f(Z) = \hP\, \dotpsiZ(f) - \tilde\gamma(Z) f(Z),
\]
where both terms are sample averages. Splitting $\hP = (\hP - \P) + \P$,
\[
M(f) = (\hP - \P)\, \dotpsiZ(f) - \tilde\gamma(Z) f(Z) + \P\, \dotpsiZ(f) - \tilde\gamma(Z) f(Z).
\]
The population first order condition at $\tilde\phi$ (minimizer of $L^\conj_{\P}$) identifies the second term: for $z^* \in \partial(\zetas \comp \rho)(\tilde\phi)$,
\begin{equation}
\label{eq:foc-population}
\P\, \dotpsiZ(g) - \tilde\gamma(Z) g(Z) = z^*(g) \qfor g \in V,
\end{equation}
so
\begin{equation}
\label{eq:residual-decomp}
M(f) = (\hP - \P)\, \dotpsiZ(f) - \tilde\gamma(Z) f(Z) + z^*(f).
\end{equation}
We now bound the subgradient term $z^*(f)$ uniformly over $f \in B_\rho - t\tilde\phi$. The chain rule decomposes $z^* = \mu_\rho \cdot w^*$ with $\mu_\rho \in \partial\zetas\{\rho(\tilde\phi)\}$ and $w^* \in \partial\rho(\tilde\phi)$. For any $f \in B_\rho - t\tilde\phi$, write $g = f + t\tilde\phi$, so that $g \in B_\rho$. Then
\[
z^*(f) = \mu_\rho\inner{w^*}{g - t\tilde\phi} = \mu_\rho\inner{w^*}{g} - \mu_\rho\, t\, \rho(\tilde\phi),
\]
using $\inner{w^*}{\tilde\phi} = \rho(\tilde\phi)$ (the norming condition for $w^* \in \partial\rho(\tilde\phi)$). Because $w^*$ is a subgradient of $\rho$, $\inner{w^*}{g} \le \rho(g)$; because $g \in B_\rho$, $\rho(g) \le 1$. So
\[
z^*(f) \le \mu_\rho - \mu_\rho\, t\, \rho(\tilde\phi) = \mu_\rho\qty{1 - t\,\rho(\tilde\phi)}.
\]
This bound is constant over $f \in B_\rho - t\tilde\phi$. Substituting \eqref{eq:residual-decomp} into the offset complexity $\xi(t)$ and pulling the constant outside the supremum gives $\xi(t) \le \overline\xi(t)$. The lemma gives $|b| \le 2\overline\xi(s)$.
\end{proof}

\begin{proof}[Refinement for linear penalties]
The bound $\inner{w^*}{f + t\tilde\phi} \le \rho(f + t\tilde\phi)$ discards the non-negative difference $\rho(f + t\tilde\phi) - \inner{w^*}{f + t\tilde\phi}$, which measures how much $w^*$ fails to norm $f + t\tilde\phi$. For linear $\zetas(\rho) = \eta\rho$, the subgradient is $z^* = \eta w^*$, so this difference enters the offset complexity scaled by $\eta$. Under strict complementarity, $\rho(g) - \inner{w^*}{g} \ge \kappa\,\rho(P_{I^c} g)$ for $g \in B_\rho$, so the difference is at least $\kappa\,\rho\{P_{I^c}(f + t\tilde\phi)\}$, penalizing off-support mass. Retaining it inside the supremum gives $\xi(s) \le \overline\xi_I(s) \le \overline\xi(s)$, and the lemma gives $|b| \le 2\overline\xi_I(s)$.
\end{proof}

\begin{lemma}[Perspective function convexity]
\label{lem:perspective-convexity}
If $g: \R^m \to \R \cup \{+\infty\}$ is convex,
then the perspective function
$(x, t) \mapsto t\, g(x/t)$
is jointly convex on $\R^m \times \R_{>0}$.
\end{lemma}
\begin{proof}
See \citet[Section~3.2.6]{boyd2004convex}.
\end{proof}

\begin{lemma}[Envelope theorem for convex programs]
\label{lem:envelope-theorem}
Let $g: V \times \R \to \R \cup \{+\infty\}$
be jointly convex
and let $v(s) = \inf_{x \in C} g(x, s)$
for a convex set $C$.
If the infimum is attained at $x^*(s)$
and $g$ is differentiable in $s$ at $(x^*(s), s)$,
then $\partial_s g(x^*(s), s) \in \partial v(s)$.
\end{lemma}
\begin{proof}
The function $\varphi(s') := g(x^*(s), s')$
satisfies $\varphi(s') \ge v(s')$ for all $s'$
(because $x^*(s) \in C$)
and $\varphi(s) = v(s)$
(because $x^*(s)$ attains the infimum).
Because $\varphi$ is convex
and touches $v$ from above at $s$,
any subgradient of $\varphi$ at $s$
is a subgradient of $v$:
$\partial_s g(x^*(s), s) \in \partial v(s)$.
\end{proof}

\begin{proof}[Proof of Lemma~\ref{lem:bias-bound}]
The proof proceeds in four steps: derive the dual problem, establish a modulus decomposition, prove properties of the offset complexity, then combine them to bound the bias.

\paragraph{Step 1: The dual problem.} The modulus \eqref{eq:modulus-primal} trades off worst-case bias against dispersion. We derive its dual by exchanging the order of optimization. Writing out the definition:
\[
\omega(s) = \inf_\gamma \sup_{f \in B_\rho} \psi(f) - \inner{\gamma}{f} + s \cdot \chi(\gamma).
\]
The objective is convex in $\gamma$ (because $\chi$ is convex and the supremum of linear functions is convex) and linear in $f$. Because $\chi$ is coercive and strictly convex, the infimum over $\gamma$ is achieved at some $\gamma^*(s)$ (\Cref{lem:coercive-minimum}). When the infimum is achieved, the minimax equals the maximin \citep[Proposition 2.63]{peypouquet2015convex}, so we may exchange inf and sup:
\[
\omega(s) = \sup_{f \in B_\rho} \psi(f) + \inf_\gamma -\inner{\gamma}{f} + s \cdot \chi(\gamma).
\]
The inner infimum over $\gamma$ can be computed via the convex conjugate, factoring out $s > 0$:
\[
\inf_\gamma -\inner{\gamma}{f} + s \cdot \chi(\gamma) = s \cdot \inf_\gamma \inner{\gamma}{-f/s} + \chi(\gamma).
\]
By definition of the convex conjugate, $\chis(\phi) = \sup_\gamma \inner{\gamma}{\phi} - \chi(\gamma)$, so $\inf_\gamma \inner{\gamma}{\phi} + \chi(\gamma) = -\chis(-\phi)$. Setting $\phi = -f/s$ gives
\[
\inf_\gamma -\inner{\gamma}{f} + s \cdot \chi(\gamma) = -s \cdot \chis(f/s).
\]
Substituting back,
\begin{equation}
\label{eq:modulus}
\omega(s) = \sup_{f \in B_\rho} \psi(f) - s \cdot \chis(f/s).
\end{equation}
This is the \emph{modulus of continuity} of the functional $\psi$ with respect to the Bregman penalty $\chis$. The scaling $f/s$ in the argument of $\chis$ reflects the bias-dispersion tradeoff: as $s \to 0$, the penalty on $\chis(f/s)$ grows, forcing $f$ toward the origin.

\paragraph{Step 2: Modulus decomposition.} We show $\omega(s) = c \cdot s + \xi(s)$ for a constant $c$ that depends only on the centering point $\phi^*$. The key is to reparametrize the supremum by centering at $s\phi^*$.

Substitute $f = s\phi^* + \Delta$ where $\Delta = f - s\phi^*$. As $f$ ranges over $B_\rho$, the offset $\Delta$ ranges over the translated set $B_\rho - s\phi^*$. The scaled argument becomes $f/s = \phi^* + \Delta/s$.

Expand $\chis$ at the centering point $\phi^*$ using the Bregman identity. For any $a, b$ in the domain of $\chis$:
\[
\chis(a) = \chis(b) + \inner{\nabla\chis(b)}{a - b} + D_{\chis}(a, b).
\]
Apply this with $a = \phi^* + \Delta/s$ and $b = \phi^*$:
\[
\chis(\phi^* + \Delta/s) = \chis(\phi^*) + \inner{\nabla\chis(\phi^*)}{\Delta/s} + D_{\chis}(\phi^* + \Delta/s, \phi^*).
\]
By the centering assumption, $\nabla\chis(\phi^*) = \tilde\gamma$. Multiply through by $-s$:
\[
-s \cdot \chis(\phi^* + \Delta/s) = -s \cdot \chis(\phi^*) - \inner{\tilde\gamma}{\Delta} - s \cdot D_{\chis}(\phi^* + \Delta/s, \phi^*).
\]

For the target functional $\psi$, use linearity and the definition $M(g) = \psi(g) - \inner{\tilde\gamma}{g}$:
\[
\psi(f) = \psi(s\phi^* + \Delta) = s \cdot \psi(\phi^*) + \psi(\Delta) = s\inner{\tilde\gamma}{\phi^*} + s \cdot M(\phi^*) + \inner{\tilde\gamma}{\Delta} + M(\Delta).
\]
The second equality uses $\psi(\phi^*) = \inner{\tilde\gamma}{\phi^*} + M(\phi^*)$ and $\psi(\Delta) = \inner{\tilde\gamma}{\Delta} + M(\Delta)$.

Combine the two expansions. The modulus objective $\psi(f) - s \cdot \chis(f/s)$ becomes:
\begin{align*}
&= \qty{ s\inner{\tilde\gamma}{\phi^*} + s \cdot M(\phi^*) + \inner{\tilde\gamma}{\Delta} + M(\Delta) } + \qty{ -s \cdot \chis(\phi^*) - \inner{\tilde\gamma}{\Delta} - s \cdot D_{\chis}(\phi^* + \Delta/s, \phi^*) }.
\end{align*}
The terms $\inner{\tilde\gamma}{\Delta}$ and $-\inner{\tilde\gamma}{\Delta}$ cancel. Group the remaining terms:
\[
= s \cdot \qty{ \inner{\tilde\gamma}{\phi^*} - \chis(\phi^*) + M(\phi^*) } + \qty{ M(\Delta) - s \cdot D_{\chis}(\phi^* + \Delta/s, \phi^*) }.
\]
Define $\check{c} = \inner{\nabla\chis(\phi^*)}{\phi^*} - \chis(\phi^*)$, the negative of the Bregman divergence $D_{\chis}(0, \phi^*)$ up to a sign, and $c = \check{c} + M(\phi^*)$. Taking the supremum over $\Delta \in B_\rho - s\phi^*$:
\[
\omega(s) = c \cdot s + \xi(s).
\]
The modulus decomposes into a linear part $c \cdot s$, which depends only on the centering point, and the offset complexity $\xi(s)$, which captures the deviation.

\paragraph{Step 3: Properties of $\xi$.} The offset complexity has two key properties: non-negativity and concavity. Neither requires the supremum to be attained. % $\xi \ge 0$ and $\xi$ concave hold by construction

\paragraph{Non-negativity.} The condition $2s \cdot \rho(\phi^*) < 1$ ensures $s\phi^* \in B_\rho$ for all $s$ in our range, so $\Delta = 0$ is feasible in $B_\rho - s\phi^*$. Evaluate the objective at $\Delta = 0$: $M(0) = \psi(0) - \inner{\tilde\gamma}{0} = 0$ by linearity of $\psi$, and $D_{\chis}(\phi^*, \phi^*) = 0$ because Bregman divergence vanishes at coincident arguments. Therefore the supremum is at least zero: $\xi(s) \ge 0$.

\paragraph{Concavity.} The objective $h(f, s) = \psi(f) - s \cdot \chis(f/s)$ is jointly concave in $(f, s)$ for $s > 0$. The map $(f, s) \mapsto s \cdot \chis(f/s)$ is the perspective function of $\chis$, which is jointly convex when $\chis$ is convex (\Cref{lem:perspective-convexity}). Because $\psi$ is linear, $h$ is the sum of a linear function and the negative of a jointly convex function, hence jointly concave.

To prove concavity of $\omega$, we use $\varepsilon$-approximate optimizers rather than assuming the supremum is attained. Let $s_1, s_2 > 0$, $t \in (0,1)$, and $\varepsilon > 0$. By definition of supremum, there exist $f_1, f_2 \in B_\rho$ with
\[
h(f_1, s_1) \ge \omega(s_1) - \varepsilon \qand h(f_2, s_2) \ge \omega(s_2) - \varepsilon.
\]
Set $f_t = t f_1 + (1-t)f_2$. Because $B_\rho$ is convex (it's a seminorm ball), $f_t \in B_\rho$. By joint concavity of $h$:
\[
h(f_t, t s_1 + (1-t)s_2) \ge t h(f_1, s_1) + (1-t)h(f_2, s_2).
\]
Because $f_t$ is feasible:
\[
\omega(t s_1 + (1-t)s_2) \ge h(f_t, t s_1 + (1-t)s_2) \ge t h(f_1, s_1) + (1-t)h(f_2, s_2) \ge t\omega(s_1) + (1-t)\omega(s_2) - \varepsilon.
\]
Because $\varepsilon > 0$ was arbitrary, $\omega$ is concave. Because $\xi(s) = \omega(s) - cs$ differs from $\omega$ by a linear function, $\xi$ is also concave.

\paragraph{Step 4: The bias bound.} We show that the bias equals $\xi(s) - s\tau$
for a subgradient $\tau$ of $\xi$ at~$s$,
then bound $|b|$ using secant inequalities
from the concavity and non-negativity of~$\xi$.
The linear part $cs$ cancels;
only the offset complexity enters.
The argument uses only concavity of $\xi$, not differentiability of the supremum.

\paragraph{The bias formula.} Let $\gamma^*$ achieve the infimum in \eqref{eq:modulus-primal}, so $\omega(s) = b(\gamma^*) + s \cdot \chi(\gamma^*)$. The bias at this optimum is $b(\gamma^*) = \omega(s) - s \cdot \chi(\gamma^*)$. By the envelope theorem (\Cref{lem:envelope-theorem}), the dispersion $\chi(\gamma^*)$ is a subgradient of $\omega$ at $s$: for $t = \chi(\gamma^*) \in \partial\omega(s)$,
\[
b = \omega(s) - s \cdot t.
\]
Because $\omega(s) = c \cdot s + \xi(s)$, any subgradient of $\omega$ at $s$ has the form $t = c + \tau$ where $\tau \in \partial\xi(s)$. Substituting:
\[
b = \{c \cdot s + \xi(s)\} - s \cdot (c + \tau) = \xi(s) - s \cdot \tau.
\]
The linear part cancels completely. The bias depends only on $\xi$ and its subgradient, not on the constant $c$.

\paragraph{Subgradient bounds from concavity.} A concave function $\xi$ has left and right derivatives at every interior point, with $\xi'_+(s) \le \xi'_-(s)$. Every subgradient $\tau \in \partial\xi(s)$ satisfies $\xi'_+(s) \le \tau \le \xi'_-(s)$. We bound these derivatives using secants.

For the upper bound, concavity implies $\xi(s/2) \le \xi(s) + \xi'_-(s)(s/2 - s)$, i.e.\ the graph lies below its tangent lines. Rearranging:
\[
\xi'_-(s) \le \frac{\xi(s) - \xi(s/2)}{s - s/2} = \frac{\xi(s) - \xi(s/2)}{s/2}.
\]
By hypothesis $\xi(s) \le \nu$, and by non-negativity $\xi(s/2) \ge 0$:
\[
\xi'_-(s) \le \frac{\nu - 0}{s/2} = \frac{2\nu}{s}.
\]

For the lower bound, concavity implies $\xi(2s) \le \xi(s) + \xi'_+(s)(2s - s)$. Rearranging:
\[
\xi'_+(s) \ge \frac{\xi(2s) - \xi(s)}{2s - s} = \frac{\xi(2s) - \xi(s)}{s}.
\]
By non-negativity $\xi(2s) \ge 0$, and by hypothesis $\xi(s) \le \nu$:
\[
\xi'_+(s) \ge \frac{0 - \nu}{s} = -\frac{\nu}{s}.
\]

\paragraph{Conclusion.} The secant bounds give
\[
2\xi(s/2) - \xi(s) \le b \le 2\xi(s) - \xi(2s).
\]
By non-negativity $\xi(s/2), \xi(2s) \ge 0$ and the hypothesis $\xi(s) \le \nu$, we have $b \le 2\nu$ and $b \ge -\nu$, so $|b| \le 2\nu$. If additionally $\delta \le \xi(s/2)$ with $\delta > \nu/2$, then $b \ge 2\delta - \nu > 0$, giving the two-sided bound $2\delta - \nu \le |b| \le 2\nu$.
\end{proof}

%% Sparsity-dependent bias subsection shelved to parked-appendices.tex (2026-02-11).

\section{Offset Complexity}
\label{sec:offset-complexity}

The central object is the offset complexity
of the unit seminorm ball.

\paragraph{Kernel balance.}
Let $K = \ker(\rho)$
and $V = K^{\perp}_{\LtP}$,
the kernel's $\LtP$-orthogonal complement.
Because $\rho(g + \alpha k) = \rho(g)$
for all $\alpha$
and $k \in K$,
the first order condition
for $\hat\phi = \argmin L_{\hP}$
is unconstrained on $K$:
\begin{equation}
\label{eq:kernel-balance}
\hP\{(\dchis \comp \hat\phi)\, k - \dotpsiZ(k)\} = 0
\qquad \text{for all } k \in K.
\end{equation}
The population minimizer $\tilde\phi$
satisfies the same condition under~$\P$.

\paragraph{Curvature-weighted norm.}
For increments $h$ around $\tilde\phi$, define
\begin{equation}\label{eq:curvature-norm}
\norm{h}_{\chis}^2 = \P (\ddchis \comp \tilde\phi) h^2.
\end{equation}
For the quadratic $\chis(\phi) = \phi^2/2$,
this reduces to $\|h\|_{\LtP}^2$.
The population Bregman divergence
$\P D_{\chis}(\tilde\phi + h, \tilde\phi)$
equals $\|h\|_{\chis}^2/2$
when $\chis$ is quadratic;
pseudo-self-concordance makes it comparable
in general.

\paragraph{Offset complexity.}
The \emph{offset complexity} at scale $s > 0$ is
\begin{equation}\label{eq:offset-complexity}
\Xi(s) = \sup_{h \in B_\rho} \qty[ (\hP - \P)\{\dotpsiZ(h) - (\dchis \comp \tilde\phi)\, h\} - \frac{1}{s}\, \hP\, D_{\chisZ}(\tilde\phi + s\,h,\, \tilde\phi) ].
\end{equation}
The Bregman offset
$(1/s)\,\hP\, D_{\chisZ}(\tilde\phi + s\,h,\, \tilde\phi)
  \approx (s/2)\, \hP (\ddchis \comp \tilde\phi)\, h^2$
penalizes large increments,
sharpening the complexity
relative to the bare empirical process.

\paragraph{Setup.}
\label{par:offset-setup}
The results below
use the following notation.
We work with the base penalty $\chis$ throughout;
by the sign-flip reduction (\Cref{sec:scope}),
$D_{\chisZ}(\phi+h,\phi)$ depends only on $\ddchis$.
Write $u_0 = \sup\{\norm{h}_{\chis} : h \in B_\rho\}$.

\begin{theorem}[Offset complexity in all regimes]
\label{thm:xi-all-regimes}
Suppose $\alpha\, \ddchis \le \dddchis \le \beta\, \ddchis$
for some $0 < \alpha \le \beta$
(Setting~E).
Suppose the interpolation inequality~\eqref{eq:interpolation}
holds with exponent $\theta \in (0,1)$
and constant $\kappa$,
the Bregman lower bound
(\Cref{thm:bregman-lower})
gives
$\hP\, D_{\chis}(\tilde\phi + h, \tilde\phi)
  \ge (1-\epsilon)\, \P\, D_{\chis}(\tilde\phi + h, \tilde\phi)$
uniformly over $\rho(h) \le s$,
$\underline{\ddchis} \le \ddchis(\tilde\phi) \le \overline{\ddchis}$
under~$\P$,
and
\[
\sup_{\substack{\rho(h) \le s \\
    \norm{h}_{\chis} \le r}}
  (\hP - \P)\qty{\dotpsiZ(h) - (\dchis \comp \tilde\phi)\,h}
  \le C_R\, n^{-1/2}\, s^{1-\theta}\, r^\theta
\]
for all $r \le s\, u_0$.
Then
\begin{equation}
\label{eq:xi-all-regimes}
\Xi(s) \le C\, n^{-1/2}\,
\begin{cases}
  (n^{1/2}\, s)^{-\theta/(2-\theta)}
    & \text{if } s \le n^{(1-\theta)/(2+8\theta)}, \\[6pt]
  \qty(\dfrac{\log n}{\alpha\, s})^{\theta/(1+\theta)}
    & \text{if } s > n^{(1-\theta)/(2+8\theta)},
\end{cases}
\end{equation}
for a constant $C$ depending on
$\kappa$, $C_R$, $\theta$,
$\beta$, $\alpha$,
$\underline{\ddchis}$, $\overline{\ddchis}$,
$\epsilon$.
\end{theorem}
Both cases vanish,
so $\Xi(s) \ll n^{-1/2}$
for every $s$ that is at most polynomial in $n$.

\begin{lemma}[Offset complexity: bounded curvature]
\label{lem:xi-bounded}
Suppose $\ddchis$ is bounded
and $|h| \le M$ for all $h \in B_\rho$.
Suppose the Bregman lower bound
(\Cref{thm:bregman-lower},
via \Cref{lem:two-pile})
gives
\begin{equation}
\label{eq:bregman-shell-bounded}
\hP\, D_{\chis}(\tilde\phi + h, \tilde\phi)
  \ge (1-\epsilon)\, \frac{r^2}{2}
\end{equation}
uniformly over $h$
with $\rho(h) \le s$ and $\norm{h}_{\chis} = r$,
with the junction tolerance included
in $\epsilon$.
For any $r_\star \le s\, u_0$ satisfying
\begin{equation}
\label{eq:critical-bounded}
\frac{r_\star^2}{2}
  \ge \frac{1}{1-\epsilon}\,
  \sup_{\substack{\rho(h) \le s \\
      \norm{h}_{\chis} \le r_\star}}
    (\hP - \P)\qty{\dotpsiZ(h) - (\dchis \comp \tilde\phi)\,h},
\end{equation}
the offset complexity satisfies
$\Xi(s) \le (1-\epsilon)\, r_\star^2/(2s)$.
\end{lemma}

\begin{proof}
The change of variables $h \leftarrow sh$
in the definition~\eqref{eq:offset-complexity}
gives
\begin{equation}
\label{eq:offset-rescaled}
s\,\Xi(s)
  = \sup_{\rho(h) \le s}
  \qty[(\hP - \P)\qty{\dotpsiZ(h) - (\dchis \comp \tilde\phi)\, h}
    - \hP\, D_{\chisZ}(\tilde\phi + h,\, \tilde\phi)].
\end{equation}
Write $\overline{M}(r)$ for the local empirical process
$\sup_{\substack{\rho(h) \le s,\,
    \norm{h}_{\chis} \le r}}
  (\hP - \P)\{\dotpsiZ(h) - (\dchis \comp \tilde\phi)\, h\}$.
Split the supremum
in~\eqref{eq:offset-rescaled}
at the critical radius $r_\star$.

\paragraph{Below $r_\star$.}
The Bregman divergence is nonnegative, so
\[
\sup_{\substack{\rho(h) \le s \\ \norm{h}_{\chis} \le r_\star}}
  \qty[(\hP - \P)\{\dotpsiZ(h) - (\dchis \comp \tilde\phi)\, h\}
    - \hP\, D_{\chisZ}(\tilde\phi + h,\, \tilde\phi)]
  \le \overline{M}(r_\star).
\]

\paragraph{Above $r_\star$.}
For $h$ with $\rho(h) \le s$
and $\norm{h}_{\chis} = r \in (r_\star, s\, u_0]$,
the hypothesis~\eqref{eq:bregman-shell-bounded}
gives a Bregman penalty of $(1-\epsilon)\, r^2/2$.
The ratio
$r^2 / \overline{M}(r)$
is non-decreasing
because $\overline{M}$ is sublinear.
To see this: $\overline{M}(r)$
is the supremum of a linear functional
over $\{\rho(h) \le s,\, \norm{h}_{\chis} \le r\}$.
This constraint set scales:
if $h$ is feasible at level $r$,
then $th$ is feasible at level $tr$
for $t \le 1$
(because $\rho(th) = t\,\rho(h) \le t\, s \le s$).
Linearity of the functional gives
$\overline{M}(tr) \le t\, \overline{M}(r)$,
so $\overline{M}(r)/r$ is non-increasing.
Because the ratio is non-decreasing
and the critical condition~\eqref{eq:critical-bounded}
says it exceeds $2/(1-\epsilon)$ at $r_\star$,
it exceeds this threshold for all $r \ge r_\star$:
$(1-\epsilon)\, r^2/2 \ge \overline{M}(r)$.
The offset is therefore non-positive:
\[
(\hP - \P)\{\dotpsiZ(h) - (\dchis \comp \tilde\phi)\, h\}
  - \hP\, D_{\chisZ}(\tilde\phi + h,\, \tilde\phi)
  \le \overline{M}(r) - (1-\epsilon)\, \frac{r^2}{2}
  \le 0.
\]

\paragraph{Combining.}
Below $r_\star$, the contribution is at most $\overline{M}(r_\star)$.
Above $r_\star$, the Bregman penalty dominates
and the contribution is non-positive.
Therefore $s\,\Xi(s) \le \overline{M}(r_\star)
  \le (1-\epsilon)\, r_\star^2/2$
by the critical condition~\eqref{eq:critical-bounded},
giving $\Xi(s) \le (1-\epsilon)\, r_\star^2/(2s)$.
\end{proof}

For entropy-like dispersions,
the curvature $\ddchis$ varies with $\phi$,
and evaluating it at the displaced point
$\tilde\phi + h$ rather than $\tilde\phi$
introduces a multiplicative discount $e^{-B_s(r)}$
in the Bregman lower bound~\eqref{eq:psc-discount-bound}.
The critical radius argument still applies,
but the monotonicity of the penalty-to-process ratio
requires a growth condition on $B_s$.

\begin{lemma}[Offset complexity: PSC]
\label{lem:xi-psc}
Let $\chis$ be pseudo-self-concordant
with constant $\beta$:
$|\dddchis| \le \beta\, \ddchis$,
and suppose $\ddchis$ is convex.
Suppose the interpolation inequality~\eqref{eq:interpolation}
holds with exponent $\theta \in (0,1)$ and constant $\kappa$.
Write $B_s(r) = \beta\kappa\, s^\theta\, r^{1-\theta}/3$
for the PSC discount~\eqref{eq:psc-discount}.
Suppose the Bregman lower bound
(\Cref{thm:bregman-lower},
via \Cref{cor:ulb-entropy})
gives
\begin{equation}
\label{eq:bregman-shell-psc}
\hP\, D_{\chis}(\tilde\phi + h, \tilde\phi)
  \ge (1-\epsilon)\, \frac{1}{2}\, e^{-B_s(r)}\, r^2
\end{equation}
uniformly over $h$
with $\rho(h) \le s$ and $\norm{h}_{\chis} = r$.
For any $r_\star \le s\, u_0$ satisfying
\begin{equation}
\label{eq:critical-psc}
\begin{aligned}
e^{-B_s(r_\star)}\, \frac{r_\star^2}{2}
  &\ge \frac{1}{1-\epsilon}\,
  \sup_{\substack{\rho(h) \le s \\
      \norm{h}_{\chis} \le r_\star}}
    (\hP - \P)\qty{\dotpsiZ(h) - (\dchis \comp \tilde\phi)\,h}, \\[4pt]
(1-\theta)\, B_s(s\, u_0) &\le 1,
\end{aligned}
\end{equation}
the offset complexity satisfies
$\Xi(s) \le (1-\epsilon)\, e^{-B_s(r_\star)}\, r_\star^2/(2s)$.
\end{lemma}

\begin{proof}
Write $\overline{M}(r)$ for the local empirical process
appearing in~\eqref{eq:critical-psc}.
The argument follows \Cref{lem:xi-bounded},
with the PSC discount modifying the monotonicity step.
As in~\eqref{eq:offset-rescaled},
$s\,\Xi(s) = \sup_{\rho(h) \le s}
  [(\hP - \P)\{\dotpsiZ(h) - (\dchis \comp \tilde\phi)\, h\}
    - \hP\, D_{\chisZ}(\tilde\phi + h,\, \tilde\phi)]$.
Below $r_\star$, the bound is identical:
the Bregman divergence is nonnegative,
so the contribution is at most $\overline{M}(r_\star)$.

\paragraph{Above $r_\star$.}
For $h$ with $\rho(h) \le s$
and $\norm{h}_{\chis} = r \in (r_\star, s\, u_0]$,
the hypothesis~\eqref{eq:bregman-shell-psc}
gives a Bregman penalty of
$(1-\epsilon)\, e^{-B_s(r)}\, r^2/2$.
The ratio
\[
\frac{e^{-B_s(r)}\, r^2}{\overline{M}(r)}
  = e^{-B_s(r)}\, r
    \cdot \frac{r}{\overline{M}(r)}
\]
is the product of two non-decreasing functions.
The second factor $r / \overline{M}(r)$
is non-decreasing because $\overline{M}$
is sublinear
(as in \Cref{lem:xi-bounded}).
The first factor $e^{-B_s(r)}\, r$
is non-decreasing when $(1-\theta)\, B_s(r) \le 1$,
by the monotonicity bound~\eqref{eq:psc-monotonicity}.
The condition $(1-\theta)\, B_s(s\, u_0) \le 1$
ensures this holds
on the full range $(r_\star, s\, u_0]$.

Because the ratio is non-decreasing
and the critical condition~\eqref{eq:critical-psc}
says it exceeds $2/(1-\epsilon)$ at $r_\star$,
it exceeds this threshold for all $r \in [r_\star, s\, u_0]$:
\[
(1-\epsilon)\, \frac{1}{2}\, e^{-B_s(r)}\, r^2
  \ge \overline{M}(r).
\]
The offset is therefore non-positive on each shell:
\[
(\hP - \P)\{\dotpsiZ(h) - (\dchis \comp \tilde\phi)\, h\}
  - \hP\, D_{\chisZ}(\tilde\phi + h,\, \tilde\phi)
  \le \overline{M}(r) - (1-\epsilon)\, \frac{1}{2}\, e^{-B_s(r)}\, r^2
  \le 0.
\]

\paragraph{Combining.}
Below $r_\star$, the contribution is at most $\overline{M}(r_\star)$.
Above $r_\star$, the Bregman penalty dominates
and the contribution is non-positive.
Therefore $s\,\Xi(s) \le \overline{M}(r_\star)
  \le (1-\epsilon)\, e^{-B_s(r_\star)}\, r_\star^2/2$
by the critical condition~\eqref{eq:critical-psc},
giving $\Xi(s) \le (1-\epsilon)\, e^{-B_s(r_\star)}\, r_\star^2/(2s)$.
\end{proof}

\begin{remark}
\Cref{lem:xi-bounded} is the special case $B_s \equiv 0$
of \Cref{lem:xi-psc},
where the growth condition is vacuous
and the monotonicity is immediate from sublinearity.
\end{remark}

\begin{lemma}[Offset complexity: exponential curvature]
\label{lem:xi-jensen}
Suppose $\chis$ satisfies
a curvature growth condition:
$\dddchis \ge \alpha\, \ddchis$
for some $\alpha > 0$.
Suppose the interpolation inequality~\eqref{eq:interpolation}
holds with exponent $\theta \in (0,1)$ and constant $\kappa$,
the Bregman lower bound
(\Cref{thm:bregman-lower})
gives
$\hP\, D_{\chis}(\tilde\phi + h, \tilde\phi)
  \ge (1-\epsilon)\, \P\, D_{\chis}(\tilde\phi + h, \tilde\phi)$
uniformly over $\rho(h) \le s$,
and $\underline{\ddchis} \le \ddchis(\tilde\phi) \le \overline{\ddchis}$
under~$\P$.
If
\[
\sup_{\substack{\rho(h) \le s \\
    \norm{h}_{\chis} \le r}}
  (\hP - \P)\qty{\dotpsiZ(h) - (\dchis \comp \tilde\phi)\,h}
  \le C_R\, n^{-1/2}\, s^{1-\theta}\, r^\theta
\]
for all $r \le s\, u_0$,
and $\alpha\, s / \log n \to \infty$,
\begin{equation}
\label{eq:xi-jensen-bound}
\Xi(s) \le C\, n^{-1/2}\,
    \qty(\frac{\log n}{\alpha\, s})^{\theta/(1+\theta)}
\end{equation}
for a constant $C$ depending on
$\kappa$, $C_R$, $\theta$, $\alpha$,
$\underline{\ddchis}$, $\overline{\ddchis}$,
$\epsilon$.
\end{lemma}

\begin{proof}
As in~\eqref{eq:offset-rescaled},
$s\,\Xi(s) = \sup_{\rho(h) \le s}
  [(\hP - \P)\{\dotpsiZ(h) - (\dchis \comp \tilde\phi)\, h\}
    - \hP\, D_{\chisZ}(\tilde\phi + h,\, \tilde\phi)]$.
Write $\overline{M}(r)$ for the local empirical process
at radius $r$, as in \Cref{lem:xi-psc}.
Split by shells of $\norm{h}_{\chis}$.
By kernel balance~\eqref{eq:kernel-balance},
the functions contributing to $\Xi(s)$
lie in $V = (\ker\rho)^{\perp}$
and satisfy $\P h = 0$.

\paragraph{Below $r_\star$.}
The Bregman divergence is nonnegative,
so the contribution is at most $\overline{M}(r_\star)$.

\paragraph{Above $r_\star$: critical radius.}
For $h \in V$ with $\rho(h) \le s$,
$\norm{h}_{\chis} = r \ge r_\star$,
$\P h = 0$,
and $\alpha\, c_1\, r^{1+\theta}/s^\theta \ge 2$
(which holds for $n$ large enough
because $\alpha\, s / \log n \to \infty$
and $r_\star^{1+\theta} \ge c\, s^\theta \log n / \alpha$
by the critical equation below),
\Cref{lem:jensen-lower}(ii)
and the Bregman lower bound give
\[
\hP\, D_{\chis}(\tilde\phi + h, \tilde\phi)
  \ge (1-\epsilon)\, \frac{c_0}{\alpha^2}\,
    \frac{r^{2\theta}}{s^{2\theta}}\,
    e^{\alpha\, c_1\, r^{1+\theta}/s^\theta}.
\]
The penalty-to-process ratio
\[
\frac{r^{2\theta}\,
  e^{\alpha\, c_1\, r^{1+\theta}/s^\theta}}
  {\overline{M}(r)}
\]
is strictly increasing in $r$:
the log-derivative is
$2\theta/r
  + \alpha\, c_1\, (1+\theta)\, r^\theta/s^\theta
  - d\log\overline{M}(r)/dr$,
and the first two terms are positive
while sublinearity of $\overline{M}$
gives $d\log\overline{M}/dr \le 1/r$.

Define $r_\star$ as the critical radius
where the penalty first matches the process:
\[
(1-\epsilon)\, \frac{c_0}{\alpha^2}\,
    \frac{r_\star^{2\theta}}{s^{2\theta}}\,
    e^{\alpha\, c_1\, r_\star^{1+\theta}/s^\theta}
  = C_R\, n^{-1/2}\, s^{1-\theta}\, r_\star^\theta.
\]
Because the ratio is strictly increasing,
the penalty dominates $\overline{M}(r)$
for all $r \ge r_\star$,
making the offset non-positive on each shell.

\paragraph{Rate of $r_\star$.}
The critical radius equation gives
\[
e^{\alpha\, c_1\, r_\star^{1+\theta}/s^\theta}
  = \frac{C_R\, \alpha^2\, s}
    {(1-\epsilon)\, c_0\, n^{1/2}\, r_\star^\theta}.
\]
Taking logarithms,
$\alpha\, c_1\, r_\star^{1+\theta}/s^\theta
  \le C''\, \log n$,
so $r_\star^{1+\theta}
  \le C''\, s^\theta\, \log n / \alpha$,
giving
\[
\Xi(s)
  \le \overline{M}(r_\star)/s
  \le C\, n^{-1/2}\,
    \qty(\frac{\log n}{\alpha\, s})^{\theta/(1+\theta)}.
\qedhere
\]
\end{proof}

\begin{proof}[Proof of \Cref{thm:xi-all-regimes}]
As in~\eqref{eq:offset-rescaled},
$s\,\Xi(s) = \sup_{\rho(h) \le s}
  [(\hP - \P)\{\dotpsiZ(h) - (\dchis \comp \tilde\phi)\, h\}
    - \hP\, D_{\chisZ}(\tilde\phi + h,\, \tilde\phi)]$.
Write $\overline{M}(r)$ for the local empirical process
at radius $r$, as in \Cref{lem:xi-psc}.
Three regimes, determined by whether
the PSC growth condition
$(1-\theta)\, B_s(s\, u_0) \le 1$ holds.
For small $s$, it holds globally
and \Cref{lem:xi-psc} applies directly.
For intermediate $s$, it fails at $s\, u_0$
but holds up to a threshold $r^\star$;
we split into three sub-shells
(empirical process below $r_\star$,
discounted Bregman penalty between $r_\star$ and $r^\star$,
Jensen quadratic floor above $r^\star$).
For large $s$,
the exponential Bregman growth of \Cref{lem:xi-jensen}
dominates.

\paragraph{Small $s$: $s \le s_0$.}
Define $s_0 = 3/\{\beta(1-\theta)\kappa\, u_0^{1-\theta}\}$.
Then $(1-\theta)\, B_s(s\, u_0)
  = (1-\theta)\beta\kappa\, s\, u_0^{1-\theta}/3 \le 1$,
so the growth condition
of \Cref{lem:xi-psc} holds.
The critical radius $r_\star$ satisfies
$e^{-B_s(r_\star)}\, r_\star^2/2
  \ge C_R\, n^{-1/2}\, s^{1-\theta}\, r_\star^\theta / (1-\epsilon)$,
giving $r_\star^{2-\theta}
  \ge c\, n^{-1/2}\, s^{1-\theta}$
(because $e^{-B_s(r_\star)} \ge e^{-1/(1-\theta)}$).
Therefore
$\Xi(s) \le \overline{M}(r_\star)/s
  \le C_R\, n^{-1/2}\, s^{-\theta}\, r_\star^\theta
  \le C\, n^{-1/2}\, (n^{1/2}\, s)^{-\theta/(2-\theta)}$.

\paragraph{Intermediate $s$:
  $s_0 < s \le C_1\, n^{(1-\theta)/(2+8\theta)}$.}
The growth condition $(1-\theta)\, B_s(s\, u_0) \le 1$
fails at $s\, u_0$ but holds at
$r^\star
  = \{3/(\beta(1-\theta)\kappa\, s^\theta)\}^{1/(1-\theta)}$,
where $(1-\theta)\, B_s(r^\star) = 1$.
Split the supremum by three shells.

\paragraph{Below $r_\star$.}
The Bregman divergence is nonnegative,
so the contribution is at most $\overline{M}(r_\star)$.

\paragraph{Between $r_\star$ and $r^\star$.}
On this range, $(1-\theta)\, B_s(r) \le 1$,
so the map $r \mapsto e^{-B_s(r)}\, r$
is non-decreasing
by the monotonicity bound~\eqref{eq:psc-monotonicity}.
The Bregman penalty
$(1-\epsilon)\, e^{-B_s(r)}\, r^2/2$
therefore grows faster than $\overline{M}(r)$
(which is sublinear),
and since the penalty dominates at $r_\star$
by the critical condition~\eqref{eq:critical-psc},
it dominates on the entire range.
The offset is non-positive.

\paragraph{Above $r^\star$.}
Kernel balance gives $\P h = 0$
for contributing $h$.
\Cref{lem:jensen-lower}(i)
gives the quadratic floor:
\[
\P\, D_{\chis}(\tilde\phi + h, \tilde\phi)
  \ge c_J\, \frac{r^{2+4\theta}}{s^{4\theta}}.
\]
The Bregman lower bound transfers this
to the empirical measure:
$\hP\, D_{\chis}(\tilde\phi + h, \tilde\phi)
  \ge (1-\epsilon)\, c_J\, r^{2+4\theta}/s^{4\theta}$.
This dominates
$\overline{M}(r) \le C_R\, n^{-1/2}\, s^{1-\theta}\, r^\theta$
when
\[
r^{2+3\theta}
  \ge \frac{C_R}{(1-\epsilon)\, c_J}\, n^{-1/2}\, s^{1+3\theta}.
\]
Because $r^{2+3\theta}$ is increasing in $r$,
the inequality holds for all $r \ge r^\star$
once it holds at $r = r^\star$.
Because $r^\star
  = \{3/(\beta(1-\theta)\kappa\, s^\theta)\}^{1/(1-\theta)}$,
\[
r^{\star\,2+3\theta}
  = C'\, s^{-\theta(2+3\theta)/(1-\theta)},
\]
and the domination condition becomes
$s^{-(1+4\theta)/(1-\theta)} \ge C\, n^{-1/2}$,
which holds when
$s \le C_1\, n^{(1-\theta)/(2+8\theta)}$.
The combined bound is
$s\,\Xi(s) \le \overline{M}(r_\star)$,
giving the same rate as the small-$s$ case.

\paragraph{Large $s$:
  $s > C_1\, n^{(1-\theta)/(2+8\theta)}$.}
Then $\alpha\, s / \log n \to \infty$
and \Cref{lem:xi-jensen}
gives~\eqref{eq:xi-jensen-bound}.
\end{proof}

\begin{remark}[Explicit constants]
\label{rmk:explicit-constants}
The constant $C$ in~\eqref{eq:xi-all-regimes}
can be computed from the proof.
In the first case
($s \le C_1\, n^{(1-\theta)/(2+8\theta)}$),
the PSC peeling argument gives
\[
C = \qty(\frac{2\, e^{1/(1-\theta)}}{1-\epsilon})^{\theta/(2-\theta)}
  \cdot C_R^{2/(2-\theta)}.
\]
The curvature parameters
$\beta$, $\kappa$, $\underline{\ddchis}$, $\overline{\ddchis}$
determine only the regime boundary~$C_1$,
not the prefactor.
In the second case,
$C$ additionally depends on the Jensen floor constant
$c_J = \underline{\ddchis}^3
  / (16\overline{\ddchis}^5\, \kappa^4)$
from~\eqref{eq:jensen-bregman-lower}
and a logarithmic correction
from the implicit critical equation
in the proof of \Cref{lem:xi-jensen}.
\end{remark}

\subsection{From offset complexity to bias}
\label{sec:offset-to-bias}

\begin{lemma}[Bias via offset complexity]
\label{lem:offset-convex}
For any $r > 2\rho(\tilde\phi)$,
\begin{equation}
\label{eq:bias-local-complexity}
|b(\hgamma)|
  \le \tfrac{5}{2}\,\Xi(r)
  + 2\dot\zetas\!\qty{\rho(\tilde\phi)}\,
  \qty(1 - \frac{\rho(\tilde\phi)}{r}).
\end{equation}
\end{lemma}

\subsubsection{Setting-specific corollaries}

\setupvariants

\theoremvariant{Q}
\begin{corollary}[Setting Q: positive-part bias bound]
\label{cor:offset-positive-part}
For $\chis(\phi) = \phi_+^2/2$,
suppose the Bregman lower bound
(\Cref{thm:bregman-lower},
via \Cref{lem:two-pile})
holds for the increment class at $\tilde\phi$
with tolerance $\epsilon$,
and
\[
\sup_{\substack{h \in B_\rho \\
    \norm{h}_{\chis} \le u}}
  (\hP - \P)\qty{\dotpsiZ(h) - (\dchis \comp \tilde\phi)\,h}
  \le C_R\, n^{-1/2}\, u^\theta.
\]
Then for any $r > 2\rho(\tilde\phi)$,
\begin{equation}
\label{eq:offset-positive-part}
|b(\hgamma)|
  \le C\, n^{-1/2}\,
  (n^{1/2}\, r)^{-\theta/(2-\theta)}
  + 2\dot\zetas\!\qty{\rho(\tilde\phi)}\,
  \qty(1 - \frac{\rho(\tilde\phi)}{r})
\end{equation}
for a constant $C$ depending on
$C_R$, $\theta$, and $\epsilon$.
\end{corollary}
\resettheorem

\theoremvariant{E}
\begin{corollary}[Setting E: unconditional bias bound]
\label{cor:unconditional-bias}
Suppose $\alpha\, \ddchis \le \dddchis \le \beta\, \ddchis$
for some $0 < \alpha \le \beta$.
Suppose the interpolation inequality~\eqref{eq:interpolation}
holds with exponent $\theta \in (0,1)$,
the Bregman lower bound
(\Cref{thm:bregman-lower})
holds for the increment class at $\tilde\phi$,
$\underline{\ddchis} \le \ddchis(\tilde\phi) \le \overline{\ddchis}$ under~$\P$,
and
\[
\sup_{\substack{h \in B_\rho \\
    \norm{h}_{\chis} \le u}}
  (\hP - \P)\qty{\dotpsiZ(h) - (\dchis \comp \tilde\phi)\,h}
  \le C_R\, n^{-1/2}\, u^\theta.
\]
Then for any $r > 2\rho(\tilde\phi)$,
\begin{equation}
\label{eq:unconditional-bias}
\begin{aligned}
|b(\hgamma)|
  &\le C\, n^{-1/2}\,
  \begin{cases}
    (n^{1/2}\, r)^{-\theta/(2-\theta)}
      & \text{if } r \le n^{(1-\theta)/(2+8\theta)}, \\[6pt]
    \qty(\dfrac{\log n}{\alpha\, r})^{\theta/(1+\theta)}
      & \text{if } r > n^{(1-\theta)/(2+8\theta)},
  \end{cases} \\
  &\quad + 2\dot\zetas\!\qty{\rho(\tilde\phi)}\,
  \qty(1 - \frac{\rho(\tilde\phi)}{r})
\end{aligned}
\end{equation}
for a constant $C$ depending on
$\kappa$, $C_R$, $\theta$,
$\beta$, $\alpha$, $\underline{\ddchis}$, $\overline{\ddchis}$,
and $\epsilon$.
Both cases vanish for any $r \ge 1$,
so $|b(\hgamma)| \ll n^{-1/2}$
without any curvature condition.
\end{corollary}
\resettheorem

\subsubsection*{Proofs}

\begin{proof}[Proof of Lemma~\ref{lem:offset-convex}]
The proof has three steps:
restrict the bias to a complement
of the kernel,
identify the correct modulus
for the restricted problem,
and reduce to offset complexity.

\paragraph{Kernel balance.}
Kernel directions are unconstrained under $\rho$,
so the first order condition~\eqref{eq:kernel-balance}
forces exact balance on $K = \ker(\rho)$:
$\hP\hgamma(Z)\,k(Z) = \hP\dotpsiZ(k)$
for all $k \in K$.
Writing $m = f + k$ with $f \in V$
and $k \in K$,
the seminorm depends only
on the $V$-component
($\rho(m) = \rho(f)$)
and kernel balance eliminates $k$
from the bias:
\[
b(\hgamma)
  = \sup_{f \in B_\rho \cap V}
    \qty{\hP\dotpsiZ(f) - \hP\hgamma(Z)\,f(Z)}.
\]

\paragraph{The $V$-restricted modulus.}
The goal is to express the bias
in terms of the offset complexity $\Xi$,
which measures how much the empirical process
can push the optimizer away from~$\tilde\phi$
net of the Bregman divergence penalty.
The difficulty is that the bias is defined
on the full seminorm ball $B_\rho$,
while the offset complexity is defined
on the $\chis$-norm shells
that the Bregman divergence naturally controls.
Kernel balance bridges these:
it eliminates the kernel directions
(which $\rho$ ignores and $D_{\chis}$ cannot control),
reducing the bias to an optimization over $V$,
where both the seminorm and the $\chis$-norm
give meaningful constraints.
The resulting $V$-restricted offset complexity $\xi_V$
uses the full centering point $\tilde\phi$
(not just $\tilde\phi_V$)
in the Bregman divergence,
because the frozen kernel component $\tilde\phi_K$
remains present in the loss.
The reduction to $\Xi$ at the end of this proof
then relaxes $\xi_V$ to the full seminorm ball,
losing only a factor of $5/4$.

Concretely:
the optimizer
$\hat\phi = \hat\phi_V + \hat\phi_K$
has its kernel component $\hat\phi_K$
pinned by the balance condition.
What remains is an optimization
over $\phi_V \in V$
with the kernel component frozen:
\[
\min_{\phi_V \in V}\;
  \P\,\chis(\phi_V + \tilde\phi_K)
  - \psi(\phi_V)
  + \zetas\!\qty{\rho(\phi_V)},
\]
where $\tilde\phi_K$ is the kernel component
of the population minimizer $\tilde\phi$.
The modulus of continuity
for this restricted problem is
\[
\omega_V(s)
  = \sup_{g \in B_\rho \cap V}
    \qty{\psi(g) - s\,\chis(\tilde\phi_K + g/s)},
\]
in which $\chis$ is evaluated
at $\tilde\phi_K + g/s$,
not at $g/s$ alone,
because the kernel component remains present.
Center at $\tilde\phi_V$
by writing $g = s\tilde\phi_V + f$.
Then the argument of $\chis$ becomes
$\tilde\phi_K + \tilde\phi_V + f/s
  = \tilde\phi + f/s$:
centering recombines the kernel
and $V$ components
into the full population minimizer.
The Bregman expansion at $\tilde\phi$
(Step~2 of the proof of
\Cref{lem:bias-bound})
gives the offset complexity,
with the full $\tilde\phi$
in both arguments of the Bregman divergence:
\[
\xi_V(s)
  = \sup_{f \in (B_\rho \cap V) - s\tilde\phi_V}
  \qty{M(f)
    - s\,D_{\chis}(\tilde\phi + f/s,\, \tilde\phi)}.
\]
This offset complexity is non-negative
($f = 0$ gives
$M(0) - s\,D_{\chis}(\tilde\phi, \tilde\phi) = 0$)
and concave
(the perspective argument
from Step~3 of the same proof),
so the bias bound~\eqref{eq:bias-local-complexity} gives
$|b(\hgamma)|
  \le 2\,\xi_V(t)
  + 2\dot\zetas\{\rho(\tilde\phi)\}\,(1 - t\,\rho(\tilde\phi))$.

\paragraph{Reduction to offset complexity.}
Set $t = 1/r$
for the given $r > 2\rho(\tilde\phi)$,
so that $2t\,\rho(\tilde\phi) < 1$.
On the shifted set
$S := (B_\rho \cap V) - t\tilde\phi_V$,
the triangle inequality gives
$\rho(f) \le 1 + t\,\rho(\tilde\phi) < 5/4$
(because $t\,\rho(\tilde\phi) < 1/2$).
Writing $f = 5g/4$
with $g = 4f/5 \in B_\rho$,
the empirical process is linear:
$M(5g/4) = (5/4)\, M(g)$.
For the Bregman offset,
convexity of $a \mapsto D_{\chis}(a, \tilde\phi)$
with $D_{\chis}(\tilde\phi, \tilde\phi) = 0$
gives
$D_{\chis}(\tilde\phi + 5g/4,\, \tilde\phi)
  \ge (5/4)\, D_{\chis}(\tilde\phi + g,\, \tilde\phi)$
pointwise,
so
\begin{equation}
\label{eq:xi-V-to-Xi}
\xi_V(t)
  \le \tfrac{5}{4}\, \sup_{g \in B_\rho}
    \qty{M(g) - t\, \hP\, D_{\chis}(\tilde\phi + g/t,\, \tilde\phi)}
  = \tfrac{5}{4}\, \Xi(r).
\end{equation}
The bias bound~\eqref{eq:bias-local-complexity} follows.
\end{proof}

\begin{proof}[Proof of Corollary~\ref{cor:offset-positive-part}]
By \Cref{lem:offset-convex} at scale $r$,
$|b(\hgamma)|
  \le \tfrac{5}{2}\,\Xi(r)
  + 2\dot\zetas\{\rho(\tilde\phi)\}$.
The positive-part dispersion has
$\ddchis(\phi) = \ind{\phi \ge 0} \le 1$,
so \Cref{lem:xi-bounded}
applies.
The critical condition~\eqref{eq:critical-bounded}
at scale $s = r$
gives a critical curvature radius $r_\star$ satisfying
$r_\star^{2-\theta}
  \ge c\, n^{-1/2}\, r^{1-\theta} / (1-\epsilon)$,
so
$\Xi(r)
  \le (1-\epsilon)\, r_\star^2/(2r)
  \le C\, n^{-1/2}\, (n^{1/2}\, r)^{-\theta/(2-\theta)}$.
\end{proof}

\begin{proof}[Proof of Corollary~\ref{cor:unconditional-bias}]
By \Cref{lem:offset-convex} at scale $r$,
$|b(\hgamma)|
  \le \tfrac{5}{2}\, \Xi(r)
  + 2\dot\zetas\{\rho(\tilde\phi)\}$.
We apply
\Cref{thm:xi-all-regimes}.
In the first case,
$r \le C_1\, n^{(1-\theta)/(2+8\theta)}$
and
\[
\Xi(r)
  \le C\, n^{-1/2}\, (n^{1/2}\, r)^{-\theta/(2-\theta)}.
\]
In the second case,
$r > C_1\, n^{(1-\theta)/(2+8\theta)}$
and
\[
\Xi(r)
  \le C\, n^{-1/2}\,
  \qty(\frac{\log n}{\alpha\, r})^{\theta/(1+\theta)}.
\]
Multiplying by $5/2$
gives~\eqref{eq:unconditional-bias}.
\end{proof}

\subsubsection{Penalty bias}
\label{sec:penalty-bias}

The offset complexity term
in~\eqref{eq:bias-local-complexity}
is $o(n^{-1/2})$
whenever $\tilde\rho \to \infty$
--- the corollaries above give the rates.
The second term,
the penalty bias $2\dot\zetas\{\rho(\tilde\phi)\}$,
is not automatically negligible:
whether it is $o(n^{-1/2})$ depends on
how the seminorm $\rho(\tilde\phi)$
of the population minimizer grows
as regularization decreases.
For quadratic regularization
$\zetas(\rho) = \eta\rho^2$,
the penalty bias is
\[
2\dot\zetas\!\qty{\rho(\tilde\phi)}
  = 4\eta\,\rho(\tilde\phi).
\]

Everything reduces to bounding
$\rho(\tilde\phi)$ as a function of~$\eta$.
There are two routes
--- the comparison bound and the balance equation ---
with different costs and different scopes.
The comparison bound works for any dispersion
but requires regularity of the unpenalized limit;
the balance equation works only
for the quadratic family
but requires only integrability
of the Riesz representer.
These are conditions on different axes
(smoothness vs.\ integrability),
and which is more natural depends on the problem.

\paragraph{The comparison bound.}
The unpenalized population minimizer
$\phi$ satisfies
$\dchis(\phi) = \griesz$
--- the inverse link applied
to the Riesz representer.
For the quadratic family,
$\phi = \griesz$:
the index \emph{is} the weights.
For entropy-like dispersions,
$\phi = \log(\griesz)/\tau$:
regularity of the index is regularity
of the log-weights,
not the weights themselves.

Because $\tilde\phi$ minimizes
the population loss,
its objective value is at most
that of any competitor~$f$:
\[
\P\,D_{\chis}(\tilde\phi, \phi)
  + \eta\,\rho^2(\tilde\phi)
  \le \P\,D_{\chis}(f, \phi)
  + \eta\,\rho^2(f).
\]
Dropping the nonnegative
Bregman term on the left
gives
\begin{equation}
\label{eq:comparison-bound}
\eta\,\rho^2(\tilde\phi)
  \le \P\,D_{\chis}(f, \phi)
  + \eta\,\rho^2(f).
\end{equation}

When $\phi$ itself has finite seminorm
--- $\rho(\phi) < \infty$ ---
we take $f = \phi$
and the Bregman term vanishes:
$\rho(\tilde\phi) \le \rho(\phi)$.
The seminorm of the regularized projection
is frozen at the seminorm
of the unpenalized limit,
independently of~$\eta$.
The penalty bias becomes
$4\eta\,\rho(\phi)$,
which is $o(n^{-1/2})$
whenever $\eta = o(n^{-1/2})$
--- condition~\ref{cond:an-A} imposes only
that $\eta$ decrease
at rate $n^{-1/2}$,
with no smoothness threshold.

When $\rho(\phi) = \infty$,
we need a competitor~$f$
that trades approximation quality
--- how small $\P\,D_{\chis}(f, \phi)$ is ---
against penalty cost
--- how large $\rho(f)$ is.
Mollification provides this tradeoff.
The penalty side is the same
for all dispersions;
the approximation side depends on
the Bregman divergence of~$\chis$.

\begin{proposition}[Mollification penalty bias, Settings~Q and~Q$_+$]
\label{prop:penalty-bias-Q}
Suppose the penalty seminorm embeds
into a $W^{s,p}$ Sobolev seminorm:
$\rho(f) \le C_\rho\,|f|_{W^{s,p}}$.
Suppose $\phi \in L_p$ and
$\phi \in W^{t,q}$
for some $t \in (0, s)$ and $q \ge 1$.
Call $t_q = t - d(1/q - 1/2)_+$
the effective $L_2$ regularity.
Then in Settings~Q and~Q$_+$
($\ddchis \le 1$),
the penalty bias satisfies
\begin{equation}
\label{eq:mollification-penalty-bias}
\eta\,\rho(\tilde\phi)
  \le c_{s,t,d}\,
  |\phi|_{W^{t,q}}^{s/(s+t_q)}\,
  (C_\rho\,\|\phi\|_{L_p})^{t_q/(s+t_q)}\,
  \eta^{(s + 2t_q)/(2(s + t_q))},
\end{equation}
where $c_{s,t,d}$ depends only on $s$, $t_q$, and $d$.
\end{proposition}

The bound~\eqref{eq:mollification-penalty-bias}
is $o(n^{-1/2})$ when
$\eta \ll n^{-(s + t_q)/(s + 2t_q)}$.

\begin{proof}
The argument has three steps:
bound the Bregman divergence,
bound the penalty of the mollifier,
and optimize the bandwidth.

\paragraph{Bregman divergence.}
The mollification competitor
$\phi_\epsilon = \phi * \omega_\epsilon$
satisfies
$D_{\chis}(\phi_\epsilon, \phi)
  \le (\phi_\epsilon - \phi)^2/2$
because $\ddchis \le 1$
in Settings~Q and~Q$_+$.
Integrating against $\P$
and applying the $L_2$ mollification estimate
for $\phi \in W^{t,q}$ gives
\begin{equation}
\label{eq:bregman-mollification-Q}
\P\,D_{\chis}(\phi_\epsilon, \phi)
  \le \tfrac{1}{2}\,
  \|\phi_\epsilon - \phi\|_{L_2(\P)}^2
  \le \tfrac{1}{2}\,\epsilon^{2t_q}\,
    |\phi|_{W^{t,q}}^2,
\end{equation}
where $t_q = t - d(1/q - 1/2)_+$:
for $q \ge 2$ the embedding
into $L_2$ is free and $t_q = t$;
for $q < 2$ it costs
$d(1/q - 1/2)$ orders of smoothness,
but the space $W^{t,q}$ is larger.

\paragraph{Penalty of the mollifier.}
Because the seminorm embeds
into $W^{s,p}$,
\begin{equation}
\label{eq:penalty-mollification}
\rho(\phi_\epsilon)
  \le C_\rho\,|\phi_\epsilon|_{W^{s,p}}
  \le C_\rho\,\epsilon^{-s}\,
  \|\phi\|_{L_p}.
\end{equation}
This requires $\phi \in L_p$
--- the entry ticket
for the comparison bound.

\paragraph{Optimizing the bandwidth.}
Substituting~\eqref{eq:bregman-mollification-Q}
and~\eqref{eq:penalty-mollification}
into the comparison bound~\eqref{eq:comparison-bound}
gives
\[
\eta\,\rho^2(\tilde\phi)
  \le \tfrac{1}{2}\,\epsilon^{2t_q}\,
    |\phi|_{W^{t,q}}^2
  + C_\rho^2\,\eta\,\epsilon^{-2s}\,
    \|\phi\|_{L_p}^2.
\]
Choosing $\epsilon$
to equalize the two terms ---
$\epsilon^{2(s+t_q)}
  = C_\rho^2\,\eta\,\|\phi\|_{L_p}^2
  / |\phi|_{W^{t,q}}^2$ up to
a factor depending on $s$ and $t_q$ ---
each term is
$\tfrac{1}{2}\,
  |\phi|_{W^{t,q}}^{2s/(s+t_q)}\,
  (C_\rho\,\|\phi\|_{L_p})^{2t_q/(s+t_q)}\,
  \eta^{t_q/(s+t_q)}$
up to the same factor.
Taking the square root
of $\eta\,\rho^2(\tilde\phi)
  \le 2 \times (\text{each term})$
yields~\eqref{eq:mollification-penalty-bias}.
\end{proof}

\begin{remark}[The role of $p$ and the entry ticket]
\label{rem:entry-ticket}
The exponent $p$ is the baseline integrability
of the penalty:
$p = 2$ for the $H^s$ penalty,
but $p$ can be smaller for weaker penalties
(at the cost of requiring $s > d/p$
for the Donsker condition).
The exponent in~\eqref{eq:mollification-penalty-bias}
depends on $t_q$ but not on $p$;
$p$ enters only through the constant
(via $\|\phi\|_{L_p}$)
and the \emph{entry ticket} $\phi \in L_p$.
When $\phi \notin L_p$,
the comparison bound cannot start;
the balance equation (below)
is the only route.

The penalty bias exponent varies
with the regularity of $\phi$:
\[
\eta\,\rho(\tilde\phi)
  = \begin{cases}
    O(\eta) & \text{if } \rho(\phi) < \infty
      \text{ (direct competitor)}, \\
    O(\eta^{3/4}) & \text{if } \phi \in H^s
      \text{ (full regularity, $t_q = s$)}, \\
    O(\eta^{(s+2t_q)/(2(s+t_q))})
      & \text{if } \phi \in W^{t,q},\;
        t < s
      \text{ (partial regularity)}.
  \end{cases}
\]
At partial regularity the exponent drops below $1$:
$\eta$ must decrease faster,
and how much faster depends on $t_q$,
which depends on both
the regularity order~$t$
and the integrability~$q$.
\end{remark}

\begin{proposition}[Mollification penalty bias, Setting~E]
\label{prop:penalty-bias-E}
In the setting
of \Cref{prop:penalty-bias-Q},
suppose additionally that
$\griesz \in L_a(\P)$
for some $a > 1$,
with $a' = a/(a-1)$ the conjugate exponent.
Call $t_{q,a'} = t - d(1/q - 1/(2a'))_+$
the effective $L_{2a'}$ regularity.
Then in Setting~E
($|\dddchis| \le \beta\,\ddchis$),
the penalty bias satisfies
\begin{equation}
\label{eq:mollification-penalty-bias-E}
\eta\,\rho(\tilde\phi)
  \le c_{s,t,d}\,
  (\tau\,\|\griesz\|_{L_a})^{s/(2(s+t_{q,a'}))}\,
  |\phi|_{W^{t,q}}^{s/(s+t_{q,a'})}\,
  (C_\rho\,\|\phi\|_{L_p})^{t_{q,a'}/(s+t_{q,a'})}\,
  \eta^{(s + 2t_{q,a'})/(2(s + t_{q,a'}))},
\end{equation}
where $c_{s,t,d}$ depends only on $s$, $t_{q,a'}$, and $d$.
\end{proposition}

The bound~\eqref{eq:mollification-penalty-bias-E}
is $o(n^{-1/2})$ when
\begin{equation}
\label{eq:lambda-upper-E}
\eta \ll n^{-(s + t_{q,a'})/(s + 2t_{q,a'})}.
\end{equation}
When $\rho(\phi) < \infty$,
the competitor $f = \phi$ avoids
the Bregman term entirely
and~\eqref{eq:mollification-penalty-bias-E}
is not needed.

\begin{proof}
The structure follows
\Cref{prop:penalty-bias-Q},
with the Bregman divergence step
replaced by a H\"older separation.

\paragraph{Bregman divergence.}
For entropy-like dispersions,
$\ddchis(\phi) = \tau\, e^{\tau\phi}$,
so the mean-value form
of the Bregman divergence gives
\[
D_{\chis}(\phi_\epsilon, \phi)
  \le \tfrac{\tau}{2}\,
  e^{\tau\max(\phi_\epsilon, \phi)}\,
  (\phi_\epsilon - \phi)^2.
\]
Because $\phi_\epsilon$ is a convex average
of translates of $\phi$,
Jensen's inequality gives
$e^{\tau\phi_\epsilon} \le (e^{\tau\phi})_\epsilon
  = (\griesz)_\epsilon$.
The maximum satisfies
$e^{\tau\max(\phi_\epsilon,\phi)}
  \le e^{\tau\phi_\epsilon} + e^{\tau\phi}
  \le (\griesz)_\epsilon + \griesz$,
so
\begin{equation}
\label{eq:bregman-mollification-E}
\P\,D_{\chis}(\phi_\epsilon, \phi)
  \le \tfrac{\tau}{2}\,
  \P\,\{(\griesz)_\epsilon + \griesz\}\,
    (\phi_\epsilon - \phi)^2.
\end{equation}

\paragraph{H\"older separation.}
Separating each weight from the approximation
by H\"older's inequality
and using Young's convolution inequality
$\|(\griesz)_\epsilon\|_{L_a} \le \|\griesz\|_{L_a}$
gives
\[
\P\,\{(\griesz)_\epsilon + \griesz\}\,
  (\phi_\epsilon - \phi)^2
  \le 2\,\|\griesz\|_{L_a(\P)}\,
  \|\phi_\epsilon - \phi\|_{L_{2a'}(\P)}^2
\]
for conjugate $1/a + 1/a' = 1$.
The $L_{2a'}$ mollification estimate
for $\phi \in W^{t,q}$ gives
$\|\phi_\epsilon - \phi\|_{L_{2a'}}
  \le c_{d,t,q}\,\epsilon^{t_{q,a'}}\,|\phi|_{W^{t,q}}$,
so the Bregman approximation term satisfies
\begin{equation}
\label{eq:bregman-approx-E}
\P\,D_{\chis}(\phi_\epsilon, \phi)
  \le \tau\,\|\griesz\|_{L_a}\,
  \epsilon^{2t_{q,a'}}\,|\phi|_{W^{t,q}}^2.
\end{equation}

\paragraph{Penalty of the mollifier.}
As in \Cref{prop:penalty-bias-Q},
the embedding $\rho \le C_\rho\,|\cdot|_{W^{s,p}}$
and Bernstein's inequality give
$\rho(\phi_\epsilon)
  \le C_\rho\,\epsilon^{-s}\,\|\phi\|_{L_p}$.

\paragraph{Optimizing the bandwidth.}
Substituting~\eqref{eq:bregman-approx-E}
and~\eqref{eq:penalty-mollification}
into the comparison bound~\eqref{eq:comparison-bound}
gives
\[
\eta\,\rho^2(\tilde\phi)
  \le \tau\,\|\griesz\|_{L_a}\,
  \epsilon^{2t_{q,a'}}\,|\phi|_{W^{t,q}}^2
  + C_\rho^2\,\eta\,\epsilon^{-2s}\,
    \|\phi\|_{L_p}^2.
\]
Choosing $\epsilon$
to equalize the two terms ---
$\epsilon^{2(s+t_{q,a'})}
  = C_\rho^2\,\eta\,\|\phi\|_{L_p}^2
  / (\tau\,\|\griesz\|_{L_a}\,
  |\phi|_{W^{t,q}}^2)$
up to a factor depending on $s$ and $t_{q,a'}$ ---
each term is
$(\tau\,\|\griesz\|_{L_a})^{s/(s+t_{q,a'})}\,
  |\phi|_{W^{t,q}}^{2s/(s+t_{q,a'})}\,
  (C_\rho\,\|\phi\|_{L_p})^{2t_{q,a'}/(s+t_{q,a'})}\,
  \eta^{t_{q,a'}/(s+t_{q,a'})}$
up to the same factor.
Taking the square root
of $\eta\,\rho^2(\tilde\phi)
  \le 2 \times (\text{each term})$
yields~\eqref{eq:mollification-penalty-bias-E}.
The additional cost relative to Setting~Q
is twofold:
the effective regularity drops
from $t_q$ to $t_{q,a'}$
(more embedding needed for $L_{2a'}$),
and the constant acquires a factor
$(\tau\,\|\griesz\|_{L_a})^{s/(2(s+t_{q,a'}))}$.
\end{proof}

\begin{proposition}[Balance equation penalty bias, Settings~Q and~Q$_+$]
\label{prop:penalty-bias-balance}
For the quadratic dispersion
$\chis(\phi) = \phi^2/2$,
suppose $\griesz \in L_2(\P)$.
Then
\begin{equation}
\label{eq:balance-penalty-bias}
\P\tilde\phi^2 \le \|\griesz\|_{\LtP}^2,
\qquad
\rho(\tilde\phi) \le \frac{\|\griesz\|_{\LtP}}{\sqrt{2\eta}},
\qquad
\eta\,\rho(\tilde\phi)
  \le \frac{\|\griesz\|_{\LtP}}{\sqrt{2}}\,
  \sqrt{\eta}.
\end{equation}
No smoothness of $\griesz$ is required
--- only $L_2$ integrability.
\end{proposition}

The penalty bias bound
is $o(n^{-1/2})$ when
\begin{equation}
\label{eq:lambda-upper-edge}
\eta \ll n^{-1}.
\end{equation}

\begin{proof}
The population minimizer $\tilde\phi$ satisfies
the first-order condition
\begin{equation}\label{eq:foc-quadratic}
\P\,\tilde\phi\,h
  + 2\eta\,\inner{\tilde\phi}{h}_\rho
  = \dotpsiZ(h)
  \qquad \text{for all } h \in \mathcal{H}.
\end{equation}
Testing at $h = \tilde\phi$
gives the balance equation
\begin{equation}
\label{eq:balance-equation}
\P\tilde\phi^2
  + 2\eta\,\rho^2(\tilde\phi)
  = \P\,\griesz\,\tilde\phi.
\end{equation}
By Cauchy--Schwarz,
$\P\,\griesz\,\tilde\phi
  \le \|\griesz\|_{\LtP}\,
  (\P\tilde\phi^2)^{1/2}$.
Write $A = \P\tilde\phi^2$
and $G = \|\griesz\|_{\LtP}$.
Then $A + 2\eta\,\rho^2(\tilde\phi)
  \le G\sqrt{A}$.
Since $A \le A + 2\eta\,\rho^2(\tilde\phi)$,
dividing by $\sqrt{A}$ gives
$\sqrt{A} \le G$,
i.e.\ $\P\tilde\phi^2 \le G^2$.
Dropping $A$ from the left
of the combined inequality
and substituting $\sqrt{A} \le G$
into the right side gives
$2\eta\,\rho^2(\tilde\phi) \le G^2$,
yielding all three
bounds~\eqref{eq:balance-penalty-bias}.
\end{proof}

\begin{remark}[Why the balance equation fails for Setting~E]
\label{rem:balance-fails-E}
For entropy-like dispersions,
$\dchis(\tilde\phi) = e^{\tau\tilde\phi}$,
and testing the first-order condition~\eqref{eq:foc-quadratic}
at $h = \tilde\phi$ gives
\[
\P\,e^{\tau\tilde\phi}\,\tilde\phi
  + 2\eta\,\rho^2(\tilde\phi)
  = \P\,\griesz\,\tilde\phi.
\]
The left side has the weighted moment
$\P\,e^{\tau\tilde\phi}\,\tilde\phi$
instead of $\P\tilde\phi^2$.
This does not control
$\|\tilde\phi\|_{L_2}$ from below
--- it can be large
when $\tilde\phi$ is large and positive,
or small (even negative)
when $\tilde\phi$ is negative ---
so the elimination argument
of \Cref{prop:penalty-bias-balance}
does not go through.
The comparison bound
is the only route for Setting~E.
\end{remark}

\begin{remark}[Positive-part quadratic]
\label{rem:positive-part-penalty-bias}
For the positive-part quadratic
$\dchis(\tilde\phi) = \tilde\phi_+$,
testing~\eqref{eq:foc-quadratic}
at $h = \tilde\phi$ gives
\[
\P\,\tilde\phi_+\,\tilde\phi
  + 2\eta\,\rho^2(\tilde\phi)
  = \P\,\griesz\,\tilde\phi.
\]
Because $\tilde\phi_+ \tilde\phi_- = 0$ pointwise,
$\P\,\tilde\phi_+\,\tilde\phi
  = \P\tilde\phi_+^2$.
The left side
has $\P\tilde\phi_+^2$
rather than $\P\tilde\phi^2$,
but the interpolation on the right side
controls $\|\tilde\phi\|_{L_q}$
(all of $\tilde\phi$),
and on the dead zone $\{\tilde\phi < 0\}$,
the FOC reduces to
$2\eta\,\inner{\tilde\phi}{h}_\rho
  = \P\,\griesz\,h$,
which controls $\tilde\phi_-$ through $\griesz$.
The balance equation
closes with the same bound
as \Cref{prop:penalty-bias-balance};
the dead zone enters
through condition~\ref{cond:an-B} and the rate,
not through condition~\ref{cond:an-A}.
\end{remark}

\paragraph{Summary.}
For the quadratic family
with the $H^s$ penalty,
the comparison bound
(\Cref{prop:penalty-bias-Q})
covers
$\griesz \in L_r$ for $r \ge 2$
(plus any regularity);
the balance equation
(\Cref{prop:penalty-bias-balance})
requires only $\griesz \in L_2$
(threshold $s > d/2$).
Under overlap,
both routes work
and the balance equation gives
the sharper bound.

For entropy-like dispersions,
the balance equation is not available
(\Cref{rem:balance-fails-E}),
so the comparison bound
(\Cref{prop:penalty-bias-E})
is the only route.
The entry ticket is
$\phi = \log(\griesz)/\tau
  \in L_p$
--- regularity of the log-representer,
a genuinely different condition
from integrability of~$\griesz$.
When $\phi$ lacks
both the entry ticket and regularity,
there is no known route
to making the penalty bias
$o(n^{-1/2})$.


\section{Rate of Convergence}
\label{sec:rate-of-convergence}

We establish rates of convergence for the sample minimizer $\hat\phi$ relative to the population minimizer $\tilde\phi$, using the offset complexity~\eqref{eq:offset-complexity}.

\subsection{Rate via offset complexity}

Consider sample and population losses of the form \eqref{eq:dual-problem}:
\begin{align*}
L_{\hP}(g) &= \hP \chisZ\{g(Z)\} - \hP \dotpsiZ(g) + \zetas\{\rho(g)\}, \\
L_{\P}(g) &= \P \chisZ\{g(Z)\} - \P \dotpsiZ(g) + \zetas\{\rho(g)\}.
\end{align*}
Let $\hat\phi$ minimize $L_{\hP}$ and $\tilde\phi$ minimize $L_{\P}$.
Define the composite penalty $\Phi = \zetas \comp \rho$
and its Bregman divergence
\[
D_\Phi(f, \tilde\phi)
  = \Phi(f) - \Phi(\tilde\phi)
  - v_{\tilde\phi}(f - \tilde\phi)
\qquad \text{for some } v_{\tilde\phi} \in \partial\Phi(\tilde\phi).
\]
Define the \emph{localization radius} $\rho_t$
as the smallest number such that
$D_\Phi(f, \tilde\phi) \le t^2$ implies
$\rho(f - \tilde\phi) \le \rho_t$.

\begin{lemma}[Rate via offset complexity]
\label{lem:generic-rate}
Let $\nu > 0$, $s^2 \ge \nu\, r^2 > 0$, and let $\rho_s$ be the localization radius at level $s$. If
\[
\rho_s\, \Xi(\rho_s/\nu) < \nu\, r^2,
\]
then $D_\Phi(\hat\phi, \tilde\phi) < \nu\, r^2$,
$\rho(\hat\phi - \tilde\phi) \le \rho_s$,
and
\begin{equation}\label{eq:surviving-bregman}
\hP\left[D_{\chisZ}\{\hat\phi, \tilde\phi\}
  - \nu\, D_{\chisZ}\!\left\{\tilde\phi
    + \frac{\hat\phi - \tilde\phi}{\nu},\,
    \tilde\phi\right\}\right]
  + D_\Phi(\hat\phi, \tilde\phi)
  \le \rho_s\, \Xi(\rho_s/\nu).
\end{equation}
When $\nu \ge 1$, the bracketed term is nonnegative
by secant monotonicity of convex $D$ at the origin,
so it can be dropped for a pure rate bound;
it becomes useful
in the variance equivalence argument
(\Cref{sec:variance-equiv}).
\end{lemma}

\begin{proof}
The loss $L_{\hP}$ is lower at $\hat\phi$ than at $\tilde\phi$, so $\ell(\hat\phi) \le 0$ for the loss difference
\begin{equation}
\label{eq:rate-loss-diff}
\begin{aligned}
\ell(f) &= L_{\hP}(f) - L_{\hP}(\tilde\phi) \\
&= \hP D_{\chisZ}\{f, \tilde\phi\} + (\hP - \P)\{(\dchisZ \comp \tilde\phi) - \dotpsiZ\}(f - \tilde\phi) + D_\Phi(f, \tilde\phi),
\end{aligned}
\end{equation}
where the decomposition uses the population first order condition $\P\{(\dchisZ \comp \tilde\phi) - \dotpsiZ\}(f - \tilde\phi) + v_{\tilde\phi}(f - \tilde\phi) = 0$.

Define the path $f_t = (1-t)\tilde\phi + t\hat\phi$ and write $h_t = f_t - \tilde\phi = t(\hat\phi - \tilde\phi)$. Convexity of $\ell$ with $\ell(\tilde\phi) = 0$ and $\ell(\hat\phi) \le 0$ gives $\ell(f_t) \le 0$ for all $t \in [0,1]$. From $\ell(f_t) \le 0$:
\begin{equation}\label{eq:path-bound}
\hP D_{\chisZ}\{f_t, \tilde\phi\} + D_\Phi(f_t, \tilde\phi) \le (\hP - \P)\{\dotpsiZ - (\dchis \comp \tilde\phi)\}(h_t).
\end{equation}
Define $T = \sup\{t \in [0,1] : D_\Phi(f_{t'}, \tilde\phi) \le s^2 \text{ for all } t' \le t\}$. At $t=0$, $D_\Phi = 0$, so $T > 0$. For $t \le T$, the localization $D_\Phi(f_t, \tilde\phi) \le s^2$ gives $\rho(h_t) \le \rho_s$ by definition of the localization radius, so $g = h_t/\rho_s \in B_\rho$ and $f_t = \tilde\phi + \rho_s\, g$. The definition of $\Xi$ at scale $\rho_s/\nu$ gives
\[
(\hP - \P)\{\dotpsiZ(g) - (\dchis \comp \tilde\phi)\, g\}
  \le \frac{\nu}{\rho_s}\, \hP\, D_{\chisZ}\!\left\{\tilde\phi + \frac{\rho_s}{\nu}\, g,\, \tilde\phi\right\}
    + \Xi(\rho_s/\nu).
\]
Because $(\rho_s/\nu)\,g = h_t/\nu$, scaling by $\rho_s$ and substituting into~\eqref{eq:path-bound}:
\[
\hP D_{\chisZ}\{f_t, \tilde\phi\} + D_\Phi(f_t, \tilde\phi)
  \le \nu\, \hP\, D_{\chisZ}\!\left\{\tilde\phi + \frac{h_t}{\nu},\, \tilde\phi\right\}
    + \rho_s\, \Xi(\rho_s/\nu).
\]
The $\hP D_{\chisZ}$ terms no longer cancel.
Because $D_{\chisZ}\{f_t, \tilde\phi\}
  \ge \nu\, D_{\chisZ}\!\left\{\tilde\phi + h_t/\nu,\, \tilde\phi\right\}$
is not guaranteed,
we keep the surviving Bregman
on the left and use only $D_\Phi \le \rho_s\,\Xi(\rho_s/\nu)$
for the continuity argument:
\[
D_\Phi(f_t, \tilde\phi)
  \le \rho_s\, \Xi(\rho_s/\nu) < \nu\, r^2 \le s^2.
\]
The strict inequality at $t = T$ and the continuity of $t \mapsto D_\Phi(f_t, \tilde\phi)$ imply $T = 1$. Setting $t = 1$ gives~\eqref{eq:surviving-bregman}, $D_\Phi(\hat\phi, \tilde\phi) < \nu\, r^2$, and $\rho(\hat\phi - \tilde\phi) \le \rho_s$.
\end{proof}

\subsection{RKHS models}
\label{sec:rkhs}

We specialize \Cref{lem:generic-rate} to RKHS settings with general regularization $\zetas$.

\begin{lemma}[Local Rademacher complexity for RKHS
  {\citep[Theorem~6.5]{bartlett2005local}}]
\label{lem:bartlett-local}
Let $\mathcal{H}$ be an RKHS
with unit ball $\mathcal{F} = \{f : \|f\|_{\mathcal{H}} \le 1\}$.
Let $T$ be the integral operator
on $L_2(\Q)$ with kernel $K$
and eigenvalues $\mu_1 \ge \mu_2 \ge \cdots > 0$.
For any $r > 0$,\footnote{%
  \Citet{bartlett2005local} state the weaker bound
  $\E_\sigma R_n \le (\ldots)^{1/2}$;
  the squared version follows
  from their proof,
  which bounds $\E_\sigma[Z^2]$
  via Cauchy--Schwarz
  in the RKHS eigenbasis
  and then derives the stated bound
  via Jensen's inequality.}
\[
\E_\sigma \qty[R_n\{f \in \mathcal{F} : \Q f^2 \le r\}]^2
  \le \frac{2}{n}
    \sum_{j=1}^\infty \min(r, \mu_j).
\]
\end{lemma}

\begin{lemma}[RKHS Rademacher complexity]
\label{lem:rkhs-rademacher}
Let $\mathcal{H}$ be an RKHS with reproducing kernel $K$
and norm $\|h\|_{\mathcal{H}}^2 = \|h\|_\chi^2 + \rho(h)^2$.
Let $T_\chi$ be the integral operator on $L_2(\ddchis \cdot \P)$
with kernel $K$,
eigenvalues $\mu_1 \ge \mu_2 \ge \cdots > 0$,
and eigenfunctions $\phi_j$
orthonormal in $L_2(\ddchis \cdot \P)$.
For independent Rademacher variables $\varepsilon_1, \ldots, \varepsilon_n$,
\begin{equation}
\label{eq:rademacher-rkhs}
\E \qty[\sup_{\rho(h) \le B,\, \|h\|_\chi \le r} \frac{1}{n} \sum_{i=1}^n \varepsilon_i h(Z_i)]^2 \le \frac{2}{n} \sum_j \min(B^2 \lambda_j, r^2).
\end{equation}
\end{lemma}
\begin{proof}
Write $\Q = \ddchis(\tilde\phi) \cdot \P$
and expand $h = \sum_j c_j \phi_j$
in the $L_2(\Q)$-orthonormal eigenbasis.
The seminorm constraint
$\rho(h)^2 = \sum_j c_j^2/\lambda_j \le B^2$
implies $c_j^2 \le B^2 \lambda_j$ for each $j$.
The curvature-norm constraint
$\|h\|_\chi^2 = \sum_j c_j^2 \le r^2$
implies $c_j^2 \le r^2$ for each $j$.
Together, $c_j^2 \le d_j := \min(B^2\lambda_j, r^2)$.
The constraint set is contained
in the unit ball of the RKHS
with eigenvalues $d_j$,
intersected with $\{\Q(h^2) \le r^2\}$.
\Cref{lem:bartlett-local} gives
$\E_\sigma[R_n]^2
  \le \frac{2}{n}\sum_j d_j\,\Q(\phi_j^2)
  = \frac{2}{n}\sum_j d_j$
(because $\Q(\phi_j^2) = 1$).
\end{proof}

\paragraph{Setup.} Let $\mathcal{H}$ be an RKHS with norm
\[
\|h\|_{\mathcal{H}}^2 = \|h\|_\chi^2 + \rho(h)^2
\]
for a seminorm $\rho$. The seminorm leaves the $\chi$-norm component unconstrained; localization through $\zetas \comp \rho$ handles this. Assume the interpolation inequality~\eqref{eq:interpolation} with exponent $\theta$ and constant $\kappa$. Let $T_\chi$ be the integral operator on $L_2(\ddchis \cdot \P)$,
\[
(T_\chi f)(x) = \int K(x,y) f(y) \ddchis\{\tilde\phi(y)\} \, d\P(y)
\]
where $K$ is the reproducing kernel of $\mathcal{H}$. Let $\mu_1 \ge \mu_2 \ge \cdots > 0$ be the eigenvalues of $T_\chi$ with eigenfunctions $\phi_j$ orthonormal in $L_2(\ddchis \cdot \P)$. By \Cref{lem:rkhs-rademacher}, the Rademacher complexity of the RKHS ball satisfies~\eqref{eq:rademacher-rkhs}.

\paragraph{Localization.} For convex increasing $\zetas: \R_{\ge 0} \to \R_{\ge 0}$, define the \emph{localization radius} $\rho_t$ as the smallest bound such that
\[
D_{\zetas \comp \rho}(f, \tilde\phi) \le t^2 \implies \rho(f - \tilde\phi) \le \rho_t.
\]

\subsubsection{Bregman lower bound for RKHS}

\begin{lemma}[Bregman lower bound in RKHS models]
\label{lem:rkhs-ulb}
Under the RKHS setup of this section,
suppose $\chis$ is twice continuously differentiable
with $\ddchis$ convex
and $|\dddchis| \le \beta\,\ddchis$,
and $0 < \underline{\ddchis}
  \le \ddchis\{\tilde\phi(X)\}
  \le \overline{\ddchis}$ a.s.
Let $M > 0$ satisfy
$\|h\|_{L_\infty} \le M$
for all $h$ with
$\P D_{\chis}(\tilde\phi + h, \tilde\phi) \le r^2$
and $\rho(h) \le \rho_t$.
Write $\Omega = e^{\beta M}$,
$\overline{D_{\chis}} = \tfrac{1}{2}\,\Omega\, \overline{\ddchis}\, M^2$,
and $C_\Omega = C\,\Omega^{10/3}\,\overline{\ddchis}^2\,M^2$
for an absolute constant $C$.
For $\delta' \ge C_\Omega/(n\,\xi\,r^2)$, suppose
\begin{equation}\label{eq:rkhs-slb-condition}
\sqrt{\frac{1}{n}
  \sum_j \min\!\qty(\rho_t^2\lambda_j,\;
      2\Omega^{1/3}r^2)}
  \le c_1^{3/2}\,\xi^{3/2}\,r^2\,
    \sqrt{\frac{n\,\underline{\ddchis}}
      {\overline{D_{\chis}}}}.
\end{equation}
Then
$\hP\, D_{\chis}(\tilde\phi+h, \tilde\phi)
  \ge (1-2\xi)\, \P\, D_{\chis}(\tilde\phi+h, \tilde\phi)$
uniformly over
$H_r^\circ = \{h \in \mathcal{H} :
  \P D_{\chis}(\tilde\phi+h,\tilde\phi) = r^2,\,
  \rho(h) \le \rho_t\}$
with probability $\ge 1 - \delta'$.
\end{lemma}

\begin{proof}
\Cref{cor:ulb-entropy} verifies
conditions~\ref{cond:shell-i}, \ref{cond:shell-iii-prime}, and~\ref{cond:shell-iv}
of \Cref{thm:bregman-lower}.
It remains to check the global Rademacher
condition~\ref{cond:shell-ii-prime}.

The Bregman ball containment gives
two constraints on $h \in H_r^\circ$:
$\|h\|_{\chis} \le r\sqrt{2\Omega^{1/3}}$
(from $\P D_{\chis} \le r^2$
and the curvature ratio)
and $\rho(h) \le \rho_t$
(from localization).
The RKHS Rademacher bound~\eqref{eq:rademacher-rkhs}
with these constraints gives
\[
\E \qty[\sup_{h \in H_r^\circ}
  \abs*{\hP \sigma\, h}]^2
  \le \frac{1}{n\,\underline{\ddchis}}
    \sum_j \min\!\qty(\rho_t^2\lambda_j,\;
      2\Omega^{1/3}r^2).
\]
Condition~\ref{cond:shell-ii-prime} requires
$\E\sup|\hP\sigma h|
  \le c_1\, \xi\, r\, \sqrt{\ell/n}/2$,
where $\ell = c_0\, n\, \xi\, r^2/(2\overline{D_{\chis}})$
from \Cref{lem:slb-bounded}
applied to $\sqrt{g}$
at tolerance $\xi/2$
(as in \Cref{cor:ulb-entropy}).
By Jensen, it suffices that
$\underline{\ddchis}^{-1}$ times the eigenvalue sum
is at most
\[
\frac{c_1^2\, \xi^2\, r^2}{4} \cdot \ell
  = \frac{c_0\, c_1^2}{8}\,
    \frac{n\, \xi^3\, r^4}{\overline{D_{\chis}}}
  = c_1^3\, \frac{n\, \xi^3\, r^4}{\overline{D_{\chis}}},
\]
using $c_0/8 = c_1$.
Squaring~\eqref{eq:rkhs-slb-condition}
gives this condition.

For the failure probability,
\Cref{cor:ulb-entropy}
with $t = 1$ gives
$\delta' \le C\,\Omega^{10/3}\,
  \overline{\ddchis}^2\, M^2 / (n\,\xi\,r^2)
  = C_\Omega/(n\,\xi\,r^2)$,
which is at most $\delta'$
by hypothesis.
The tolerance is
$\xi_{1,k} = 2\xi + (2-\xi)e^{-k} \le 2\xi + 1/k$,
which is at most $2\xi + \delta'$
since $k \ge 1/\delta'$.
\end{proof}

\subsubsection{Rate for RKHS models}

\begin{corollary}[RKHS model]
\label{cor:rkhs}
With the setup above,
suppose $\chis$ is twice continuously differentiable
with $\ddchis$ convex, $|\dddchis| \le \beta\,\ddchis$,
and $0 < \underline{\ddchis}
  \le \ddchis\{\tilde\phi(X)\}
  \le \overline{\ddchis}$ a.s.
Suppose
\[
\P\qty{(\dchis \comp \tilde\phi)(X)\, h(X) - \dotpsiZ(h)}^2
  \le \bar\nu \|h\|_{\chis}^2
  \quad \text{for all } h \in \mathcal{H}.
\]
For a given analysis radius $r > 0$,
let $M$ be a uniform bound
on $\|h\|_{L_\infty}$
for all $h$ with
$\P D_{\chis}(\tilde\phi+h,\tilde\phi) \le r^2$
and $\rho(h) \le \rho_t$;
write $H_r$ for this set.
Write $\Omega = e^{\beta M}$,
$\overline{D_{\chis}}
  = \tfrac{1}{2}\,\Omega\, \overline{\ddchis}\, M^2$,
$C_\Omega = C\,\Omega^{10/3}\,\overline{\ddchis}^2\,M^2$,
and $\nu_{\min} = \min(\nu, 1-2\xi)$.
Suppose
\begin{equation}\label{eq:rkhs-condition}
\sqrt{\frac{1}{n}
  \sum_j \min\!\qty(\rho_t^2\lambda_j,\;
      2\Omega^{1/3}r^2)}
  \le \frac{\nu_{\min}\sqrt\delta}{2\sqrt{2\bar\nu}}
    \cdot r^2
\qqtext{for}
\frac{2C_\Omega}{n\,\xi\,r^2}
  \le \delta
  \le \frac{8\bar\nu\,c_1^3\,\xi^3}
    {\nu_{\min}^2\,\overline{D_{\chis}}}.
\end{equation}
Then with probability $\ge 1 - \delta$,
\[
\P D_{\chis}(\hat\phi, \tilde\phi) < r^2
\qqtext{and}
D_{\zetas \comp \rho}(\hat\phi, \tilde\phi) \le \nu\, r^2,
\]
and the surviving Bregman bound~\eqref{eq:surviving-bregman}
holds with $\nu \ge 1$ from \Cref{lem:generic-rate}.

For the positive-part quadratic,
see \Cref{cor:sobolev-two-pile}.
\end{corollary}

\begin{proof}[Proof of Corollary~\ref{cor:rkhs}]
The argument has five steps:
bound $\|h\|_{L_\infty}$ for $h \in H_r$,
verify the Bregman lower bound,
bound the offset complexity,
apply the rate lemma,
and show by contradiction
that the population Bregman divergence stays inside $r^2$.
We split the failure budget $\delta$
into $\delta/2$ for the Bregman lower bound
and $\delta/2$ for the offset complexity.

\paragraph{Boundedness.}
For $h$ with
$\P D_{\chis}(\tilde\phi+h,\tilde\phi) \le r^2$
and $\rho(h) \le \rho_t$,
we need a uniform bound on $\|h\|_{L_\infty}$.
The chain runs through the curvature-weighted norm:
Bregman divergence controls $\|h\|_{\chis}$,
which controls $\|h\|_{L_\infty}$ via interpolation.

For the first link,
\Cref{lem:quadratic-bregman}
and the curvature ratio~\eqref{eq:curvature-ratio-psc}
at displacement $h/3$ give
\[
D_{\chis}(\tilde\phi+h,\tilde\phi)
  \ge \tfrac{1}{2}\,
    \ddchis(\tilde\phi + h/3)\, h^2
  \ge \tfrac{1}{2}\,
    e^{-\beta|h|/3}\, \ddchis(\tilde\phi)\, h^2.
\]
Because $|h| \le M$
and $\Omega = e^{\beta M}$,
the exponential is at least $\Omega^{-1/3}$.
Integrating against $\P$:
$\P D_{\chis}(\tilde\phi+h,\tilde\phi)
  \ge \frac{1}{2}\, \Omega^{-1/3}\, \|h\|_{\chis}^2$,
so
$\|h\|_{\chis}^2 \le 2\, \Omega^{1/3}\, r^2$.
The exponent $1/3$ (rather than $1$)
comes from evaluating the curvature ratio
at the intermediate point $h/3$
in \Cref{lem:quadratic-bregman},
which is sharper than
the sandwich bound~\eqref{eq:bregman-sandwich}
evaluated at $h$.

For the second link,
the interpolation inequality~\eqref{eq:interpolation}
with $\rho(h) \le \rho_t$
and $\|h\|_{\chis}^2 \le 2\Omega^{1/3} r^2$ gives
\[
\|h\|_{L_\infty}
  \le \kappa\,
    \rho_t^\theta\,
    (2\Omega^{1/3})^{(1-\theta)/2}\,
    r^{1-\theta}
  = M.
\]
This defines $M$---and hence $\Omega = e^{\beta M}$---implicitly.
Writing $F(M) = \kappa\, \rho_t^\theta\,
  (2e^{\beta M/3})^{(1-\theta)/2}\, r^{1-\theta}$,
$F$ is continuous with $F(0) > 0$;
the condition $M \le r$ restricts
to a compact domain
on which $F(M) = M$ has a solution
by the intermediate value theorem.
When $\beta M \ll 1$
(the analysis radius is small
relative to $1/\beta$),
$\Omega \approx 1$
and $M \approx \kappa\,\rho_t^\theta\, r^{1-\theta}$,
recovering the quadratic case.

\paragraph{Bregman lower bound.}
We verify the hypotheses of \Cref{lem:rkhs-ulb}
with failure probability $\delta/2$.
The lower bound on $\delta$
in~\eqref{eq:rkhs-condition} gives
$\delta/2 \ge C_\Omega/(n\,\xi\,r^2)$,
which is the sample-size condition
of \Cref{lem:rkhs-ulb}.
For the eigenvalue condition~\eqref{eq:rkhs-slb-condition},
we square~\eqref{eq:rkhs-condition}
and use the upper bound on $\delta$:
\[
\frac{1}{n}\sum_j \min(\cdots)
  \le \frac{\nu_{\min}^2\,\delta}{8\bar\nu}
    \cdot r^4
  \le \frac{c_1^3\,\xi^3\,r^4}
    {\overline{D_{\chis}}},
\]
recovering~\eqref{eq:rkhs-slb-condition};
the second inequality uses
$\delta \le 8\bar\nu\,c_1^3\,\xi^3
  /(\nu_{\min}^2\,\overline{D_{\chis}})$.
\Cref{lem:rkhs-ulb} concludes that
$\hP\, D_{\chis}(\tilde\phi+h, \tilde\phi)
  \ge (1-2\xi)\, \P\, D_{\chis}(\tilde\phi+h, \tilde\phi)$
uniformly over $H_r^\circ$,
with probability $\ge 1 - \delta/2$.

\paragraph{Offset complexity.}
We bound the centered empirical process
that defines the offset complexity,
using the remaining failure budget $\delta/2$.
Write $w(Z) = (\dchis \comp \tilde\phi)(X) - \griesz(Z)$
for the noise process weight:
$(\hP - \P)\{(\dchis \comp \tilde\phi)\, h
  - \dotpsiZ(h)\}
= (\hP - \P)(w \cdot h)$.

The $\bar\nu$ hypothesis says
$\P(w h)^2 \le \bar\nu\, \|h\|_{\chis}^2$
for all $h \in \mathcal{H}$.
This is an operator dominance condition:
the integral operator $T_w$
with kernel $w(x)\, K(x,y)\, w(y)$
satisfies $T_w \preceq \bar\nu\, T_\chi$
(the hypothesis says exactly this
in quadratic-form language),
so $\nu_j \le \bar\nu\, \lambda_j$
by the min-max principle.
The noise process inherits
the spectral decay of the RKHS,
scaled by $\bar\nu$.

Symmetrization (\Cref{lem:symmetrization}
with $\Psi(x) = x^2$)
replaces the centered empirical process
with a Rademacher process,
and the RKHS Rademacher
bound~\eqref{eq:rademacher-rkhs}
applied to $\{w \cdot h : h \in H_r\}$
with eigenvalues $\nu_j$ gives
\[
\E \qty[\sup_{h \in H_r}
  (\hP - \P)(w \cdot h)]^2
\le 4\bar\nu \cdot
  \frac{1}{n}
  \sum_j \min\!\qty(\rho_t^2\lambda_j,\;
    2\Omega^{1/3}r^2).
\]
Squaring the eigenvalue
condition~\eqref{eq:rkhs-condition}
bounds the right side by
$\nu_{\min}^2\,\delta\, r^4/2$.
By Markov's inequality,
the supremum exceeds $\nu_{\min}\, r^2$
with probability at most
$\nu_{\min}^2\,\delta\,r^4
  /(2\,\nu_{\min}^2\,r^4)
  = \delta/2$.
Because $\nu_{\min} \le \nu$
and $\nu_{\min} \le 1-2\xi$,
this single bound handles
both the rate application
(which needs the supremum below $\nu\, r^2$)
and the population Bregman contradiction
(which needs it below $(1-2\xi)\, r^2$).

\paragraph{Rate bound.}
The offset complexity bound
$\sup_{h \in H_r}
  (\hP - \P)(w \cdot h)
  \le \nu_{\min}\, r^2$
restricts to each $\chis$-norm shell:
$\overline{M}(u) \le \nu_{\min}\, r^2$
for all $u$.
\Cref{lem:xi-bounded}
at scale $\rho_s/\nu$
gives $\rho_s\, \Xi(\rho_s/\nu) < \nu\, r^2$.
\Cref{lem:generic-rate} gives
$D_\Phi(\hat\phi,\tilde\phi) < \nu\, r^2$,
$\rho(\hat\phi-\tilde\phi) \le \rho_s$,
and the surviving Bregman
bound~\eqref{eq:surviving-bregman}.

\paragraph{Population Bregman bound.}
It remains to show
$\P D_{\chis}(\hat\phi,\tilde\phi) < r^2$.
Suppose for contradiction
that $\P D_{\chis}(f_{t_0}, \tilde\phi) = r^2$
at some $t_0$ along the convex path
$f_t = (1-t)\tilde\phi + t\hat\phi$.
Then $f_{t_0} \in H_r^\circ$
and $\rho(f_{t_0} - \tilde\phi) \le \rho_s$,
so the Bregman lower bound
(\Cref{lem:rkhs-ulb}) gives
$\hP D_{\chis}(f_{t_0},\tilde\phi)
  \ge (1-2\xi)\, r^2$.
On the other hand,
$\ell(f_{t_0}) \le 0$,
so the path bound~\eqref{eq:path-bound} gives
\[
\hP D_{\chis}(f_{t_0},\tilde\phi)
  \le (\hP - \P)\{\dotpsiZ
    - (\dchis \comp \tilde\phi)\}(h_{t_0}).
\]
The empirical process bound
from the offset complexity paragraph---$\sup_{h \in H_r}
  (\hP - \P)(w \cdot h)
  \le \nu_{\min}\, r^2
  \le (1-2\xi)\, r^2$---shows
the right side is less than
$(1-2\xi)\, r^2$
on the same event.
This contradicts the lower bound,
so $T = 1$.
\end{proof}

\begin{remark}[Localization radii for common $\zetas$]
\label{rmk:localization}
The localization radius $\rho_t$ depends on $\zetas$:
\begin{itemize}
\item \emph{Quadratic}: $\zetas(\rho) = \eta\rho^2$ gives $\rho_t = t/\sqrt{\eta} + 2\rho(\tilde\phi)$. When $\sqrt{\eta}\rho(\tilde\phi) \ll t$, we have $\rho_t \asymp t/\sqrt{\eta}$.
\item \emph{Hard constraint}: $\zetas(\rho) = \mathcal{I}\{\rho \le a\}$ forces $\rho(f) \le a$, so $\rho_t = a + \rho(\tilde\phi)$ independent of $t$.
\item \emph{Linear}: $\zetas(\rho) = \eta\rho$ gives $\rho_t = t^2/\eta + 2\rho(\tilde\phi)$.
\end{itemize}
\end{remark}

\begin{remark}[Sobolev spaces]
\label{rmk:sobolev}
Sobolev spaces satisfy the interpolation inequality \eqref{eq:interpolation} with $\rho(g) = \|(-\Delta)^{s/2} g\|_{L_2(\mathcal{X})}$ and $\theta = d/(2s)$ when $s > d/2$. The constant $\kappa$ depends on $s$, $d$, and the domain $\mathcal{X}$. The eigenvalues satisfy $\lambda_j \asymp j^{-2s/d}$, giving $d_\tau \asymp \tau^{d/s}$ for small $\tau$.
\end{remark}

\subsection{Sobolev Models via Two-Pile Accounting}
\label{sec:sobolev-two-pile}

We now give a Sobolev application for the positive-part quadratic $\chis(x) = x_+^2/2$ using \Cref{lem:two-pile}. This approach handles the ``dead zone'' where $\tilde\phi + h$ crosses zero by accepting a tolerance degradation that depends on the boundedness parameter $M$ and the mass $\omega(\epsilon) = \P(\tilde\phi \in [0,\epsilon])$ near the dead-zone boundary.

\begin{corollary}[Sobolev models with positive-part quadratic via two-pile]
\label{cor:sobolev-two-pile}
Let $\chis(x) = x_+^2/2$
and let $\zetas: \R_{\ge 0} \to \R_{\ge 0}$
be convex and increasing.
Let $\rho(g) = \|(-\Delta)^{s/2} g\|_{L_2}$
be the Sobolev seminorm with $s > d/2$,
and let $X$ have density $p \ge \pmin > 0$
on $[0,1]^d$.
Suppose:
\begin{enumerate}
\item[(i)] \textbf{(Dead-zone control)}
  The mass $\omega(\epsilon) = \P(\tilde\phi \in [0,\epsilon])$ satisfies $\omega(\epsilon) \to 0$ as $\epsilon \to 0$.
\end{enumerate}

Let $\nu \ge 1$.
For increments $h = f - \tilde\phi$,
define $\tilde h = \ind{\tilde\phi \ge 0} \cdot h$,
$r = \|\tilde h\|_{\LtP}$,
The boundedness parameter is
\[
M = C_{\mathrm{emb}}\, \rho_{\sqrt{\nu}r} + C_K\, r,
\]
where $C_{\mathrm{emb}} = (1 + C_K)\kappa_{\mathrm{emb}}$
combines the Sobolev embedding constant $\kappa_{\mathrm{emb}}$
(depending on $s$, $d$, $\pmin$)
with the kernel norm equivalence constant $C_K$
(depending on $\dim(\ker\rho)$ and $\P(\tilde\phi \ge 0)$),
and $\rho_t$ is the localization radius
(\Cref{rmk:localization}).
The two-pile transfer (\Cref{lem:two-pile,lem:dead-zone-concentration}) gives combined tolerance~\eqref{eq:two-pile-tolerance} $\xi = \xi_0 + (\mu + 4MR + t)/r^2$ where the dead-zone cost is $\mu = \omega(\epsilon)\, M^2 + r^2 M^2/\epsilon^2$, optimized over $\epsilon > 0$.
Because $M \ge C_K r$, the ratio $M/r \ge C_K$
is bounded below; when $\rho_{\sqrt{\nu}r} \ll r$,
$M/r \to C_K$.

Fix $\xi_0 \in (0,1)$.
Then with probability at least
$1 - 2e^{-c n \xi_0^2 r^2/M^2}$,
\[
\|\hat\phi - \tilde\phi\|_{\chis}^2 < r^2
\qqtext{and}
D_{\zetas \comp \rho}(\hat\phi, \tilde\phi) \le \nu\, r^2,
\]
and the surviving Bregman bound~\eqref{eq:surviving-bregman}
holds with $\nu \ge 1$ from \Cref{lem:generic-rate}.
\end{corollary}

\begin{proof}
We verify the conditions
of \Cref{thm:bregman-lower}
for the positive-part quadratic.

\paragraph{Boundedness.}
Decompose $h = h_0 + k$
with $h_0 \in \ker(\rho)^\perp$
and $k \in \ker(\rho)$.
Sobolev embedding gives
$\|h_0\|_{L_\infty}
  \le \kappa_{\mathrm{emb}}\, \rho(h)
  \le \kappa_{\mathrm{emb}}\, \rho_{\sqrt{\nu}r}$.
For the kernel component:
$\tilde k = \ind{\tilde\phi \ge 0} \cdot k$ satisfies
$\|\tilde k\|_{\LtP}
  \le \|\tilde h\|_{\LtP} + \|\tilde h_0\|_{\LtP}
  \le r + \kappa_{\mathrm{emb}}\, \rho_{\sqrt{\nu}r}$,
and norm equivalence on $\ker(\rho)$ gives
$\|k\|_{L_\infty}
  \le C_K\, \|\tilde k\|_{\LtP}$,
where $C_K$ depends on $\dim(\ker\rho)$
and $\P(\tilde\phi \ge 0)$.
(For constants, $C_K = \P(\tilde\phi \ge 0)^{-1/2}$.)
On the active zone,
$|\tilde h|
  \le \|h_0\|_{L_\infty} + \|k\|_{L_\infty}
  \le (1+C_K)\kappa_{\mathrm{emb}}\, \rho_{\sqrt{\nu}r}
    + C_K\, r
  = M$.

\paragraph{Bregman SLB\@.}
We need condition~\ref{cond:shell-i}
of \Cref{lem:bregman-shell}:
a stable lower bound
for $D_{\chis}(\tilde\phi + h, \tilde\phi)$.
\Cref{lem:two-pile} gives a pointwise SLB
for $D_{\chis}(\tilde\phi + h, \tilde\phi)$
with tolerance depending on
the empirical junction error $\hP f_h$.
\Cref{lem:dead-zone-concentration},
with boundedness $M$ from above
and the Rademacher bound
from~\eqref{eq:rademacher-rkhs},
concentrates $\sup_h \hP f_h$
on an event $E$
that is independent of $h$.
The combined tolerance~\eqref{eq:two-pile-tolerance} includes
the dead-zone cost:
$\mu = \omega(\epsilon)\, M^2 + r^2 M^2/\epsilon^2$,
optimized over $\epsilon > 0$.
Because $E$ is $h$-independent,
it does not inflate the net union bound
in \Cref{thm:bregman-lower};
$\P(E^c)$ enters the probability
as a single exponential term.

\paragraph{Derivative equicontinuity.}
For $\chis(x) = x_+^2/2$,
$\dchis(x) = x_+$ is $1$-Lipschitz,
so condition~\ref{cond:shell-iv}
of \Cref{lem:bregman-shell}
holds with $L = 1$ and $\Q = \P$.

\paragraph{Complexity and rate.}
Condition~\ref{cond:shell-ii-prime}
requires a Rademacher bound
for the active-zone functions
$\tilde h = \ind{\tilde\phi \ge 0} \cdot h$.
Because $\Q = \P$ for Q$_+$
(the derivative equicontinuity
uses the unweighted norm),
the eigenvalue bound~\eqref{eq:rademacher-rkhs}
applies with the $L_2(\P)$ eigenvalues
of the active-zone kernel
$K_+(x,y) = \ind{\tilde\phi(x) \ge 0}\,K(x,y)\,\ind{\tilde\phi(y) \ge 0}$.
\Cref{thm:bregman-lower}
then gives the uniform lower bound
$\hP D_{\chis}(\tilde\phi+h, \tilde\phi)
  \ge (1 - \xi)\, r^2$
over all radii in $[\underline{r}, \overline{r}]$.
The probability terms
$p_{\mathrm{net}}$, $p_{\mathrm{trunc}}$,
$p_{\mathrm{fluct}}$,
and $\P(E^c)$ are all exponential
in $n\, \xi_0^2\, r^2/M^2$,
giving the stated bound.
\Cref{lem:xi-bounded}
at scale $s = \rho_s/\nu$
gives $\Xi(\rho_s/\nu) \le (1-\epsilon)\, r_\star^2\, \nu/(2\rho_s)$;
the eigenvalue condition verifies~\eqref{eq:critical-bounded}
and gives
$\rho_s\, \Xi(\rho_s/\nu) < \nu\, r^2$.
\Cref{lem:generic-rate}
gives the rate.
\end{proof}


\begin{remark}[Rate summary for common penalties]
\label{rmk:rate-table}
For Sobolev spaces with smoothness $s > d/2$, let $q = s/(2s-d)$.
The seminorm embedding~\eqref{eq:seminorm-embedding} gives
\[
M = C_{\mathrm{emb}}\,\rho + C_K\, r,
\qquad
M/r = C_{\mathrm{emb}}\,\rho/r + C_K,
\]
which bounds how large $|\tilde h(x)|$ can be on the active zone, uniformly over the localized class. The following table summarizes $r$, $\rho$, and $M/r$ for two penalty types, suppressing dimension-dependent constants.

\begin{center}
\begin{tabular}{lccc}
\hline
\textbf{Penalty} $\zetas(\rho)$ & $r$ ($L_2$ rate) & $\rho$ (seminorm) & $M/r$ \\
\hline
$\eta \rho^2$ (ridge) & $n^{-q} \eta^{-q/2}$ & $n^{-q} \eta^{-q}$ & $\eta^{-q/2} + C_K$ \\[0.5em]
$\mathcal{I}\{\rho \le a\}$ (constraint) & $n^{-1/2}$ & $a$ & $n^{1/2} a + C_K$ \\
\hline
\end{tabular}
\end{center}

\noindent The ratio $M/r$ governs the two-pile tolerance degradation. Under \emph{undersmoothing} ($\eta \to 0$ or $a \to \infty$), $\rho/r$ diverges:
\begin{itemize}
\item \emph{Ridge}: $M/r \asymp \eta^{-q/2} \to \infty$ as $\eta \to 0$.
\item \emph{Constraint}: $M/r \asymp n^{1/2} a \to \infty$ as $a \to \infty$.
\end{itemize}
Because the dead-zone cost includes the term $M^2/\epsilon^2$, the two-pile tolerance degrades when $M/r$ is large; the rate of degradation depends on the anticoncentration modulus $\omega$. The seminorm embedding~\eqref{eq:seminorm-embedding} requires no curvature structure and applies to the positive-part quadratic where the interpolation inequality~\eqref{eq:interpolation} does not.
\end{remark}



\section{Variance Equivalence}
\label{sec:variance-equiv}

\begin{lemma}[Setting E: variance equivalence]\label{lem:variance-equiv}
Under Setting~E,
write $\tilde\gamma = \dchis(\tilde\phi)$,
$\hgamma = \dchis(\hat\phi)$,
$h = \hat\phi - \tilde\phi$.
Let $v: \mathcal{Z} \to \R$ satisfy $\underline{v} \le v \le \bar{v}$.
Suppose $\hP D_{\chis}(\hat\phi, \tilde\phi) \le \nu r^2$,
$\|h\|_\infty \le M$,
and $\|\tilde\gamma\|_{L_3(\hP)}^3 := \hP \tilde\gamma^3 < \infty$.
Then
\begin{equation}\label{eq:variance-equiv-bound}
\abs{\hP \hgamma^2 v - \hP \tilde\gamma^2 v} \le \beta_1 e^{\beta M}(2 + \beta_1 e^{\beta M} M)\, \bar{v} \|\tilde\gamma\|_{L_3}^{3/2}\, \sqrt{2 e^{\beta M}\nu/\alpha_1}\, r.
\end{equation}
\end{lemma}

The difference $\hP \hgamma^2 v - \hP \tilde\gamma^2 v$ decomposes into a linear term $2\,\hP\, \tilde\gamma\, \ddchis(\tilde\phi)\, h\, v$ and a nonlinear remainder. The key step is \emph{Bregman absorption}: the integral representation $D_{\chis}(\tilde\phi+h,\tilde\phi) = h^2 \int_0^1 (1-t)\,\ddchis(\tilde\phi+th)\,dt$ and pseudo-self-concordance ($|\dddchis| \le \beta\,\ddchis$) convert the rate bound into a weighted $L^2$ bound $\hP\, \ddchis(\tilde\phi)\, h^2 \le 2e^{\beta M}\, \nu r^2$. The Level~1 bound ($\alpha_1\, \dchis \le \ddchis \le \beta_1\, \dchis$) converts this to $\hP\, \tilde\gamma h^2 \le (2 e^{\beta M}/\alpha_1)\,\nu r^2$ and bounds the linear term by $\beta_1\,\hP\,\tilde\gamma^2|h|$. Cauchy--Schwarz with the split $\tilde\gamma^2|h| = \tilde\gamma^{3/2} \cdot \tilde\gamma^{1/2}|h|$ introduces the third-moment condition.

\begin{proof}[Proof of Lemma~\ref{lem:variance-equiv}]
Write $\Delta\gamma = \hgamma - \tilde\gamma
  = \int_0^1 \ddchis(\tilde\phi + th)\, dt \cdot h$.
Decompose
\begin{equation}\label{eq:var-decomp}
\hP \hgamma^2 v - \hP \tilde\gamma^2 v
  = 2\, \hP\, \tilde\gamma\, \ddchis(\tilde\phi)\, h\, v
  + \hP\, R\, v,
\end{equation}
where
$R = 2\tilde\gamma \int_0^1
  [\ddchis(\tilde\phi + th) - \ddchis(\tilde\phi)]\, dt \cdot h
  + (\Delta\gamma)^2$.

\paragraph{Bregman absorption.}
We convert the rate bound $\hP D_{\chis} \le \nu r^2$
into a weighted $L_2$ bound
$\hP\, \ddchis(\tilde\phi)\, h^2 \le 2\Omega\, \nu r^2$
using the Bregman sandwich.
By~\eqref{eq:bregman-sandwich}
with $|h| \le M$
and $\Omega = e^{\beta M}$:
\begin{equation}\label{eq:var-bregman-absorb}
\ddchis(\tilde\phi)\, h^2/(2\Omega)
  \le D_{\chis}(\tilde\phi+h, \tilde\phi)
  \le \Omega\, \ddchis(\tilde\phi)\, h^2/2.
\end{equation}
By~\eqref{eq:weighted-l2-from-bregman} and the rate bound,
$\hP\, \ddchis(\tilde\phi)\, h^2 \le 2\Omega\, \nu r^2$.
The link--curvature bound
from \Cref{lem:e-type-bounds}(iv)
gives
\begin{equation}\label{eq:var-weighted-l2}
\hP\, \tilde\gamma\, h^2
  \le \frac{2\Omega}{\alpha_1}\, \nu r^2.
\end{equation}

\paragraph{Linear term.}
The Level~1 bound
$\ddchis \le \beta_1\, \dchis$
gives
$\tilde\gamma\, \ddchis(\tilde\phi)\, |h|
  \le \beta_1\, \tilde\gamma^2\, |h|$.
To separate pure $\tilde\gamma$-moments
from the $h$-dependent factor,
split
$\tilde\gamma^2 |h|
  = \tilde\gamma^{3/2} \cdot \tilde\gamma^{1/2}|h|$
and apply Cauchy--Schwarz:
\[
\abs{\hP\, \tilde\gamma\, \ddchis(\tilde\phi)\, h\, v}
  \le \beta_1\, \bar{v} \sqrt{\hP\, \tilde\gamma^3}\,
    \sqrt{\hP\, \tilde\gamma\, h^2}
  \le \beta_1\, \bar{v} \|\tilde\gamma\|_{L_3}^{3/2}\,
    \sqrt{2 e^{\beta M}\nu/\alpha_1}\, r
\]
by the third-moment bound and~\eqref{eq:var-weighted-l2}.
The linear term in~\eqref{eq:var-decomp}
is therefore bounded by
$2\beta_1\, \bar{v} \|\tilde\gamma\|_{L_3}^{3/2}\,
  \sqrt{2 e^{\beta M}\nu/\alpha_1}\, r$.

\paragraph{Remainder.}
The remainder $R$ has two parts:
a curvature variation term,
first-order in $\Omega - 1$,
and the squared increment $(\Delta\gamma)^2$,
second-order in $h$.

\emph{Curvature variation.}
The curvature ratio~\eqref{eq:curvature-ratio-psc} bounds
the integrand in $R$:
\[
|\ddchis(\tilde\phi+th) - \ddchis(\tilde\phi)|
  \le \ddchis(\tilde\phi)(\Omega - 1).
\]
The Level~1 bound $\ddchis \le \beta_1\, \dchis$
converts this to a $\tilde\gamma$-weighted bound.
Together with the integral representation of $R$,
\[
|R_{\text{curv}}|
  := 2\tilde\gamma
    \int_0^1 |\ddchis(\tilde\phi + th) - \ddchis(\tilde\phi)|\, dt
    \cdot |h|
  \le 2\beta_1(\Omega - 1)\, \tilde\gamma^2\, |h|.
\]
The same Cauchy--Schwarz split
$\tilde\gamma^2 |h| = \tilde\gamma^{3/2} \cdot \tilde\gamma^{1/2} |h|$
as in the linear term gives
\[
\hP\, |R_{\text{curv}}|\, v
  \le 2\beta_1(e^{\beta M}-1)\, \bar{v}
    \|\tilde\gamma\|_{L_3}^{3/2}\,
    \sqrt{2 e^{\beta M}\nu/\alpha_1}\, r.
\]

\emph{Squared increment.}
The derivative Lipschitz
bound~\eqref{eq:derivative-lipschitz}
and Level~1 give
\[
|\Delta\gamma|
  \le \Omega\, \ddchis(\tilde\phi)\, |h|
  \le \beta_1\, \Omega\, \tilde\gamma\, |h|,
\]
so
$\hP\, (\Delta\gamma)^2\, v
  \le \beta_1^2 \Omega^2\, \bar{v}\, \hP\, \tilde\gamma^2 h^2$.
To bound $\hP\, \tilde\gamma^2 h^2$,
apply Cauchy--Schwarz:
\[
\hP\, \tilde\gamma^2 h^2
  \le \sqrt{\hP\, \tilde\gamma^3}\,
    \sqrt{\hP\, \tilde\gamma\, h^4}
  \le \|\tilde\gamma\|_{L_3}^{3/2}\, M\,
    \sqrt{2 e^{\beta M}\nu/\alpha_1}\, r,
\]
where the second inequality uses
$h^4 \le M^2 h^2$,
the weighted bound
$\hP\, \tilde\gamma\, h^4
  \le M^2\, \hP\, \tilde\gamma\, h^2
  \le 2 M^2 e^{\beta M}\, \nu r^2/\alpha_1$
from~\eqref{eq:var-weighted-l2},
and the third-moment bound.
The squared-increment contribution is therefore
\[
\hP\, (\Delta\gamma)^2\, v
  \le \beta_1^2 e^{2\beta M} M\, \bar{v}
    \|\tilde\gamma\|_{L_3}^{3/2}\,
    \sqrt{2 e^{\beta M}\nu/\alpha_1}\, r.
\]

\paragraph{Conclusion.}
Combining: linear ($2\beta_1$),
curvature variation ($2\beta_1(e^{\beta M} - 1)$),
and quadratic ($\beta_1^2 e^{2\beta M} M$),
all multiplied by
$\bar{v} \|\tilde\gamma\|_{L_3}^{3/2}\,
  \sqrt{2 e^{\beta M}\nu/\alpha_1}\, r$.
The sum of coefficients is
$2\beta_1 e^{\beta M} + \beta_1^2 e^{2\beta M} M
  = \beta_1 e^{\beta M}(2 + \beta_1 e^{\beta M} M)$,
giving~\eqref{eq:variance-equiv-bound}.
\end{proof}

The rate bound is the output of \Cref{cor:rkhs}. The sup-norm bound follows from Sobolev interpolation when $\rho = |\cdot|_{\dot H^s}$ with $s > d/2$: the population norm bound $\|h\|_{\chis} < r$ from \Cref{cor:rkhs} and the penalty Bregman bound $D_{\zetas \comp \rho}(\hat\phi, \tilde\phi) \le \nu r^2$ together control $\|h\|_{L_2}$ and $\rho(h)$, and the embedding $\dot H^s(\R^d) \hookrightarrow L^\infty(\R^d)$ gives $\|h\|_\infty \le C_s \|h\|_{L_2}^{1-\theta} \rho(h)^\theta$ where $\theta = d/(2s)$.

The third-moment bound is a population property: $\hP \tilde\gamma^3 \to \P \tilde\gamma^3$ by concentration whenever $\P \tilde\gamma^6 < \infty$.

\setupvariants

\theoremvariant{Q}
\begin{corollary}[Setting Q: variance equivalence]\label{cor:variance-equiv-quadratic}
For $\chis(\phi) = \phi^2/2$ or $\chis(\phi) = \phi_+^2/2$,
write $\tilde\gamma = \dchis(\tilde\phi)$,
$\hgamma = \dchis(\hat\phi)$,
$h = \hgamma - \tilde\gamma$.
Let $v: \mathcal{Z} \to \R$ satisfy $\underline{v} \le v \le \bar{v}$.
If $\hP D_{\chis}(\hat\phi, \tilde\phi) \le \nu r^2$,
then $\hP h^2 \le 2\nu r^2$ and
\begin{equation}\label{eq:variance-equiv-quadratic}
\abs{\hP \hgamma^2 v - \hP \tilde\gamma^2 v}
  \le 2\bar{v}\, \|\tilde\gamma\|_{L_2(\hP)}\, \sqrt{2\nu}\, r
  + 2\bar{v}\, \nu r^2.
\end{equation}
\end{corollary}

\begin{proof}
The pointwise bound
$(\hgamma - \tilde\gamma)^2
  \le 2\, D_{\chis}(\hat\phi, \tilde\phi)$
holds in both cases.
For $\chis = \phi^2/2$
this is $h^2 = 2 \cdot h^2/2$.
For $\chis = \phi_+^2/2$,
check the four cases:
when $\tilde\phi \ge 0$ and $\tilde\phi + h \ge 0$,
the weight difference is $h$
and $D = h^2/2$;
when both are negative,
both sides vanish;
when $\tilde\phi \ge 0$
and $\tilde\phi + h < 0$,
the weight difference is $\tilde\phi$
and $D = \tilde\phi|h| - \tilde\phi^2/2
  \ge \tilde\phi^2/2$
(because $|h| > \tilde\phi$);
when $\tilde\phi < 0$
and $\tilde\phi + h > 0$,
the weight difference is $\tilde\phi + h$
and $D = (\tilde\phi + h)^2/2$.
Averaging gives
$\hP h^2 \le 2\nu r^2$.
Decompose
$\hP \hgamma^2 v - \hP \tilde\gamma^2 v
  = 2\,\hP\,\tilde\gamma\, h\, v
  + \hP h^2 v$.
The cross term is at most
$2\bar{v}\, \|\tilde\gamma\|_{L_2(\hP)}\, \sqrt{2\nu}\, r$
by Cauchy--Schwarz;
the squared term is at most
$2\bar{v}\, \nu r^2$.
\end{proof}
\resettheorem

\theoremvariant{E}
\begin{corollary}[Setting E: variance equivalence under interpolation]\label{cor:variance-interpolation}
Under Setting~E,
suppose $\rho$ satisfies
the interpolation inequality~\eqref{eq:interpolation}
for some $\theta \in (0,1)$.
Let $v: \mathcal{Z} \to \R$ satisfy $\underline{v} \le v \le \bar{v}$.
On any event where
\[
\|h\|_{\chis} < r, \quad
\hP D_{\chis}(\hat\phi, \tilde\phi) \le \nu r^2, \quad
\rho(h) \le \rho_r, \qqtext{and}
\hP \tilde\gamma^3 < \infty,
\]
the bound~\eqref{eq:variance-equiv-bound} holds with
$M = \kappa\, \rho_r^\theta\, r^{1-\theta}$:
\begin{equation}\label{eq:var-ratio-interpolation}
\abs{\hP \hgamma^2 v - \hP \tilde\gamma^2 v}
  \le \beta_1 e^{\beta M}(2 + \beta_1 e^{\beta M} M)\,
     \bar{v} \|\tilde\gamma\|_{L_3}^{3/2}\, \sqrt{2 e^{\beta M}\nu/\alpha_1}\, r.
\end{equation}
The right side vanishes as $r \to 0$
provided $M = \kappa\, \rho_r^\theta r^{1-\theta} \to 0$.
\end{corollary}

For Sobolev models with quadratic regularization $\zetas(\rho) = \eta\rho^2$, \Cref{cor:sobolev-main} gives $\rho_r = \sqrt{\nu}\, r/\sqrt{\eta}$, so $M = \kappa\, \nu^{\theta/2}\, \eta^{-\theta/2}\, r$, which vanishes at any regularization $\eta \ge c_0 n^{-2s/(2s+d)}$ when $s > d/2$.

\begin{proof}[Proof of Corollary~\ref{cor:variance-interpolation}]
The rate bound and third-moment bound hold by hypothesis.
The interpolation inequality~\eqref{eq:interpolation}
and the bounds $\|h\|_{\chis} < r$, $\rho(h) \le \rho_r$
give $\|h\|_\infty \le \kappa\, \rho_r^\theta\, r^{1-\theta} = M$,
verifying the sup-norm bound.
\Cref{lem:variance-equiv} gives~\eqref{eq:var-ratio-interpolation}.
\end{proof}
\resettheorem


\section{Asymptotics}
\label{sec:asymptotics}

\subsection{Abstract conditions}
\label{subsec:abstract-conditions}

\paragraph{Regularized influence function.}
The goal is to show that the one-step estimator
is asymptotically linear in a regularized influence function
determined by the population regularized weights
$\tilde\gamma = \dchis \comp \tilde\phi$:
\begin{equation}\label{eq:regularized-influence}
\tilde\varphi(Y, Z)
  = \tilde\gamma(Z)\{Y - \mu(Z)\}
  + \psiZ(\mu)(Z) - \psi.
\end{equation}
\Cref{thm:asymptotic-normality} shows that
$(\hat\psi - \psi)/\hat\sigma \to_d N(0,1)$
with $\hat\sigma^2 / \tilde\sigma^2 \to_p 1$,
for the second moment
\begin{equation}\label{eq:tilde-sigma}
\begin{aligned}
\tilde\sigma^2 = \P\tilde\varphi^2 = \P\tilde\gamma^2\nu + \P\tilde m^2
\qqtext{where} &\tilde m  = \tilde\gamma \mu_\perp + \psiZ(\mu) - \psi \\
\qfor          &\mu_\perp = \E(Y \mid Z) - \mu.
\end{aligned}
\end{equation}

\paragraph{Background assumptions.}
We assume throughout:
\begin{enumerate}
\item[\textup{(i)}]
\textit{Conditional variance bounded away from zero.}
$\nu(Z)\ind{\nu(Z) > 0} \ge \underline{\nu} > 0$.
\item[\textup{(ii)}]
\textit{Uniform integrability.}
\begin{equation}\label{eq:uniform-integrability}
\max_{i \le n}
\E\qty[\varepsilon_i^2
  \ind{\varepsilon_i^2 \ge M}
  \;\big|\; Z_1 \ldots Z_n]
\to_p 0
\qqtext{as} M \to \infty
\qqtext{for} \varepsilon_i = Y_i - \E(Y_i \mid Z_i).
\end{equation}
\end{enumerate}

% o_p/O_p notation is appropriate here: this is a genuine asymptotic result
\begin{theorem}[Asymptotic normality]
\label{thm:asymptotic-normality}
Let $\hat\psi$ be the one-step estimator~\eqref{eq:one-step}
with $\sigma(Z_1, \ldots, Z_n)$-measurable weights $\hgamma$
and $\hat\Gamma = \hP\hgamma^2\nu$.
Then
$\sqrt{n}(\hat\psi - \psi) \to_d N(0, \tilde\sigma^2)$
if the following conditions hold.
\begin{enumerate}[label=\textup{(\Alph*)}]
\item\label{cond:an-A}
\textit{Negligible imbalance.}
$\hP\hgamma(\hat\mu - \mu)
  - \dpsiZ(\hat\mu - \mu)
  = o_p(n^{-1/2}\hat\Gamma^{1/2})$.

\item\label{cond:an-B}
\textit{Infinitesimal weights.}
$\max_i \hgamma^2(Z_i)\nu(Z_i) / (n\hat\Gamma) \to_p 0$.

\item\label{cond:an-C}
\textit{Variance equivalence.}
$\hat\Gamma / \P\tilde\gamma^2\nu \to_p 1$.

\item\label{cond:an-D}
\textit{Projection coherence.}
$\P(\hgamma - \tilde\gamma)\mu_\perp
  = o_p(n^{-1/2}\hat\Gamma^{1/2})$
and
$\hP(\hgamma - \tilde\gamma)^2 \mu_\perp^2 \to_p 0$.

\item\label{cond:an-E}
\textit{Centering.}
$\P\tilde\gamma\, \mu_\perp
  = o_p(n^{-1/2}\hat\Gamma^{1/2})$.

\item\label{cond:an-F}
\textit{Functional stability.}
$(\hP - \P)\,\psiZ(\hat\mu) - \psiZ(\mu)
= o_p(n^{-1/2}\hat\Gamma^{1/2})$.

\item\label{cond:an-G}
\textit{Linearization.}
$\P\psiZ(\hat\mu) + \dpsiZ(\mu - \hat\mu) - \psi
  = o_p(n^{-1/2}\hat\Gamma^{1/2})$.
\end{enumerate}
\end{theorem}

Conditions~\ref{cond:an-A}--\ref{cond:an-C}
govern the weights:
negligible imbalance,
infinitesimal individual contributions,
and variance stability.
Conditions~\ref{cond:an-D}--\ref{cond:an-E}
handle misspecification:
coherence between estimated and population weights,
and centering.
Conditions~\ref{cond:an-F}--\ref{cond:an-G}
handle the outcome estimator:
empirical process stability
and linearization error.
For linear functionals
(Examples~\ref{ex:treatment-specific-mean}
and~\ref{ex:average-linear-regression}),
conditions~\ref{cond:an-F}--\ref{cond:an-G} are vacuous.

The CLT is not pivotal:
the limiting variance $\tilde\sigma^2$
depends on the population regularized weights
$\tilde\gamma$ and the misspecification $\mu_\perp$.
To construct confidence intervals,
we need a consistent variance estimate.
The plug-in estimator~\eqref{eq:variance-estimate}
achieves this under a mild additional condition.

\begin{proposition}[Variance consistency]
\label{prop:variance-consistency}
Under the conditions of
\Cref{thm:asymptotic-normality},
if
$\hP\hgamma^2(\hat\mu^\sigma - \mu)^2
  + \{\psiZ(\hat\mu^\sigma) - \psiZ(\mu)\}^2
  \to_p 0$
for the nuisance estimate $\hat\mu^\sigma$
used in~\eqref{eq:variance-estimate},
then $\hat\sigma^2 / \tilde\sigma^2 \to_p 1$.
\end{proposition}

\paragraph{Efficiency.}
The variance $\tilde\sigma^2$
depends on the regularization
through $\tilde\gamma$.
When $\griesz \in L_2(\P)$
and $\tilde\gamma \to \griesz$ in $\LtP$,
the regularized influence function
converges to the efficient influence function
\begin{equation}\label{eq:influence-function}
\varphi(Y, Z)
  = \griesz(Z)\{Y - \mu(Z)\}
  + \psiZ(\mu)(Z) - \psi
\end{equation}
and $\tilde\sigma^2 \to \sigma^2 = \P\varphi^2$,
the semiparametric efficiency bound.

\begin{corollary}[Efficient variance]
\label{cor:efficient-variance}
Under the conditions of
\Cref{thm:asymptotic-normality},
if $\griesz \in L_2(\P)$
and $\tilde\gamma \to \griesz$ in $L_2(\P)$,
then $\tilde\sigma^2 \to \sigma^2 = \P\varphi^2$.
\end{corollary}

When $\griesz \notin L_2(\P)$---as
in the treatment-specific mean
without overlap---no square-integrable
efficient influence function exists.
The CLT still holds
with the regularization-dependent
variance $\tilde\sigma^2$:
the confidence interval is valid,
but wider than what overlap would give.

\subsubsection{Proofs}

\begin{proof}[Proof of Corollary~\ref{cor:efficient-variance}]
$L_2$ convergence
$\tilde\gamma \to \griesz$ gives
$\P\tilde\gamma^2\nu \to \P\griesz^2\nu$
and $\P\tilde m^2 \to \P m^2$
for $m = \griesz\mu_\perp + \psiZ(\mu) - \psi$,
so $\tilde\sigma^2 \to \sigma^2$.
\end{proof}

\begin{proof}[Proof of Proposition~\ref{prop:variance-consistency}]
Write $U_i
= \hgamma_i\{Y_i - \hat\mu^\sigma(Z_i)\}
  + \psiZ(\hat\mu^\sigma)(Z_i)$
and $\bar U = \hat\psi$,
so $n\hat\sigma^2 = \hP(U - \bar U)^2$.
Define the discrepancy
$\Delta_i = U_i - (V_i + \psi)$; explicitly,
\[
\Delta_i
= \{\hgamma(Z_i) - \tilde\gamma(Z_i)\}\, \mu_\perp(Z_i)
  - \hgamma(Z_i)\,\delta\mu^\sigma(Z_i)
  + \psiZ(\hat\mu^\sigma)(Z_i) - \psiZ(\mu)(Z_i)
\]
for $\delta\mu^\sigma = \hat\mu^\sigma - \mu$.
We show $\hP\Delta^2 = o_p(1)$.
By the triangle inequality
$\hP\Delta^2 \le 3(I + II + III)$
where
$I = \hP(\hgamma - \tilde\gamma)^2 \mu_\perp^2 = o_p(1)$
by the second part of~\ref{cond:an-D},
and $II + III
= \hP\hgamma^2(\delta\mu^\sigma)^2
  + \hP\{\psiZ(\hat\mu^\sigma) - \psiZ(\mu)\}^2
= o_p(1)$
by hypothesis.

Because $U_i = V_i + \psi + \Delta_i$
and $\bar U = \hP U$,
\[
\hat\sigma^2
= \hP(V - \hP V)^2
+ \hP(\Delta - \hP\Delta)^2
+ 2\hP(V - \hP V)(\Delta - \hP\Delta).
\]
The second term is $o_p(1)$
because $\hP(\Delta - \hP\Delta)^2 \le \hP\Delta^2$.
The cross-term is $o_p(1)$
by Cauchy--Schwarz.

It remains to show
$\hP(V - \hP V)^2 \to_p \tilde\sigma^2$.
Because $\sqrt{n}\,\hP V$
is bounded in distribution by
\Cref{thm:asymptotic-normality},
$(\hP V)^2 \to_p 0$,
so it suffices to show
$\hP V^2 \to_p \tilde\sigma^2$.
Expand
\[
\hP V^2
= \hP\hgamma^2\varepsilon^2
+ 2\,\hP\hgamma\varepsilon \tilde m
+ \hP \tilde m^2.
\]
The first term satisfies
$\hP\hgamma^2\varepsilon^2
= (\hP\hgamma^2\varepsilon^2/\hat\Gamma)
  \cdot \hat\Gamma
\to_p 1 \cdot \P\tilde\gamma^2\nu$
by~\eqref{eq:sum-of-squares} and~\ref{cond:an-C}.
The cross-term vanishes:
conditional on $Z_1, \ldots, Z_n$,
each summand has mean zero,
and the conditional variance
$n^{-2}\sum_i \hgamma^2(Z_i)\, \tilde m^2(Z_i)\, \nu(Z_i)
\le n^{-1}\max_i\{\hgamma^2(Z_i)\nu(Z_i)\}
  \cdot \hP \tilde m^2$
vanishes by~\ref{cond:an-B}
and the law of large numbers.
The third term satisfies
$\hP \tilde m^2 \to_p \P \tilde m^2$
by the law of large numbers.
Combining gives
$\hP V^2 \to_p \P\tilde\gamma^2\nu + \P\tilde m^2
= \tilde\sigma^2$.
\end{proof}

\begin{lemma}[Conditional martingale CLT]
\label{lem:conditional-clt}
Let $\xi_{1,n}, \ldots, \xi_{n,n}$ be a triangular array
of martingale differences
adapted to a filtration $\mathcal{F}_{i,n}$,
with $\sum_{i=1}^n \E[\xi_{i,n}^2 \mid \mathcal{F}_{0,n}] = 1$.
If the conditional Lindeberg condition holds:
\[
\sum_{i=1}^n
  \E[\xi_{i,n}^2\,\ind{\xi_{i,n}^2 \ge \epsilon}
    \mid \mathcal{F}_{0,n}] \to_p 0
\qfor \epsilon > 0,
\]
then $\sum_{i=1}^n \xi_{i,n} \to_d N(0,1)$
conditionally on $\mathcal{F}_{0,n}$.
Furthermore,
the Lindeberg condition implies convergence
of the sum of conditional squares:
$\sum_{i=1}^n \xi_{i,n}^2 \to_p 1$.
\end{lemma}
\begin{proof}
\Citet[Theorem~3.2]{hall2014martingale}
gives the martingale CLT
for arrays with deterministic norming.
The random-norming extension
(conditioning on $\mathcal{F}_{0,n}$
and normalizing by $\hat\Gamma^{1/2}$)
follows from the argument
in \citet[\S3.2(vi)]{hall2014martingale}:
$\{\xi_{i,n}/\hat\Gamma^{1/2},\, \mathcal{F}_{i,n}\}$
is itself a martingale array
satisfying the conditions of Theorem~3.2
conditionally on $\mathcal{F}_{0,n}$.
Convergence of the sum of squares
is \citet[Theorem~2.23]{hall2014martingale}.
\end{proof}

\begin{proof}[Proof of Theorem~\ref{thm:asymptotic-normality}]
Write $\varepsilon_i = Y_i - \E(Y_i \mid Z_i)$
for the centered noise,
$\delta\mu = \hat\mu - \mu$,
and
$\tilde m(Z)
  = \tilde\gamma(Z)\, \mu_\perp(Z)
  + \psiZ(\mu)(Z) - \psi$.

\paragraph{Step 1: reduction to leading terms.}
Expand
\begin{align*}
\hat\psi - \psi
&= R_{\mathrm{lin}}
+ \bigl[\hP\hgamma(\mu - \hat\mu)
  - \dpsiZ(\mu - \hat\mu)\bigr]
+ \hP\hgamma\varepsilon
+ \hP\hgamma \mu_\perp \\
&\quad
+ (\hP - \P)\psiZ(\mu)
+ (\hP - \P)[\psiZ(\hat\mu)
  - \psiZ(\mu)].
\end{align*}
The linearization remainder is
$o_p(n^{-1/2}\hat\Gamma^{1/2})$ by~\ref{cond:an-G},
the imbalance term by~\ref{cond:an-A},
and the empirical process term
$(\hP - \P)[\psiZ(\hat\mu) - \psiZ(\mu)]$
by~\ref{cond:an-F}.
For the projection term,
decompose
$\hP\hgamma\mu_\perp
  = \hP\tilde\gamma\mu_\perp
  + \hP(\hgamma - \tilde\gamma)\mu_\perp$.
The first piece is
$(\hP - \P)\tilde\gamma\mu_\perp
  + \P\tilde\gamma\mu_\perp$,
where $\P\tilde\gamma\mu_\perp
  = o(n^{-1/2}\hat\Gamma^{1/2})$
by~\ref{cond:an-E}.
For the second piece,
$\P(\hgamma - \tilde\gamma)\mu_\perp
  = o_p(n^{-1/2}\hat\Gamma^{1/2})$
by the first part of~\ref{cond:an-D},
and $\hP(\hgamma - \tilde\gamma)^2\mu_\perp^2 \to_p 0$
by the second part of~\ref{cond:an-D},
so $(\hP - \P)(\hgamma - \tilde\gamma)\mu_\perp
  = o_p(n^{-1/2}\hat\Gamma^{1/2})$ by Chebyshev.
In total,
$\hP\hgamma\mu_\perp
  = (\hP - \P)[\tilde\gamma\mu_\perp]
  + o_p(n^{-1/2}\hat\Gamma^{1/2})$.
Combining with
$(\hP - \P)\psiZ(\mu) = (\hP - \P)[\psiZ(\mu) - \psi]$
and defining
$V_i = \hgamma(Z_i)\varepsilon_i + \tilde m(Z_i)$,
we have
$\hP \tilde m
  = \P\tilde\gamma\mu_\perp
  + (\hP - \P)\tilde m$.
The first term is
$o_p(n^{-1/2}\hat\Gamma^{1/2})$ by~\ref{cond:an-E}, so
\begin{equation}\label{eq:an-leading}
\hat\psi - \psi
= \hP V + o_p(n^{-1/2}\hat\Gamma^{1/2}).
\end{equation}

\paragraph{Step 2: asymptotic normality of $\sqrt{n}\,\hP V$.}
Conditional on $Z_1, \ldots, Z_n$,
the weights $\hgamma(Z_i)$ are fixed
and $\varepsilon_i = Y_i - \E(Y_i \mid Z_i)$
are independent with mean zero
and variance $\nu(Z_i)$.
So $V_i$ has conditional mean $\tilde m_i = \tilde m(Z_i)$
and conditional variance $s_i^2 = \hgamma^2(Z_i)\nu(Z_i)$.
Define the standardized array
\[
\xi_{i,n}
= \frac{\hgamma(Z_i)\varepsilon_i}
       {(n\hat\Gamma)^{1/2}},
\]
a martingale difference array
adapted to
$\mathcal{F}_{i,n}
= \sigma\{Z_1, \ldots, Z_n,
  \varepsilon_1, \ldots, \varepsilon_i\}$,
standardized:
$\sum_{i=1}^n \E[\xi_{i,n}^2 \mid Z_1, \ldots, Z_n]
= 1$.

\textit{Lindeberg condition.}
We verify
$\sum_{i=1}^n
  \E[\xi_{i,n}^2\,\ind{\xi_{i,n}^2 \ge \epsilon}
    \mid Z_1, \ldots, Z_n] \to_p 0$
for all $\epsilon > 0$.
Split the indicator at
$\varepsilon_i^2 = M$
using $\ind{\hgamma_i^2\varepsilon_i^2
  \ge \epsilon n\hat\Gamma}
  \le \ind{\varepsilon_i^2 \ge M}
  + \ind{\hgamma_i^2 M \ge \epsilon n\hat\Gamma}$:
\begin{align*}
&\sum_{i=1}^n
  \E[\xi_{i,n}^2\,\ind{\xi_{i,n}^2 \ge \epsilon}
    \mid Z_1, \ldots, Z_n] \\
&\quad= \frac{1}{n\hat\Gamma}\sum_{i=1}^n
  \hgamma_i^2\,
  \E\qty[\varepsilon_i^2\,
    \ind{\hgamma_i^2\varepsilon_i^2
      \ge \epsilon\, n\hat\Gamma}
    \mid Z_1, \ldots, Z_n] \\
&\quad\le \frac{\hP\hgamma^2}{\hat\Gamma}
  \cdot \max_{i \le n}
  \E\qty[\varepsilon_i^2\,
    \ind{\varepsilon_i^2 \ge M}
    \mid Z_1, \ldots, Z_n]
  \;+\;
  \frac{M\,\hP\hgamma^2}{\hat\Gamma}
  \cdot \max_{i \le n}
  \ind{\hgamma_i^2 M \ge \epsilon\, n\hat\Gamma}.
\end{align*}
The leading ratio
$\hP\hgamma^2 / \hat\Gamma \le 1/\underline\nu$
is bounded,
where $\underline\nu = \inf\{\nu(Z) : Z \in \mathrm{supp}(\hgamma^2)\} > 0$
(the conditional variance is positive
wherever the weights are nonzero).
The first bound's second factor
vanishes for large $M$
by~\eqref{eq:uniform-integrability}.
For fixed $M$,
the second bound vanishes
because
$\max_i \hgamma_i^2 / (n\hat\Gamma)
\le \max_i \hgamma_i^2\nu_i
  / (n\hat\Gamma\underline\nu) \to_p 0$
by~\ref{cond:an-B}.

The conditional Lindeberg CLT
(\Cref{lem:conditional-clt})
gives
\begin{equation}\label{eq:conditional-clt}
\sum_{i=1}^n \xi_{i,n}
\;\Big|\; Z_1, \ldots, Z_n
\;\to_d\; N(0,1),
\end{equation}
i.e.\
$\sqrt{n}\,\hP\hgamma\varepsilon
/ \hat\Gamma^{1/2}
\to_d N(0,1)$
conditional on $Z_1, \ldots, Z_n$.
The Lindeberg condition
also implies convergence
of the sum of squares
(\Cref{lem:conditional-clt}):
\begin{equation}\label{eq:sum-of-squares}
\hP\hgamma^2\varepsilon^2 / \hat\Gamma
\to_p 1.
\end{equation}

\textit{Unconditional conclusion.}
Because the conditional convergence
in~\eqref{eq:conditional-clt}
holds in probability over $Z_1, \ldots, Z_n$,
the unconditional convergence holds as well.
The shift
$\sqrt{n}\,\hP \tilde m
= \sqrt{n}\,\P\tilde\gamma\mu_\perp
+ \sqrt{n}\,(\hP - \P)\tilde m$
is measurable with respect to $Z_1, \ldots, Z_n$.
The bias $\sqrt{n}\,\P\tilde\gamma\mu_\perp$
is $o(\hat\Gamma^{1/2})$ by~\ref{cond:an-E},
and by the CLT,
$\sqrt{n}\,(\hP - \P)\tilde m
  \to_d N(0, \Var(\tilde m))$.
Because $\P\tilde m = \P\tilde\gamma\mu_\perp
  = o(n^{-1/2})$ by~\ref{cond:an-E},
Slutsky's lemma gives
$\sqrt{n}\,\hP\tilde m
  \to_d N(0, \P\tilde m^2)$.
Conditional on $Z_1, \ldots, Z_n$,
$\sqrt{n}\,\hP V
= \sqrt{n}\,\hP\hgamma\varepsilon
+ \sqrt{n}\,\hP \tilde m$
is the noise
(conditionally normal)
plus a fixed shift.
By~\ref{cond:an-C}, $\hat\Gamma \to_p \P\tilde\gamma^2\nu$.
Because
$\E[\varepsilon_i \mid Z_i] = 0$,
the noise and shift
are uncorrelated---$\E[\hgamma_i\varepsilon_i \mid Z_1, \ldots, Z_n] = 0$---so
the conditional characteristic function
of the sum factorizes,
and the unconditional limit is
\[
N(0,\, \P\tilde\gamma^2\nu + \P\tilde m^2)
= N(0,\, \tilde\sigma^2). \qedhere
\]
\end{proof}

\subsection{Penalized Bregman weights}
\label{subsec:penalized-bregman}

We now specialize to the setting
of the rest of the paper:
the penalty seminorm $\rho$
admits a Sobolev interpolation inequality
(\Cref{lem:sobolev-interpolation})
with exponent $\theta = d/(2s)$,
the regularization is quadratic
$\zetas(\rho) = \eta\rho^2$,
and the dispersion $\chis$
is one of Settings~Q, Q$_+$, or~E.
The abstract conditions~\ref{cond:an-A}--\ref{cond:an-C}
of \Cref{thm:asymptotic-normality}
involve the empirical weights $\hgamma$ directly;
here we translate them
into conditions on the regularization parameter~$\eta$
and the population projection~$\tilde\phi$,
which the analyst controls
and the penalty bias machinery
of \Cref{sec:penalty-bias} bounds.
The remaining conditions~\ref{cond:an-D}--\ref{cond:an-G}
depend on the structure
of the functional and the model
and are verified application by application.

\Cref{cor:sobolev-main}
gives the convergence rate
\[
\hP\, D_{\chis}(\hat\phi, \tilde\phi)
  + \eta\,\rho^2(\hat\phi - \tilde\phi)
  \le \frac{C\, d_\eta}{\delta\, n}.
\]
The constant $C$ depends on
the curvature ratio
$\overline{\ddchis}/\underline{\ddchis}$,
the PSC parameter~$\beta$,
and the interpolation constant~$\kappa$;
$d_\eta \le C_s\,\eta^{-\theta}$
for Sobolev eigenvalues.

\begin{lemma}[Bias rate]
\label{lem:bias-rate}
For the quadratic dispersion
$\chis(\phi) = \phi^2/2$,
suppose $\griesz \in L_2(\P)$.
For any competitor $f$
with $\rho(f) < \infty$,
write
$r_f = 2\sqrt{\|f - \griesz\|_{\LtP}^2/(2\eta)
  + \rho^2(f)} + 1$.
Then $r_f > 2\rho(\tilde\phi)$ and
\begin{equation}
\label{eq:bias-rate}
|b(\hgamma)|
  \le \tfrac{5}{2}\,\Xi(r_f)
  + 2\sqrt{2}\,
  \|f - \griesz\|_{\LtP}\sqrt{\eta}
  + 4\eta\,\rho(f).
\end{equation}
At $f = 0$,
$r_0 = \|\griesz\|_{\LtP}\sqrt{2/\eta} + 1$
and the bound reduces to
$\tfrac{5}{2}\,\Xi(r_0)
  + 2\sqrt{2}\,\|\griesz\|_{\LtP}\sqrt{\eta}$.
When $\griesz \in \mathcal{H}$,
taking $f = \griesz$ gives
$\tfrac{5}{2}\,\Xi(2\rho(\griesz)+1)
  + 4\eta\,\rho(\griesz)
  = O(\eta)$.
\end{lemma}

\begin{proof}
The comparison bound~\eqref{eq:comparison-bound}
at competitor~$f$ gives
$\eta\,\rho^2(\tilde\phi)
  \le \tfrac{1}{2}\,\|f - \griesz\|_{\LtP}^2
  + \eta\,\rho^2(f)$,
so $\rho(\tilde\phi)
  \le \sqrt{\|f-\griesz\|_{\LtP}^2/(2\eta) + \rho^2(f)}$
and $r_f > 2\rho(\tilde\phi)$.
\Cref{lem:offset-convex} at $r = r_f$ gives
$|b(\hgamma)|
  \le \tfrac{5}{2}\,\Xi(r_f)
  + 4\eta\,\rho(\tilde\phi)$.
The penalty bias:
$4\eta\,\rho(\tilde\phi)
  \le 4\sqrt{\eta\,\|f-\griesz\|_{\LtP}^2/2
    + \eta^2\,\rho^2(f)}
  \le 2\sqrt{2}\,\|f-\griesz\|_{\LtP}\sqrt{\eta}
  + 4\eta\,\rho(f)$.
\end{proof}

\begin{theorem}[Asymptotic normality
  for penalized Bregman weights]
\label{thm:general-clt}
In the setting above,
suppose $\griesz \in L_2(\P)$:
\begin{enumerate}[label=\textup{(\alph*)}]
\item\label{cond:clt-a}
\textit{Imbalance.}
For some competitor $f$ with $\rho(f) < \infty$
and some decomposition
$\hat\mu - \mu = g_n + r_n$
with $\rho(g_n) < \infty$,
\begin{equation}
\label{eq:clt-a}
[\|f - \griesz\|_{\LtP}\sqrt{\eta}
  + \eta\,\rho(f)]
  \cdot \rho(g_n)
  + [\|f - \griesz\|_{\LtP}
  + \sqrt{\eta}\,\rho(f)]
  \cdot \|r_n\|_{\LtP}
  + \frac{\|\griesz\,r_n\|_{\LtP}}{\sqrt{n}}
  = o_p(n^{-1/2}).
\end{equation}
When $\hat\mu - \mu \in \operatorname{dom}(\rho)$,
take $g_n = \hat\mu - \mu$, $r_n = 0$.
\item\label{cond:clt-c}
\textit{Infinitesimal weights.}
$\|\tilde\gamma\|_{L_\infty}^2
  = o(n\,\P\tilde\gamma^2\nu)$.
Combined with~\ref{cond:clt-d},
this holds whenever
$\|\tilde\gamma\|_{L_\infty}^2 = o(n)$.
\item\label{cond:clt-d}
\textit{Denominator.}
$\P\tilde\gamma^2\nu \ge c > 0$
uniformly in $\eta$.
\item\label{cond:clt-e}
Conditions~\ref{cond:an-D}--\ref{cond:an-G}
of \Cref{thm:asymptotic-normality}
hold.
\end{enumerate}
Then
$\sqrt{n}\,(\hat\psi - \psi)/\hat\sigma
  \to_d N(0,1)$.
If additionally
$\tilde\gamma \to \griesz$ in $\LtP$,
then $\hat\sigma^2/\sigma^2 \to_p 1$
for the semiparametric efficiency bound
$\sigma^2 = \P\varphi^2$.
\end{theorem}

\begin{proof}
Condition~\ref{cond:an-A} of
\Cref{thm:asymptotic-normality}
requires
$\hP\hgamma(\hat\mu - \mu)
  - \dpsiZ(\hat\mu - \mu)
  = o_p(n^{-1/2}\hat\Gamma^{1/2})$.
Split $\hat\mu - \mu = g_n + r_n$
as in~\ref{cond:clt-a}
and bound the left side
by its smooth and rough parts.

\paragraph{Smooth part.}
Because $g_n \in \operatorname{dom}(\rho)$,
$g_n/\rho(g_n) \in B_\rho$
and the bias functional gives
$|\hP\hgamma\,g_n - \dpsiZ(g_n)|
  \le |b(\hgamma)|\,\rho(g_n)$.
\Cref{lem:bias-rate}
at the competitor $f$ from~\ref{cond:clt-a} gives
$|b(\hgamma)|
  \le \tfrac{5}{2}\,\Xi(r_f)
  + C[\|f - \griesz\|_{\LtP}\sqrt{\eta}
    + \eta\,\rho(f)]$.
The penalty bias contributes
$C[\|f-\griesz\|_{\LtP}\sqrt{\eta}
  + \eta\,\rho(f)]
  \,\rho(g_n)
  = o_p(n^{-1/2})$
by the first term of~\eqref{eq:clt-a}.
The offset complexity contributes
$\tfrac{5}{2}\,\Xi(r_f)\,\rho(g_n)$;
since $r_f \ge 1$,
$\Xi(r_f) = o(n^{-1/2})$,
and $\rho(g_n)$ grows at most polynomially
in $n$ (from~\eqref{eq:clt-a}
with nonvanishing coefficient),
the product is $o_p(n^{-1/2})$.

\paragraph{Rough part.}
Write
$\hP\hgamma\,r_n - \dpsiZ(r_n)
  = \hP(\hgamma - \griesz)\,r_n
  + (\hP - \P)\griesz\,r_n$.
By Cauchy--Schwarz on $\hP$,
$|\hP(\hgamma - \griesz)\,r_n|
  \le \|\hgamma - \griesz\|_{L_2(\hP)}\,
  \|r_n\|_{L_2(\hP)}$.
The comparison bound~\eqref{eq:comparison-bound}
at competitor $f$ gives
$\|\tilde\gamma - \griesz\|_{\LtP}
  \le \|f - \griesz\|_{\LtP}
  + \sqrt{2\eta}\,\rho(f)$,
and $\|\hgamma - \tilde\gamma\|_{\LtP} \to 0$
by the Bregman rate,
so
$\|\hgamma - \griesz\|_{L_2(\hP)}
  \le C[\|f - \griesz\|_{\LtP}
  + \sqrt{\eta}\,\rho(f)]
  + o_p(1)$.
The product with $\|r_n\|_{L_2(\hP)}$
is $o_p(n^{-1/2})$
by the second term of~\eqref{eq:clt-a}.
The centering term
$(\hP - \P)\griesz\,r_n
  = O_p(\|\griesz\,r_n\|_{\LtP}/\sqrt{n})
  = o_p(n^{-1/2})$
by the third term of~\eqref{eq:clt-a}.
So~\ref{cond:an-A} follows from~\ref{cond:clt-a}.

Condition~\ref{cond:an-B}
requires
$\max_i \hgamma^2(Z_i)\,\nu(Z_i)
  / (n\,\hat\Gamma) \to_p 0$.
Because $\hat\Gamma \to_p \P\tilde\gamma^2\nu$
by condition~\ref{cond:an-C} below,
and $\max_i \hgamma^2(Z_i)
  \le (1 + o_p(1))\,\|\tilde\gamma\|_{L_\infty}^2$
(by the rate bound,
$\rho(\hat\phi - \tilde\phi) \le \rho_s$
and $\|\hat\phi - \tilde\phi\|_{\chis} < r$;
the Sobolev interpolation
$\|f\|_{L_\infty}
  \le \kappa\,\rho(f)^\theta\,\|f\|_{L_2}^{1-\theta}$
applied to $f = \hat\phi - \tilde\phi$
gives $\|\hat\phi - \tilde\phi\|_{L_\infty} \to 0$;
for the quadratic family
$\|\hgamma - \tilde\gamma\|_{L_\infty}
  \le \|\hat\phi - \tilde\phi\|_{L_\infty}$;
the empirical max $\max_i \hgamma^2(Z_i)
  \le \|\hgamma\|_{L_\infty}^2
  \le (\|\tilde\gamma\|_{L_\infty} + o(1))^2$),
\ref{cond:an-B} follows from~\ref{cond:clt-c} and~\ref{cond:clt-d}.

Condition~\ref{cond:an-C}
requires
$\hat\Gamma / \P\tilde\gamma^2\nu \to_p 1$.
For the quadratic and positive-part
quadratic dispersions ($\hgamma = \hat\phi$
or $\hat\phi_+$),
the Bregman convergence rate
gives $\hP(\hgamma - \tilde\gamma)^2 \to_p 0$,
so~\ref{cond:an-C} follows from
the expansion
$\hat\Gamma
  = \hP\tilde\gamma^2\nu
  + 2\,\hP\tilde\gamma(\hgamma - \tilde\gamma)\nu
  + \hP(\hgamma - \tilde\gamma)^2\nu$
and Cauchy--Schwarz.
For entropy-like dispersions,
\Cref{cor:variance-interpolation}
gives~\ref{cond:an-C}
under the interpolation inequality~\eqref{eq:interpolation}
and pseudo-self-concordance.

The CLT follows from
\Cref{thm:asymptotic-normality}.
The efficiency claim follows from
\Cref{cor:efficient-variance}.
\end{proof}

For Setting~Q
($\tilde\gamma = \tilde\phi$),
condition~\ref{cond:clt-c} reduces to a growth rate
on the seminorm:
the Sobolev interpolation inequality
$\|\tilde\phi\|_{L_\infty}
  \le \kappa\,\rho(\tilde\phi)^\theta\,
    (\P\tilde\phi^2)^{(1-\theta)/2}$
gives
$\|\tilde\phi\|_{L_\infty}^2
  / (n\,\P\tilde\phi^2)
  \le \kappa^2\,
    \rho(\tilde\phi)^{2\theta}\,
    (\P\tilde\phi^2)^{-\theta} / n$,
so condition~\ref{cond:clt-c} holds when
$\rho(\tilde\phi)^{2\theta} = o(n)$
--- the slower the seminorm grows,
the weaker the requirement.
For Setting~E
($\tilde\gamma = e^{\tau\tilde\phi}$),
condition~\ref{cond:clt-c} becomes
$e^{2\tau\|\tilde\phi\|_{L_\infty}}
  = o(n\,\P e^{2\tau\tilde\phi}\,\nu)$,
an exponential growth constraint.

\begin{corollary}[Penalized least-squares outcome estimator]
\label{cor:outcome-estimator}
Let $\mu = \E(Y \mid Z)$
be the conditional mean,
and let $\hat\mu$ be
the penalized least-squares estimator
\begin{equation}\label{eq:general-outcome}
\hat\mu
  = \argmin_{m}
    \hP (Y - m(Z))^2/2
    + \eta_\mu\, \rho^2(m)/2
\end{equation}
with $\eta_\mu = n^{-b}$
for $b < 1/(1+\theta)$.
If $\rho(\mu) < \infty$,
then $\rho(\hat\mu - \mu)$
is bounded in probability.
\end{corollary}

\begin{proof}
The argument has two steps:
the population comparison bound
controls the regularized population minimizer,
and the Bregman convergence bound
controls the empirical fluctuation.

\paragraph{Population comparison.}
The population loss is
$L(m) = \P(Y - m)^2/2
  + \eta_\mu\,\rho^2(m)/2$.
Because $\tilde\mu$ minimizes $L$,
$L(\tilde\mu) \le L(\mu)$.
The tower property
$\P(Y - \mu)(m - \mu)
  = \P\,\E(Y - \mu \mid Z)\cdot(m - \mu)
  = 0$
gives the decomposition
$\P(Y - m)^2
  = \P(Y - \mu)^2 + \P(\mu - m)^2$,
so the comparison reduces to
\[
\P(\tilde\mu - \mu)^2/2
  + \eta_\mu\,\rho^2(\tilde\mu)/2
  \le \eta_\mu\,\rho^2(\mu)/2.
\]
Both terms on the left are nonnegative,
so in particular
$\rho(\tilde\mu) \le \rho(\mu)$:
the regularized population minimizer
is no rougher than the truth.

\paragraph{Empirical fluctuation.}
\Cref{cor:sobolev-main}
applied to the outcome problem
--- with quadratic dispersion
($\overline{\ddchis}/\underline{\ddchis} = 1$) ---
gives
\[
\eta_\mu\,\rho^2(\hat\mu - \tilde\mu)/2
  \le \frac{C\, d_{\eta_\mu}}{\delta\, n}
  \le \frac{C\, C_s\, \eta_\mu^{-\theta}}
       {\delta\, n}
\]
with probability at least $1 - \delta$.
With $\eta_\mu = n^{-b}$, the right side is
$C'\,n^{b\theta - 1}/\delta$,
so
$\rho^2(\hat\mu - \tilde\mu)
  \le C''\,n^{b(1+\theta) - 1}/\delta$,
which is bounded for any fixed $\delta$
when $b(1+\theta) < 1$.

\paragraph{Conclusion.}
By the triangle inequality,
$\rho(\hat\mu - \mu)
  \le \rho(\hat\mu - \tilde\mu)
    + \rho(\tilde\mu - \mu)$.
The first term is bounded in probability
by the empirical fluctuation bound.
The second is at most $2\rho(\mu)$
by the population comparison.
Both are bounded,
so $\rho(\hat\mu - \mu)$
is bounded in probability.
\end{proof}

\begin{remark}
\label{rem:arm-specific-outcome}
For arm-specific estimation ---
$\hat\mu_w$ estimated from
the subsample with $W = w$ ---
the same argument applies with
$n$ replaced by $n_w$.
Because $n_w/n \to p_w > 0$
by the law of large numbers,
the condition
$b < 1/(1+\theta)$
is unchanged.
\end{remark}

\begin{lemma}[Denominator anchoring]
\label{lem:denominator-anchoring}
Suppose the quadratic dispersion
$\chis(\phi) = \phi^2/2$,
so $\tilde\gamma = \tilde\phi$.
If there exists
$h_0 \in \ker(\rho)$
with $\dotpsiZ(h_0) \ne 0$,
then
\begin{equation}
\label{eq:denominator-anchoring}
\P\tilde\phi^2
  \ge \frac{\dotpsiZ(h_0)^2}
           {\|h_0\|_{L_\infty}^2}
  > 0,
\end{equation}
uniformly in $\eta$.
\end{lemma}

\begin{proof}
Because $h_0 \in \ker(\rho)$,
the penalty term vanishes
in the first-order condition~\eqref{eq:foc-quadratic}
for $\tilde\phi$,
giving
$\P\,\tilde\phi\, h_0
  = \dotpsiZ(h_0)$.
By Cauchy--Schwarz,
$\dotpsiZ(h_0)^2
  = (\P\,\tilde\phi\, h_0)^2
  \le \P\tilde\phi^2\,
      \|h_0\|_{L_\infty}^2$,
which rearranges
to~\eqref{eq:denominator-anchoring}.
\end{proof}

For the treatment-specific mean,
$h_0 = 1 \in \ker(\rho)$
and $\dotpsiZ(1) = p_w$,
giving $\P\tilde\phi^2 \ge p_w^2$.
For average linear regression
with components $h(W,X) = a(X) + W\,b(X)$,
the function $h_0 = (0,1)$
lies in $\ker(\rho)$
and $\dotpsiZ((0,1)) = 1$.
For the CATE variance,
$\dotpsiZ^\mu(h)(Z)
  = 2\tau(X)\{h(1,X) - h(0,X)\}$
kills constants,
but $h_0(w,x) = w$
lies in $\ker(\rho)$
(constant in $x$)
with $\dotpsiZ^\mu(h_0)
  = 2\E\tau(X)$.
When $\E\tau \ne 0$,
the denominator is anchored
by $\P\tilde\phi^2
  \ge (2\E\tau)^2$;
this requires
heterogeneous effects,
not overlap.
When $\E\tau = 0$,
the influence function for $\Var(\tau(X))$
degenerates
and all efficient estimators
may fail to be asymptotically normal
\citep{sanchez-becerra2023};
the comparison bound remains a fallback
when $\griesz \in L_2(\P)$
and $\eta$ is small enough.

For the quadratic family,
the balance equation~\eqref{eq:balance-equation}
(\Cref{prop:penalty-bias-balance})
and the Sobolev interpolation inequality
together determine a two-sided window
for~$\eta$:
the lower edge comes
from the imbalance condition~\ref{cond:clt-a},
and the upper edge
from the sup-norm growth condition~\ref{cond:clt-c}.

\begin{corollary}[Tuning window, quadratic family]
\label{cor:tuning-window}
For the quadratic dispersion
$\chis(\phi) = \phi^2/2$
with Sobolev penalty,
suppose $\griesz \in L_2(\P)$.
Suppose the denominator is anchored:
$\P\tilde\phi^2 \ge c > 0$
uniformly in $\eta$.
If additionally
$\rho(\hat\mu - \mu) = O_p(1)$
and conditions~\ref{cond:an-D}--\ref{cond:an-G} hold,
then writing $\eta = n^{-a}$,
the CLT of
\Cref{thm:general-clt}
holds whenever
\begin{equation}
\label{eq:general-lambda-window}
1 < a < \frac{1}{\theta},
\end{equation}
which is nonempty for all $\theta < 1$
(equivalently $s > d/2$).

Under overlap
--- $\griesz \in \mathcal{H}$
and $\|\tilde\gamma\|_{L_\infty}
  \le \|\griesz\|_{L_\infty}$ ---
taking $f = \griesz$,
$g_n = \hat\mu - \mu$, $r_n = 0$
in~\ref{cond:clt-a}:
$\|f - \griesz\|_{\LtP} = 0$,
so~\eqref{eq:clt-a} reduces to
$\eta\,\rho(\griesz)\,\rho(\hat\mu-\mu)
  = o_p(n^{-1/2})$,
which holds whenever $\eta \to 0$.
Condition~(b) is automatic,
and the window
becomes $\eta \to 0$ with $a\theta < 1$.
\Cref{cor:efficient-variance}
gives $\tilde\sigma^2 \to \sigma^2$:
the CLT is both valid and efficient.
\end{corollary}

\begin{proof}
For condition~\ref{cond:clt-a},
take $f = 0$, $g_n = \hat\mu - \mu$, $r_n = 0$.
Write $G = \|\griesz\|_{\LtP}$.
The condition~\eqref{eq:clt-a} reduces to
$G\sqrt{\eta}\,\rho(\hat\mu - \mu)
  = o_p(n^{-1/2})$.
Because $\rho(\hat\mu - \mu) = O_p(1)$,
this holds
when $\sqrt{\eta} = o(n^{-1/2})$,
i.e.\ $a > 1$.

For condition~\ref{cond:clt-c},
the Sobolev interpolation inequality gives
$\|\tilde\phi\|_{L_\infty}^2
  \le \kappa^2\,\rho(\tilde\phi)^{2\theta}\,
      (\P\tilde\phi^2)^{1-\theta}$.
From \Cref{prop:penalty-bias-balance},
$\rho(\tilde\phi)^{2\theta}
  \le C\,\eta^{-\theta}$,
so
$\|\tilde\phi\|_{L_\infty}^2
  / (n\,\P\tilde\phi^2)
  \le C\,\eta^{-\theta}
       / (n\,c^\theta)$.
Condition~\ref{cond:clt-c} holds when
$n\,\eta^{\theta} \to \infty$,
i.e.\ $a < 1/\theta$.

The window $1 < a < 1/\theta$
is nonempty when $\theta < 1$,
i.e.\ $s > d/2$.

Under overlap ($\griesz \in \mathcal{H}$),
taking $f = \griesz$ in~\ref{cond:clt-a}:
the approximation term vanishes
and the penalty term is
$\eta\,\rho(\griesz)\,\rho(\hat\mu-\mu) = o_p(n^{-1/2})$
whenever $\eta \to 0$.
The sup-norm $\|\tilde\gamma\|_{L_\infty}$
is bounded,
so~\ref{cond:clt-c} holds with $C/n$.
The window reduces to $\eta \to 0$
with $a\theta < 1$
(from the effective dimension
$d_\eta \ll n$).
\end{proof}


\begin{figure}[ht]
\centering
\setlength{\tabcolsep}{4pt}
\renewcommand{\arraystretch}{1.6}
\begin{tabular}{r@{\quad}l@{\quad}l}
& \textbf{Window for $a$} & \textbf{Condition} \\[2pt]
\hline\\[-8pt]
\textup{(i)} balance equation
  & $1 < a < 1/\theta$
  & $\griesz \in L_2$ \\[4pt]
\textup{(ii)} comparison
  & $\dfrac{s}{s+s_\gamma}
    < a <
    \dfrac{s}{\theta(s - s_\gamma)}$
  & $\griesz \in H^{s_\gamma}$, $s_\gamma < s$ \\[6pt]
\textup{(iii)} full
  & $1/2 < a < 1/\theta$
  & $\griesz \in H^s$ \\[2pt]
\end{tabular}
\caption{Nested tuning windows for the penalty exponent~$a$
  ($\eta = n^{-a}$)
  in the quadratic-dispersion setting,
  with $\theta = d/(2s)$.
  Each level of regularity on the Riesz representer
  widens the window.
  Window~\textup{(i)} uses
  the balance equation
  (\Cref{prop:penalty-bias-balance})
  and requires only $\griesz \in L_2$.}
\label{fig:tuning-windows}
\end{figure}

\subsection{Self-contained Sobolev CLT}
\label{sec:sobolev-clt}

The following theorem assembles
the rate (\Cref{cor:sobolev-main}),
bias (\Cref{cor:offset-positive-part}),
variance equivalence (\Cref{sec:variance-equiv}),
and asymptotic normality (\Cref{thm:general-clt})
into a single usable statement
for Sobolev models
with quadratic regularization.

\begin{theorem}[Sobolev CLT for penalized Bregman weights]
\label{thm:sobolev-clt}
Let $\rho(h) = \|(-\Delta)^{s/2} h\|_{L_2}$
for integer $s > d/2$,
$\zetas(\rho) = \eta\rho^2$
with $\eta = n^{-a}$,
$\theta = d/(2s)$,
and $X$ with density $p \ge \pmin > 0$
on $[0,1]^d$.
Let $\chis$ be one of
Settings~Q, Q$_+$, or~E.
Suppose:
\begin{enumerate}[label=\textup{(\Alph*)}]
\item\label{cond:sob-overlap}
\textit{Overlap.}
The derivative functional is $L_2$-continuous:
$|\dotpsiZ(h)| \le C_\psi\,\|h\|_{L_2}$
for all $h \in L_2$,
equivalently $\griesz \in L_2(\P)$.
\item\label{cond:sob-denom}
\textit{Denominator anchoring.}
$\dotpsiZ(1) \neq 0$.
\item\label{cond:sob-dispersion}
For Setting~E:
$\log(\griesz)/\tau \in H^s$
(the unpenalized index has finite Sobolev seminorm).
\item\label{cond:sob-misc}
For nonlinear or misspecified functionals:
conditions~\ref{cond:an-D}--\ref{cond:an-G}
of \Cref{thm:asymptotic-normality}.
(For linear functionals with well-specified models,
these are vacuous.)
\end{enumerate}

Then the tuning window
\begin{equation}\label{eq:sobolev-clt-window}
1 < a < \frac{1}{\theta}
\end{equation}
is nonempty when $\theta < 1$,
equivalently $s > d/2$.
Suppose the outcome estimator $\hat\mu$
uses the same Sobolev seminorm $\rho$,
with penalty $\eta_\mu = n^{-b}$
and $\mu \in H^{s_\mu}$ for some $s_\mu > 0$.
For $\eta = n^{-a}$
in this window,
with $b$ satisfying the conditions
of \Cref{prop:outcome-clt-condition}
(an estimation-rate condition
$b(1+\theta) < a$
and a double-robustness condition
on $s_\mu$ and the regularity of $\griesz$):
\begin{enumerate}[label=\textup{(\roman*)}]
\item \textit{Rate.}
$\hP\,D_{\chis}(\hat\phi, \tilde\phi)
  + \eta\,\rho^2(\hat\phi - \tilde\phi)
  \le C\,d_\eta / (\delta\,n)$
with probability $\ge 1 - \delta$,
where $d_\eta = \sum_j \lambda_j/\max(\lambda_j, \eta)
  \asymp \eta^{-\theta}$.
\item \textit{Bias.}
$|b(\hgamma)| = o(n^{-1/2})$.
\item \textit{CLT.}
$\sqrt{n}\,(\hat\psi - \psi)/\hat\sigma
  \to_d N(0,1)$.
\item \textit{Efficiency.}
Under overlap,
$\tilde\gamma \to \griesz$ in $\LtP$
(by the comparison bound;
see the proof below),
and $\hat\sigma^2/\sigma^2 \to_p 1$
for the semiparametric efficiency bound.
\end{enumerate}
\end{theorem}

\begin{proof}
The argument has four stages:
denominator anchoring,
weight convergence,
bias control, and the CLT\@.

\paragraph{Denominator anchoring.}
Since $1 \in \ker(\rho)$,
the first-order condition at $h = 1$ gives
$\P\dchis(\tilde\phi) = \dotpsiZ(1) \neq 0$
by~\ref{cond:sob-denom}.
For Setting~Q,
$\dchis(\tilde\phi) = \tilde\phi$,
so $\P\tilde\phi = \dotpsiZ(1)$
and by Jensen's inequality
$\P\tilde\phi^2 \ge (\P\tilde\phi)^2
  = (\dotpsiZ(1))^2 > 0$
uniformly in $\eta$
(\Cref{lem:denominator-anchoring}).

\paragraph{Weight convergence.}
Sobolev eigenvalues $\lambda_j \asymp j^{-2s/d}$
give effective dimension
$d_\eta \asymp \eta^{-\theta}$.
\Cref{cor:sobolev-main} gives the rate
$\hP D_{\chis}(\hat\phi, \tilde\phi)
  + \eta\,\rho^2(\hat\phi - \tilde\phi)
  \le C\,d_\eta/(\delta n)$
with probability $\ge 1 - \delta$,
provided $n\pmin^{1-\theta}\eta^\theta$
exceeds a constant
times $(\overline{\ddchis}/\underline{\ddchis})\,
  \log(2/\delta)$.
For the quadratic family
$\overline{\ddchis}/\underline{\ddchis} = 1$.
For Setting~E,
the comparison bound
$\rho(\tilde\phi) \le \rho(\phi_\infty) < \infty$
(by~\ref{cond:sob-dispersion})
and the denominator bound
$\P\tilde\phi^2 \ge c > 0$
give $\|\tilde\phi\|_{L_2}$ bounded;
Sobolev embedding then gives
$\|\tilde\phi\|_{L_\infty}
  \le \kappa\,\rho(\tilde\phi)^\theta\,
  \|\tilde\phi\|_{L_2}^{1-\theta}
  \le C$
uniformly in $\eta$,
so the curvature ratio
$e^{2\tau\|\tilde\phi\|_\infty}$ is bounded.

\paragraph{Bias control and the tuning window.}
The bias splits
(\Cref{lem:offset-convex})
into an offset complexity term
and a penalty bias term.
The offset complexity
$\Xi(r) \le C\, n^{-1/2}
  (n^{1/2} r)^{-\theta/(2-\theta)}$
is $o(n^{-1/2})$ for any $r > 2\rho(\tilde\phi)$,
with no condition on $\eta$.

The penalty bias is where
the tuning window enters.
\Cref{prop:penalty-bias-balance}
(applied to~\ref{cond:sob-overlap})
gives
\begin{equation}\label{eq:sobolev-clt-balance}
\P\tilde\phi^2 \le \|\griesz\|_{\LtP}^2,
\qquad
\eta\,\rho(\tilde\phi) \le \frac{\|\griesz\|_{\LtP}}{\sqrt{2}}\,\sqrt{\eta}.
\end{equation}
(For Q$_+$, $\P\tilde\phi_+\tilde\phi
  = \P\tilde\phi_+^2$
plays the same role;
see \Cref{rem:positive-part-penalty-bias}.)
Writing $G = \|\griesz\|_{\LtP}$,
\begin{equation}\label{eq:sobolev-penalty-bias}
\eta\,\rho(\tilde\phi)
  \le \frac{G}{\sqrt{2}}\,\sqrt{\eta}.
\end{equation}
This is $o(n^{-1/2})$ when
$a > 1$:
the lower edge of the window.
For Setting~E
the comparison bound gives
$\rho(\tilde\phi) \le \rho(\phi_\infty)$,
so the penalty bias is $O(\eta)$,
which is sharper.
(For Setting~E,
\Cref{rem:balance-fails-E} explains
why the balance equation does not apply;
the comparison bound replaces it.)

For the upper edge:
condition~\ref{cond:clt-c} requires
$\|\tilde\gamma\|_{L_\infty}^2
  = o(n\,\P\tilde\gamma^2\nu)$.
The interpolation inequality gives
$\|\tilde\phi\|_{L_\infty}^2
  \le C\,\rho(\tilde\phi)^{2\theta}\,
  (\P\tilde\phi^2)^{1-\theta}$,
so
\[
\frac{\|\tilde\gamma\|_{L_\infty}^2}
     {n\,\P\tilde\gamma^2\nu}
  \le \frac{C\,\rho(\tilde\phi)^{2\theta}}
           {n\,(\P\tilde\phi^2)^\theta}.
\]
The penalty bias bound~\eqref{eq:sobolev-penalty-bias}
gives $\rho(\tilde\phi)
  \le G\,\eta^{-1/2}/\sqrt{2}$,
so $\rho(\tilde\phi)^{2\theta}
  \le C\,\eta^{-\theta}$,
and denominator anchoring gives
$(\P\tilde\phi^2)^\theta \ge c^\theta > 0$.
The ratio is therefore at most
$C\,\eta^{-\theta}
  / (n\,c^\theta)$,
which vanishes when
$a < 1/\theta$.
The window~\eqref{eq:sobolev-clt-window}
is nonempty when $\theta < 1$,
i.e.\ $s > d/2$.

\paragraph{CLT and efficiency.}
The conditions of \Cref{thm:general-clt}
are satisfied.
The imbalance condition~\ref{cond:clt-a}
is verified by \Cref{prop:outcome-clt-condition},
which handles $\mu \in H^{s_\mu}$
by choosing the decomposition
$g_n = \hat\mu - \tilde\mu$,
$r_n = \tilde\mu - \mu$
in~\eqref{eq:clt-a},
without requiring $\rho(\mu) < \infty$.
Condition~\ref{cond:clt-c}
--- that individual weights
are infinitesimal ---
holds by the interpolation bound above.
Condition~\ref{cond:clt-d}
--- that the denominator is bounded away from zero ---
holds by denominator anchoring.
For Q$_+$,
the two-pile tolerance~\eqref{eq:two-pile-tolerance}
depends on the dead-zone mass
$\omega(\varepsilon) = \P(\tilde\phi \in [0,\varepsilon])$;
the control $\omega(\varepsilon) \to 0$
is verified for each application below
using \Cref{prop:anticonc-sobolev}
and the approximation rate~\eqref{eq:approx-competitor}.

\paragraph{Efficiency.}
When $\griesz \in \mathcal{H}$,
the comparison bound~\eqref{eq:comparison-bound}
$L(\tilde\phi) \le L(\griesz)$
gives $\P(\tilde\gamma - \griesz)^2
  \le 2\eta\,\rho^2(\griesz) \to 0$.
When $\griesz \in L_2(\P)$
but $\griesz \notin \mathcal{H}$,
approximate:
for $f_\epsilon \in \mathcal{H}$
with $\|f_\epsilon - \griesz\|_{\LtP} < \epsilon$,
the comparison bound at $f_\epsilon$ gives
$\P(\tilde\phi - \griesz)^2
  \le 4\eta\,\rho^2(f_\epsilon) + 2\epsilon^2
  \to 2\epsilon^2$;
letting $\epsilon \to 0$
gives $\tilde\phi \to \griesz$ in $\LtP$.
For Setting~E,
the comparison bound and Sobolev embedding
give $\|\tilde\phi\|_{L_\infty} \le C$
uniformly in $\eta$,
so $\ddchis$ is bounded below
on the relevant set,
converting Bregman convergence
to $\LtP$ convergence of the index;
Lipschitz continuity of the link
$\phi \mapsto e^{\tau\phi}$
on bounded sets
then gives $\tilde\gamma \to \griesz$ in $\LtP$.
In all cases,
$\tilde\sigma^2 \to \sigma^2$
by \Cref{cor:efficient-variance}.
\end{proof}

\begin{proposition}[Outcome estimator
  sufficient condition for~\ref{cond:clt-a}]
\label{prop:outcome-clt-condition}
In the setting of \Cref{thm:sobolev-clt},
suppose the outcome estimator $\hat\mu$
is a penalized method-of-moments estimator
with the same Sobolev seminorm $\rho$,
penalty $\eta_\mu = n^{-b}$,
and $\mu \in H^{s_\mu}$ for some $s_\mu > 0$.
Then the imbalance condition~\ref{cond:clt-a}
is satisfied whenever
\begin{equation}\label{eq:outcome-clt-condition}
b(1+\theta) < a
\end{equation}
and
\begin{equation}\label{eq:outcome-dr-condition}
c_\gamma + \frac{b\,s_\mu}{2s} > \frac{1}{2},
\end{equation}
where $c_\gamma
  = \min\!\qty{
    \frac{1-a\theta}{2},\,
    \frac{a s_\gamma}{2s}}$
is the rate of weight convergence
in $\LtP$,
with $s_\gamma > 0$ the Sobolev regularity
of the Riesz representer
($\griesz \in H^{s_\gamma}$).
In particular,
when $s_\gamma + s_\mu > s$
a suitable choice of $a$ and $b$
exists throughout the window.
\end{proposition}

\begin{proof}
We verify condition~\eqref{eq:clt-a}
with competitor $f = 0$
and decomposition
$g_n = \hat\mu - \tilde\mu$,
$r_n = \tilde\mu - \mu$.
Because both $\hat\mu$ and $\tilde\mu$
are penalized estimators in $H^s$,
$g_n \in H^s$ and $\rho(g_n) < \infty$.
Because $\mu \in H^{s_\mu}$ with $s_\mu < s$,
$r_n = \tilde\mu - \mu \notin H^s$ in general,
and we bound $\|r_n\|_{\LtP}$
using approximation theory.

\paragraph{Rates.}
\Cref{cor:sobolev-main}
applied to the outcome problem
(quadratic dispersion,
$\overline{\ddchis}/\underline{\ddchis} = 1$)
gives the estimation rate
\begin{equation}
\label{eq:outcome-seminorm-rate}
\rho^2(g_n)
  = \rho^2(\hat\mu - \tilde\mu)
  \le C\,n^{b(1+\theta)-1}/\delta.
\end{equation}
For the approximation error,
write $\mu = \mu_J + r_J$
where $\mu_J = \sum_{j \le J} \mu_j\,e_j$
and $r_J = \mu - \mu_J$.
Because $\mu \in H^{s_\mu}$
and $\lambda_j \asymp j^{-2s/d}$,
$\|r_J\|_{L_2}
  \le C\,J^{-s_\mu/d}$.
The comparison bound at $\mu_J$ gives
$\|\tilde\mu - \mu\|_{L_2}^2
  \le \|r_J\|_{L_2}^2
  + \eta_\mu\,\rho^2(\mu_J)
  \le C\,J^{-2s_\mu/d}
  + \eta_\mu\,C\,J^{2(s-s_\mu)/d}$.
Taking $J = n^{b\theta}$,
both terms are $O(n^{-bs_\mu/s})$,
so
\begin{equation}
\label{eq:approx-l2-rate}
\|r_n\|_{\LtP}
  = \|\tilde\mu - \mu\|_{\LtP}
  = O(n^{-bs_\mu/(2s)}).
\end{equation}

\paragraph{Term~1 (smooth part).}
With $f = 0$, the first term
of~\eqref{eq:clt-a} is
$\|\griesz\|_{\LtP}\sqrt{\eta}\,\rho(g_n)
  = O_p(n^{-a/2 + (b(1+\theta)-1)/2})$,
which is $o_p(n^{-1/2})$
when $b(1+\theta) < a$.

\paragraph{Term~2 (rough part).}
With $f = 0$,
the second term is
$\|\griesz\|_{\LtP}\,\|r_n\|_{\LtP}
  = O(n^{-bs_\mu/(2s)})$.
This is $o(n^{-1/2})$
whenever $bs_\mu/s > 1$,
which is implied
by~\eqref{eq:outcome-dr-condition}
(the weight convergence rate
$c_\gamma > 0$ allows the full burden
to fall on the outcome side).

A better bound uses the weight
convergence rate:
\Cref{cor:sobolev-main}
gives $\|\hgamma - \tilde\gamma\|_{\LtP}
  = O_p(n^{(a\theta-1)/2})$,
and the comparison bound gives
$\|\tilde\gamma - \griesz\|_{\LtP}
  = O(n^{-as_\gamma/(2s)})$
when $\griesz \in H^{s_\gamma}$,
so
$\|\hgamma - \griesz\|_{\LtP}
  = O_p(n^{-c_\gamma})$.
Multiplying:
$n^{-c_\gamma}\,n^{-bs_\mu/(2s)}
  = n^{-c_\gamma - bs_\mu/(2s)}
  = o(n^{-1/2})$
when~\eqref{eq:outcome-dr-condition} holds.
(The proof of \Cref{thm:general-clt}
uses both the comparison-bound rate
from~\eqref{eq:clt-a}
and the Bregman rate;
the combined bound is
$n^{-c_\gamma}$
with $c_\gamma
  = \min\{(1-a\theta)/2,\,as_\gamma/(2s)\}$.)

\paragraph{Term~3 (centering).}
Because $\griesz \in H^{s_\gamma}$
with $s_\gamma > 0$
and the Sobolev embedding
$H^{s_\gamma} \hookrightarrow L_p$
for $p > 2$,
$\|\griesz\,r_n\|_{\LtP}
  \le \|\griesz\|_{L_p}\,\|r_n\|_{L_q}$
by H\"older
with $1/2 = 1/p + 1/q$.
Because $\|r_n\|_{L_2} \to 0$
and $\|r_n\|_{L_{q'}}$
is bounded for some $q' > 2$
(Sobolev embedding of $\tilde\mu$ and $\mu$),
interpolation gives $\|r_n\|_{L_q} \to 0$.
So $\|\griesz\,r_n\|_{\LtP}/\sqrt{n}
  = o(n^{-1/2})$.
\end{proof}

\begin{remark}[Double robustness]
\label{rem:double-robustness}
Condition~\eqref{eq:outcome-dr-condition}
is the double-robustness boundary
of \citet{hirshberg2025retargeted}:
in interpolation-space notation,
$\kappa_\mu + \kappa_\gamma > 1$
where $\kappa_\mu = s_\mu/s$
and $\kappa_\gamma = s_\gamma/s$.
Under strong overlap
($\griesz$ bounded, $s_\gamma$ large),
any $s_\mu > 0$ suffices.
Under limited overlap,
the outcome regression
must compensate
with greater smoothness.
At $s = d/2 + \epsilon$,
the condition approaches $s_\mu + s_\gamma > d/2$,
the semiparametric minimax boundary
of \citet{robins2009semiparametric}.
\end{remark}

Under overlap ($\griesz \in L_2$),
the window~\eqref{eq:sobolev-clt-window}
is $1 < a < 1/\theta$,
nonempty when $s > d/2$,
and efficiency follows.
The three applications below
instantiate this theorem
by verifying the overlap
condition~\ref{cond:sob-overlap}
for each functional.

\section{Treatment-specific mean}
\label{sec:app-tsm}

We verify the conditions of \Cref{thm:asymptotic-normality}
for the treatment-specific mean
$\psi^w(\mu) = \E \mu(w, X)$
of Example~\ref{ex:treatment-specific-mean},
with Sobolev seminorm $\rho(h) = \|(-\Delta)^{s/2} h\|_{L_2}$
for $s > d/2$ and quadratic regularization
$\zetas(\rho) = \eta \rho^2$.

\paragraph{Common setup.}
The derivative functional is
$\dotpsiZ(h) = h(w, X)$ and the Riesz representer is
$\griesz(W,X) = \ind{W=w}/\pi_w(X)$
for propensity score $\pi_w(x) = \P(W=w \mid X=x)$.
We assume $X$ has density $p \ge \pmin > 0$
on $[0,1]^d$.
The interpolation exponent is $\theta = d/(2s)$
(\Cref{lem:sobolev-interpolation}),
and the effective dimension is
$d_\eta = \sum_j \lambda_j / \max(\lambda_j, \eta)
  \asymp \eta^{-\theta}$
for Sobolev eigenvalues $\lambda_j = j^{-2s/d}$.

\paragraph{Conditions common to all dispersions.}
Because $\mathcal{H}$ contains $\E(Y \mid Z)$
in Example~\ref{ex:treatment-specific-mean},
$\mu_\perp = 0$
and~\ref{cond:an-D} and~\ref{cond:an-E} are vacuous.
Because $\psiZ$ is linear,
the linearization error $R_{\mathrm{lin}} = 0$,
so \ref{cond:an-F} and~\ref{cond:an-G} are vacuous.
It remains to verify~\ref{cond:an-A},
\ref{cond:an-B}, and~\ref{cond:an-C} for each dispersion.

\subsection{Normality}

\paragraph{Setting~Q: denominator anchoring and the balance equation.}
Here $\dchis(\phi) = \phi$,
so $\hgamma = \hat\phi$,
$\tilde\gamma = \tilde\phi$,
and $\ddchis = 1$ everywhere.
Pseudo-self-concordance holds with $\beta = 0$,
so $\Omega = 1$.

Constant functions have zero Sobolev seminorm,
so the first-order condition~\eqref{eq:foc-quadratic} at $h = 1$
gives $\P\tilde\phi = \dotpsiZ(1) = p_w$
with no contribution from the penalty:
the penalized projection integrates
to the treatment probability.
\Cref{lem:denominator-anchoring} then gives
\begin{equation}
\label{eq:tsm-variance-lower}
\P\tilde\phi^2 \ge p_w^2 > 0,
\end{equation}
uniformly in $\eta$ and without overlap.
This anchors the denominator
$\hat\Gamma \ge \underline\nu\, \hP\hat\phi^2
  \to_p \underline\nu\, \P\tilde\phi^2
  \ge \underline\nu\, p_w^2 > 0$.

Under overlap ($\Var(W \mid X) \ge \underline{v} > 0$),
the Riesz representer $\griesz = \ind{W=w}/\pi_w(X)$
is in $L_2(\P)$
with $\|\griesz\|_{\LtP}^2 = \E(1/\pi_w(X))$.
\Cref{prop:penalty-bias-balance} gives
\begin{equation}\label{eq:tsm-balance-bounds}
\eta\,\rho(\tilde\phi)
  \le \frac{\|\griesz\|_{\LtP}}{\sqrt{2}}\,\sqrt{\eta},
\qquad
\rho(\tilde\phi)^{2\theta}
  \le C\,\eta^{-\theta}.
\end{equation}
\Cref{cor:tuning-window}
translates these into the two-sided window
\begin{equation}
\label{eq:tsm-lambda-window}
n^{-1/\theta} \ll \eta
  \ll n^{-1},
\end{equation}
nonempty when $\theta < 1$,
equivalently $s > d/2$.


Under overlap ($\griesz \in \mathcal{H}$),
taking $f = \griesz$ in~\ref{cond:clt-a}:
the approximation term vanishes
and $\eta\,\rho(\griesz)\,\rho(\hat\mu-\mu)
  = o_p(n^{-1/2})$
whenever $\eta \to 0$.
Condition~\ref{cond:clt-c} holds with $C/n$
because $\|\tilde\gamma\|_\infty \le 1/\pi_{\min}$.
In this regime,
$\tilde\gamma \to \griesz$ in $\LtP$,
and \Cref{cor:efficient-variance} gives
$\tilde\sigma^2 \to \sigma^2$:
the CLT is both valid and efficient.

The practitioner need not check
which regime applies.
Choosing $\eta$
in the window~\eqref{eq:tsm-lambda-window}
guarantees a valid CLT
with variance $\tilde\sigma^2$.
If the problem has overlap,
efficiency follows automatically;
if not, the confidence interval
is valid but conservative.

\paragraph{Positive-part quadratic (Setting~Q with dead zone).}
Here $\dchis(\phi) = \phi_+$,
so the weights are $\hgamma = \hat\phi_+$
and $\tilde\gamma = \tilde\phi_+$.
The curvature is $\ddchis(\phi) = \ind{\phi \ge 0}$,
giving $\underline{\ddchis} = 0$, $\overline{\ddchis} = 1$.
The dead zone $\{\phi < 0\}$
means the Bregman divergence
\[
D_{\chis}(\hat\phi, \tilde\phi)
  = (\hat\phi_+ - \tilde\phi_+)^2/2
\]
controls only the positive parts:
convergence of $\hat\phi$ to $\tilde\phi$
on $\{\tilde\phi < 0\}$ is lost.
But the weights are $\phi_+$,
and the Lipschitz property
$|\hat\phi_+ - \tilde\phi_+| \le |\hat\phi - \tilde\phi|$
ensures this doesn't affect conditions~\ref{cond:an-B} or~\ref{cond:an-C}.

The dead-zone cost enters condition~\ref{cond:an-A}
through the tolerance degradation
in \Cref{lem:two-pile}.
On $\{W = w\}$,
$\griesz = 1/\pi_w(X) \ge 1$,
so $\tilde\phi \in [0,\varepsilon]$
with $\varepsilon < 1/2$
gives $D_{\chis}(\tilde\phi, \phi)
  \ge (1 - \varepsilon)^2/2 > \varepsilon^2/2$;
on $\{W \ne w\}$,
the dispersion forces $\tilde\phi \le 0$.
Therefore $\omega(\varepsilon)
  \le 2\,\P D_{\chis}(\tilde\phi, \phi)/\varepsilon^2$,
which is controlled
by the approximation rate~\eqref{eq:approx-competitor}:
neither strict overlap
nor an integer-smoothness condition is needed.

\paragraph{Setting~E: the comparison bound.}
Here $\dchis(\phi) = e^{\tau\phi}$,
so $\hgamma = e^{\tau\hat\phi}$
and $\tilde\gamma = e^{\tau\tilde\phi}$.
All consecutive derivative ratios equal $\tau$,
so pseudo-self-concordance holds with $\beta = \tau$
and $\Omega = e^{\tau\|h\|_{L_\infty}}$.

The first-order condition~\eqref{eq:foc-quadratic} at $h = 1$ gives
$\P e^{\tau\tilde\phi} = p_w$:
a normalization of the exponential weights.
Cauchy--Schwarz anchors the denominator:
\begin{equation}\label{eq:tsm-entropy-denominator}
\P\tilde\gamma^2
  = \P e^{2\tau\tilde\phi}
  \ge (\P e^{\tau\tilde\phi})^2
  = p_w^2 > 0.
\end{equation}

The balance-equation argument
from Setting~Q does not carry over.
Testing the first-order condition~\eqref{eq:foc-quadratic}
at $h = \tilde\phi$ gives
$\P\,e^{\tau\tilde\phi}\,\tilde\phi
  + 2\eta\rho^2(\tilde\phi)
  = \dotpsiZ(\tilde\phi)$,
but the weighted moment
$\P\,\tilde\gamma\,\tilde\phi$
does not bound $\rho(\tilde\phi)$
the way $\P\tilde\phi^2$ did
for quadratic.
The exponential link entangles
weights and penalty,
and no choice of test direction
disentangles them.

Instead,
write $\phi_\infty = \log(1/\pi_w)/\tau$
for the unpenalized population minimizer
and decompose the loss
via the Bregman divergence
from $\phi_\infty$.
Comparing $L(\tilde\phi) \le L(\phi_\infty)$
and using $D_{\chis} \ge 0$ gives
\begin{equation}
\label{eq:entropy-comparison}
\rho(\tilde\phi) \le \rho(\phi_\infty)
\end{equation}
whenever $\rho(\phi_\infty) < \infty$:
the log-propensity score
$\log(1/\pi_w)$
has finite Sobolev seminorm.
Write $B = \rho(\phi_\infty)$ for this bound.
The Poincar\'e inequality gives
$\|\tilde\phi - \P\tilde\phi\|_{\LtP}
  \le C_P^{1/2}\, \rho(\tilde\phi)
  \le C_P^{1/2}\, B$,
and the interpolation inequality~\eqref{eq:interpolation}
(applied to the mean-zero part
$\tilde\phi - \P\tilde\phi$,
which has the same seminorm) gives
\[
\|\tilde\phi - \P\tilde\phi\|_{L_\infty}
  \le \kappa\, B^\theta\,
    (C_P^{1/2}\, B)^{1-\theta}
  = \kappa\, C_P^{(1-\theta)/2}\, B.
\]
Write $M_0 = \kappa\, C_P^{(1-\theta)/2}\, B$
for this bound.
The normalization pins the mean:
$\P e^{\tau\tilde\phi} = p_w$
with $|\tilde\phi - \P\tilde\phi| \le M_0$
gives
$e^{\tau(\P\tilde\phi - M_0)}
  \le \P e^{\tau\tilde\phi}
  \le e^{\tau(\P\tilde\phi + M_0)}$,
so $\log(p_w)/\tau - M_0
  \le \P\tilde\phi
  \le \log(p_w)/\tau$.
By the triangle inequality,
\begin{equation}\label{eq:tsm-entropy-supnorm}
\|\tilde\phi\|_{L_\infty}
  \le |\P\tilde\phi| + M_0
  \le |\log p_w|/\tau + 2M_0,
\end{equation}
so
$\|\tilde\gamma\|_\infty
  = e^{\tau\|\tilde\phi\|_\infty}$
is bounded
by a constant depending on $B$, $p_w$,
and $\tau$
but not on $\eta$.
This is the key simplification:
the comparison bound
freezes both the seminorm and the sup-norm,
making conditions~\ref{cond:an-A} and~\ref{cond:an-B} immediate.
The resulting window
\begin{equation}\label{eq:tsm-entropy-window}
n^{-1/\theta} \ll \eta \ll n^{-1/2}
\end{equation}
is one-sided ---
only condition~\ref{cond:an-C} constrains $\eta$
from below ---
and nonempty for all $s > d/2$,
unlike the quadratic case
which requires $s > (1+\sqrt{2})\,d/2$.

\emph{The threshold is sharp.}
If $\rho(\phi_\infty) = \infty$
(the log-propensity score
lacks the penalty's native smoothness),
the Fourier-truncation comparison
gives $\rho(\tilde\phi)
  \le C\,\eta^{(s'-s)/(2s)}$
for $\log(1/\pi_w) \in H^{s'}$
with $s' < s$,
and the penalty bias
$\eta\,\rho(\tilde\phi) \le C\,\eta^{(s+s')/(2s)}$
is tractable for~\ref{cond:an-A}.
But the sup-norm
\[
\|\tilde\phi\|_{L_\infty}
  \le C\,\eta^{\theta(s'-s)/(2s)}
\]
now grows as $\eta \to 0$,
and condition~\ref{cond:an-B} requires
$e^{C\eta^{-\delta}}/n \to 0$
for $\delta = \theta(s-s')/(2s) > 0$---a
logarithmic rate incompatible
with the polynomial rate from~\ref{cond:an-A}.
The exponential link converts
polynomial sup-norm growth
into super-polynomial weight growth;
for quadratic,
the linear link keeps both polynomial,
and partial regularity
widens the $\eta$ window
rather than closing it.

Under overlap,
$\rho(\phi_\infty)
  = \rho(\log(1/\pi_w)/\tau) < \infty$
when the propensity score
is smooth
(as it is whenever
the Sobolev RKHS
contains the conditional model),
and $\tilde\gamma \to \griesz$ in $\LtP$.
\Cref{cor:efficient-variance} gives
$\tilde\sigma^2 \to \sigma^2$:
the CLT is both valid and efficient,
as in Setting~Q.

\paragraph{Outcome estimator.}
The condition~\ref{cond:an-A} verification
requires that the offset function
$f = \hat\mu - \mu$
have bounded Sobolev seminorm:
this is how the offset enters
the imbalance bound
of \Cref{lem:offset-convex}.
The pure Riesz-weighting estimator
uses $\hat\mu = 0$,
giving $f = -\mu$
and requiring $\rho(\mu) < \infty$.
But the same penalized-projection framework
that estimates the weights
can estimate the outcome.
The conditional mean
$\mu_w(\cdot) = \E(Y \mid W{=}w, \cdot)$
is the Riesz representer
for the functional
$h \mapsto \P_w[Y\, h]$,
so penalized least squares
with the Sobolev penalty
targets $\mu_w$ directly:
\begin{equation}\label{eq:tsm-outcome}
\hat\mu_w
  = \argmin_{m}
    \hP_w (Y - m(X))^2/2
    + \eta_\mu\, \rho^2(m)/2.
\end{equation}
The comparison bound
gives $\rho(\tilde\mu_w) \le \rho(\mu_w)$:
the population minimizer
inherits the regularity of the truth.
The cross term that makes this work ---
$\P(Y - \mu_w)(m - \mu_w) = 0$
for all $m$ ---
vanishes by the tower property,
because $\mu_w$ is the conditional mean.

\begin{corollary}[Treatment-specific mean]
\label{cor:tsm}
Assume $X$ has density $p \ge \pmin > 0$
on $[0,1]^d$,
$\rho(\mu) < \infty$,
and $\rho(\hat\mu - \mu)$ bounded in probability.
Write $\theta = d/(2s)$ and $\eta = n^{-a}$.
\begin{enumerate}
\item[\textup{(Q)}]
\textit{Quadratic.}
If $1/(1-\theta) < a < 1/(\theta(1+\theta))$
--- nonempty when $s > (1+\sqrt{2})\,d/2$ ---
then
$\sqrt{n}(\hat\psi^w - \psi^w)/\hat\sigma \to_d N(0,1)$.
No overlap condition enters.
If additionally $\griesz \in L_p(\P)$
for some $p \in (1,2)$,
the window widens
(the threshold relaxes to $\theta < \theta_p$).
Under overlap,
$\rho(\tilde\phi) \le \rho(\griesz) < \infty$
and the window becomes $\eta \to 0$.
\item[\textup{(Q$_+$)}]
\textit{Positive-part quadratic.}
The same window as~\textup{(Q)} applies.
The dead-zone cost~(i) of \Cref{cor:sobolev-two-pile}
is controlled
by the approximation rate~\eqref{eq:approx-competitor}.
\item[\textup{(E)}]
\textit{Entropy-like.}
If $\rho(\log(1/\pi_w)/\tau) < \infty$
and $1/2 < a < 1/\theta$
--- nonempty for all $s > d/2$ ---
then
$\sqrt{n}(\hat\psi^w - \psi^w)/\hat\sigma \to_d N(0,1)$.
\end{enumerate}
In all cases,
$\hat\sigma^2 / \tilde\sigma^2 \to_p 1$.
Under overlap,
$\tilde\sigma^2 \to \sigma^2$,
the semiparametric efficiency bound.
\end{corollary}

\begin{proof}
Conditions~\ref{cond:an-D}--\ref{cond:an-G} are vacuous:
$\psiZ$ is linear
($R_{\mathrm{lin}} = 0$, so~\ref{cond:an-F} and~\ref{cond:an-G} hold)
and $\mu_\perp = 0$
(so~\ref{cond:an-D} and~\ref{cond:an-E} hold).
We verify~\ref{cond:an-A}--\ref{cond:an-C}
for each dispersion.

\emph{Setting~Q.}
We apply \Cref{thm:general-clt}.
Condition~\ref{cond:clt-d} is~\eqref{eq:tsm-variance-lower}.
Taking $f = 0$ in~\ref{cond:clt-a}:
$\|\griesz\|_{\LtP}\sqrt{\eta}\,\rho(\hat\mu-\mu)
  = o_p(n^{-1/2})$
when $\rho(\hat\mu - \mu) = O_p(1)$
and $a > 1$.
For condition~\ref{cond:clt-c},
the interpolation inequality~\eqref{eq:interpolation}
and the Jensen bound $\P\tilde\phi^2 \ge p_w^2$
give
$\|\tilde\phi\|_{L_\infty}^2
  / (n\,\P\tilde\phi^2)
  \le C\,\rho(\tilde\phi)^{2\theta}
      / (n\, p_w^{2\theta})
  \le C / (n\,\eta^{\theta(1+\theta)})$,
which vanishes when $a < 1/(\theta(1+\theta))$.
Because $1/(\theta(1+\theta)) < 1/\theta$,
\ref{cond:clt-c} is more restrictive
than the effective-dimension condition $d_\eta \ll n$,
and the window is~\eqref{eq:tsm-lambda-window}.

\emph{Setting~Q$_+$: imbalance~\ref{cond:an-A}.}
\Cref{cor:sobolev-two-pile}
gives the offset complexity rate,
with probability bound
$1 - 2e^{-cn\xi_0^2 r^2/M^2}$
where $M = C_{\mathrm{emb}}\,\rho_{\sqrt{\nu}r} + C_K\, r$.
\Cref{cor:offset-positive-part}
gives the bias bound
$|b(\hgamma)|
  \le C\, n^{-1/2}\, (n^{1/2}\, \tilde\rho)^{-\theta/(2-\theta)}
  + 4\eta\,\rho(\tilde\phi)$,
the same leading rate as the quadratic case.
The dead-zone cost enters through
the tolerance degradation
in \Cref{lem:two-pile}:
$\mu = \omega(\epsilon)\, M^2 + r^2 M^2/\epsilon^2$,
optimized over $\epsilon > 0$.
Taking $\epsilon = r^a$ for any $a \in (0,1)$
gives $\mu/r^2 \le (M/r)^2[\omega(r^a) + r^{2-2a}]$,
which vanishes as $r \to 0$
because $M/r$ is bounded
and $\omega(\varepsilon)
  \le 2\,\P D_{\chis}(\tilde\phi, \phi)/\varepsilon^2 \to 0$
by the approximation rate~\eqref{eq:approx-competitor}
(the first term of \Cref{prop:anticonc-sobolev}
vanishes because $\griesz \ge 1$ on $\{W = w\}$
and $\tilde\phi \le 0$ on $\{W \ne w\}$).
Condition~\ref{cond:an-A} holds as in the quadratic case.

\emph{Setting~Q$_+$: infinitesimal weights~\ref{cond:an-B}.}
The weight map $\phi \mapsto \phi_+$
is $1$-Lipschitz,
so $|\hgamma - \tilde\gamma|
  = |\hat\phi_+ - \tilde\phi_+|
  \le |\hat\phi - \tilde\phi|$.
The same sup-norm bound as Setting~Q applies:
$\max_i \hgamma^2(Z_i)\nu(Z_i)/(n\hat\Gamma)
  \le C/(n\eta^{\theta(1+\theta)})$,
and the condition is the same.

\emph{Setting~Q$_+$: variance equivalence~\ref{cond:an-C}.}
Because $D_{\chis}(\hat\phi, \tilde\phi)
  \ge (\hat\phi_+ - \tilde\phi_+)^2/2$,
Bregman convergence gives
$\hP(\hgamma - \tilde\gamma)^2 \le 2\nu r^2$
directly.
The Lipschitz property ensures
the dead zone does not affect variance:
$\hP(\hgamma - \tilde\gamma)^2\nu \to_p 0$
by the same expansion
$\hat\Gamma
  = \hP\tilde\gamma^2\nu
  + 2\hP\tilde\gamma(\hgamma - \tilde\gamma)\nu
  + \hP(\hgamma - \tilde\gamma)^2\nu$
and Cauchy--Schwarz as in Setting~Q.

\emph{Setting~E: imbalance~\ref{cond:an-A}.}
\Cref{thm:xi-all-regimes}
controls the offset complexity,
with $\Omega = e^{\tau\|h\|_{L_\infty}} \to 1$
as $r \to 0$.
Because $\rho(\tilde\phi) \le B$
by the comparison bound~\eqref{eq:entropy-comparison},
the penalty bias satisfies
$\eta\,\rho(\tilde\phi) \le \eta B$,
which is $o(n^{-1/2})$
whenever $\eta \ll n^{-1/2}$,
i.e.\ $a > 1/2$.
The offset complexity term
is $o(n^{-1/2})$
by \Cref{cor:unconditional-bias}
at any $r > 2\rho(\tilde\phi)$.

\emph{Setting~E: infinitesimal weights~\ref{cond:an-B}.}
Because $\|\tilde\gamma\|_\infty
  = e^{\tau\|\tilde\phi\|_\infty}$
is bounded independently of $\eta$
by~\eqref{eq:tsm-entropy-supnorm},
\[
\frac{\max_i \hgamma^2(Z_i)\,\nu(Z_i)}
     {n\,\hat\Gamma}
  \le \frac{\|\tilde\gamma\|_\infty^2}{n\,\underline\nu\,p_w^2}
  \le \frac{C}{n} \to 0.
\]
Condition~\ref{cond:an-B} holds automatically,
with no constraint on $\eta$.
This is the key advantage of the comparison bound:
because $\rho(\tilde\phi) \le B$ regardless of $\eta$,
the sup-norm and therefore the weights
are bounded independently of $\eta$,
and \ref{cond:an-B} imposes no lower bound
on the undersmoothing rate.

\emph{Setting~E: variance equivalence~\ref{cond:an-C}.}
\Cref{cor:variance-interpolation}
with $\alpha = \alpha_1 = \beta = \beta_1 = \tau$ gives
$\hP\hgamma^2\nu / \hP\tilde\gamma^2\nu \to 1$
directly,
without requiring $L_2$ weight convergence.
The point is that for entropy-like dispersion,
the Bregman divergence
$D_{\chis}(\hat\phi, \tilde\phi)$
does not directly control
$(\hgamma - \tilde\gamma)^2
  = (e^{\tau\hat\phi} - e^{\tau\tilde\phi})^2$
the way it does for quadratic.
The interpolation inequality
of \Cref{cor:variance-interpolation}
bridges this gap,
using the PSC structure
to convert Bregman convergence
into variance-ratio convergence.
Condition~\ref{cond:an-C} holds whenever $d_\eta \ll n$,
i.e.\ $a < 1/\theta$.

\emph{Setting~E: tuning window.}
Collecting:
\ref{cond:an-A} requires $a > 1/2$,
\ref{cond:an-B} is automatic,
\ref{cond:an-C} requires $a\theta < 1$.
The window is given by~\eqref{eq:tsm-entropy-window},
nonempty for all $s > d/2$.

\emph{Efficiency.}
Under overlap,
$\tilde\gamma \to \griesz$ in $\LtP$
in all three settings,
and \Cref{cor:efficient-variance}
gives $\tilde\sigma^2 \to \sigma^2$.
\end{proof}

\begin{corollary}[Penalized LS outcome estimator]
\label{cor:tsm-outcome}
In the setting of \Cref{cor:tsm},
let $\hat\mu_w$ be given by~\eqref{eq:tsm-outcome}
with $\eta_\mu = n^{-b}$
for $b < 1/(1+\theta)$.
Then $\rho(\hat\mu - \mu)$
is bounded in probability
and the conclusions of
\Cref{cor:tsm} hold.
\end{corollary}

\begin{proof}
\Cref{cor:outcome-estimator}
with \Cref{rem:arm-specific-outcome}
($n \to n_w$).
\end{proof}

\section{Average linear regression}
\label{sec:app-regression}

We verify the conditions of \Cref{thm:asymptotic-normality}
for the average linear regression function
$\psi(\mu) = \E b_0(X)$
of Example~\ref{ex:average-linear-regression},
with quadratic dispersion $\chis(\phi) = \phi^2/2$
and Sobolev regularization
$\zetas(\rho) = \eta\rho^2$
for $s > d/2$.

\paragraph{Setup.}
The tangent space $\mathcal{H}$
is the space of functions linear in $W$:
$h(W,X) = \alpha(X) + W \cdot c(X)$.
The Riesz representer
$\griesz(W,X)
  = \{W - \E(W \mid X)\}/\Var(W \mid X)$
lies in $\mathcal{H}$.
The conditional linear predictor
$\mu(W,X) = a_0(X) + W \cdot b_0(X)$
is the $L_2(\P)$ projection of $\E(Y \mid Z)$
onto $\mathcal{H}$,
with residual
$\mu_\perp = \E(Y \mid Z) - \mu$.
The projection property gives
$\P h\,\mu_\perp = 0$
for all $h \in \mathcal{H}$;
in particular, $\P\griesz\,\mu_\perp = 0$.

Only the quadratic dispersion is natural here:
the weights $\hgamma = \hat\phi$
estimate a contrast that takes both signs,
so positivity constraints
or sign-definite links would be inappropriate.

\paragraph{Penalty.}
Each $m \in \mathcal{H}$
decomposes uniquely as
$m(W,X) = a(X) + W \cdot b(X)$.
We penalize the Sobolev smoothness
of the coefficient functions:
\begin{equation}\label{eq:regression-penalty}
\rho^2(m) = \rho_s^2(a) + \rho_s^2(b),
\end{equation}
where $\rho_s$ is the Sobolev seminorm
on functions of $X \in [0,1]^d$
with $d = \dim(X)$.
The nonparametric estimation problem
is the smoothness of the coefficient functions
in $X$---not the linearity in $W$,
which is structural.
The effective dimension satisfies
$d_\eta \le C_s\, \eta^{-\theta}$
with $\theta = d/(2s)$,
where $C_s$ depends on $s$, $d$,
and the domain.

\paragraph{Misspecification and coherence.}
Unlike the treatment-specific mean,
the nuisance $\mu$ is a projection
of $\E(Y \mid Z)$ onto $\mathcal{H}$,
not the conditional mean itself.
The residual $\mu_\perp = \E(Y \mid Z) - \mu$
is generally nonzero,
so conditions~\ref{cond:an-D} and~\ref{cond:an-E}
are no longer vacuous.
The projection property is what saves us:
because $\mathcal{H}$ is a linear subspace,
both $\hat\phi$ and $\tilde\phi$ lie in $\mathcal{H}$,
and $\P h\,\mu_\perp = 0$ for all $h \in \mathcal{H}$
kills the leading term of~\ref{cond:an-D} exactly.
For~\ref{cond:an-E},
the same projection property gives
$\P\tilde\gamma\,\mu_\perp
  = \P\tilde\phi\,\mu_\perp = 0$
(using $\tilde\gamma = \tilde\phi$
for the quadratic dispersion),
so the centering condition
holds exactly.
Conditions~\ref{cond:an-A}--\ref{cond:an-C}
follow from the same machinery as the TSM,
with the Cauchy--Schwarz balance equation
replacing Jensen's inequality
(because the test direction $h = 1$
contributes nothing useful here).

\paragraph{Outcome estimator.}
The condition~\ref{cond:an-A} verification
requires that the offset function
$f = \hat\mu - \mu$
have bounded Sobolev seminorm.
Because $\mu \in \mathcal{H}$,
the same penalized-projection framework
that estimates the weights
can estimate the outcome:
\begin{equation}\label{eq:regression-outcome}
\hat\mu
  = \argmin_{m \in \mathcal{H}}
    \hP (Y - m(W,X))^2/2
    + \eta_\mu\, \rho^2(m)/2.
\end{equation}
The cross term
$\P(Y - \mu)(m - \mu)$
vanishes:
$\P\epsilon(m - \mu) = 0$
by the tower property
($\epsilon = Y - \E(Y \mid Z)$
has conditional mean zero),
and $\P\mu_\perp(m - \mu) = 0$
by the projection property
($m - \mu \in \mathcal{H}$).
The comparison bound
gives $\rho(\tilde\mu) \le \rho(\mu)$:
the population minimizer
inherits the regularity of the truth.

\begin{corollary}[Average linear regression]
\label{cor:regression}
Assume $X$ has density $p \ge \pmin > 0$
on $[0,1]^d$,
bounded $W$ and $\E(Y \mid Z)$,
$\rho(\mu_\perp) < \infty$,
and $\rho(\hat\mu - \mu)$ bounded in probability.
Write $\theta = d/(2s)$ and $\eta = n^{-a}$.
\begin{enumerate}[label=\textup{(\roman*)}]
\item\label{case:reg-structural}
\textit{$L_2$ integrability.}
Under overlap ($\griesz \in L_2(\P)$),
if $1 < a < 1/\theta$
--- nonempty when $s > d/2$ ---
then
$\sqrt{n}(\hat\psi - \psi)/\hat\sigma \to_d N(0,1)$.
\item\label{case:reg-comparison}
\textit{Partial regularity.}
If additionally $\griesz \in H^{s_\gamma}$
for some $s_\gamma \in (0,s)$,
the window widens to
$s/(s+s_\gamma) < a < s/(\theta(s-s_\gamma))$,
nonempty whenever $\theta < (s+s_\gamma)/(s - s_\gamma)$.
\item\label{case:reg-full}
\textit{Full regularity.}
If $\griesz \in H^s$,
the window widens to
$1/2 < a < 1/\theta$.
\end{enumerate}
In all cases,
$\hat\sigma^2 / \tilde\sigma^2 \to_p 1$.
Under overlap
($\Var(W \mid X) \ge \underline{v} > 0$),
$\tilde\sigma^2 \to \sigma^2$,
the semiparametric efficiency bound.
\end{corollary}

\begin{proof}
We verify conditions~\ref{cond:an-A}--\ref{cond:an-G} of
\Cref{thm:asymptotic-normality}.

\emph{Vacuous conditions.}
The functional $\psiZ$ is linear in $\mu$,
so $R_{\mathrm{lin}} = 0$
and conditions~\ref{cond:an-F} and~\ref{cond:an-G} hold.

\emph{Projection coherence~\ref{cond:an-D}.}
For the first part:
$\hat\phi$ and $\tilde\phi$
are both in $\mathcal{H}$,
so $\hat\phi - \tilde\phi \in \mathcal{H}$,
and the projection property
$\P h\,\mu_\perp = 0$ for $h \in \mathcal{H}$
gives
$\P(\hgamma - \tilde\gamma)\mu_\perp
  = \P(\hat\phi - \tilde\phi)\mu_\perp
  = 0$.
For the second part:
$\hP(\hat\phi - \tilde\phi)^2\mu_\perp^2
  \le \|\mu_\perp\|_{L_\infty}^2\,
      \hP(\hat\phi - \tilde\phi)^2
  \to_p 0$
by Bregman convergence
($\mu_\perp$ is bounded
because $\E(Y \mid Z)$ and $\mu$ are).

\emph{Balance equation.}
The first-order condition~\eqref{eq:foc-quadratic}
at $h = \tilde\phi$ gives
\begin{equation}\label{eq:regression-balance}
\P\tilde\phi^2 + 2\eta\rho^2(\tilde\phi)
  = \dotpsiZ(\tilde\phi).
\end{equation}
For $h(W,X) = \alpha(X) + W\, c(X)$,
the derivative functional satisfies
$\dotpsiZ(h) = \E[c(X)]$:
differentiating $\psi(\mu) = \E\, b_0(X)$
at $\mu + th$ gives
$\tfrac{d}{dt}\E[b_0(X) + t\, c(X)]|_{t=0}
  = \E[c(X)]$.
Under overlap ($\Var(W \mid X) \ge \underline{v} > 0$),
$G := \|\griesz\|_{\LtP} < \infty$,
and \Cref{prop:penalty-bias-balance} gives
\begin{equation}\label{eq:regression-balance-bounds}
\P\tilde\phi^2 \le G^2,
\qquad
\rho(\tilde\phi) \le G/\sqrt{2\eta},
\qquad
\eta\,\rho(\tilde\phi) \le G\sqrt{\eta/2}.
\end{equation}

\emph{Centering~\ref{cond:an-E}.}
The projection property
$\P h\,\mu_\perp = 0$
for $h \in \mathcal{H}$,
applied to $h = \tilde\phi \in \mathcal{H}$,
gives $\P\tilde\gamma\,\mu_\perp
  = \P\tilde\phi\,\mu_\perp = 0$.
Condition~\ref{cond:an-E} holds exactly.

\emph{Conditions~\ref{cond:an-A}--\ref{cond:an-C}
  via \Cref{thm:general-clt}.}
Condition~\ref{cond:clt-d} holds:
$\P\tilde\phi^2 > 0$
from \Cref{lem:denominator-anchoring}
(with $h_0 = (0, 1) \in \ker(\rho)$
and $\dotpsiZ(h_0) = 1$).

For case~\ref{case:reg-structural},
taking $f = 0$ in~\ref{cond:clt-a}
and using $\rho(\hat\mu-\mu) = O_p(1)$
gives condition~\ref{cond:clt-a} when $a > 1$.
Condition~\ref{cond:clt-c} holds when $a < 1/\theta$.
\Cref{cor:tuning-window} gives
conditions~\ref{cond:clt-a}--\ref{cond:clt-c}
in the window $1 < a < 1/\theta$.

\emph{Comparison bound.}
For cases~\ref{case:reg-comparison}
and~\ref{case:reg-full},
the loss decomposes via
$L(\tilde\phi) \le L(\phi_J)$
for the $J$-term Fourier truncation
$\phi_J$ of $\griesz$.
When $\griesz \in H^{s_\gamma}$,
optimizing over $J$ gives
$\rho(\tilde\phi)
  \le C\,\eta^{(s_\gamma-s)/(2s)}$.
This replaces the balance-equation bound
whenever $s_\gamma > 0$.
When $s_\gamma = s$,
$\rho(\tilde\phi) \le \rho(\griesz)$
and all bounds involving $\rho(\tilde\phi)$
become constants.

\emph{Efficiency.}
Under overlap ($\Var(W \mid X) \ge \underline{v} > 0$),
$\griesz \in \mathcal{H} \subset L_2(\P)$,
$\tilde\phi \to \griesz$ in $\LtP$
as $\eta \to 0$,
and \Cref{cor:efficient-variance}
gives $\tilde\sigma^2 \to \sigma^2$.
\end{proof}

\begin{corollary}[Penalized LS outcome estimator, regression]
\label{cor:regression-outcome}
In the setting of \Cref{cor:regression},
let $\hat\mu$ be given by~\eqref{eq:regression-outcome}
with $\eta_\mu = n^{-b}$
for $b < 1/(1+\theta)$.
Then $\rho(\hat\mu - \mu)$
is bounded in probability
and the conclusions of
\Cref{cor:regression} hold.
\end{corollary}

\begin{proof}
The cross term $\P(Y - \mu)(m - \mu)$
vanishes by both the tower property
($\P\epsilon(m - \mu) = 0$)
and the projection property
($\P\mu_\perp(m - \mu) = 0$
for $m - \mu \in \mathcal{H}$).
The rest of the argument
is identical to
\Cref{cor:outcome-estimator}.
\end{proof}

\section{Variance of the CATE}
\label{sec:app-vcate}

We verify the conditions of \Cref{thm:asymptotic-normality}
for the variance of the CATE
$V(\mu) = \Var(\tau(X))$
of Example~\ref{ex:vcate},
with Sobolev seminorm $\rho(h) = \|(-\Delta)^{s/2} h\|_{L_2}$
for $s > d/2$ and quadratic regularization
$\zetas(\rho) = \eta \rho^2$.
We work with the nonlinear component
$\P\psiZ(\mu) = \E[\tau(X)^2]$,
where $\psiZ(m) = \{m(1,X) - m(0,X)\}^2$;
the linear component $(\E\tau)^2$
is handled by \Cref{sec:app-tsm}.

\paragraph{Setup.}
The derivative at $\mu$ in direction $h$ is
$\dpsiZ^\mu(h) = 2\,\P\,\tau(X)\{h(1,X) - h(0,X)\}$
with Riesz representer
\[
\griesz^\mu(W,X)
  = 2\tau(X)
    \qty\Big{\frac{\ind{W=1}}{\pi_1(X)}
         - \frac{\ind{W=0}}{\pi_0(X)}}.
\]
We assume $X$ has density $p \ge \pmin > 0$
on $[0,1]^d$,
bounded outcomes $\|\E(Y \mid Z)\|_{L_\infty} < \infty$
(so $\|\tau\|_{L_\infty} \le T$ for some $T < \infty$),
heterogeneous effects $\Var(\tau(X)) > 0$,
$\E\tau(X) \ne 0$,
and marginal treatment probabilities $\P(W = w) > 0$.

The interpolation exponent is $\theta = d/(2s)$
and the effective dimension
$d_\eta \asymp \eta^{-\theta}$,
as in \Cref{sec:app-tsm}.

\paragraph{Outcome estimator.}
We estimate $\mu(w, \cdot)$ in each arm
by penalized least squares:
\begin{equation}\label{eq:vcate-outcome}
\hat\mu_w
  = \argmin_{m}
    \hP_w (Y - m(X))^2/2
    + \eta_\mu\, \rho^2(m)/2,
\end{equation}
where $\hP_w$ averages over the subsample
$\{i : W_i = w\}$.
The outcome estimation problem
fits the penalized-projection framework
with quadratic dispersion $\chis(m) = m^2/2$,
derivative $\dpsiZ(m) = \hP_w[Y\, m(X)]$,
and Riesz representer $\mu(w, \cdot)$ itself.
The Bregman convergence bounds of
\Cref{sec:bias} give
\begin{equation}\label{eq:vcate-outcome-rate}
\hP_w (\hat\mu_w - \tilde\mu_w)^2
  \lesssim d_{\eta_\mu}/n_w,
\end{equation}
and the regularization bias satisfies
$\|\tilde\mu_w - \mu_w\|_{\LtP}^2
  \le \eta_\mu\, \rho^2(\mu_w)$.
Together,
$\|\hat\mu_w - \mu_w\|_{\LtP}^2
  \lesssim d_{\eta_\mu}/n + \eta_\mu\, \rho^2(\mu_w)
  \asymp \eta_\mu^{-\theta}/n + \eta_\mu$.
The CATE estimate $\hat\tau = \hat\mu_1 - \hat\mu_0$
satisfies
\begin{equation}\label{eq:vcate-cate-rate}
\|\hat\tau - \tau\|_{\LtP}^2
  \lesssim \eta_\mu^{-\theta}/n + \eta_\mu.
\end{equation}
This is $o(n^{-1/2})$
when $\eta_\mu^{-\theta}/n = o(n^{-1/2})$
and $\eta_\mu = o(n^{-1/2})$,
i.e.\ $n^{-1/(2\theta)} \ll \eta_\mu \ll n^{-1/2}$.
This window is nonempty when $1/(2\theta) > 1/2$,
equivalently $\theta < 1$,
equivalently $s > d/2$.
Writing $\eta_\mu = n^{-b}$:
\begin{equation}\label{eq:vcate-outcome-window}
1/2 < b < 1/(2\theta).
\end{equation}

\paragraph{Balance equation and denominator.}
Because $\mathcal{H}$ is all bounded functions
of $(W,X)$,
$\mu_\perp = 0$
and conditions~\ref{cond:an-D} and~\ref{cond:an-E}
are vacuous.
The first-order condition~\eqref{eq:foc-quadratic}
at $h = 1$ gives
$\P\dchis(\tilde\phi) = \dpsiZ^\mu(1) = 0$:
the derivative kills constant functions
because $\psiZ$ depends on $m$
only through differences $m(1,X) - m(0,X)$.

But $h_0(w,x) = w$
lies in $\ker(\rho)$
(constant in $x$)
and satisfies $\dpsiZ^\mu(h_0) = 2\E\tau(X)$.
When $\E\tau \ne 0$,
\Cref{lem:denominator-anchoring} gives
\begin{equation}\label{eq:vcate-denominator-anchor}
\P\tilde\phi^2
  \ge (2\E\tau)^2 > 0,
\end{equation}
uniformly in $\eta$,
with no overlap condition.

The balance equation satisfies
\begin{equation}\label{eq:vcate-balance}
|\dpsiZ^\mu(\tilde\phi)|
  = 2\,|\P\,\tau(X)\{\tilde\phi(1,X) - \tilde\phi(0,X)\}|
  \le 4T\,\|\tilde\phi\|_{L_\infty}.
\end{equation}

Under overlap ($\pi_w \ge \pi_{\min} > 0$),
$\griesz^\mu \in L_2(\P)$
with $G := \|\griesz^\mu\|_{\LtP}
  \le 2T/\sqrt{\pi_{\min}}$,
and the first-order condition
at $h = \tilde\phi$ gives
$\P\tilde\phi^2 + 2\eta\rho^2(\tilde\phi)
  = \dpsiZ^\mu(\tilde\phi)
  = \P\griesz^\mu \tilde\phi
  \le G\,(\P\tilde\phi^2)^{1/2}$.
By \Cref{prop:penalty-bias-balance},
\begin{equation}\label{eq:vcate-balance-bounds}
\P\tilde\phi^2 \le G^2,
\qquad
\rho(\tilde\phi) \le G/\sqrt{2\eta},
\qquad
\eta\,\rho(\tilde\phi) \le G\sqrt{\eta/2}.
\end{equation}
The comparison bound~$L(\tilde\phi) \le L(\griesz^\mu)$
gives
\begin{equation}\label{eq:vcate-lower-anchor}
\P(\tilde\phi - \griesz^\mu)^2
  \le 2\eta\,\rho^2(\griesz^\mu),
\end{equation}
so by the triangle inequality in $\LtP$,
$\P\tilde\phi^2
  \ge (G
    - \sqrt{2\eta}\,\rho(\griesz^\mu))^2
  \ge G^2/4$
for $\eta \le G^2/(8\rho^2(\griesz^\mu))$.

\paragraph{Nonlinearity.}
The previous two examples had linear $\psiZ$,
making conditions~\ref{cond:an-F} and~\ref{cond:an-G} vacuous.
Here $\psiZ(m) = \{m(1,X) - m(0,X)\}^2$
is quadratic in $m$,
and both conditions require control of the outcome estimator.

Because $\psiZ$ is quadratic,
the linearization remainder is exact ---
the remainder equals the squared $L_2$ error
of the CATE estimate $\hat\tau$,
and condition~\ref{cond:an-G} asks that this be $o(n^{-1/2})$.
Condition~\ref{cond:an-F} asks that $\hat\tau^2 - \tau^2$
have small empirical process fluctuation;
because $\hat\tau$ is data-dependent,
this requires sample splitting.

The representer $\griesz^\mu$
depends on $\mu$ through $\tau$,
so the penalized projection targets
$\dpsiZ^{\hat\mu}$ rather than $\dpsiZ^\mu$.
This introduces a derivative error
proportional to $R_{\mathrm{lin}}$,
again controlled by the outcome rate.
The population minimizers
$\tilde\phi^{\hat\mu}$ and $\tilde\phi$
coincide in the limit
because the first-order conditions~\eqref{eq:foc-quadratic}
are linear in the derivative functional.

Only the quadratic dispersion is natural here:
the representer $\griesz^\mu$
takes both signs
($>0$ when $W=1$ and $\tau(X) > 0$,
$<0$ when $W=0$ and $\tau(X) > 0$),
so positivity constraints
or sign-definite links would be inappropriate---the
same situation as for
average linear regression
(\Cref{sec:app-regression}).

\begin{corollary}[Variance of the CATE]
\label{cor:vcate}
In the setting of Example~\ref{ex:vcate},
assume density $p \ge \pmin > 0$ on $[0,1]^d$,
bounded outcomes $\|\E(Y \mid Z)\|_{L_\infty} \le T$,
heterogeneous effects $\Var(\tau(X)) > 0$,
$\E\tau(X) \ne 0$,
and Sobolev smoothness $s > d/2$
with $\mu_w = \mu(w, \cdot) \in H^s$.
Estimate the outcome by~\eqref{eq:vcate-outcome}
with $\eta_\mu = n^{-b}$
for $b$ satisfying~\eqref{eq:vcate-outcome-window},
on an independent sample,
and the weights by quadratic penalized projection
with $\eta = n^{-a}$.
Write $\theta = d/(2s)$.
\begin{enumerate}[label=\textup{(\roman*)}]
\item\label{case:vcate-structural}
\textit{$L_2$ integrability.}
Under overlap ($\griesz^\mu \in L_2(\P)$),
if $1 < a < 1/\theta$
--- nonempty when $s > d/2$ ---
then
$\sqrt{n}(\hat V - V)/\hat\sigma \to_d N(0,1)$.
\item\label{case:vcate-comparison}
\textit{Partial regularity.}
If additionally $\griesz^\mu \in H^{s_\gamma}$
for some $s_\gamma \in (0,s)$,
the window widens to
$s/(s+s_\gamma) < a < s/(\theta(s-s_\gamma))$.
\item\label{case:vcate-full}
\textit{Full regularity.}
If $\griesz^\mu \in H^s$,
the window widens to
$1/2 < a < 1/\theta$.
\end{enumerate}
In all cases,
$\hat\sigma^2 / \tilde\sigma^2 \to_p 1$.
Under overlap
($\pi_w \ge \pi_{\min} > 0$),
$\tilde\sigma^2 \to \sigma^2_{\mathrm{eff}}$,
the semiparametric efficiency bound for $V$.
\end{corollary}

\begin{proof}
We verify conditions~\ref{cond:an-A}--\ref{cond:an-G} of
\Cref{thm:asymptotic-normality}.
Because the outcome model $\mu$ is unrestricted
($\mu_\perp = 0$),
conditions~\ref{cond:an-D} and~\ref{cond:an-E} hold trivially.

\emph{Conditions~\ref{cond:an-A}--\ref{cond:an-C}
  via \Cref{thm:general-clt}.}
For case~\ref{case:vcate-structural},
the denominator is anchored
by~\eqref{eq:vcate-denominator-anchor}:
$\P\tilde\phi^2 \ge (2\E\tau)^2 > 0$.
Under overlap,
taking $f = 0$ in~\ref{cond:clt-a}
and using $\rho(\hat\mu - \mu) = O_p(1)$
gives condition~\ref{cond:clt-a} when $a > 1$;
condition~\ref{cond:clt-c} holds when $a < 1/\theta$.
\Cref{cor:tuning-window} gives
conditions~\ref{cond:clt-a}--\ref{cond:clt-c}
in the window $1 < a < 1/\theta$.

For cases~\ref{case:vcate-comparison}
and~\ref{case:vcate-full},
the comparison bound
$L(\tilde\phi) \le L(\phi_J)$
for the Fourier truncation of $\griesz^\mu$
gives $\rho(\tilde\phi) \le C\,\eta^{(s'-s)/(2s)}$.
Under overlap,
write $G = \|\griesz^\mu\|_{\LtP}
  \le 2T/\sqrt{\pi_{\min}}$.
The denominator is anchored
by the comparison bound~\eqref{eq:vcate-lower-anchor}:
$\P\tilde\phi^2 \ge G^2/4$
for small $\eta$.
Taking $f = 0$ in~\ref{cond:clt-a}
and using $\rho(\hat\mu - \mu) = O_p(1)$
gives condition~\ref{cond:clt-a} when $a > 1$;
condition~\ref{cond:clt-c} holds when $a\theta < 1$.
The window is
\begin{equation}\label{eq:vcate-lambda-window}
1 < a < 1/\theta,
\end{equation}
nonempty when $\theta < 1$,
equivalently $s > d/2$.

\emph{Functional stability~\ref{cond:an-F}.}
We need
$(\hP - \P)[\psiZ(\hat\mu) - \psiZ(\mu)]
  = (\hP - \P)[\hat\tau^2 - \tau^2]
  = o_p(n^{-1/2})$.
Because $\hat\tau$ is data-dependent,
we estimate $\hat\mu$ on one half of the sample
and evaluate $\hP$ on the other.
Conditional on the estimation sample,
$\hat\tau^2 - \tau^2$ is a fixed function
and the conditional CLT gives
$\sqrt{n}\,(\hP - \P)[g]
  = O_p(\|g\|_{\LtP})$
for any fixed $g \in L_2(\P)$.
It remains to show
$\|\hat\tau^2 - \tau^2\|_{\LtP} \to 0$.

Decompose
$\hat\tau^2 - \tau^2
  = 2\tau(\hat\tau - \tau)
  + (\hat\tau - \tau)^2$.
The first term satisfies
$\|\tau(\hat\tau - \tau)\|_{\LtP}
  \le \|\tau\|_{L_\infty}\,
      \|\hat\tau - \tau\|_{\LtP}
  \le T\,\|\hat\tau - \tau\|_{\LtP}
  \to 0$
by~\eqref{eq:vcate-cate-rate}.
For the second,
$\|(\hat\tau - \tau)^2\|_{\LtP}
  = \|\hat\tau - \tau\|_{L_4}^2
  \le \|\hat\tau - \tau\|_{L_\infty}\,
      \|\hat\tau - \tau\|_{\LtP}$,
so it suffices to control
the sup-norm growth.
The interpolation inequality gives
$\|\hat\tau - \tau\|_{L_\infty}
  \le \kappa\,
      \rho(\hat\tau - \tau)^\theta\,
      \|\hat\tau - \tau\|_{\LtP}^{1-\theta}$,
and therefore
\begin{equation}\label{eq:vcate-L4-bound}
\|\hat\tau - \tau\|_{L_4}^2
  \le \kappa\,
      \rho(\hat\tau - \tau)^\theta\,
      \|\hat\tau - \tau\|_{\LtP}^{2-\theta}.
\end{equation}
The seminorm factor grows:
$\rho(\hat\mu_w - \mu_w)
  \le \rho(\hat\mu_w - \tilde\mu_w)
    + 2\rho(\mu_w)$
by the comparison bound~\eqref{eq:comparison-bound}
$\rho(\tilde\mu_w) \le \rho(\mu_w)$,
and the Bregman convergence bound gives
$\rho^2(\hat\mu_w - \tilde\mu_w)
  \le C\,\eta_\mu^{-(1+\theta)}/n
  = C\, n^{b(1+\theta) - 1}$,
so $\rho(\hat\tau - \tau)^\theta
  \le C\, n^{\theta(b(1+\theta)-1)/2}$
when $b(1+\theta) > 1$.
The $L_2$ factor decays:
by~\eqref{eq:vcate-cate-rate},
$\|\hat\tau - \tau\|_{\LtP}^{2-\theta}
  \le C\, n^{-(2-\theta)(1-b\theta)/2}$.
The product in~\eqref{eq:vcate-L4-bound}
has exponent
\[
\frac{\theta(b(1+\theta) - 1)
  - (2-\theta)(1-b\theta)}{2}
  = \frac{3b\theta - 2}{2},
\]
which is negative
when $b < 2/(3\theta)$.
Because $4 > 3$,
this threshold exceeds $1/(2\theta)$,
so the quadratic term vanishes
throughout the outcome window~\eqref{eq:vcate-outcome-window}.

\emph{Linearization~\ref{cond:an-G}.}
By direct computation,
\begin{align*}
R_{\mathrm{lin}}
&= \P\psiZ(\hat\mu) + \dpsiZ^\mu(\mu - \hat\mu) - \P\psiZ(\mu) \\
&= \P\hat\tau^2
  - 2\,\P\,\tau(\hat\tau - \tau)
  - \P\tau^2 \\
&= \P(\hat\tau - \tau)^2
  = \|\hat\tau - \tau\|_{\LtP}^2.
\end{align*}
By~\eqref{eq:vcate-cate-rate},
this is $o(n^{-1/2})$
when $b$ satisfies~\eqref{eq:vcate-outcome-window}.

\emph{The cost of using $\hat\mu$.}
The penalized projection targets
$\dpsiZ^{\hat\mu}$, not $\dpsiZ^\mu$.
The imbalance~\ref{cond:an-A} decomposes as
\[
\hP\hgamma(\hat\mu - \mu) - \dpsiZ^\mu(\hat\mu - \mu)
  = \underbrace{\hP\hgamma(\hat\mu - \mu)
    - \dpsiZ^{\hat\mu}(\hat\mu - \mu)}_{\text{imbalance at }\hat\mu}
  + \underbrace{(\dpsiZ^{\hat\mu} - \dpsiZ^\mu)
    (\hat\mu - \mu)}_{\text{derivative error}}.
\]
The first term is the imbalance
of the penalized projection
formed at $\hat\mu$,
controlled by the offset complexity
exactly as in~\ref{cond:an-A}.
The second satisfies
\[
|(\dpsiZ^{\hat\mu} - \dpsiZ^\mu)(\hat\mu - \mu)|
  = 2\,|\P(\hat\tau - \tau)^2|
  = 2\,\|\hat\tau - \tau\|_{\LtP}^2
  = 2\, R_{\mathrm{lin}},
\]
which is $o_p(n^{-1/2})$ by~\ref{cond:an-G}.

\emph{Population limit under $\hat\mu$.}
The convergence analysis for~\ref{cond:an-B} and~\ref{cond:an-C}
is stated in terms of the population minimizer
$\tilde\phi$ targeting $\dpsiZ^\mu$.
In practice, $\hat\phi$ targets $\dpsiZ^{\hat\mu}$,
with population limit $\tilde\phi^{\hat\mu}$.
Subtracting the first-order conditions~\eqref{eq:foc-quadratic}
for $\tilde\phi^{\hat\mu}$ and $\tilde\phi$
and testing at $h = \tilde\phi^{\hat\mu} - \tilde\phi$ gives
\[
\P(\tilde\phi^{\hat\mu} - \tilde\phi)^2
  + 2\eta\,\rho^2(\tilde\phi^{\hat\mu} - \tilde\phi)
  = (\dpsiZ^{\hat\mu} - \dpsiZ^\mu)
    (\tilde\phi^{\hat\mu} - \tilde\phi).
\]
The left side is at least
$\P(\tilde\phi^{\hat\mu} - \tilde\phi)^2$.
The right side satisfies
$|(\dpsiZ^{\hat\mu} - \dpsiZ^\mu)(h)|
  = 2\,|\P(\hat\tau - \tau)\{h(1,X) - h(0,X)\}|
  \le C\,\|\hat\tau - \tau\|_{\LtP}\,
      (\P h^2)^{1/2}$
by Cauchy--Schwarz.
Setting $h = \tilde\phi^{\hat\mu} - \tilde\phi$
and dividing by
$(\P(\tilde\phi^{\hat\mu} - \tilde\phi)^2)^{1/2}$
gives
$\P(\tilde\phi^{\hat\mu} - \tilde\phi)^2
  \le C\,\|\hat\tau - \tau\|_{\LtP}^2 \to 0$:
the two population limits coincide
as $\hat\mu \to \mu$.

\emph{Efficiency.}
Under overlap,
$\tilde\phi \to \griesz^\mu$ in $\LtP$
by~\eqref{eq:vcate-lower-anchor},
and \Cref{cor:efficient-variance}
gives $\tilde\sigma^2 \to \sigma^2_{\mathrm{eff}}$.
\end{proof}

\end{document}
